\xchapter{Introduction}

CARDINAL
NEWMAN in his Church of England
days always read his sermons. He discontinued this practice, except on
very special occasions, after his conversion. At both periods he was
following what happened to be the more general custom in the Communion
to which he belonged.

It is not likely that his preaching suffered much
by the change. If it had, he would have gone back to his old practice.
He would have preferred, it is true, to do what most other priests
did, for he hated singularity in any shape, but this was not a
sufficient reason for running the risk of even partial failure before \emph{mixed
congregations} in a
town where he was a stranger, and at a time when from various causes
the no-Popery feeling was particularly strong. Neither does the
change, though he was past middle life when he made it, seem to have
been a difficulty to him. Apparently he soon discovered that the
thoughts that he had in his mind when he entered the pulpit developed
themselves and took new shapes while he was speaking; for the Notes
which are now being published were for the most part written out not
before, but \emph{after},
the sermon.

These Notes were given by the late Father William
Neville, the Cardinal's literary executor, to Father Henry Bellasis,
then a member of the Birmingham, now of the Roman Oratory. 'The Sermon
Notes,' writes Father Bellasis, 'were given me by Father William as a
present one Christmas. The only thing I remember his saying was that
the Cardinal had the practice of going to his room after preaching,
and writing down in the form of notes what he had said. This is how
they come to be in these books. You may remember that the Cardinal
went on preaching in his turn, till he complained that he sometimes
forgot what he was going to say—he must have been eighty-three when
this (not uncommon) complaint attacked him ... However, he felt it at
eighty-three, and then, if you remember, he took to reading some of
his Parochial Sermons from MS.—touched up, I suppose, a little
bit—and very nice it was for those who had never heard him read in
his younger days. I have a recollection of his doing this several
times ... As to his preaching, I was too young to remember much about
it, except that he always had a small Bible in his hands, and quoted a
good deal, always reading the texts, after finding them, and not
(quoting) from memory \footnote{He wished to quote from the \emph{Rheims and
Douay} Version, which was not familiar to him as was the \emph{Authorised
Version}.}. He
used to be rung down by the M.C. \footnote{The Master of Ceremonies at the Sunday High Mass, more often than not
one of the bigger boys of the Oratory School. Needless to say this
functionary used not his discretion but a watch.}, usen't he? I always heard that \emph{he was very obedient to the
bell}, and stopped shortly afterwards.'

'I was too young to remember' is what all who now
survive and used to hear the Cardinal preach, before he had come to
extreme old age, have to repeat. Still it may be worth while setting
down the little that can be recalled. He held the Bible which was in
his hands while he was preaching rather close to his face, for the
print was small and he was short-sighted. Memory pictures him as
constantly turning over its leaves, after the rather fumbling manner
of an old man, while he was speaking, presumably in order to find the
next passage he intended to quote. It is impossible to say whether on
the whole he spoke quickly or slowly, for there was no appearance
either of haste or deliberation. His manner of speaking was the same
in the pulpit as on ordinary occasions; in fact, he was not preaching
but conversing, very thoughtfully and earnestly, but still conversing.
His voice, with its gentleness, the trueness of every note in it, its
haunting tone of (if sadness be too strong a word) patient enduring
and pity, has often been described by those who heard it at St. Mary's
in the old Oxford days, and, judging from their descriptions, it seems
to have been the same in old age as it was then. Probably the
initial impression on one who heard it for the first time would be
that it varied very little. This, however, was certainly not the case.
Changes of expression or feeling were constantly coming over it, but
so naturally and in such perfect unison with what was being said at
the moment, that they were hardly noted at the time. It was only
afterwards, if something had struck home and kept coming back to the
mind, that one realised that it was not the words only, but something
in the tone of the voice in which they were said, that haunted the
memory.

What is the kind of impression that Newman's
sermons were likely to make on a boy or very young man who listened to
them? It would probably not be long before he felt that the preacher
had the power of making things seem very real. He would also be rather
surprised, and perhaps half puzzled, as if it was something a little
incongruous, that a man who seemed so aloof from everyday life should
speak even more plainly and simply than ordinary men. From time to
time he might almost be startled at some change in the preacher's
voice and the words which accompanied it. Take, for example, in the
present volume (p. 301) the way in which the ignominy of the
Crucifixion is described—'as we fix a noxious bird up.' Only those
who have heard Newman can imagine the distress which would have come
over his voice in uttering these words, and the kind of haste, as
if to get them out and done with; followed by a quick return to the
calm with which he had been speaking a few seconds before. Very
probably some of his hearers were, without knowing why, almost as
shocked as if they had now heard of the Crucifixion for the first
time. These little outbreaks came and went as a flash of lightning.
They seemed like a momentary loss of the perfect self-restraint
habitual to the speaker followed by an instant recovery. An
extraordinary thing about them was the very slight change in the voice
which they seemed to entail. It was like a mere breath of wind passing
over the surface of perfectly still water. The Cardinal's voice, as is
well known, was not a strong one. It was of the low and gentle kind,
which in ninety-nine cases out of a hundred tends to be monotonous and
even flat. But in his case the lack of volume or compass capable of
changes on a large scale was made up for by a purity of tone upon
which the faintest modulations told. Those who knew the Cardinal hear
him while they read him, and often a passage, whether in a sermon or a
lecture, comes home to them in a way it can hardly do to others,
because they have the music of his voice in their memory as well as
the printed words before their eyes. Take in the present volume (p.
50) the account of a man's enjoyment of his life. One hears the almost
genial tone which has come over the preacher's rather sad voice
while he is speaking of pleasures in themselves innocent, the summer
stroll, the warm fire, and so forth, then the note of sorrow when he
comes to pleasures not harmless.

Newman's power of entering into minds quite alien
to his own has often been pointed out. 'We,' wrote James Mozley in
1846, 'have often been struck by the keen way in which he enters into
a regular tradesman's vice—avarice, fortune-getting, and so on. This
is not a temper to which we can imagine Mr. Newman ever having felt in
his own mind even the temptation but he understood it ... No man of
business could express it more pungently, more \emph{ex animo}.' From
time to time he would show this power in conversation even when he was
very old, and the voice played its part. One evening he was talking
quietly about the progress of unbelief. He anticipated a time when the
world at large would assume that Christianity had been disproved.
Those who persisted in believing in it would neither be listened to
nor reasoned with. What would be said to them amounted to this: 'It
has been disproved, we cannot disprove it again.' The tone of anger
and impatience he put into his voice just for the moment it took to
say these words, is the reason why a not otherwise remarkable
conversation is remembered by one person who was present, nearly a
quarter of a century afterwards.

Perhaps the most vividly remembered
characteristic of Newman's preaching was his marvellous reading of
Holy Scripture. This was the same in his Anglican sermons, and has so
often been described by those who heard him at St. Mary's that it
would be hardly worth alluding to here except for a reason which will
be presently given. At the Sunday High Mass, when it was his turn to
preach (the boys never heard him in the evening), the Church notices
were usually read for him from the altar steps by one of the other
priests. There must have been some difficulty over them in the past to
make him submit, old though he was, to this. When these were disposed
of, his voice was heard like a soft piece of music from a distance
reading the Epistle. He read it, so far as can be remembered, with
very little variation in his voice, except perhaps a barely
perceptible lingering over the last words, in which they seemed to die
away. There was also in this just a touch of the weariness of old age
which, when recalled in later years, reminds one of the \emph{Expectans
expectavi} of the psalmist. His manner of reading the Gospel was
different. There were of course the pauses required to mark off the
purely narrative portions from the words of different speakers,
accompanied by some slight changes in the voice. But the marked thing,
which cannot be described, was the increased reverence in the reader's
voice which culminated when he came to words of our Saviour.
Before and after these there was a kind of hush. A most wonderful
thing about it all was the complete elimination of the personality of
the reader. He seemed to be listening as much as reading. The words
were the living agent, he but their instrument.

Most striking was the contrast between the humble
pleading way in which he spoke his own words, and the reverence with
which he read such passages of Scripture as he might quote in the
course of his sermon. This contrast must almost necessarily have been
greater in his Catholic than his Anglican sermons, which were written
out beforehand and \emph{read} from the pulpit. It made every sermon a
sermon on the objectivity of Revealed Truth. This is one of the most
ineffaceable impressions left on the mind after nearly all memory of
details has passed away.

Father Neville's statement that these Notes were
taken down after the sermon is borne out by scraps of paper, found
between the leaves of the books in which they are written down,
containing short memoranda, apparently used in the pulpit. But there
were exceptions to the rule, for against the headings of some of the
sermons is written 'not preached yet.' It is quite evident the books
themselves were never taken into the pulpit. They are too large to be
held in the hands, and the writing is much too minute (a
magnifying-glass has often been necessary in order to decipher
it) for them to have been of any use lying on a desk in front of the
preacher. The passages of Scripture referred to in the text of the
Sermons have as a rule been quoted in full in the footnotes. With
regard to the text, the chief difficulty of the editors was in the
case of sermons which the writer seems to have gone over a second time
and added to. He wrote his additions between the lines, or anywhere
where there was a vacant space on the page, without always making it
clear how and where they were to be inserted. It must be remembered
that these Notes were written down by the Cardinal for his own use,
and his own use alone. When, therefore, a subject or some particular
train of thought was familiar to him, a word or two was enough to
remind him of it. If the reader will turn to the notes at the end of
the book, he will find some illustrations of this in the passages
quoted from the Cardinal's published works, where a thought expressed
in the Sermon Notes so briefly as to be hardly intelligible, is found
drawn out at considerable length. The number of these illustrative
passages might have been considerably increased had space permitted.

THE ORATORY,
BIRMINGHAM

\emph{February} 21, 1913.

\xchapter{Sermon Notes}

\xsection{July 1, 1849 Purity and Love—Love in the Innocent and the Penitent}

\footnote{See Note 1, p. 334.}

1. INTROD.—All
the saints of God, all holy persons, all the faithful have, each in
his measure, purity and love—the two graces go together. What is
purity but the having the heart fixed [on], the loving (\emph{al}. an
affection for) things unseen? What turns away the soul from God so
much as impurity? What is impurity but loving \emph{(al}. having affection for) what is sinful? We have such an affection
for money, etc.

2. Yet, though this be so, here as in other
graces, some [saints] are instances of one grace, some of another; and
therefore it may be said without exaggeration that there [are] two
kinds of saints, the saints of purity, the saints of love—the lily
and the rose.

3. St. John the Baptist, sanctified from his
mother's womb, living in the desert. St. John the Evangelist, the
virgin saint—perhaps he never voluntarily yielded to venial sin, and
hence so much favoured—lying in our Lord's bosom. How different these two, yet agreeing in this that they lived \emph{out of} the
world. It is the characteristic of the virgin saints that their love
[is] contemplative, tranquil, etc.—nay, can hardly be called love;
so intimately one with the Supreme that it is with Him rather than it
loves Him. It does not approach towards Him so much as already have
Him; it \emph{is} heaven rather than loves heaven; it is a partaker of
the divine Nature rather than a lover of it. Therefore we talk of such
for their purity.

4. Such above all the Virgin Mother. She
therefore, as coming so near to God, is associated with His titles—\emph{sedes
sapientiae, janua coeli, vita, dulcedo}, etc.

5. But on the other hand, when a soul has given
itself to sin, when it has lost its first estate or never had
grace—when it is in \emph{bonds}, whether to be converted or to be
reclaimed, what is to counteract and antagonise to pride and pleasure?
[to compete] with formed habits? What but a superior attraction? (St.
Augustine in Pentecost \footnote{The reference is to a
passage of St. Augustine [\emph{Tract. In Joann. XXVI}.] read in the
matins of Wednesday in Whit-week. It is quoted in \emph{Sermons to Mixed
Congregations}, p. 70.}.)

Hence \emph{love} is the great instrument of
conversion: (therefore as purity is the emblem of the one [the
Innocent], \emph{i.e}. it shows itself more, takes a prominent
place—so love of the other [the Penitent]), therefore, as the love
of the pure is tranquil, so the love of the penitent is energetic,
zealous, active, belligerent, [full] of emotion, of work, of
passion—the one the love which is of peace, the other which is of
warfare. [Examples of the love of the penitent.] (1) David in the
Psalms; (2) St. Mary Magdalene—her energy, thrusting herself
into the room, tears in the room and in the garden; (3) St.
Peter—loving more than the rest—walking on the sea—John xxi.
(contrast St. John's tranquil 'It is the Lord') rushed
forward—weeping for his denial—crucified head downwards; (4) St.
Paul, 'the life I now live in the flesh,' \footnote{Gal. ii. 20.} 'the love of Christ constraineth us,' \footnote{2 Cor. v.14.} etc.; (5) St. Augustine represented in pictures as loving; (6)
thus in the Confiteor 3, [Our Lady, St. Michael, St. John the Baptist,
the Innocent] then 2, [SS. Peter and Paul—Penitents].

\xsection{July 8 Power of Prayer}

1. INTROD.—If
there is anything which distinguishes religion at all, which is meant
by the very word, it is the power of prayer. Yet wonderful at first
[sight] that prayer should have an effect.

2. The \emph{order} of the world seems to forbid
it—sun rising and setting—everything uniform—laws (causes and
effects) \footnote{Over these words is
written an incomplete sentence: 'and the more you inquire the more ...
'}—winds, indeed,
and atmosphere and sea irregular, but a regular irregularity. All
things cause and effect, bound together like a steam-engine; would
prayer to its Maker make it stop?

3. Hence at all times the wise in this world have
laughed prayer to scorn: they have thought it was a superstition \footnote{This sermon was not
completed. Underneath it is pasted a small scrap of paper with notes,
apparently used in the pulpit.}.

\xsection{July 22 Causes which Keep Men from Catholicity}

\footnote{Note by the writer—'Not
preached yet.'}

1. INTROD.—Catholics
often surprised that every one does not become a Catholic. And they
have a difficulty, how any one can see the Church without
acknowledging it.

2. Now this often arises from \emph{invincible
ignorance}, as almost all would admit; viz. when the reasons for
Catholicity have never been brought home to a man. Born in another
religion—not come across the Church—never been a practical
question. They are Protestant simply from Protestantism being in
possession.

3. Prejudice, violent—yet some danger here that
not invincible ignorance, because they cannot but hear other side in
part.

4. The reason I shall here mention is a main one,
viz. not liking to belong to a body. (1) Contrast of Church of
England—taking a pew, etc.;—they can believe what they please,
they are not bound. (2) This applies in part to dissenting bodies. You
will say 'I grant it in part.' Still a great difference, for the
Catholic Church [is] a body, a society such as no body of men is.

Enumerate particulars.—(1) Not able to believe
what they please, not knowing what they are pledged to; strictness of
confession, the Church having a \emph{hold} upon them. (2) Pride—not
liking to condescend. (3) Human respect—not liking to be laughed
at. It involves outward profession.

5. But this is the \emph{strength} as well as the
\emph{difficulty}, for it is a body or society which has privilege
from Christ. Men come for the sacraments—for pardon of sin—hence
they feel peace on joining it. Communion of saints, of merits, of
indulgences.

\xsection{July 22 On Human Respect}

1. On Magdalene's entrance to our Lord at the
feast—circumstances of it; the way the guests sat—custom of lying
at table. The Pharisee would start from a bad woman coming in \footnote{The writer corrected this
statement in note written in pencil and much of which is illegible:
'No, probably they mixed with them without thinking of converting
them, but despising them; only solicitous they should not touch them.'}. Not strange a woman should come in, but that she should anoint
feet, not head, and weep. She chose a banquet, the very place where
she might have been seen in splendour, not weeping—Maldonatus on
Matt. xxv., p. 286 \footnote{See Note
2, p. 334.}. How
she was able to get in—\emph{ib}., \emph{in loc}. [Luke vii. 38], p.
167; neglectful of self—she did not think what others would say,
because she saw Christ. Just allude to people not going to church with
bad clothes.

2. Now I suppose no more urgent motive than 'what
people will think of us'—fear of the world and human
respect—extends from high to low. The feeling is
tyrannical.—fashion.

3. I am far from meaning that it is bad or
useless; it is not \emph{in itself} bad because it is natural. Nothing
natural bad, for from God, except under circumstances—in excess, not
so as God wills it, or to the exclusion of God \footnote{Written over in pencil
'perfectly becoming we should have a regard for each other—``provide
things honest in the sight of men.'''—Rom. xii. 17.}.

4. Nor useless—the contrary. A number of good
things are done which would not else be [done], and bad avoided which
would not else be [avoided]. We cannot go right without a \emph{feeling
of responsibility}. Now the mass of mankind do not realise God, and
therefore human opinion makes them responsible, and makes them act \emph{rightly}—a
present \emph{visible judgment} useful as far as true. \emph{E.g}. (1)
No public office would go on well without responsibility—abuses
coming in. (2) When upper classes have been shielded from
responsibility, we know what excesses. (3) In colonies and abroad
society sinks down to an immoral level—even ministers of religion.

5. But while we acknowledge it \emph{good on a large
scale, we should be very jealous of it in our own case}. Insist on
this contrast—it \emph{usurps the place} of God—'Loving the praise
of men rather than God' [John xii. 42]. (1) Instances with what I
began with—Magdalene; not going to church with bad clothes—a
little thing, but most expressive; it extends to all communions and
religions, and to all but the higher classes who wear the same clothes
on Sundays and other days. (2) A greater matter by contrast making
propriety of appearance taking [take] place of virtue —unchastity
nothing so that you are not found out; infanticide; Spartan boy
stealing fox; being ridiculed for religion; not daring to obey God;
obliged to take part in bad discourse; ashamed of being known to
pray—I do not wish to be hard upon them—St. Augustine before his
conversion, \emph{pudebat me}, etc.; \emph{St. Alfonso's Sermons}, p.
172.

6. Hence saints have been so set against human
respect. \emph{St. Francis Borgia} carrying a vessel of broth to some
prisoners met his son on horseback, etc., vide \emph{St. Alfonso's
Sermons}, pp. 175-6—hence strange penances of S. Filippo Neri.

7. We should not go out of our way, without
direction, to do strange things, but one thing we \emph{should} aim
at—to substitute the presence of God for dependence on the world.
Act in thought of God. How many things we do in private which it would
be the greatest punishment possible we can conceive to do in public.
Well, act as if God's eye were on you—fear Him more than man, etc.

CONCLUSION.—And,
O my friends, if any Protestants are here, are you sure that \emph{you}
would not become Catholics but for the fear of men? Are you prepared
to say, 'I will follow wherever God leads me. I do not, indeed, see my
way to be a Catholic, but I will become one, in spite of the world, if
I find it is my duty to be one, and if I \emph{suspect} it is my duty,
I will inquire and not give over.' The world passes; in a little time
those only are blessed who, putting aside (\emph{al}. thinking nothing
of) the world's opinion (He that is ashamed of Me and of My words,
etc.) like Magdalene, see Christ alone and gain His favour.

\emph{Note
appended by the preacher to this sermon}

QUERY.—This
sermon would be more complete, if not so short, if confined to the
following:

1. As above (1) circumstances of St. Mary M.
coming to our Lord.

2. On appearing before the Lord in the particular judgment—longing
to see Him though He punishes.

3. Looking on our Lord in benediction which follows the sermon.

\xsection{July 29 On the Gospel for Pentecost IX. [Christ Weeping over Jerusalem]}

1. INTROD.
Wonderful union of mercy and severity in God, as in our Saviour
weeping over Jerusalem. Who can have hard thoughts of Him? Yet who can
presume?

2. Take the case of the Jews. St. Paul says
'these things happened in figure' \footnote{Alluding to the epistle
of the Sunday, 1 Cor. x. 6-13.}: they are a figure of God's dealings with every soul. Consider
the frequent judgments mentioned in the epistle—how often He had to
punish before He gave them up. At last wrath came without remedy—He
wept while He denounced—it was over and there was no hope. He dried
His tears—He rose up—He executed wrath—He rejected them and
burned up their city.

3. And so in every age. Consider what a wonderful
patience—the same thing acted over and over again—how \emph{weary}
the angels must get of the \emph{history} of the world—every
generation beginning with sinners, and then some turning to
repentance—looking at individual souls, seeing them plunge into sin
fearlessly—yet they are afterwards to repent—they must feel
indignation that God should be trifled with. The very same poor souls
who now sin will repent as the generation before them: they even take
their fill of sin before they turn to God.

4. Then they see repentance—all so promising,
such a good start; yet God sees that those very persons, who are
beginning so well, are again to fall from Him—to profane and
ungratefully treat all His gifts.

5. But so the world goes on. Numbers never coming
back to God [at] all, numbers coming back then falling again, numbers
repenting only in the end of life, numbers sinning against light and
warning, again and again, till they are cast off without remedy.
Observe how stern the words, '\emph{for now they are hidden from thine
eyes}.'

6. Yet \emph{how beautiful} the temple looked.
Describe—goodly stones—how unlikely that it should be destroyed,
yet it was doomed.

7. Blessed they who do not sin; next blessed they
who consider God's wrath and mercy—\emph{not one of them} only, lest
they despair or presume.

\xsection{August 5 On Private Judgment}

SUBJECT.—Why
such opposition to Catholicity?

1. Many reasons may be given, but this is one of
the chief: viz. (1) popularly and rhetorically (for they cannot
speak calmly) Catholicity a system which conspires against the peace
and liberty of man; a tyrannical system which imposes a load of things
upon the conscience, which terrifies the weak; anathemas; a grasping
secret system; and so they go on—and worse, till they rise to
priestcraft, Babylon, Antichrist, man of sin, etc., etc.

(2) Really this (for almost all exaggeration is
founded on truth: now the question is what truth there is in this):
viz. that it [Catholicity] is intrusive—interferes between a man and
God. Religion [they say] a private matter; every one has a right to
judge for himself—I am quite able to teach myself; I will never
allow dictation; I am a freeborn Briton; Britons never were slaves,
and such like vulgar swaggering, etc., etc. \footnote{The following sentences
are jotted down in pencil between (1) and (2): 'Let us sift this, for
it is spoken rhetorically. For nothing like prose—all rhetoric [and]
declamation. An inspiring subject and controversial, it always rises
into oratory.'}

2. How much truth? We must sift it still further:
viz. there is interference, but not by an individual, as if the clergy
might bind the laity by their private judgment, but by a system, a
system of laws, etc. This must be cleared up. (1) Mistake to suppose
the Pope can order what he will to be believed. You say, Suppose the
Pope were to say [this or that preposterous thing] we must believe it.
Well, but suppose there were no God, etc.? or again, what would you do
with two and two making five? (2) The confessor cannot do what he
will: (i) a penitent chooses his confessor; (ii) he need not be known to him, \emph{e.g}. extraordinary missions on purpose; (iii) the
confessor goes by rule as a judge does; (iv) the penitent may appeal
to another confessor; (v) the confessor may not speak out of
confession even to the penitent; (vi) a penitent may not introduce the
name of others—detraction; (vii) the confessor knows his penitent
too little, not too much; (viii) two confessors may not talk over a
penitent. Account of Protestant who said the confessor looked like a
God—how absurd! a father and child [is nearer the mark]. How the
penitent may tell, [while] the confessor can't defend himself, can't
employ his knowledge in any way, not even to defend himself against
poisoning; therefore in this individual arbitrary sense, the reproach
against Catholicity, as stated above, not true. Yet it is true that
Catholicity, the Church, interposes between man and his God, teaching
him, warning him, and judging for him. Now I have brought it down to
what is true, and here I join issue with it.

3. Now the Protestant view [of no interference]
is unnatural, irrational, unscriptural, etc. Now to show this—why
Protestants cannot carry it out themselves! If they were consistent
they would not educate their children; yet so eager for education, and
against Catholics educating! If every one should form an opinion for
himself, children should be let alone. Some people have done so. But \emph{you}
take, not a common time, but the first time, and the most
impressionable; it is half, three-fourths of the battle to educate
children [as Protestants virtually admit].

4. But you will say it can't be helped. If we
don't educate, others will—the devil will. Well, [that shows
that] the nature, the state of things, is against you. This is just
what I said, that your way [view of non-interference] was \emph{unnatural}.

If then children, why not grown men?

5. But you will say this is absurd, because they
can judge for themselves. No. No proof that they \emph{can} judge for
themselves, only that they \emph{will}. They can't in worldly matters:
wouldn't it be better if they didn't in worldly matters? Why then in
spiritual? Men are all their lives children as regards religion; they
can't in spiritual [things] judge. You know they can't in business and
cares of life—this is what I meant by irrational—and in matter of
fact they don't. You know they are influenced first by their
education, next by the persons they meet; to say they go by [their own
judgment is] simply absurd—they go by counsels. Why not right of
private judgment in children? Simply because they are weak and depend
on you.

6. Well then, you make every one a prey of every
confident talker, as you do. God has provided the Church to prevent
this very evil.

7. Now think how else you would get out of the
difficulty, viz. the fact that human nature requires guidance, and
will take the first that comes. Does not God's goodness point to a
church?

8. O my brethren, which is the more scriptural?
One church in heaven and earth—the saints—souls in
purgatory—communication of merits—each depends on each—hand and
foot—(explain). This why Englishmen so unamiable—coldness.

\xsection{August 19 On the Fitness of Our Lady's Assumption}

1. Recollect Luke xxiv. 26, 'Ought not Christ,'
etc.; Hebrews ii. 10, 'It became him,' etc.; Rom. xii., 'Analogy of
Faith.' \footnote{See Note
3, p. 334.}

2. In like manner it \emph{became} our Lord to
raise His mother, and her so sinless. Let us think of this.

3. Doctrine from the first—that God was her
Son, lay in her womb, was suckled by her, etc. She enjoyed His voice,
smile, etc.

4. Esther vi. 6, 'What should be done to the man
whom the king is desirous to honour?' 'He ought to be clothed in royal
robes,' etc. And so of our Lady—she should be the Mirror of Justice,
the Mystical Rose, etc. Thus has King Solomon risen up to meet his
mother.

5. Now go into details—sanctity and spiritual
office or work go commonly together: (1) the angels; (2) seraphim.
Exceptions: (i) Balaam, Caiaphas, overruled; (ii) many shall say in
that day, 'Lord, Lord, have we not prophesied in Thy Name,' and He
shall answer, 'I never knew you'; (3) They may have fallen away; (4) \emph{gifts}
are [imparted] separately from sanctity, but gifts are not offices. On
the contrary (5) Enoch, Noe, Moses, Samuel, David, etc.—except
Judas. [These are instances where, according to the general rule,
'sanctity and spiritual office go together.']

6. If such to whom the word was \emph{made}, much
more [does sanctity go with spiritual office] in [her, in] whom He was
born. Was it not fitting? Do human parents otherwise? Do they give
their children to suckle to common persons? \emph{Nature} says that
the \emph{fount} of truth should be \emph{holy}. Here is the
difference [as between] miracles and sacraments \footnote{Speaking generally, a
miracle is a testimony to some truth. Its worker therefore is a \emph{fount
of truth}, and therefore presumably holy. A like presumption does
not exist in the case of the administrator of a sacrament.}. Prophets receive, beget, and bear the Divine Word.
Scripture-writers different from each other, and so the Fathers \footnote{The fact that the
inspired writer, or the Father, preserves his individuality, shows
that he is something else than 'a mere instrument.'}. As the tree, so the fruit. 'Beware of false prophets' [says
our Lord, and then He adds, 'from their fruits ye shall know them'].
Mary not a mere \emph{instrument}—as the first-fruit, so the mass.
The Word did not pass through [her]. He \emph{took} a body from her,
therefore she was worthy of the Creator—\emph{full of grace}.

7. Hence doctrine of Immaculate
Conception—grace before Gabriel [\emph{i.e}. before the
Annunciation], before vow of virginity, before Temple, before birth,
before St. John the Baptist [\emph{i.e}. at an earlier period in her
existence than that in which it was bestowed upon the Baptist]. \emph{She
must surpass all saints}.

8. Again with her \emph{co-operation}—this her
merit. She was peccable; she \emph{grew} in grace, etc. Enoch merited,
Noe merited, Abraham merited, Levites merited, David, Daniel
[merited]; how much more Mary since her reward was such!

9. Her glories were not simply from her being
Mother of God: \emph{it} implies something before it. The feast of
Annunciation implies feast of Conception and of Assumption.

10. Come then (I would not weary you) to
Assumption—more difficult \emph{not} to believe Assumption than to
believe it after the Incarnation. Human Sons sustain their mothers.
She died; she saw no corruption, for she had no original sin. 'Dust
thou art,' etc.

11. Therefore she died in private. Give history
of Assumption.

12. What is \emph{it} fitting that we should be
with such a mother?

\xsection{August 28 On Want of Faith}

\footnote{See Note
4, p. 335.}

1. INTROD.—Many,
as I said lately, wonder the beauty of Catholicity does not attract
multitudes of people 'to see that great sight' \footnote{Exod. iii. 3.} 'come and see' \footnote{John i. 39.}—and
join the Church.

2. They do not become Catholics because they have
not faith. This no truism; faith is a certain faculty, like justice,
etc.

3. Describe faith—assent on the word of
another—as coming from God, not by sight or reason, not as word of
man, which is a kind of faith, but not firm. We take man's word for
what it is worth, but [divine] faith is most certain.

4. This was faith in the apostles'
time—certainty. The converts entered the Church in order to
learn. Could they have disputed with an apostle? or separated from
them? or believed them not inwardly? or wait for further proof? No,
there was no private judgment then; they either believed or did not
believe. They could not say, I will choose, I will believe a little,
as I please, etc., etc.

5. And this is plain from Scripture texts—1
Thess. ii. 13, \emph{ib}. iv. 8, Luke x. 16—'know most certainly' \footnote{Acts ii. 36.} or 'believe and be saved,' the 'word of hearing.' \footnote{1 Thess. ii. 13.}

6. What a contrast between this or deducing from
a book to master it! The one is submitting, the other is judging.
Faith, then, in the apostles' age consisted in submitting.

7. Plainly not a temper of the world now; they
have not what the apostles meant by faith. Men change to and fro now:
this is opinion, not faith.

8. Again they laugh at faith as servile, the work
of priestcraft. Would they not have rejected the apostles? Would they
not have died pagans?

9. Object—the pagan did so laugh—quote 1 Cor.
i. 23 \footnote{' ... Christ crucified
... unto the Gentiles foolishness.'}; Matt. xi. 25 \footnote{' ... Thou hast hid
these things from the wise and prudent.'}.

10. They have not faith. How then do they believe
the Scriptures? They don't; it is a nursery habit. When they think of
their contents they begin to doubt.

11. What faith [was] in the apostles' days [it
is] now. I have proved, then, the world has not faith. Though many men
may admire, may encourage, yet they will say, 'O that we were
Catholics!' and get no further.

12. Deplorable state! for faith so necessary; \emph{e.g}.
Heb. xi. 6 \footnote{'Without faith it is
impossible to please God.'}; Mark xvi.
16 \footnote{'He that believeth and
is baptized shall be saved.'}; John iii. 18 \footnote{'He that doth not
believe is already judged.'}. Christ might have saved by sight, but He saves by faith.

13. Let them try to put faith elsewhere if they
can. Faith through grace.

14. Exhortation to Protestants to use existing
grace.

\xsection{September 2 Prejudice as a Cause why Men are not Catholics}

1. INTROD.—It
may surprise those Catholics who live to themselves why so many
Protestants are not Catholics, but it will not surprise those who go
into the world. They will there find what will account for it, viz. a
prejudice about Catholicism such that the wonder is not why men do not
become Catholics, but why any do at all.

2. Such is the power of prejudice. What is
prejudice? It is forming a judgment without sufficient grounds. We
cannot help being prejudiced, because there are ten thousand things
about which we can but have an opinion; but the fault is—and this is
what we truly mean by prejudice—when we stick to it in spite of
better information, or will not listen to other information.

3. Now as regards Catholicism. Men in childhood
have always heard Catholics abused. They are considered to be
cruel; stories of torments inflicted by them are circulated: (1)
exaggerated stories of individuals; (2) a horror of some particular
practice; (3) viewing things separately, not as a whole. One thing
fixes in their mind, and seems to justify all anticipations, for
persons take fright at some one particular doctrine or fact, \emph{e.g}.
about the confessional, and at one bad priest, etc., etc. This
especially inconsistent in those who profess to go by private
judgment. Why do they go by what they hear, and that casually perhaps?

4. Not believing [that] (1) priests believe what
they say; (2) are continent; (3) [that] converts are
satisfied—looking out for some change in them. The consequence of
this deep prejudice is that from the nature of the case there are no
ways of overcoming it. If Catholics are particular, devout, or
charitable, etc., they are said to be hypocrites; if all things
apparently simple, they think there is something in the background;
they call them plausible; if nothing can be found against them, how
well they conceal things; if they argue well, what clever sophists; if
charitable, they have vast wealth; if they succeed, not of God's
blessing, but of craft. I wish we had half the cleverness they impute
to us.

5. Hence they circulate lies about us, not
inquiring the authority, and when they are disproved, instead of
giving over, circulate others which can't be. When any particular lie
is put out, they embrace it at once as being so likely, \emph{i.e}.
like their prejudice. They take not this age and place, but a
thousand miles away and two hundred years ago. Catholics alone can
suffer this, because they are in all times and places; they could not,
\emph{e.g}., treat Quakers so. Explain how far honest. They say to
themselves, if this is not true, yet something else is true quite as
bad.

6. But this cannot last, \emph{i.e}. prejudice
(whether they become Catholics or not); as the ice goes in Canadian
sea in a week, so when prejudice once begins to thaw it will go
quickly. Now the remedy for all this is to see us—(enlarge on this).
They cannot keep up their theories against us, but they are afraid to
be puzzled with something on our side. They have a sort of feeling
that if they were to see us we should contradict their prejudices, so
they do all they can to keep us out of sight.

7. Hence no person hardly who has been much
abroad and lived with the people can keep up their prejudices; no one
who has read much history: the strength of prejudice is with those who
are not informed.

\xsection{September 30 Faith and Doubt}

\footnote{See Note
5, p. 335.}

1. INTROD.—Those
who are curious ask, 'May I doubt when I am a Catholic?' Those who
object say it is a tyranny, violence on the mind, immoral, etc.; \emph{i.e}.
they ought to hold that it is wrong to make up the mind on any
religious subject whatever, however sacred. A liberty to doubt [is
what they ask].

2. First, doubt is incompatible with faith. Who
would say a man believed the apostles' mission who added that perhaps
he should one day doubt about it? A real but latent doubt. 'I perhaps
am excited,' or 'in a delusion,' or 'everything may turn out.' What
men object to is \emph{faith}. If the thing is \emph{true}—that God
became man—why must it be \emph{doubted}? Either they have faith
now, and then, etc., or not, and then, etc. I may love and obey by
halves, not believe.

3. And so when a Catholic [doubts] he has already
lost it (\emph{i.e}. faith). Persons converted to Protestantism by
reading the Bible. Protestants only show by their objections that they
do not know what faith is.

4. Secondly, love [rejects doubt] as well as
faith. What would you think of a friend who bargained not to trust
you? who said he should be trifling with truth if he did not [so
bargain]? May I never have such a friend—jealous minds, etc.—[Give
me] cordial, openhearted ones, etc. And so of God. If a man thought
God might be unjust, or bargained to believe in, worship Him, only
while his reason told him, he would be worshipping himself. And so of
the Church. Fetters! [Yes,] cords of Adam.

5. The world thinks faith a \emph{burden}—cannot
understand joy of believing. [It imagines] confession [to be] chiefly
of doubts. On the contrary, \emph{popule meus, quid feci tibi}, etc.

6. Third view. (Doubt does not destroy
intellectual conviction, but) faith a gift of God. I may see I ought
to believe, yet cannot. Conviction and acting, conviction and
faith—faith not of necessity, [\emph{i.e}. the will not
coerced] but of will—merit in faith.

7. Conviction may remain without faith. If we
listen to objections without cause—case of those who fall
away—they cannot answer arguments. Thus they either linger about the
Church or go into atheism.

8. Fourth reason [why doubt is not permissible].
Inconsistent in the Church to do otherwise. St. Paul, St. John,
Eliseus \footnote{Eliseus prohibiting the
search for Elias, 4 Kings ii. 16.}.

9. No other body can demand faith—not
Dissenters, not Church of England.

10. Be sure, before you join us. You must come to
learn. Do not distress yourselves \emph{whether} your faith will last.

11. Get \emph{conviction}. Act when it comes; it
comes differently to different persons. We are anxious about persons,
not as wishing them to act without conviction, but because perhaps the
time is past.

12. Oh the misery if you \emph{have} not become
Catholics! Oh the misery if we \emph{had} not!

\xsection{October 14 Maternity of Mary}

\footnote{See Note
6, p. 336}

1. INTROD.—When
we look upon earth and sky we find everything connected together in a
wonderful way—everything answering to each. Nothing could be
altered—if there were one star less or more—and so of animal
power—atmosphere and sea. The like happens in the Catholic
religion, and the more a person examines the more he will find it,
though people have no time for examination. But concerning the
doctrines of our Lady, apropos of the Maternity. (It seems fair to
say, that if God would restore the world, she must be without sin. Yet
[again] one truth follows from another.)

2. When God intended to take flesh, He might have
taken a body like Adam or Eve (or from the sky), but either men would
not have believed He was man, or not so readily. The notion of God
becoming man is so hard that the human mind will evade it if it can.
Hence, in order to \emph{seal the doctrine}, He took a human mother.

3. Hence the great doctrine that Mary is His
mother: the Mother of God has ever been the bulwark of our Lord's
divinity. And it is that which heretics have ever opposed, for it is
the great witness that the doctrine of God being man is true. The
making \emph{much} of, the \emph{prominence} of the doctrine is the
bulwark, hence she had her gifts—(1) to erect her as a Turris
Davidica, lest she should be forgotten; (2) to prepare her fitly, as a
temple (no unclean thing can enter heaven); (3) lest she should be
puffed up as Satan. She said, '\emph{Ecce ancilla Domini}'; thus she
ministers as a creature, and does glory to God.

4. The \emph{first} mode heretics took was saying
that our Lord's body came down from heaven, or that He was an apparent
man, etc. They affected to be reverent, and said the idea was shocking
that a born man should be God. Now you see how the doctrine excludes \emph{this}
idea; hence, in the creed, 'born of the Virgin Mary.'

The \emph{second} mode was to say that Mary bore
a man, and not God—mother of our Lord's manhood, etc.—but that God
was in a particular way in Him, as He was in the prophets and good
men.

The Council of Ephesus, about four hundred years
after Christ, [decreed] the title of \emph{Mother of God}.

The \emph{third} ground was at the
Reformation—bolder—that it [\emph{i.e}. Catholic teaching about
our Lady] was idolatry, etc., Satan hoping so to destroy the belief in
our Lord's divinity. Here again false reverence, so they abolished the
honour of our Lady out of tenderness to Christ's divinity! Look at the
issue. The truth is, the doctrine of our Lady keeps us from a
dreaming, unreal way. If no mother, no history, how did He come here,
etc? He is from heaven. It startles us and makes us think what we say
when we say Christ is God; not merely like God, inhabited by, sent by
God, but really God; so really, that she is the mother of God because
His mother.

\emph{Fourth}, the tendency of this age to
depress man, fancying the stars inhabited, etc. Why, then, should God
think of us? why should His son be incarnate? Not so much meant. Now
this doctrine \emph{fixes} that so much is meant, \emph{coelum animatum}.

\xsection{November 11 Purgatory}

1. INTROD.—Not
wonderful if God, who might have saved us without, yet saved us with
Christ's passion—though He might have saved us without pain, [yet]
saves us with pain.

2. Hence 'through many tribulations we must
enter,' etc., either in this life or next. As Christ [suffered]
without sin, so we for our own sins. Suffering in next life is in
purgatory.

3. Now I shall best describe purgatory by first,
'He descended into hell.' What is meant by hell? Plainly the place to
which \emph{souls} go—as His body to the grave, so His soul to hell.
So 'thou shalt not leave My soul in hell,' etc., Ps. xv. 10. Yet hell
cannot mean the home of the devil, therefore it means something short
of hell, though called hell.

4. Remarkable it should be so called. It cannot
be called so without reason. A joyful place would not be called hell.
Yet so also in the Mass, etc., 'de porta inferi,' 'de ore leonis et de
profundo lacu,' '\emph{ne absorbeat}eum'; and so again Phil. ii. 10 \footnote{'That in the name of
Jesus every knee should bow, of those that are in heaven, on earth and
under the earth.'} and Rom. x. 6-7 \footnote{'Say not in thy heart,
Who shall ascend into heaven? that is, to bring Christ down. Or, Who
shall descend into the deep? that is, to bring up Christ again from
the dead.'},
and Samuel, 'from the earth.'

5. Evidently then near hell, or in some respects
like hell—\emph{absorbeat}, like a whirlpool. Such is purgatory, and
it is not wonderful that it should be a place of great punishment.
Hence it is that, being near hell, the holy fathers say the flames are
hell flames, like being scorched by a house on fire.

6. Still at least it cannot be a happy place, for
it \emph{is not heaven}. Pain of loss—(describe). God our good. We
manage in this life to lean on creatures—our friends, etc.; our
comforts, etc. The soul downcast, miserable, dreary, as being \emph{hungry},
like the feeling of sinking—fainting to the body. Such is
purgatory, at worst flame, at best and always, desolation.

7. Different mansions in purgatory (as in
heaven). This shown in the vision of St. Felicitas, and of St.
Malachi.

8. Hence a received, or at least a pious, opinion
that there is a region where there is no pain of sense at all. St.
Bede speaks, as St. Felicitas suggests, of a meadow. St. Gertrude; St.
Mechtilda. Such the place of the old patriarchs such as Samuel. Hence
Abraham's bosom, though in sight of hell—and, 'with me in
paradise'—a garden.

9. But since it is the place which, with all our
penances and satisfactions here, we cannot escape from, I will add
some consolations. First, they \emph{do not sin}—no ruffling or
impatience; they are the \emph{holy} souls in purgatory. (1) They hate
their sin so much that they have greater pleasure in suffering than in
not suffering with the feeling of sin. (2) No impatience; they \emph{will}
to suffer, for it is God's will. Thus every consolation—full
resignation. A holy soul plunging into the place where it sent
itself—rather feeling the pains of hell than the least sin.

10. Secondly, they \emph{know} they have done
with sin; it was a phantom haunting them all through life, night and
morning. Weariness—all at an end. Resignation in the storm; ecstatic
feeling; nay, he rejoices to combat with the antagonist trial.

11. Assurance of salvation, as they know that
each hour brings them nearer to the end.

12. Consoled by angels, etc. St. Francis de
Sales. Transcendental state; not single bliss and single pain;
nor mixed as in this life, but both together, pure and antagonistic \footnote{See Note
7, p. 336.}.

\xsubsection{St. Francis de Sales on Purgatory }

It is true that its torments are so great that
the most extreme pains of this life cannot be compared with it; yet,
on the other hand, the internal satisfactions there are such that
there is no prosperity or contentment on earth which can equal them.
(1) The souls there are in a continual union with God. (2) They are
perfectly resigned to His will, or, to speak more exactly, their will
is so transformed into His will that they cannot will otherwise than
God wills; so that if paradise were opened to them, they would rather
precipitate themselves into hell than appear before God with the
defilements which they still recognise in themselves. (3) They are
purified there voluntarily and lovingly, since such is the divine good
pleasure. (4) They wish to be there in the manner which pleases God,
and for so long as pleases Him. (5) They cannot sin, and cannot
experience the least motion of impatience, nor commit the least
imperfection. (6) They love God more than themselves, or anything,
with a full love, pure and disinterested. (7) They are there consoled
by angels. (8) They are there assured of their salvation, with a hope
which cannot be confounded in its expectation. (9) Their bitterness,
most bitter as it is, is in the midst of peace most profound. (10)
Though purgatory be a sort of hell as regards the pain, yet it is
a paradise as regards the sweetness which charity spreads abroad in
the heart; charity more strong than death, more powerful than hell,
the lights of which are all fire and flame. (11) Happy state, more
desirable than formidable, since the flames are flames of love and
charity. (12) Formidable nevertheless, since they retard the end of
all consummation, which consists in seeing God and loving Him, and by
that sight and that love, in praising and glorifying Him through the
whole extent of eternity.

\xsection{December 9 On Man as Disobedient by Sin as Contrasted with Mary}

1. INTROD.—Our
Saviour came at this time of year to bring peace on earth.

2. Prince of peace—leopard and lamb \footnote{'The leopard shall lie
down with the \emph{kid}.'—Isaias xi. 6.}—'on earth peace'; hence, as type, peace in Roman Empire.

3. He reconciled man to man, God to man, but
especially the soul to itself. He made peace within—this the great
gift.

4. Man created at unity with himself; his
different powers, irascible and concupiscible—how are they to be
brought together? Only by God's grace. He is not sufficient for his
own happiness. Stoics have tried to subject the passions to the
reason, without subjecting the reason to God. Sin is self-destructive.

5. Such the case, about eternal punishment—it
is not religion brings in the doctrine; it is a fact in prospect
before us—for suppose no God, and man immortal, he would be his own
eternal torment, and could not free himself.

6. Give a person riches, health, name, power,
ability, let him live centuries here, would that be a gain, or the
contrary? Would not the very time show that these things had failed?

7. Two great principles, the irascible and the
concupiscible. Solomon in Eccles. ii.—indulgence of sensuality; what
does mirth and grasping profit? Tired out—sated—the same dishes
daily; the same faces; the same servants behind chairs—Lord
Byron—the man who killed himself because he had to get up and go to
bed. When such men get to the end of life they would not live longer;
they want rest, as the man in 'The Siege of Corinth'—'The Giaour.'

8. Satiety would make way for gloom—ill-temper.
The misery of ill-temper—gloom; eating the heart out. On kings with
unrestrained power, what brutes they become! Their furious passions.
Youth is gay, age is crabbed—vain regret of first youthful feelings,
gone for ever. Why, such feelings would tend to madness. Oh the awful
misery of a man living an eternity in this world!

9. Yet they do not live on, but die. And then,
what the additional agony of a soul left to itself! with nothing
corporeal; no means of communicating with others; thrown on itself;
voluntarily cut off from God, who is our only stay, comfort,
then—and so for eternity.

10. Pain of the body great, but pain of the mind
worse, though we do not know much of it here. Scaring, bad
dreams, hair turning white—what when it comes in its fulness? The
wicked is like the troubled sea. Here is your portion, my brethren, if
you will not turn to God.

11. Oh, what dreadful thoughts for the future!
This is how man, then, will appear before his Maker—covered with
wounds, etc. Suppose at the judgment God, without positive infliction,
merely left a man to himself.

12. What a contrast our Lady to this—our
Saviour is God and cannot afford the contrast. Immaculate in her
conception—so sweet, so musical, etc. She holds up to us what man is
intended to be, as a type, the most perfect submission of his powers
to grace.

13. Instinctive feeling in the Church that it is
so.

14. Christ the source, Mary the work of grace.

\xsection{December 16 On the Last Times of the World}

1. INTROD.—Two
Advents of Christ.

2. The difference between them: (1) the latter
sudden; the former, a long course of preparation, so that He could not
have come sooner than He did—the latter, hardly any
preparation—Antichrist alone—else it may come any day.

3. (2) An apostasy before the second—quote 2
Thess. ii. 3-4 \footnote{'Let no man deceive you
by any means, for ... unless there come a revolt first, and that man
of sin be revealed, the son of perdition; who opposeth and is lifted
up above all that is called God,' etc.}; yet
this no infringement on its suddenness, for the apostasy began
working even in apostolic times \footnote{Cf. 2 Thess. ii. 7; 1
John iv. 3.}.

4. On the contrary, since it is always working,
the contemplation of it may be useful to us.

5. \emph{Characteristic} of the apostasy—\emph{not
idolatry}, not presumption, as 'the Temple of the Lord,' etc. \footnote{'Trust ye not in lying
words, saying, The temple of the Lord, The temple of the Lord, It is
the temple of the Lord.'—Jer. vii. 4.} Not despair, as 'why should I wait on the Lord,' 2 Kings vi.
33, but infidelity—quote 2 Thess. ii. 4, 'Shewing himself,' etc. \footnote{'Shewing himself as if
he were God.'}

6. Particular sins have \emph{particular}
punishments, as fire for Sodom and Gomorrah. (1) Parallel of flood,
first destruction of the world: (i) An \emph{apostasy}—\emph{filii Dei
ad filias hominum}; a new state of things followed; a sort of
perfection—\emph{viri famosi}; (ii) St. Peter called it the world of
the impious; (iii) St. Paul, by \emph{faith} Enoch and Noe \emph{endured
the world}; (iv) St. Jude, Enoch's prophecy against the \emph{impious};
(v) remarkable; Tubal-cain and Jubal—useful and fine arts \footnote{'Jubal the father of
them that play upon the harp and the organs ... Tubal-cain ... a
hammerer and artificer in every work of brass and iron.'—Gen. iv.
21, 22.}; and so [\emph{anomos}], 2 Thess. ii., with \emph{iniquitas},
Gen. vi. 13 \footnote{'The end of all flesh is
come before me; for the earth is filled with iniquity through them;
... I will destroy them with the earth.'}.

7. (2) Description of the last apostasy in the
New Testament: (i) St. Paul, 'depart (apostatise) from the \emph{faith}';
(ii) St. Paul, 'wax worse and worse' ; (iii) St. Jude, 'mockers,' etc.
A still more remarkable passage, 2 Peter ii. 4-9 \footnote{'For if God spared not
the angels that sinned, but delivered them, drawn down by infernal
ropes to the lower hell unto torments, to be reserved unto judgment.
And spared not the original world, but preserved Noe, the eighth
person, the preacher of justice, bringing in the flood upon the world
for the ungodly ... The Lord knoweth how to deliver the godly from
temptation, but to reserve the unjust unto the day of judgment to be
tormented.'}, where the state of antediluvian and last days (unbelief) are
connected. They thought nature must go on as hitherto,—'Where is the
promise of the coming?' etc. Nature all-sufficient, all in all, that
it should come to nought—\emph{an idle tale}.

8. This further illustrated by the miracles of
Antichrist, in whom the apostasy will terminate, 2 Thess. ii., Apoc.
xiii. 13. Now the devil cannot do \emph{real} miracles, therefore they
are miracles of knowledge. Knowledge is power—parallel of Tubal-cain
above—and they say power is but knowledge, \emph{i.e}. the revealed
miracles are not real ones.

9. Hence so plausible, that even the elect might
be deceived by the sophism.

10. Such the apostasy, and while it is brought
before us by the season, it concerns us because St. Paul says, 'It
already worketh.' It is in all ages, and surely not the least in
this—open infidelity, specious objections, various kinds of argument
from long ages, geology, history of civilisation, antiquities, etc.

11. You may say it doesn't concern us; it
does—specious objections. But let us ask our \emph{hearts}, do not
they speak for religion in spite of these?

12. It is all founded on \emph{pride}. Pride is
dependence on \emph{nature} without grace, thinking the \emph{supernatural}
impossible. Eating the forbidden fruit was pride and unbelief; thus
the world will end with the sin with which it begins.

Continue

Back

\xsection{January 6, 1850 On the Catholic Church}

1. INTROD.—Today
the birthday of the Catholic Church, for the Gentiles came to it.

2. From eternity in the councils of God. At
length in time it began to be; it was conceived and lay in the womb.
Its vital principle faith, therefore with Abraham especially it began.
It remained in the womb of former dispensations its due time; long
expectations; burstings of hope, till the time came; and was born when
Christ came.

3. In the fulness and consummation of time. OBJECTION.—Why
so late? True answer, because unmerited. God may choose His time and
place. Again, because He had to \emph{work through} human wills, and
therefore, so to say, \emph{under} the present order \emph{could} not
choose His time. But here I say fulness and consummation of time, \emph{i.e}.
man is born after months in the womb. He is born in due time, not an
abortion. So of the church.

4. When born, a robust and perfect offspring,
fulfilling its promise—its promise that it was to be everywhere, and
was \emph{to be able} to be everywhere.

5. \emph{Able to be}, for this is the difficulty
which no other religion ever attempted. None but the Catholic has \emph{been
able} to be everywhere. \emph{Local} religions—whether
Eastern mythologies or Protestantism.

6. But even earthly empires do not spread over
the world so \emph{widely} \emph{as} this and so diversely—now from
east to west, now from north to south. Mahomet by the sword.

7. But even empires of this world gained by the \emph{sword}
do not \emph{last}. Not only is this a single religious empire, but it
has lasted out earthly empires, and now shows as little decay as ever.

8. And in such tumults—the whole world broken
up so many times—present revolutions nothing to former. The deluge;
describe waters—whirlpools, waterspouts, currents, rush of waters,
cataracts, waves, yet the ark on them. This, the ark, \emph{the greatest
of miracles}. Well, it is but the acknowledged type of the Church:
as this was the miracle (as we all confess) of the deluge, such that
morally of the Church.

9. A house \emph{not} divided against itself \emph{does}
stand—other religions specimens of the reverse. House at this moment
less divided than ever. Protestants have looked: they felt the
question was, whether we were in extremities? not whether the Pope was
alive, but whether nations acknowledged him? (1) No jealousy about
Pope's power. Pope never so powerful as now—perfectly good
understanding; jealousy at an end. (2) No heresies now. (3) Nay, \emph{schools}
at an end, [\emph{e.g}.] Immaculate Conception \footnote{All opposition to this doctrine had
disappeared though it was not yet (in 1850) defined. It no longer
divided the \emph{schools}.}.

\xsection{January 27 On Labour and Rest}

1. INTROD.—On
Septuagesima, beginning a time of penance and penitential work. No
more Alleluias. The colour purple.

2. Labour is the lot, the punishment of man. Bad
and good labour, nay, evil labours and virtue labours.

3. It is otherwise as God made things. There is
motion and activity in Nature, but it is \emph{without effort}; all
creation is as it were hung upon wheels, and moves noiselessly and
gracefully—the sun, the stream, the breeze, life.

4. And so in paradise. Adam's tending the flowers
was but a specimen of divine labour without effort; such, too, was his
service of God; such the angels' service—without effort.

5. But sin has made things otherwise. Henceforth
labour changes its character. It is no longer Eden, but that vineyard
into which the labourers were sent in today's gospel—to pull out
stones, to destroy the weeds, worms, blights—and a wall round
it—for there is a warfare. Labour is a war and aims at conquest.

6. Take bodily labour, labour of the
field—preparing the earth, felling trees, making roads,
canals—then building houses; it is all penitential, all the
punishment of sin—the mind does not come in, but a weariness.

7. And much more with \emph{intellectual} labour
and the labour of the mind—the mere wear and tear of business;
the necessity of providing for a family; anxiety, suspense, fear,
failure, dreariness and hopelessness. But even when successful, one
enterprise leads to another, till the mind is overburdened and
overwrought, and is sucked into a vortex. Most engrossing; no time for
the thought of religion; religion must take its chance, and that they
feel.

8. Much more sin; the bondage and service of the
devil most wearisome—the drunkard, the sensualist. (I knew one who
was tempted to fatalism.) Wearing, restless feeling, even when they
call themselves happy.

9. Nay, virtue here is too a toil, because there
is war between good and evil. Read the saints' lives. Such is labour,
and it wearies soul and body. The body shows it. whether it is manual,
mental, or intellectual.

10. Oh! if we must labour, let us labour in the
service of the Great Master of the Vineyard—that only pays, that
only has hire. Then we shall labour that we may rest, then only. Sin
never rests; there is no rest in hell. This is that penny which they
one and all received, because nothing better or higher.

11. When the evening of life comes, then shall we
know most fully the meaning of labour by being freed from it.

12. The blessedness of rest, of freedom from sin
and toil, even though in purgatory. Purgatory is rest compared with
this life.

13. And much more in heaven, where we see the
face of God.

\xsection{February 24 On Grace, the Principle of Eternal Life}

1. INTROD.—God,
who had been the sole life from eternity, is the life of all things.
He did not lose His prerogative or give to others or creation what He
is Himself.

2. Nothing lives without Him; nothing is.
Animated nature, vegetables, nay, the very material substances, have
their life, if it may be so called, their motion and activity in
Him—the elements. What is called Nature, a principle of life, is
from Him.

3. Moreover, the life He has given to Nature is
but transient and fleeting. It is beautiful while it lasts, but it
comes to an end. Nay, it is self-destructive; thus the water and the
fire, which are the conservation, have been and shall be the
destruction of the earth. And so growth tends to decay. It is the same
process; all things grow to an end.

4. Thus this earth, as I have said, will be
consumed. Thus the year, too, comes to an end—how beautiful spring,
yet it doesn't last. The year runs a reckless course, like a
spendthrift; it cannot help going on till it is nothing. So it is with
bodily health—'dust thou art,' etc. We see it again in animals,
which are sportive and playful when young, but get old and miserable
and sullen. Thus in Nature the best is first.

5. Nature, then, has no immortal principle in it.
All natural things run a course; and this is true of the soul, of the
natural soul. The soul as it is by Nature, by original creation,
has no principle of permanent life in it. The soul grows old as
anything else.

6. Describe the engaging manners of the
young—fascinating, light-heartedness, cheerfulness; affections warm;
imagination, conversation, wit; all pain shaken off—what can be
better? Why is not Nature enough? Wait awhile.

7. Wait awhile, for the soul grows old as
anything else—as the leaves turn yellow, as the animal frame grows
stiff, so wait on a few years, the natural soul too grows old; the
beauty decays as beauty of person; the soul contracts, stiffens,
hardens, instead of being supple and versatile, and elastic and
vigorous; its limbs are cramped; everything is a burden; it is a fear
to it to be pulled out into new positions; it cannot take pleasure in
what once pleased—not in poetry or works of fiction, not in
friendship; it cannot form new friends; it is bereaved [of the old
ones], and does not replace them; it cannot laugh; disappointment
breaks it; it cannot recover. Hence relapsing into natural
imperfections (as crabbedness, ill-nature, etc.) which a man had
seemed to overcome, having ever struggled against them.

8. Oh terrible! old people hard-hearted, without
affections, careless of the loss of friends—not from high
motives—they have no faith—virtue seems a fancy—with hearts like
stone, etc.

9. Follow such a one into the next world. What is
to be his happiness for eternity?—immortal, yet dead, eternal death.
Life of the soul is in the affections; he has no affections—a closed
heart. The devil cannot love God—\emph{vide} St. Catherine of
Genoa in \emph{St. Alfonso's Sermons}, p. 335.

10. Such is the course of Nature in the soul as
in the body. Nature ages; it has in it no principle of life. No, grace
is the only principle of immortality. We must go beyond nature; we
must go to something higher. Here, then, is one characteristic
difference between Nature and Grace.

11. EXHORTATION,
Eccles. xii. All saints have lived by love—the martyrs, confessors,
etc., etc. St. Valentine \footnote{Not the St. Valentine
whose name is found in the Kalendar, but a martyr of the same name
whose relics were found in the catacombs, and given to Newman by Pius IX.},
who connects us with the first age, shows that the Church has kept,
not lost, her first bloom.

\xsection{March 24 (Palm Sunday) On Our Lord's Agony}

1. INTROD.—We
naturally seek to be told something of the death and the deathbeds of
those we know and love. We are drawn to the deathbed of the saints and
holy people; and much more if anything remarkable about it, and much
more if a man be our benefactor, parent, etc. How much more the death
of the great God?

2. Thus, above all, our Lord's death—how sudden
it was! One day brought into the city in triumph, the next plotted
against, betrayed and seized.

3. God from eternity—the Holy Trinity. Each
person all God; the Son the only God, as if only Person.

4. God most happy; Son all happy—bliss, peace,
calmness, glory, beauty, perfection from all eternity.

5. And now look at that one only God, as we
contemplate Him at this time of year. He is still one, sole, and
alone. He was one in heaven; He is one in the garden, one on the tree.
He trod the winepress \emph{alone}. When He went into the garden He
took but a few with Him, and separated Himself from them; and
afterwards the disciples 'left Him alone,' and fled. Easy for the
traitor to take Him, for He was alone.

6. But though one and alone, how different! He
who was glorious is become a leper; He who was so peaceful has lost
His rest.

7. It is said that nothing is so fearful as the
overwhelming sorrow of man as contrasted with woman, of a hero or
great and firm man overcome by adversity or bereavement; for it being
more difficult, it bursts more [violently]; it is like a storm rending
and shattering. What, then, in the most peaceful and serene? What a
conflict in the sinless!—(enlarge).

8. It is said that 'the wicked are like a boiling
sea'; what means this in the innocent? Yet so it is. He began to grow
weary, sad, frightened. (Explain.) On the devil, who was foiled in the
wilderness, to his surprise finding our Lord in the garden agitated as
a sinner. He had gained his point—his eternal enemy vanquished. On
the apostles sleeping for sorrow, but Christ praying more earnestly.

9. Pain of mind greater than that of body, though
we are more conversant in bodily pain—grief, fear, anxiety, terror,
despair, disappointment—\emph{poena damni} of the lost greater than \emph{poena
sensus}. On the effect of mental pain—hair turning white;
Nabal \footnote{But early in the morning
when Nabal had digested his wine, his wife told him those words and
his heart died within him, and he became as a stone.'—1 Samuel xxv.
37.}. So effect on
Christ—agony of blood.

10. Let us gather round and look at Him whom God
has punished; but in no idle way, for His pain is from our sins.
Address to sinners.

\xsection{July 14 On the Particular Judgment}

1. INTROD.—Give
an account of thy stewardship, from the Gospel of Pentecost viii.

2. (A judgment will take place directly.) It
would be well if we could realise what our actual position is. We \emph{happen}
at this moment to be in this world, but \emph{any moment} we may find
ourselves in the world unseen. We are now talking to each other; we
see each other, etc. Yet just as walking we may cross over a street,
so suddenly we may cross to the next world—'Thou fool, this night
thy soul will be required of thee,' etc.—a veil drawn across.

3. It is difficult even for a Catholic who
believes it to realise. Thus a person who never was at York could not
realise that this time tomorrow he will be there. Still, the more we
meditate the more we shall realise it; and it is our duty to realise
it more and more. Saints realise it.

4. But a Protestant really has no notion of it.
This is proved by any \emph{sudden} death—sudden deaths throw
them off their balance and detect them. They at once betray by their
words that, whatever they may say or wish, they do not really believe.
They call it an \emph{unknown} state; but though it is unseen it is
not unknown. But a Protestant does not know \emph{whither} he is going
more than Adrian with his \emph{anima blandula}, etc. He in his heart
confesses it. He says, 'After all, we know nothing.' Whether he will
lose consciousness, or be asleep; whether in heaven, or what is
heaven. In a word, he is all abroad; the question is new to him, and
he has not one idea about it, no more than a pagan. What has he more
than a pagan?

I am not talking of heaven, or eternity, but of
what will happen to him personally directly after death.

5. But a Catholic knows—particular judgment and
purgatory—judgment on the very spot and time, expeditious like an
inquest, as necessarily [following death] as an undertaker. 'After
death the judgment': (1) Philip the Second to his courtiers, \emph{St.
Alphonso}, p. 249; (2) St. Mary Magdalene of Pazzi trembled in her
sickness, \emph{St. Alphonso}, p. 248.

6. Yet true as this in general, still none of us
understand as we should what is meant by \emph{stewardship}. What,
have all stewardship?

7. [Bound to make our lives] conformable to the
life of Christ—'if the just scarcely be saved.' Account [to be
rendered] of (1) sins; (2) blessings and graces; (3) idle word—\emph{ictus
oculi}; (4) thoughts of heart; (5) \emph{time}, recreations.

8. Let us live ever in the thought of judgment.

\xsection{August 11 On the Doctrine of Prayer as Reconciling Us to the Catholic Teaching about Our Blessed Lady}

1. In the course of Nature everything proceeds in
order; system of cause and effect proceeds illimitably, so that we do
not know where it stops. It is a vast web. Hence all things seem
fated. This is what unassisted reason seems to teach:

2. and hence no religion proper, for God cannot
act upon us \emph{directly}, but only through a system; and therefore
it is a system only which acts. OBJECTION.—The
laws of God will go on whether we strive or not: \emph{e.g}. how can
prayer save from trouble, give health, cut off a persecutor?

3. Yet conscience, feeling, and the religious
sense, which are part of us, speak contrariwise, viz. of particular
providence.

4. Revelation confirms this. 'Not a sparrow falls
without ... hairs of head,' etc. And specially it reveals the power \emph{of
prayer}.

5. Prayer in its effect, though the idea is so
familiar to us, is one of the greatest of mysteries and miracles, yet
it is the clear doctrine of Scripture.

6. Now I am going to use this without reference
to the subjects, which will be brought before us in a few days in the
Assumption.

7. Much is said by Protestants against our Lady's
power, but our Lady's power is nothing else than the greatest
exemplification of the power of prayer. We don't give her power of
atonement, etc., but simply prayer, as we give ourselves; we in a
degree, she in fulness. Now I can understand persons scrupling at
the power of prayer \emph{altogether}; but why, that there should be
one instance [\emph{i.e}. great exemplification] of it? We do not
introduce a mystery, but realise it. The great mystery is that prayer
should have influence. When once we get ourselves to believe the power
of prayer, etc.

8. \emph{E.g}. even Protestants say the strongest
things about prayer—'prayer of faith'—'Satan trembles'—'faith
moves mountains'—'faith can do all things,' \footnote{\emph{Written in pencil above},
'faith can do. Prayer constrains God: well, this is the very thing
Protestants think so shocking when we say it of our Lady.'} that is, is omnipotent—this is just what we say about our
Lady—omnipotent through prayer—'Let me alone' [Exod. xxxii.
9-11]—Amalec \footnote{'And when Moses lifted up
his hands Israel overcame: but if he let them down a little, Amalec
overcame. And Moses' hands were heavy: so they took a stone, and put
under him and he sat on it: and Aaron and Hur stayed up his hands on
both sides.'—Exod. xvii. 11, 12.}—Jacob's
wrestling \footnote{'He remained alone, and
behold a man wrestled with him till morning.'—Gen. xxxii. 24.}—'I cannot
until thou come thither' [Gen. xix. 21, 22]—Luke xviii. 5—violent
carry it away by force.

9. Perseverance. Again Amalec, Jacob's wrestling;
Luke xviii. 2-7 \footnote{'There was a judge in a
certain city who feared not God nor regarded man. And there was a
certain widow in that city; and she came to him, saying, Avenge me of
my adversary. And he would not for a long time: but afterwards he said
within himself, Although I fear not God, nor regard man, yet because
this widow is troublesome to me, I will avenge her, lest continually
coming she weary me. And the Lord said, hear what the unjust judge
saith. And will not God revenge his elect, who cry to him day and
night?'}, the
woman gives him no rest.

10. Sanctity. 'God heareth not sinners.'

11. Now who can persevere as our Lady? Who is as
holy as she? What wonder that in her [the power of prayer] should be
fulfilled most perfectly?

12. The more we pray the more we shall be
reconciled to the doctrine.

13. Let us at this time make her both the example
of prayer and its object.

\xsection{October 6 On the Necessity of Securing Our Election}

1. INTROD.—There
are mysteries of revelation (most) which are beyond our experience; we
receive them only on faith, \emph{e.g}. Trinity, Incarnation, etc.,
but there is one which we see, viz. the great mystery of election and
predestination.

2. We see before our eyes the astonishing fact,
that all are not in the same condition as regards religious truth.
Some are born in heathen countries—of bad parents—in heresy.
Again, of two men who sin, one is cut off, one lives to repent. Again,
men who go on well for years suddenly fall away—how all are mixed
together. The world and Catholics not distinguishable; they dress
alike, etc.

3. Now Scripture recognises this awful fact. It
speaks about the elect being few, the flock being little. It says much
of God's grace, of a choice, etc. This certainly is most wonderful,
for it was a prophecy at the time, which every age has confirmed since, and a curious \emph{combination}, viz. that Christ's religion
should at once surround and subdue the world, yet be thus small, and
weak, and despised.

4. Such is the doctrine of Scripture, and it is
put even more strongly. There is an awful text, 'If the mighty works,'
Matt. xi. 20-24 \footnote{'Then began he to upbraid
the cities wherein were done the most of his miracles, for that they
had not done penance: Woe to thee, Chorozain! woe to thee, Bethsaida!
for if in Tyre and Sidon had been wrought the miracles that have been
wrought in you, they had long ago done penance in sackcloth and ashes.
But I say unto you, It shall be more tolerable for Tyre and Sidon in
the day of judgment than for you. And thou, Capernaum, shalt thou be
exalted up to heaven. Thou shalt go down even unto hell. For if in
Sodom had been wrought the miracles that have been wrought in thee,
perhaps it had remained unto this day. But I say unto you, That it
shall be more tolerable for Sodom in the day of judgment than for
thee.'}.

5. Now in saying this, it is not (as I have many
times urged) as if God did not give enough to all, but He gives more
to one than another. Why, we know not. (Do not think I put it as a
speculative mystery merely; it is most instructive; it is not only
awful and mysterious, but on the other hand, a most profitable fact to
consider.)

6. It arises from God's self-dependence,
self-sufficiency. Eternally happy from everlasting. Creation did not
make Him dependent. What is it to Him if thou art virtuous? You have
no claims on Him except what He has given by pledging Himself. Beware
of pride. He does not want you, etc.—you can do Him no good. In your
own nature you are indefinitely removed from Him; it is only by
superabundant grace that you come near Him. He is not bound to
give grace, but He does give it to all. And as He is not bound to give
at all, He is not bound in measure. He has full right to give so much
as He pleases, more to one than to another.

7. Hence we cannot argue for certain that because
He has forgiven our sins already, He will in future, if we sin. We
cannot count that He will a second time give us the grace of
repentance. Nor because He has forgiven others, therefore He will us.
And so again, if we are external to the Church, we cannot rely on His
always giving us the grace which He has given so far as He has, solely
to bring us into the Church. And so again we have no confidence,
because we are in a good way now, that therefore we shall persevere.

8. Do not think I am putting this as a harsh
speculative mystery; it is as a practical consideration. Beware of
quenching grace. The grace of conversion is rare; the grace of
illumination is precious. You do not know but this may be the last
grace given you, if you resist it. You have claim on nothing if you
are external to the Church. You may have good feelings, dare not rely
on them; you may at present be in God's grace, do not conclude too
easily that you will persevere. Watch against sin; for what you know,
the least wilful venial sin may act upon your deathbed, and subtract
from the aid which would then have been given you.

9. The need of prayer. God sovereign, but prayer
almighty. God has given to us as a means to overcome, as I may say,
Himself. Let us never be satisfied with getting prayers for our
perseverance.

10. Here is the special office of our Lady, and
its bearing on us. She does not predestinate, she does not give
grace, she does not merit grace for us, but she gains it by prayer;
she gains perseverance by prayer. Thus she overcomes God, as I may
say.

11. Suitable on Rosary Sunday. \emph{Nunc et in hora
mortis nostrae}.

12. May we die in peace.

\xsection{October 18 On External Religion}

(\emph{For St. Peter's, on the
opening} \footnote{Preached at St. Peter's,
Birmingham.})

1. INTROD.—Gospel
of the day, Pentecost xiii.—'Glory to God.'

2. What is meant by glory? We unprofitable, but
we can \emph{show} our worship, etc.

3. I am led to this, the natural subject of the
day, now that the chapel, embellished, etc., is reopened.

4. How natural, you see from what every one of
you does—children forming little altars, etc.

5. So the first Christians, even in caves \footnote{See Note
8, p. 337.} (which are most alien to Christianity), in catacombs, adorned
them in times of persecution.

6. Whereas when an unreligious movement, the
first business to destroy these [embellishments]—the Danes, the
Reformers, the Huguenots, the French Revolution. As the devil
delighted to destroy our Lord's beauty, so the beauty of His Church
(even organ and surplice).

7. But it is exemplified in all religions. Look
out into the first ages, patriarchal times, before religion was
corrupted—a mountain top (beautiful prospect), or grove or rivers,
where sweet smells and sounds of birds. No matter if afterwards
corrupted—Garden of Eden.

8. Particularly in south, where scenery beautiful
and weather fine—out-of-doors worship. Next, the first artificial
part was processions, vestments, and music and statues.

9. Then came flowers and incense. Describe a
procession—children with garments; the victim; no matter that
superstition afterwards corrupted it to false gods.

10. Then they took it out of the open
air—Jewish tabernacle. Then we come to furniture, as you read in
Exodus—and so jewels, etc., marble, pictures; then painting,
sculpture and music.

11. Lastly came architecture, which has to do
with form. With it [is] so great a part of the beautiful. The dome
represents the heaven, the arch the wood.

12. Thus at length all things—the eye, the ear,
the scent—form, colour, music, incense.

13. One more characteristic in all
this—costliness. Sacrifices. David spoke of 'that which doth cost me
nothing.' \footnote{'And the king answered
him and said, Nay; I will buy it of thee at a price. I will not offer
to the Lord my God holocausts free-cost.'—2 Samuel xxiv. 24.} So also of
the widow, 'She hath done what she could'—'two mites, which make a
farthing,' Mark xii. 42, 43 \footnote{'And there came a
certain poor widow, and she cast in two mites, which make a farthing.
And calling his disciples together, he said to them, Amen I say to
you, this poor widow has cast in more than all they who have cast into
the treasury: For all they did cast in of their abundance; but she of
her want cast in all she had, even her whole living.'}

14. 'The poor you have always with you; not me,'
Matt. xxvi. 11; occasional call, as now, 'a stranger, and you took me
in.' St. Peter and St. Paul—'shall receive a prophet's reward.'

\xsection{December 1 (Advent Sunday) On Death}

1. INTROD.—Again
Advent. Christmas!—the day darkens; the year dies; all things tend
to dissolution. It is the end; we have to think of death and all
connected with it.

2. We are going on right to death; a truism, yet
not felt. We are on a stream, rushing towards the ocean; every morning
we rise nearer to death; every meal we take; every time we see our
friends, etc.; nearer the time when we shall lose them. We rise, we
work, we eat; all such acts are as milestones. As the clock ticks, we
are under sentence of death. The sands of the glass run out; we are
executed; we die.

3. And when it comes, what happens? We all know.
This happens—we are no longer here. We see not indeed whither we go,
but this we know full well, we are not here. The body which was ours
is no longer ours; we have slipped it off; no longer a part of
us. It is a mask, as a dress; but it is not our instrument or organ.
We who think, feel, speak, etc., are not here. Where we are, nothing
that is here tells us; but this we know full well, we are not in the
body. We are cut off from all here. This minute here, the next a wall
impenetrable has grown up; we are as utterly cut off as if we had
never been here; as if we had never known any one here. We don't go by
degrees—we do not (as it were) lessen in perspective and disappear
in the horizon—we go at once and for all.

4. Where it is we see not; what it is we know
not; but what it is not we know, as we know where it is not. The man
is not what he was. He took pleasure; he depended on this world. He
depended for its enjoyment on the senses. That life was not a burden;
that it was dear to him; that he enjoyed it; that he was unwilling to
quit it, was because he saw, he heard, etc., his amusements, his
pleasures; he went to his club, or to business, with his friends; he
liked the warm fire, the light; he liked his family, home comforts,
his dinner; he strolled out in summer, or he went to places of
merry-making and enjoyed the gratifications of sin—nothing
supernatural: how many we have known such! Why are people unwilling to
die? What is the one reason? There is no pain in it. Because they
leave what is known; they go to what is unknown. They leave the sun,
etc.; they leave their families, their schemes, their wealth.

5. Oh, how much is implied in this! Men witness
against themselves. They are afraid to leave this life; they own
they are going to the unknown, yet they are unwilling to make that
unknown known. Do lay this to heart;—you are going to the unknown.

6. Now I will tell you what you are going
to—not to creatures as here, but to God. Oh the dreadful state of
the soul when this step is over! Another world is close to us. It has
taken the step, and is in that other world. Have you any relations
with God? Do you know aught about Him? Do you know what He is like?
Have you tried to make Him your friend? Have you made your peace with
Him? What madness! If men are going on a voyage they take letters of
introduction; they inquire about the country; they try to make friends
beforehand; they take money with them, etc. Yet you do not try to
disperse the thick darkness; on the contrary, you learn to be content,
because you do not know.

7. Yet that acquiescence is an additional alarm,
for it shows God is angry with you. Men lightly say: 'It is a matter
of opinion.' No, it is a matter of punishment. This very discordance
of sects is a sign of God's displeasure.

8. The longest life comes to an end. You may be
young, you may be vigorous, but you must die. When it is over, the
longest life is short.

9. Seek the Lord therefore; this is the
conclusion I come to; this world is nothingness. Seek Him where He can
be found, \emph{i.e}. in the Catholic Church. He is here in the same
sense in which we are.

\xsection{December 29 On the Office of the Church—St. Thomas the Martyr}

1. INTROD.—This
is the birthday of a great saint, one of the greatest of English
saints, whose fame has gone out, etc.; a saint of the universal
Church, especially known in France, North Italy, Roman States, etc.;
nay, whose feast is embodied in the octave of Christmas.

2. His manner of death.

3. What did he die for? If you ask a Protestant
history, it will mention some minute ground, some question of detail,
of course; but, if examined, for \emph{that} which is ever the cause
of battle between the world and the Church. Parallel it to the early
age, a grain of incense; the present moment, calling bishops bishops
of sees, etc. But all these are accidents—the ground, one and the
same.

4. To explain this I must go into the subject.
State of the world before Christ came—the world left to itself;
doubt and inquiry; philosophers; pagans; yet no known truth.
Philosophers felt it impossible to throw truth into a popular form.
Hence they were tempted to believe there was no truth. Great
difference between religious truth and scientific, etc. We can get to
sciences of geology, etc., because we start from what we see, but who
shall tell the designs of the Divine Mind?

5. Prophecy of a Teacher—voice behind thee
[?]—a law—a light (Isaias ix.)—Isaias xxv.—Thus a master, or
guide, or monitor to be set up.

6. Such is the one province which Christianity
was to fulfil. Now Protestants think this fulfilled in the \emph{Bible}.
But the Bible has not \emph{in fact} been the means. (1) The majority
have not been able to read. (2) And now fifty years' experience shows
it is not God's way \footnote{Does the preacher refer
to the British and Foreign Bible Society founded in 1804?}.
(3) Nor can it be, for a book does not speak; it is shut till it is
opened. A law cannot enforce itself; it implies an executive; not a
book instead of a physician, etc. (4) It is nowhere said in Scripture
that Scripture was to be the guide, but it is said what is to fight
with the gates of hell, viz.

7. the Church—texts. This is set up, and did
exist before, etc., in all lands to appeal to high and low, to all
ranks and callings—(enlarge). To moderate, and in a certain sense to
interfere, viz. with the conscience—on the misery of princes being
made so much of from youth—to give the law and to teach the faith.

8. This is the quarrel—the world does not like
to be taught. (The Jewish kings did not like prophets.) The Church
interferes with it; she lifts up a witness. Men regret the old pagan
times when each could say and think what he pleased. Kings and
ministers, etc., etc., don't like to be interfered with.

9. This, then, was the world's quarrel with St.
Thomas. Henry II. felt the
Protestant ground just as the meetings now held do—it is the same
spirit—therefore does the world persecute us now. When, then, men
object that we interfere with conscience, etc., etc., we say 'yes.'
And if we did not, we should not be the Church; if we did not,
there would be no good in a Church.

10. And you may be sure that the Church will
never betray its trust.

\xsection{January 19, 1851 On the Name of Jesus}

1. It has been from the beginning the order of
Providence—nay, even \emph{verbum}—not to create without giving a
name. As grace is necessary to keep things together lest they
dissipate, so a name is, as it were, the crown of the work, as giving
it a meaning and description, and, as it were, registering it before
Him. Henceforth it lives in His sight, as being in His \emph{catalogue}.

2. Thus 'day' and 'night,' 'earth' and 'seas.'
Hence Adam named his wife and the beasts, etc. Hence Abraham's name
changed; Jacob's, Sarah's, Isaac's; Isaac's given, Jacob's changed;
St. John Baptist; St. Peter and St. Paul. These names are descriptive.

3. Hence anxiety of men to know God's name. They
are born in ignorance. They have a sense there is a God, but what is
He? The heavens and earth do not condense and concentrate His manifold
attributes, etc. They give hints, glimpses, snatches, but what is He?
Hence He is the unknown God, and men are but 'feeling after Him' by
what they see. They are in God; He surrounds them, but they want to
gaze on Him objectively.

4. Thus Jacob about the angel, 'What is thy
name?' And to Manue, 'Why askest thou my name, which is \emph{mirabile}?'
Judg. xiii. 18. Moses bolder. God had been called 'God of Abraham,'
etc. \emph{Adonai}.

5. Hence you see a meaning why the Eternal Son
would reveal this, that the Name of that Son was of consequence; it
was a manifestation of the nature and attributes of God—Admirabilis,
Isa. ix. 6 \footnote{'For a child is born to
us ... and his name shall be called Wonderful, Counsellor, God the
mighty, the Father of the world to come, the Prince of Peace.'}; Emmanuel,
Isa. vii. 14 \footnote{'For behold a virgin
shall conceive, and bear a son, and his name shall be called
Emmanuel.'}. Still,
however, the name was not told. At length Gabriel said it, Luke i. 31
\footnote{'Behold, thou shalt
conceive in thy womb, and shall bring forth a son, and thou shalt call
his name Jesus.'}; circumcision, Luke
ii. 21; angel to Joseph, Matt. i. 21 \footnote{'She shall bring forth a
son, and thou shalt call his name Jesus.'}, His name was called Jesus. And hence the devils: 'Jesus the
Son of God'; 'I know thee who thou art.' On the cross \footnote{'And the writing was,
Jesus of Nazareth, the King of the Jews.'—John xix. 19.}. The first miracle of St. Peter and St. John, Acts iii.—'in
the name,' 'this name,' 'no other name'—and St. Paul in Phil. ii.
8-11 \footnote{'He humbled himself,
becoming obedient unto death, even to the death of the cross. For
which cause God also hath exalted him, and hath given him a name which
is above all names. That in the name of Jesus every knee should bow,
of those that are in heaven, on earth, and under the earth; And that
every tongue should confess that the Lord Jesus Christ is in the glory
of God the Father.'}. The two great
apostles, the angels from Gabriel, devils from the possessed, and
men from the circumcision.

6. For in this the whole history of salvation,
the whole creed—how God would save men, how He loved them, etc.,
recounting the Christian doctrine \footnote{\emph{I.e}. the Holy
Name sums up in Itself the history of salvation.}. Thus when we would know who God is, we answer, Jesus. We see
God in the clouds, in the mountains, etc., and who is He? Jesus. Who
then rules? Who is looking, the ruler of bad men? Who is looking, the
guardian of the virtuous? Who, etc.? and we answer, Jesus. He is the
one word containing in itself all power, etc., because in it we
thereby have in our minds the full description of Almighty God.

7. And in it an answer to all objections and
difficulties. It surpasses all (this is the point of the sermon):
whatever difficulties, whatever mysteries in religion, this
comprehends and protects them. What is more wonderful than that God
should become man. Real Presence, power of Mary, purgatory, eternal
punishment, intercession of saints, election, original sin. The whole
Catholic system bound up in it.

8. Hence, and since Protestants have the name of
Jesus on their lips, it is the test whether or not they understand it,
\emph{i.e}. their taking Catholic doctrine or not. If they don't, if
they stumble at it, they don't understand Jesus. On invincible
ignorance, as alone hindering Catholicism.

9. Let us then rejoice in the fulness of this
Name. Let us use it as the Name of virtue against devils, bad
thoughts, evil men, the world, dangers and frights. It is our banner.

\xsection{January 26 (Third Epiphany) On Disease as the Type of Sin}

1. INTROD.—When
our Lord came His chief miracles were, not like Moses', etc., on
elements, but on men, on diseases.

2. Why? Because He was the Redeemer. The physical
world had not to be redeemed, but men, and disease was a defect;
whereas the physical world was perfect after its kind, very good. He
did work some miracles on the elements, to show He was the Creator;
most on the infirmities of human nature, to show He was its Redeemer.

3. He had to do with sin; and bodily diseases are
at once its symptoms and its representations (they represent sins both
in their intensity and variety). When man fell, the grace which
covered his soul and body was like a skin torn off and leaving him
raw—(enlarge). Men would forget sin, but they cannot. Hence it is
that all false views of religion fail—the views of the day; a bright
careless religion does in the sunshine, not in the shade. Here it is
that Christ spoke to the heart. He comes to do that which false
religions and infidelity \emph{ignore}—to cure sin.

4. Nothing more awful than bodily pain, except
mental; but mental is a private matter, and can be denied, can be put
off; bodily is before us. Disease represents sin in its \emph{intensity}
and \emph{variety}. Go through bodily complaints—fever, ague,
sinking from weakness, oppression of breath, etc., cholera,
restlessness, etc., paralysis, leprosy—here you have sin in its
various forms.

5. And it suggests to us future punishment. No
dreams of God's mercy can overcome the fact—and our Lord, most
merciful though He is, requires it. 'Jerusalem, Jerusalem,' etc.,
'Behold your house shall be left to you desolate,' Matt. xxiii. 37-38,
Luke xiii. 34-35, and see His tears over Lazarus' grave, John xi. 35 \footnote{'Where have you laid
him? They say to him, Lord, come and see. And Jesus wept. The Jews
therefore said, behold how he loved him.'—John xi. 34-36.} 'An enemy has done it,' Matt. xiii. 25-28 \footnote{'The kingdom of heaven
is like to a man who sowed good seed in his field, but while men were
asleep his enemy came and oversowed cockle ... And the servants of the
goodman of the house coming, said to him, Sir, didst thou not sow good
seed in thy field? Whence then hath it cockle? And he said to them, An
enemy hath done this.'}. It is necessary by the immutable laws of truth.

6. Nor can you say it is merely \emph{remedial},
for why allowed, \emph{i.e}. sin? God could have hindered sin. Again,
He need not have died for it, and yet might have pardoned it—'I
will, be thou clean.' No, it is the beginning, not the no-ending of
pain, etc., which is the marvel. When once it comes in, there is no
reason why it should not continue, etc.

7. Come to the Physician of Souls, etc.

\xsection{February 9 (Fifth Epiphany) On the Descent into Egypt}

1. INTROD.—There
is one subject not much thought of, the descent into Egypt, though
belonging to this season, and it may be accounted one of the
Epiphanies of Christ. (1) Magi; (2) purification; (3) idol-breaking in
Egypt, for as the Ark levelling the walls of Jericho, and
[overthrowing] Dagon, so much more Christ.

2. Circumstances. How one journey to
Bethlehem—then they got home to Nazareth—not enough, must set out
again. It is remarkable that these first years were spent in a heathen
country. There till seven years old.

3. Now why? I will give you a reason. He would
undergo every suffering; He would be in a heathen country to share the
trials of His apostles and missioners. In Jerusalem was the Temple of
God, in the Holy Land His religion; but even there He chose not the
Temple, but Nazareth—and the first years of His life Heliopolis, in
heathen Egypt.

4. Now it must not be supposed that our Lord was
too young to have a trial. (Explain.) Ignorance came from the fall. He
was as sensitive [in childhood] as [we] when we are grown. Thus He saw
all the evil of the place; and as His body made Him feel in the
crucifixion, so His soul was exposed to moral sufferings from the
first.

5. And the suffering was greater than we
conceive. To live among heathens is a misery, the greater, the purer
the mind—Lot in Sodom, St. Paul at Athens—the world is everywhere,
and we can understand from a country which is not heathen, such as
this, how evil it is, though it would be a great deal worse [among
heathens].

6. Even in this country, I say, which is not
heathen, the misery of being in the world is great to any holy
mind. Take \emph{e.g}. a city like this, and fancy the thoughts of an
apostle in it. Could he go about it freely? A continual service of the
devil here. How? By sins of the tongue; not like the seven Catholic
Hours coming at intervals, but incessantly; a continual light talk in
a thousand places, from morning to night, with scarce breaks. Who is
honoured like the devil? Blasphemy and immodesty, so that most men's
mouths and all men's ears are polluted from year to year's end. And
are not their hearts too? Then imagination. Alas! \emph{this} is why
the devil loves the bad talk; it is the \emph{pabulum}, the \emph{silva}
of corruption; it sets the heart on fire, as shavings round the wood
and coal for a fire. I don't know anything more awful. Other sins men
commit from time to time, but this one now. The evil concupiscence
boils over and burns without exhaustion, and involves every one, so
that religious people are like the Three Children [in the fiery
furnace].—and how many, many fall!

7. Well this, bad as it is, is not so bad as
Egypt, as heathen Heliopolis, for this country has been
Catholic—remains of good, which have soaked in. Grant that a modern
city is a furnace of sin—yet it [sin] was deified in Egypt. Vices
canonised in animals—heathen idolatry—all vices made gods—the
world lieth in wickedness, etc. \footnote{2 Cor. iv. 4.}—the god of this world, etc. \footnote{1 John v. 19.}.—the prince of the power of this air, of the spirit that now
worketh on the children of unbelief \footnote{Eph. ii. 2.}. O misery of the infant Jesus walking in the streets! St.
Aloysius fainting at the mention of a mortal sin—smell—saints
detecting mortal sin [by its smell]. As sick men cannot bear
strong scent or sound, so purity here. What a living martyrdom, etc.

8. I have said He did this for our sakes, to
taste every trial, to sanctify every state, to sanctify the state of
those who live in the world.

9. You who live in the world, resist evil. On
confraternities—third order of St. Francis, and so allude to the
Oratorium Parvum. Your confessors may, or may not, from not liking to
put burdens on you, speak of these.

\xsection{February 16 (Septuagesima) On Labour—Our Work Here}

1. INTROD.—Before
Lent the Church begins by setting before us work as an introduction.

2. Epistle and gospel—beginning of Genesis.
Even before the fall, and much more after—thorns and thistles.

3. This the contrast between before the fall and
after. The ground typifies our hearts—and now we have labour.

4. And this will show us the heinousness of the
fall, for before it, the labour, the effort, was to sin—before as
difficult to sin as now to be a hero. Grace was so great.

5. But grace being gone, the lower nature rose
against the upper \footnote{Written above in pencil,
'Is this right?'} as
the upper against God.

6. This then, I say, our work—labour of one
kind or another. It has different names—self-discipline,
self-denial, penance, reformation, mortification—all meaning the
bringing under of ourselves. Don't think it hard if you find a thing
difficult; it is your work.

7. This implied in the subduing our 'ruling
passion,' so called.

8. Also exemplified in particular examination.

9. Also done in suffering. Suffering is a work.
On satisfaction and \emph{satispassio} \footnote{\emph{Satispassio} is
paying the full penalty—'the last farthing.'}; on bearing pain with sweetness or patience, with sweet faces,
ways, voice, etc., etc. On the discipline when associated with the
thought of Christ's sufferings, more meritorious; for the mind goes
with it and is not otiose.

10. Thus let us begin this sacred time.

\xsection{February 23 (Sexagesima) On St. Paul the Type of the Church as Missionarising}

1. INTROD.—This
day seems especially set apart for the consideration of the apostle
St. Paul, in collect, in epistle, in gospel, for he is the sower.

2. How he sowed in all places. How he preached.
He fought. The great soldier. David goes out against a giant, but he
against the world. What a great ideal! Patriots, Joan of Arc, etc.,
etc., but this, not in one country only. To east and west, north and
south, he goes and forms a kingdom.

3. But this great portent is completed by the
history of the Church after him. It is not solitary, not an
accident, not like one great man, as Buonaparte; but he \emph{intended}
and the work has \emph{lasted} eighteen hundred years, going on the
same way.

4. Now the warfare goes on just the same, and
with the same enemies. This again most extraordinary. The view of the
battle is just the same. As a shadow may move onwards and presents the
same outline over hills and dales, so as time has gone, this one
grouping has gone on for eighteen centuries.

5. Look at it in St. Paul's day—zealots and
indifferents, statesmen and philosophers. Describe them.

6. Zealots—Jews and pagans. Pagan, Acts xix.,
[tumult of silversmiths at Ephesus]: Jewish, Acts xxiii. 12 [forty men
bound under a great curse neither to eat nor drink till they killed
Paul].

7. Indifferents—magistrates. Gallio, Acts xviii.
12; Festus, 'Paul, thou art beside thyself,' Acts xxvi. 24;
philosophers at Athens, Acts xvii. 18.

8. Application to the present times. Furious
evangelicals and statesmen. Their different ground. The first call
Rome Antichrist.

The second profess to care nothing for doctrine,
but only go to political grounds.

9. Nay, our blessed Lord. Pharisees
furious—'The Son of God.' Then they come to Pilate (What is truth?)
with a different plea. 'Thou art not Caesar's friend,' etc. The
emperor's supremacy, etc., denied.

10. This awful unity of the Church is our
consolation. While it proves the Church comes from God, it proves
nothing comes strange and new to her.

11. No, our business to sow and to fight, and to
leave the rest to God. It is never to be supposed we shall not go on
doing the same as before.

\xsection{March 9 (First Lent) On the Accepted Time}

1. INTROD.—Lent
an apostolical observance.

2. And well did it become the Divine Mercy to
appoint a time for repentance, who had in the fulness of time died for
our redemption. For what is every one's business is no one's; what is
for all times is for no time.

3. And even those who will not take God's time,
feel a time there must be. They always profess a time; they quiet
their conscience by naming a time; but when?

4. 'Go thy way for this time; when I have a
convenient season,' etc., Acts xxiv. 24-25 \footnote{'And after some days
Felix coming with Drusilla his wife, who was a Jewess, sent for Paul,
and heard of him the faith that is in Christ Jesus. And as he treated
of justice and chastity and of the judgment to come, Felix being
terrified, answered, For this time go thy way, but when I have a
convenient time I will send for thee.'}. When the present temptation is out of the way. When the
present business or trouble is got through. When they have enjoyed
life a little more.

5. When 'a little more,' for there is no
satisfaction in sin, each sin is the last. But the thirst comes
again; there is no term at which we can quit it; it is like drinking
salt water—horizon recedes.

6. End of life, time of retirement. The
seriousness will come as a matter of course; passions will naturally
burn out—\emph{otium cum dignitate}—alas, the change of nature is
not the coming of grace. We may change, but we shall not be nearer
heaven. To near heaven is not a natural change, but a specific work,
as much as building a house. It is not a growth till there is
something to grow from.

7. Feeling then there must be a time, and having
the conscience of men on this point with her, the Church appoints a
time and says, 'Now is the appointed time.' She blows the trumpet;
proclaims forgiveness; an indulgence—scattering gifts—inviting all
to come and claim. Not sternly, but most lovingly and persuasively she
does it.

8. Oh for those who have neglected the summons
hitherto, year after year, conscience pleading!

9. Or perhaps we have repented just through Lent
and then relapsed and undone, and more than undone, all.

10. And so we get older, older, and farther from
heaven every year, till we come to our last Lent, and we do not keep
it a bit the better.

11. Then we come near death, yet won't believe
that death is near. Set thy house in order—packing up, and how many
things left out. We cannot realise it. All hurry and confusion.
Between illness, delirium, weakness, relations, worldly affairs, etc.,
we shall be able to recollect nothing—all in disorder. No real
contrition. And so we die.

12. Ah! then in that very moment of death we
shall recollect everything; all things will come before us. We shall
wish to speak; it will be too late. We shall have passed from this
life; the accepted time will have passed by.

\xsection{March 23 (Third Lent) On the Strong Man of Sin and Unbelief}

1. INTROD.—'The
strong man' represents the sinner in his strength and security. It
represents him fortified by his three friends—the world, the flesh,
and the devil. 'The old man,' Eph. iv. 22 \footnote{'For our wrestling is
not against flesh and blood, but against principalities and powers,
against the rulers of the darkness of this world, against the spirits
of wickedness in the high places.'—Eph. vi. 12.}, the old Adam, the evil spirit who has taken possession of
him.

2. He has a 'house.' It is a castle: nor is it
the work of a day. How long it takes to build a castle! and buildings
grow up about it, fort after fort, treasure house after treasure
house, viz. by habits. (Explain about habits.) No one remains without
them; they are intended to be a defence for the good. They also become
a defence in wickedness. Supernatural habits and natural habits.

3. Absence of faith—'The light that is in
them.' His standard of things—scoffs at things supernatural; does
not think himself a bad man because he does not pray; is in 'peace';
perfectly satisfied with his standard; may not come up to it; is
firmly seated. He may be educated, learned, able, etc.; this only
increases the evil.

4. Enormous strength of a bad man. His \emph{vis
inertiae}, his \emph{momentum}. In his black panoply, armed \emph{cap-à-pie}
like a knight in story, such the bad man. Then fancy a host of them,
the rulers of this world, like a bodyguard of Satan, or his 'guards.'

5. Such are the enemies of Christ, described in
the Gospel: 'We wrestle not against,' etc., Eph. vi. 12. Then Christ's
grace more powerful: 'A stronger than he,' etc.

6. No one can come up to the strength of God's
grace—stronger than the elements; stronger than miracles. It bears
up against anything; it overcomes everything. On the wonderful way in
which Christianity overthrew the establishment of paganism (\emph{vide}
Döllinger).

7. Let this be your comfort if you feel afraid,
and have to do a great work. God's grace can convert; it has converted
from sin of whatever kind.

\xsection{March 30 (Fourth Lent) On Bearing Mockery}

1. INTROD.—Laetare
Sunday. Joy, like a flower springing out of desolation and
mortification, as Christ goes along the desert.

2. What flower shall be our offering? We cannot
do much in the way of fasting, or other bodily mortifications. Why,
the time supplies one, and which the epistle suggests, viz. our
bearing reviling, etc. On the epistle of the day—Hagar and
Sarah. It was a strong boy bullying a small child—cowardly and
ungenerous. This animal nature. (Describe Ishmael.) Sarah childless
till Isaac. \emph{Laetare}—the mocking. 'Even so it is now,' says
St. Paul. It is the mark of the true Church, and the form of its
warfare—mockery. And so it is at this minute.

3. The scoffings, etc., which surround us not
exactly violence or suffering, but slander, etc. The huge
Protestantism of this land cannot keep from grinning, scoffing, etc.

4. Now this has ever been the case with the
Church, as I have said, \emph{e.g}. Isaac.

5. Joseph, Job, David, Jeremias, Daniel.—Heb.
xi. 36.

6. Our Lord—(particulars)—bowing the knee,
etc. Christ's sensitiveness.

7. Something very irritating in mockery, irony,
etc. Indignation and anger natural, and not sinful, yet to be
restrained lest they become sinful. Slander, misrepresentation, abuse
of the good, blasphemy of things sacred, ludicrous views, pictures,
etc. Nay, the people who throng the doors of a chapel like this, with
persons going to and fro, and insult them.

8. All painful, yet \emph{laetare}. Rejoice in
your desolation; let it be your Lent. Rejoice and leap for joy, for
great is your reward in heaven.

9. Rejoice if you are made like Christ and His
saints.

10. Rejoice, for it is a proof of your real
strength. \emph{Quare fremuerunt gentes}. Our Lord's whisper terrifies
this great country. His vicar, a feeble old man, by a bit of
paper frightens it—\emph{vox Domini super aquas}. Can Wesleyans,
etc., do so? When did Protestantism ever raise a whole state as a
small act of the vicar of Christ has done? You see how the devils
fear. Tall Ishmael is mocking in our streets; a strong boy beating a
small one.

11. Rejoice, for it is an augury for the future.
The desolate has many more children than she that has a husband. So
Protestantism is married to the state. Rejoice not against me, O my
enemy, etc. \footnote{'Rejoice not thou, my
enemy, over me, because I am fallen: I shall arise; when I sit in
darkness, the Lord is my light.'—Mic. vii. 8.}

\xsection{April 6 (Passion Sunday) On the Priesthood of Christ}

1. INTROD.—Go
through the gospel of the day, showing the strangeness of our Lord's
doctrine, and the surprise and contempt of the Jews, in detail—modes
of expression, ideas, objects, different.

2. So it was: it was a different system. If the
world was true, He was not; if He, the world not.

3. They felt it obscurely and in detail, though
He did not speak openly. How would they have felt if our Lord had said
openly, 'I am the priest of the world'? What a great expression! But
this is the truth, as forced on us by today's epistle. What the gospel
says obscurely the epistle speaks out.

4. What is a priest? See how much it implies:
first the need of reconciliation—it has at once to do with sin;
it presupposes sin. When then our Lord is known to come as a priest,
see how the whole face of the world is changed. Describe the world,
how it goes on, buying and selling, etc.; then the light thrown on it
that it is responsible to God, and has ill acquitted itself of that
responsibility.

5. Again, it implies one the highest in rank. The
head of the family was a priest—primogeniture. Hence Christ the Son
of God.

6. Christ then, the Son of God, offers for the
whole world, and that offering is Himself. He who is high as eternity,
whose arms stretch through infinity, is lifted up on the cross for the
sins of the world.

7. And He is a priest for ever. 'Thou art a
priest \emph{for ever} according to the order of Melchisedec.' The
offering of the Mass. Say not it is an \emph{historical} religion,
done and over; it lasts.

8. And as, for ever, so \emph{all things} with
blood. Why? Grace of Christ, and Adam's grace before the fall. Men
'washed their robes in the blood of the Lamb'; 'the blood of Christ
cleanseth,' 1 John i. 7 \footnote{'We have fellowship one
with another, and the blood of Jesus Christ his Son cleanseth us from
all sin.'}.

9. Now turn back and see how different from what
we see—need of \emph{faith}, so says our Lord in the gospel of the
day.

10. And this awful addition, 'He that heareth the
word of God is of God,' etc., John viii. 47 \footnote{'He that is of God
heareth the words of God: therefore you hear him not, because you are
not of God.'}.

11. This a reason for these yearly
commemorations, to bring on us the thought of the unseen world.

\xsection{April 13 (Palm Sunday) On Christ as Hidden}

1. INTROD.—At
this season we veil our images. Why? because the light of our eyes has
gone from us. God is hidden. He showed Himself. He manifested Himself,
but He is gone.

2. This is especially referred to in the gospels
of the past week. Go through them; there is only one, that on
Thursday, which does not obviously refer to it.

3. He had shown Himself through His ministry for
three years as all beautiful—'Blessed is the womb that bore Thee,
and the breasts that gave Thee suck,' etc., Luke xi. 27; 'He has done
all things well,' Mark vii. 37. But now a change. Hidden, bloody
sweat, indignities, blows, etc., called a deceiver, Isa. liii. 3-4 \footnote{'Despised and most
abject of men; a man of sorrows, and acquainted with infirmity. And
his look was as it were hidden and despised, whereupon we esteemed him
not. Surely he hath borne our infirmities, and carried our sorrows:
and we have thought him as it were a leper, and as one struck by God,
and afflicted.'}. Ps. xxii. 6-7 \footnote{'But I am a worm, and no
man; the reproach of men, and the outcast of the people. All they that
saw me have laughed me to scorn: they have spoken with the lips and
wagged the head.'};
even the disciples doubting, etc.

4. Epiphany the beginning, Palm Sunday the end.
From Passion Sunday till now He had been hidden.

5. He did not show himself after the resurrection
'to all the people.' Ascension, and now the Holy Eucharist, a hidden
manna. John xiv. 19 \footnote{'Yet a little while, and
the world seeth me no more; but you see me: because I live and you
shall live.'}.

6. Difference in the mode in which He has been
hidden before and after—'Verily thou art a God that hidest Thyself,'
Isa. xlv. 15—before, all men in ignorance; now, 'the people that sat
in darkness,' etc.; but still, before, He would not be found; now, men
will not seek. John xvi. 16 \footnote{'A little while, and now
you shall not see me: and again, a little while, and you shall see
me.'}.

7. 'The light shineth in the darkness.' Go
through John i. and thus explain the gospel for Thursday [in Passion
week]. \footnote{Luke vii. 36-50.} Magdalene saw
what the Pharisee did not see.

8. Hence He is at once hid and not hid, John
xiv.-xvi.: xiv. 19-23: xvi. 16. Plenty of Catholics in this country,
yet how little they are known. Falsehoods circulated against them.

9. But let us beware how we refuse the light when
it comes. State of the Jews on Palm Sunday. Today's Mass implies they
were visited by grace. A sudden great grace illuminated that
day—(enlarge on the palms, procession, etc.). Alas, how soon it
went!

10. Alas, it was like the 'stronger than he'
taking possession of His house, and the evil spirit returning.

11. O may that not be the case of any of us. We
look up at the cross now, and cannot see Christ's face. A veil, a
thick veil is over it. O let us say, 'My Saviour, let it not be
so really with my soul. I know I cannot always enjoy Thy consolations,
but let not Thy face really be hidden from me. It is my eternal joy.'

Continue

Back

\xsection{April 15 (Tuesday in Holy Week) Maria Addolorata}

1. INTROD.—It
is often said that men in trial act well or ill, according to their
previous life, which then is brought out. They cannot work themselves
up to be martyrs. This applies also to the contemplation of the
sufferings of the saints, and of our Lord and His blessed Mother. I
should like \emph{e.g}. to bring before you the subject of the Mater
Addolorata. But how am I to do so? It depends on yourselves. Are you
familiar with her image? Is she a household word? If so, you will
meditate well; if not, ill.

2. This comes on us at this time of year, when we
wish so much to meditate, and find it so difficult. We shall keep this
time well, according as we have kept the year well. As we have
meditated through the year, so shall we celebrate this season. We
cannot force our minds into love, compassion, gratitude, etc.

3. So as to matters of this world. We hear of
deaths, losses, accidents, etc., with emotion or not, according as we
know the persons, according as the name is familiar to us.

4. I cannot impress this knowledge upon you or
myself, and this makes me almost loth to discourse on these great
topics. The very sight of a crucifix or holy picture, such as we have
in our chambers, should be enough. It is not a matter of words, but of
heart.

5. Think then of her, first as she was, as she
had been, and you will understand what she was in her grief. Go
through her character—so lovely, so perfect, so glorious; the ideal
of painters and poets; yet superhuman, the flower of human
nature—the soul so beaming through her that you could not tell her
features, etc.; so gentle, winning, harmonious, attractive; so loving
towards others; so pained at sorrow and pain; so modest, so retiring:
her voice, her eyes—yet still so chaste and holy that she inspired
holiness. Hence the fulness of the sanctity of St. Joseph: it was
inspired by her.

6. And she had lived with a Son who cannot be
described in this way only, because He is God; who surpassed her
infinitely, but in another order. In the one the attributes of the
Creator, in the other the most perfect work. What a picture! what a
vision! Mother and Son.

7. Next, that Son has left her. And now the news
comes to her that He is to die, to be tortured; that He is to die a
criminal's death of shame and torment; His limbs to be torn to pieces,
etc., and He so innocent. Why, it is worse than killing and torturing
the innocent babe.

8. Under those circumstances, remarkable boldness
in coming to see Him die. Does a mother commonly so act? Here the
perfection of Mary's character. Hagar, 'Let me not see the death
of my son,' Gen. xxi. 15-16 \footnote{'And when the water in the bottle was spent,
she cast the boy under one of the trees that were there. And she went
her way, and sat over against him a great way off, as far as a bow can
carry: for she said, I will not see the boy die. And sitting over
against him, she lifted up her voice and wept.'}.

9. [She saw] Christ bearing the cross. Then at
cross.

10. Our distress at seeing mother's grief, which
we cannot help.

11. On mental pain. Greatest. Christ's mental
pain would have swallowed up even His bodily, had He not willed to
feel it.

\xsection{April 27 (Low Sunday) Faith the Basis of the Christian Empire}

1. INTROD.—Our
Lord came to form a kingdom all over the earth unto the end of time.
And to this end the commission to preach the Gospel, etc. [Mark xvi.
15.]

2. Now observe what a great problem is this. It
had never been done before; it has never been done since, except in
the instance of that kingdom. Why, a large empire extending over many
countries, \emph{mole ruit sua}! On the four empires and others—but
were by an effort and ephemeral.

3. Even the Jews, small as they were, could not
keep together in one. Divisions of Reuben— Benjamin
slaughtered; ten tribes; various sects, Pharisees, etc.

4. (So in Protestants, though fain would be one,
but cannot), but the Catholic Church has lasted one through all time,
and is as much or more one now than ever she has been.

5. Now today's Mass tells us in the epistle and
gospel how it is, and what means God took. It was by means of faith,
which is not only the beginning of all acceptable service, but is the
binding principle of the Church. John xx. 29, St. Thomas, 'Blessed are
they that have not seen, and have believed.' 1 John v. 4, 'This is the
victory which overcometh the world, our faith.' Now consider this
attentively. It is a problem which has never been solved before. He
did not take self-interest, worldly benefit, etc., because they would
not last; and it is what the world proposes to mankind.

6. When He would make a universal empire, He did
not take a book or law for the basis. Some would have said the Bible,
but the event, the divisions of Protestants show it would not do.

7. Not a law, nor a polity, nor episcopacy (as
Anglo-Catholics say). \emph{Quis custodiet}, etc. What shall make
bishops obeyed, etc.?

8. Nor reason (as Liberals and Latitudinarians
will say), for it only arrives at opinion.

9. Nor love (as religious persons may think),
quoting: 'By this shall all men know that you are my disciples,
because you have love one for another'; for concupiscence overcomes
love \footnote{Written over these
words—'love the heart, not the whole body.'}, and the good will never be the many. Some principle must be taken which all can
have.

10. Therefore He took faith—a supernatural
gift. Faith may be possessed by good and bad, and is most influential;
even the bad are made to serve His glory and praise. And it is the
bond, for thus all have common objects. Faith is not easily lost.

11. Hence 1 Peter ii. 9, 'a royal priesthood' as
in yesterday's epistle—hence Jeremias xxxi. [33-34 \footnote{But this shall be the
covenant that I will make with the house of Israel; After those days,
saith the Lord, I will put my law in their bowels, and I will write it
in their hearts; and I will be their God, and they shall be my people.
And they shall teach no more every man his neighbour, and every man
his brother, saying, Know the Lord: for all shall know me, from the
least of them even to the greatest.'}]—and 'our hearts enlarge.' They obey because they believe. It
is not the Church enforces on them faith, but faith obliges them to
take the Church—1 John ii. [20], 'know all things'; [\emph{ib}. 27],
'no one to teach you.'

12. Hence the people never wrong (individuals
indeed, and sometimes nations, may apostatise), but I mean the whole
body. Unlike the Jewish Church. Aaron and calf. Pilate 'willing to
content the people.' But the Christian people cannot be wrong. \emph{Vox
populi}, etc. Hence 'when the Son of Man cometh shall faith be
found,' etc., because of the obscuration under Antichrist.

13. This is our consolation at all times. Our
very sins do not overcome the Church, for faith is independent of sin.

\xsection{May 1 (Month of Mary 1) On Mary as the Pattern of the Natural World}

1. INTROD.—Why
May the month of Mary?

2. Consider what May denotes. It is the youth of
the year; its beauty, grace and purity. Next is its fertility; all
things bud forth. The virgin and mother.

3. See how the ecclesiastical year answers to it.
Our Lord passed His time in the winter—born at Christmas, etc. He
struggles on. We sympathise with Him. We fast in Lent—the rough
weather continues. He comes to His death and burial when the weather
is still bad, yet with promise—fits of better anticipations. He
rises; the weather mends; but, as He was not known as risen, not all
at once. But at length it is not doubtful. He is a risen king, and,
still the weather gets warmer. As a climax May comes, and He gives His
mother.

4. Such is the comparison. Nothing so beautiful
in the natural world as the season when it opens. Nothing so beautiful
in the supernatural as Mary. The more you know of this world the more
beautiful you would know it to be—in other climates—beauty of
scenery, etc., etc.

5. But this is not all. Alas, the world is so
beautiful as to tempt us to idolatry. St. Peter said, 'It is good to
be here' [on Mt. Thabor], but 'It is not good to be in the world.' Say
'Hast thou tracked a traveller round,' etc.; all that is so beautiful
tempts us. Hence all Nature tends to sin (not in itself), etc.

6. Here then a further reason why the month is
given to Mary, viz. in order that we may sanctify the year.

And thus she is a better Eve. Eve, too, in the
beginning may be called the May of the year. She was the first-fruits
of God's beautiful creation. She was the type of all beauty; but alas!
she represented the world also in its fragility. She stayed not in her
original creation. Mary comes as a second and holier Eve, having the
grace of indefectibility and the gift of perseverance from the first,
and teaching us how to use God's gifts without abusing them.

\xsection{May 4 (Second Easter) On the Good Shepherd and Lost Sheep}

1. INTROD.—God
is from eternity and ever blessed in Himself, and needs nothing.

2. On His, being such, taking part in things of
time.

3. An office of ministration—one towards things
physical; a further towards things moral, \emph{i.e}. which have free
will.

4. A further still towards man fallen—on his
waywardness, arising from concupiscence and ignorance—and even the
just [not exempt]—of which ignorance remains more fully in all.
Ignorance is the best estate. This is portrayed in sheep. Other
animals \footnote{Written above—'animals,'
'or swine.'} are fearful,
etc., and represent sinners, but the innocent sheep, ignorant and
helpless, is the fit type of the just. What a picture this gives
us! We are tempted to laugh at sheep, who will not go the right way,
start at every noise, do not know the meaning of anything, and are
obliged to be forced by terror, as by the dog; yet it is our best
image. Our Lord, the Good Shepherd, is obliged to frighten us, etc.,
etc. Yet so patient.

5. O how patient towards us! But more than
patient—the lost sheep, and His laying down His life for it—the
wolf \footnote{Killing the shepherd? \emph{or},
running away with one sheep?}—nay, and that a
one, though one.

6. What is meant by \emph{one}? Because any one
must consider Himself the one. Every one is \emph{worst} to himself:
he alone knows himself.

7. On St. Augustine, this day St. Monica's day.

8. Does the Church lament over you, O \emph{one}
sinner! Hero we are in the happy time of the year—Christ risen and
the month of May come—yet you have not been to your duties, or have
not got absolution, or have fallen again into sin. Mater Ecclesia
deplores you, our blessed Lord deplores you, etc.

\xsection{May 8 (Month of Mary 2) On Mary as Our Mother}

1. INTROD.—Our
Lord from the cross said, 'Behold thy mother.' These words, spoken to
St. John, have been considered by the Church to apply to us all.

2. When our Lord went up on high, He supplied us
with all those relations in a spiritual way which we have in a natural
way. He is all of them—our physician, our teacher, our ruler or
pastor, our father and our mother. Explain how our mother—as bearing
us in pain. 'Shall a mother forget her sucking child?' and in
nourishing us with the milk of the Holy Eucharist.

3. And as St. Peter the one pastor, as St. John,
etc., and the prophets and doctors [as teachers], as priests His
physicians, so He has left His own mother to be our mother.

4. 'Behold thy mother,' etc. Month of Mary.

5. Now consider what is meant by this—a mother's
special gift—fostering care, tenderness, compassion, unfailing love,
so that whenever we would express what is home, and a refuge, and a
retreat, and a school of love, we call it our mother. Our country is
our mother: our schools, colleges, universities, etc., etc. \footnote{\emph{E.g}. Alma mater.} Hence the Church.

6. This is what Mary fulfils to all who seek her
care, and in a far higher degree than any mother can do; for,

7. First, many lose their mothers, or have unkind
mothers, etc. Everything of earth fades.

8. Second, a human mother's standard of things
may be wrong: it may lead from God, hence human affections keep so
many from the Church.

9. Everything human has a chance of fostering
idolatry. What is always present hides the unseen.

10. Our heavenly mother cannot fail and cannot
err, cannot obscure her Son and Lord, but reminds of Him.

11. Let us try to get this filial feeling, though
we can only learn it by degrees, and cannot force ourselves into it.

\xsection{May 11 (Third Easter) Oratory of Brothers— On the General Scope of the Institute}

1. INTROD.—Perhaps
some of you do not know what it is we are offering you in this
association.

2. In one word, which, vague as it is, still is
true, we are meeting together to do something towards saving our
souls.

3. Difficulty of saving the soul. St. Philip's
saying that no one could be expected to get to heaven who had not
feared hell. Scripture texts, 'narrow is the way,' etc.

4. Grace most abundant. Till the last day, we
shall not know how much.

5. But there is a most unaccountable waywardness
in man. It is needless to speculate on it. Every one feels it. He
cannot steady, command, direct himself—inefficacious desires. He is
beaten about here and there at the mercy of the waves. Sloth,
cowardice, anger, fretfulness, sullenness, vanity, curiosity,
concupiscence, ever lead him astray.

6. Hence all serious men look out for a rule of
life to defend them against themselves.

7. This leads many into religion for assistance,
for sympathy, for guidance.

8. The Oratorium Parvum is a slight bond of
sympathy and of mutual assistance.

9. Hence it matters not what we do, or whether
you have anything definite in it beyond this end, if you secure it.

\xsection{May 18 (Fourth Easter) On the World Hating the Catholic Church}

1. INTROD.
In the discourse of which the gospel is part, our Lord speaks of the
world hating us.

2. This remarkable, viz. that we should be hated.
That the Catholic faith is difficult and a stumbling-block is
intelligible—but hateful! Difficult to realise, for we are drawn to
all, and cannot believe they hate us.

3. Consider its beauty—acknowledged by
intellectual men—of its services; of its rites; of its majesty;
doctrine of our lady, etc., etc. Its connection with art, etc., etc.
Paley on Romans xii., in \emph{Evidences}.

4. Yet so our Lord has said—quote John xv.
18-19 \footnote{'If the world hate you, ye
know that it hath hated me before you. If you had been of the world,
the world would love its own: but because you are not of the world,
but I have chosen you out of the world, therefore the world hateth
you.'}, John xvii. 14
\footnote{'I have given them thy
word; and the world hath hated them, because they are not of the
world, as I also am not of the world.'}, and 1 John iii. 1 \footnote{'Behold, what manner of
charity the Father hath bestowed on us, that we should be called and
should be the Sons of God. Therefore the world knoweth not us because
it knew not Him.'}. 'Wonder not if the world hate you.'

5. And what is remarkable further, it is a
prophecy. It has been fulfilled and is fulfilled to this day; it is
literal honest hate. The world is not merely deceived; it has an
instinct, and hates.

6. But more than this, or again, it is a note of
the Church in every age; in the Middle Ages, when religion was
established as much as now.

7. And none but the Church thus hated. So that
our Lord's prophecy falls on us, and connects us with the apostles.

8. Others, indeed, by an accident and for a time.

9. For sects have (1) something true and good in
them; (2) are extravagant; and these two things make them persecuted.

10. But it is for a time. The truth goes off, and
the extravagance—they tame down; thus the Methodists and the
Quakers.

11. But Catholics, nothing of this—sober—by
token men of the world get on with us.

12. Yet the suspicion, irritability, impatience,
etc., etc.—Demoniacs, and it is the devil's work.

13. This must not make us misanthropic, but cast
us on the unseen world and purify our motives. This one benefit of the
present agitation.

\xsection{June 8 (Whitsunday) The Life-Giving Spirit}

1. INTROD.—We
have what we have waited for. Paschal time is not only a time of
rejoicing, but of waiting for a gift. The \emph{whole creation groaning},
etc. Hence, now being the end, we go no further, but date our time
from Pentecost.

2. The gift of today set up the Church, hence it
is said to be a vehement wind filling the house. Solomon's temple
filled with the glory, as the sweet nard filled the house.

3. For up to this date the Church was not formed.
The multitude who followed Christ was but matter \footnote{I.e. \emph{materia sine
forma}. A bold figure or comparison which must not be taken too
literally.}. They were not a body filled with Christ. Christ was with
them, but external \footnote{'External' is followed
by some words which are nearly illegible. They look like 'for was not
a form.'};
they were not confirmed. They were all scattered abroad as sheep.
Hence as an individual may have first actual, then habitual grace—so
the \emph{multitudo fidelium} all Paschal time is begging to be the
bride of Christ.

4. Now then the Spirit came down, to gather
together the children of God, etc., all those who had fled away, etc.;
returned—3000-5000 \footnote{'There were added in
that day about three thousand' (Acts ii. 41). 'The number of the men
was made five thousand' (Acts iv. 4).}.

5. Like the resurrection of dry bones, Ezech.
xxxvii.

6. Such is the power, the manifestation, of the
Spirit; thus sudden, thus gentle, thus silent. It is life from
death—what health is after sickness. It makes young. Oh what a gift
is this! Who would not wonder if a physician could make an old
man young? See him, unable to do more than grope about, his limbs
stiff, his face withered, etc., etc. But the physician comes, and
health and comeliness and vigour return, etc. This is what is
fulfilled by the power of the Spirit, in a measure in individuals,
certainly in the body.

7. And is it possible such is in store for
England?—(explain). Nothing unexpected, nothing too difficult. It is
grace, yet spreading not at once.

8. Prayer for it. Never so much prayer as now.

\xsection{June 29 The Rock of the Church—St. Peter and St. Paul}

1. INTROD.—If
nothing else could be said for our holy religion than the topic of
this day suggests, I should think it abundantly proved.

2. At present we see a vast body with vast power
all over the earth. We know how great the British power. Such (I don't
say with the same weapons) is the Catholic Roman Church, nay, far more
fully, because it reigns more directly—not through other powers, as
the British in India, etc.

3. Now look at the British Empire. What is its
peculiarity? It has grown, as it happens, in the course of a century;
but never mind that. The Catholic Church has never grown; it always
has been [what it now is].

4. Now one point is the great youth of all other
powers compared with the Catholic Church, but I won't dwell on that.

5. What I wish to dwell on is, that whether they
be young or old, they have had a growth—a beginning, a progress, and
an ending like a tree—(enlarge). Look at the great Roman Empire;
Gibbon has written its decline and fall.

6. No one can write, I will not say the decline
and fall, but the growth of the Catholic Church—(explain). I don't
say it has not developed in many respects; in consolidation, in
temporal power, in definition of doctrine, in experience; but it is
stationary.

7. Look back five centuries. Just the
same—stationary. Look back ten, etc. No, it expanded at once in the
apostles, and has ever since possessed the earth. 'Blessed are the
meek,' etc.

8. But further. Suppose not only the British
Empire had lasted long, that not only it was stationary, being just
what it was in Alfred's time; but supposing Alfred declared it should
last; suppose all the kings who ever were declared it would
last—moreover, in consequence of an old prophecy in Julius Caesar's
time, etc.

9. This fulfilled in the Church—St. Leo 1400
years ago. Our Lord's test—the rock—how exactly it fulfils it.
'The house upon the sand'—Protestantism.

\xsection{August 10 (Ninth Pentecost) On the Death of the Sinner}

1. INTROD.—The
gospel—our Lord weeping over Jerusalem. Particulars of it. The Jews
so little aware. They thought a great conqueror was coming to
them. Their great infatuation. They had a vast future (they thought)
before them. The Temple rebuilt. Our Lord saw through it all.

2. Application to the soul of the individual.
Type of sinner in death. Our Lord looking and prophesying
ill—(particulars). 'Cast a trench,' 'hedge them in.'

3. 'Hedge them in.' Yes, Satan will take
possession of him; keep God out; keep him all to himself. What a
portentous thought!

4. Christ foresees it, weeps over the man, but He
leaves him.

5. But does He not give grace? Yes, but it is
ineffectual.

6. Why does He not give more? What is that to the
purpose? He does not.

7. We cannot change things by asking questions.
Why does He punish him? Can you change it by disputing? Your wisdom is
to take things as they are, and submit and improve them. Is not this
the way you do with this world? You do not quarrel with the wind, the
flame, etc., but use them. Our Lord with Judas. His denunciations of
eternal woe. His own sufferings [are as if He said], 'I say not why,
but I suffer.'

8. Well, then, the fact is this. The sinner
generally is thus 'walled in.' \emph{Vide} St. Alfonso on this day.

9. Saul. Antiochus.

10. Encircled—wild beasts. Sins as faithful
friends who encircle you in their arms.

11. The priest's prayers in vain.

12. The sacraments in vain.

13. Our Lady not. Ave Maria! St. Andrew Avellino!

14. Let us ask her to intercede for us.

\xsection{August 31 (Twelfth Pentecost) On Christ the Good Samaritan}

1. INTROD.—Go
through the parable briefly, applying it in a secondary sense to the
sinner and Christ.

2. In the parable the traveller was robbed
against his will, the sinner with his will. Satan cannot conquer us
against ourselves. Eve—temptation, etc.; it is a bargain.

3. Thus he gets from us justice, habitual grace,
etc., nay, part of our mere nature, for he leaves wounds. Thus he may
be said to suck the blood from us. A vampire bat sucking the blood
out. All terrible stories of ghosts, etc., etc., are fulfilled in him
who is the archetype of evil.

4. He has the best of the bargain, as is evident.
What have we to show for it?—there are improvident spendthrifts who
anticipate their money, and get nothing for it. What have we to show
if we have given ourselves to Satan?

5. (1) Those who commit frauds—ill gains go.
(2) Anger, swearing and blasphemy—what remains? (3) Sensuality is
more rational, because men get something.

6. Yet in a few years where is it all? Let a man
enjoy life, let him be rich, but he gets old, and then! Wisdom
[v. 8]. 'What hath pride profited us?'

7. Thus Satan has the best of the bargain, and we
lie like the traveller.

8. Nothing of this world can help us—priest or
Levite: there we should lie for ever, etc.

9. Christ alone, by His sacraments.

10. Mind He is a Samaritan—so
Nazareth—because the Catholic Church is hated. She is the good
Samaritan to Protestants. Observe again the text, 'He who showed mercy
to him.' Has the Catholic Church or Protestantism done this for us?

\xsection{September 28 (Sixteenth Pentecost) On the M. Addolorata—the Seven Dolours}

1. INTROD.—The
usual representation which painters make of our Lord and His mother is
that of virgin and child. Describe the peaceful virgin, secure because
she has Him, and He the Life and Light. Hence she the Seat of Wisdom,
etc., etc.

2. But let thirty years pass, and there is a
great change come over the picture. It melts into something different.
He is taken up from her soft arms. He is lifted aloft. Something else
embraces Him. He is in the arms of the cross. There He lies not
easily, etc. He has grown to man's estate. He has been scourged, etc.
And she is standing still, but it is at His feet. She can be of no use
to Him; she can only lament. How the group is changed! He is
covered with wounds; she is almost killed with grief.

Such is the picture which the Church puts before
us today, and that because, we may suppose, Easter is so long past.

3. Well, as to the sufferings of the Son of God,
they are awful mysteries; but they need not surprise us, for He comes
to suffer. He indeed might have saved us without suffering, but it was
in fact bound up in His coming. He was a combatant—combatants
suffer. He was prophesied as a warrior and man of blood. He fought
with the devil. He fought with sin, not indeed His own, but sin was
imputed to Him. He came in the place and character of a sinner: no
wonder He should suffer.

4. But there was one who neither sinned nor took
on her the character of a sinner. What had she to do with blood, or
wounds, or grief? She had ever lived in private; she bore Him without
pain; she had never come forward. She had on the whole been sheltered
from the world, yet she suffered. This makes Mary's suffering so
peculiar. She is the queen of martyrs.

5. Yet she too was to suffer. She is innocent, so
harmless, not provoking the devil, etc. She was to suffer, and be the
queen of martyrs. Joseph was taken away; she remained.

6. It is true she was not to undergo that bodily
pain and violent death which literally makes a martyr. He alone
suffered all who died for all. He alone suffered bodily and mentally.
Her tender flesh was not scourged, but His was; her virginal form was
not rudely exposed, but His was. All this would have been
unseemly and unnecessary. He was to save us by that body and blood
which she furnished; not she. He was to be made a \emph{sacrament} for
us as well as a sacrifice.

7. Yet she was privileged to share the acutest
part of His sufferings, the mental, once she came into the midst, at
His crucifixion.

8. Mental pain all in a moment, like a spear;
despondency, sinking of nerves; no support.

9. Yet she stood.

10. Surely it quite changed her outward
appearance to the end of her life.

\xsection{October 26 (Twentieth Pentecost) On the Patrocinium B.V.M.}

1. INTROD.—This
festival of our Lady [is] more immediately interesting to us than any,
because by it we are made over to her and she to us. [In] the
Incarnation, the Assumption, etc. [we celebrate more immediately her
relations to Almighty God], but [in] this [feast we call to mind
particularly her relations to ourselves].

2. It is like the divine works to turn things to \emph{account}.
Thus, though she subserved the Redeemer, she also subserves the
redeemed. Hers is a \emph{ministry} to us, and it was to Him
originally.

3. As a pope makes a congregation over to a
cardinal, or a king gives some one a ring, etc., saying, 'Whatever you
want, send the ring and you shall have it.'

4. Thus she is the fount of mercy, as a
magistrate of justice, etc.

5. Hence Protestant absurdity of saying [that] we
rate her more merciful than Christ. Christ is the judge also. Show
what is meant by it. Can a ring be merciful?

6. As this [is] the feast most intimately
interesting to us, so we hear much of this character and office in
Scripture, in the Holy Fathers.

7. Gen. iii., Apoc. xii.—\emph{Advocata} with
clients; mother of all living. 'Behold thy mother,' John xix. 27.

8. Hence first instances in history represent her
in this character—St. Gregory Thaumaturgus—St. Justina—against
unbelief, against impurity respectively \footnote{For the stories here
referred to, see \emph{Development of Christian Doctrine}, pp. 417,
418.}.

9. St. Gregory Thaumaturgus has a creed given
him—St. Ignatius—St. Philip.

10. Experience of all saints.

11. Let us use it, for living, for dead, for
young, for old. The two first instances [given] above are [of] a young
man, a young woman.

\xsection{1851 The Immaculate Conception the Antagonist of an Impure Age}

1. INTROD.—The
world always the same, and its history the same. It is always sinning,
always going on to punishment. Judgments and visitations always
[coming] upon it. Christ always coming [in judgment].

\emph{Sin provoking wrath}.

2. This is seen in the judgments on cities for
their crimes—Nineveh, Babylon, etc., and above all, Sodom and
Gomorrah—all figures of the end of the world.

3. And especially eras—the deluge—the
Christian era \footnote{The allusion must be to
the destruction of the Temple and the rejection of the Jews.}—the
end of the world. And they are compared together in Scripture, Matt.
xxiv., etc.

4. \emph{What} sin (provoking wrath)? Sensuality.

As the loss of vital powers brings on dissolution
of [the] body, so when passion emancipates itself from conscience, the
death of the world.

5. The truth is, that the flesh is so strong, it
is always struggling against conscience. It is like a wild beast in a
cage, ever trying to get out, and but slowly subdued. Heavy things
fall; steam rises up. So with concupiscence; and hence St. Peter
[speaks of] 'The corruption of that concupiscence which is in the
world,' 2 Peter i. 4.

6. Now as this goes on in a state, reason becomes
infidel and the conscience goes, and then there is nothing to restrain
concupiscence.

7. Hence we are sure (\emph{exceptis excipiendis})
that wherever there is not religion there is immorality. What is to
keep a man from indulgence?

8. Statesmen see this so well that they advocate
religion.

9. Hence [came the] deluge—[the] Christian era
\footnote{Cf. Romans i.}. [hence will come
the] end of the world, [\emph{i.e}. when] infidelity [has] brought in
sensuality.

10. This age [is] an impure age.

11. Hence [the] B[lessed] V[irgin] M[ary] [is]
attacked.

12. Hence [the devotion to] the Immaculate
Conception is so apposite.

\xsection{1851 The Special Charm of Christmas}

1. INTROD.—[The]
two chief festivals [of the Church are] Easter and Christmas; [of
these] Easter [is] the greater.

2. Yet somehow we adorn our churches more
brightly and spontaneously, now than then. There is more of heart,
apparently, in what we do. And there is an inexpressible charm over
all. The midnight Mass, the three Masses. The special representations,
whether the Stable or the Infant. [Again, the singing of] carols.

3. Why is this? Christmas is easier to understand
to the mass of men; it comes home to them more readily, and imposes an
easier duty on our worship.

4. It is the difference between coming and going.
The apostles felt that sorrow filled their hearts [at the going of the
Lord]. \emph{Mane nobiscum Domine}.

Easter is the feast of the perfect. If we were
perfect, we should rejoice in Easter the more [of the two festivals].
In the one Christ comes to us, in the other we go to Christ.

5. All our human feelings are soothed by
Christmas— Abraham had to leave his country.—We naturally do
not like to move. We are allowed to remain at home: Christ comes to us
as our guest.

6. And coming, He \emph{brightens} everything. He
does not take away, He adds. He adds grace to Nature. If at any time
we might love the world, it is now. If at any time, [it is when He is
come to be our Emmanuel].

7. He makes the world our home, for he deigns to
be the light of it. He sanctifies families with the image of Mary and
Jesus. And where there is no home in a family, then He brings us all
together in one family in church. The midnight Mass is our holy
celebration [of Christmas], eclipsing the world's merrymaking.

8. And we think of Him who put off all His glory,
of which our celebrations are but a type. The priestly vestments a
type of His glory, [which He put off in order] to come into this bleak
prison and suffer for us.

9. Let us rejoice \emph{in Him}.

\xsection{December 28 (Sunday in Octave of Christmas) On Christian Peace}

1. INTROD.—Peace
is, as we all know, the special promise of the Gospel.

2. Isa. xl., Rom. xi., Isa. ix., 'Peace on
earth.' 'Peace I leave with you' [John xiv. 27]. 'Peace be with you,'
and St. Paul 'making peace' [Rom. xii. 18].

3. This is the great want of human nature. It is
what all men are seeking; they are restless because they have not
peace. They always think the time will come when they shall be happy,
yet it never comes.

The schoolboy—the young man—the soul in
disorder.

4. Hence it forms to itself notions of peace and
happiness, [such notions as we find in] novels, tales, poems; [notions
which are] imaginary.

And above all, [notions of] religion. It attempts
to make religions for itself, where everything shall be beautiful,
etc.

5. Thus it goes on, and then it looks down on
Christianity. Christ Jesus (they say) does not bring peace.

This is the way of so many infidels now. They say
they want a religion more beautiful, more comfortable than the Gospel.
They point to the gloominess of Catholicity—nothing sunny and
bright—confession, penance, mortifications of the senses and the
will; monks, etc., etc.; and they say this is a dreary religion, and
they could form a better one. They say they could form a better god
than the Father of Jesus Christ—a god of their own dreams; [they
could form] a religion without sin and without punishment.

6. Thus they go on; but what is this but to say,
'Peace, peace, where there is no peace'?

7. The more haste, the worse speed. Shrubs
putting out their leaves too soon—the hare and the tortoise. 'The
end is the trial.'

8. The truth is, once beautifulness and peace \emph{did}
come first, viz. in the Garden of Eden. Since then there has been
a fall. There must be a restoration, and it is painful.

9. Contrasting pantheism with true religion,
recollect we are only in process, etc., and therefore we look to
disadvantage.

Hence religion gloomy, because it is an
intermediate state.

10. But we look forward for peace to the next
world.

\xsection{January 11, 1852 On the Epiphany, as Christ's Reign Manifested to Faith}

1. INTROD.—On
the peculiarity of this octave.

2. Viz. no saint's day in it. Contrast Christmas.
Contrast Easter and Whitsun as not \emph{perfect} \footnote{'The octave of Christmas
is full of saints' days—St. Stephen, St. John, etc. Those of Easter
and Pentecost are cut short by Low Sunday and Trinity Sunday
respectively.}, [the latter containing] fast days. Contrast [octaves of the]
Ascension, Corpus Christi, [the] Assumption.

3. Why? Christ [is] a king, and we anticipate His
reign. It is the season most nearly typical of heaven.

4. Now, how was this fulfilled? His palace a
stable, His throne a manger—(enlarge).

5. Here it was the three kings came. They came a
long way to see, what? The poor child of a poor woman—(describe).
They entered. Mary drew off the covering cast over the sleeping Child.
They gazed, etc.; they offered gifts; they adored.

6. What a remarkable scene! And this was the manifestation of His glory! For this they had travelled their weary
way!

7. Describe what they had to go through—the
wonder of their people—why were they setting off?—Then, they did
not know whither they were going, etc.

8. Describe their state of mind. They \emph{knew}
they ought to go; they \emph{knew} there was something to find.

9. Enlarge on faith and reason, and explain.

10. This is that faith which is the beginning of
salvation in every age, and the greatest specimen [of it]. It is like
St. Thomas's, with less evidence, 'My Lord and my God.'

11. Greater than, yet like that in the Holy
Eucharist.

\xsection{(No date) Self-Denial in Comforts}

\footnote{'Not used as yet.'}

1. INTROD.—Contrast
between men and other animals, that they [the latter] are sufficient
for themselves.

2. The Creator has so ordained things that
everything is there, where it can flourish. External nature and the
nature of animals correspond.

3. Thus warmth and air, abode and food given to
all; and when external nature is likely to press hard, [there are
given] internal means of meeting it, \emph{e.g}. furs, or hardiness,
or instincts, etc., etc.

4. But man an exception. Strange to say, if born
in a state of simple nature, he would die. His delicate frame
ill-suited to the elements, etc. He needs clothes, a house, etc.

5. Revelation tells us it was [not] always so,
not in his creation, for he was in Paradise; but it is one of the
consequences of the fall.

6. Hence man is ever striving to get out of this
state of fallen nature (so far [as concerns the needs of his body]). \emph{Curis
acuens mortalia corda} \footnote{Virgil, \emph{Georgics},
i. 123.}.
Hence his arts, etc. Hence his loom and his carpentering, etc., etc. I
may say the whole course of life is escaping from this state of fallen
nature, i.e. \emph{as regards the body}: for the worst penalties, viz.
the wounds of the soul, he leaves untouched.

7. Till at length he surrounds himself with
comforts. They are called comforts, and make the whole world minister
to him, and make his home and his rest here.

8. Now it is startling how our Lord took just the
reverse course. He threw away comforts—born in a stable, carried
into Egypt, not a place to lay His head, etc.

9. WHAT AN
AWFUL CONTRAST
between Him and us—(enlarge).

10. Let us take a lesson from it. We have here no
abiding city, etc.

\xsection{January 25 On the Character of the Christian Election— St. Paul's Conversion}

1. INTROD.—A
great principle—not many mighty, noble, wise, called.

2. St. Paul—\emph{exceptio probat regulam}.

3. Still, such is the awful phenomenon in every
age. When Catholicism [is] national, then indeed \emph{all} Catholics.
But when the Church acts freely, then the same characteristic as at
the first.

4. \emph{E.g}. the Church now [is] what it was in
the apostles' time—few learned, etc.

5. It is a most wonderful phenomenon how it goes
on. Why it does not fall to pieces, [seeing there are but] just enough
of learned, etc., men to keep it going.

6. And here we see the reason, viz. that it may
be manifestly God's doing.

7. This [is] set forth in Epistle to Corinthians
\footnote{'For see your vocation,
brethren, that there are not many wise according to the flesh, not
many mighty, not many noble: but the foolish things of the world hath
God chosen, that he may confound the wise,' etc.—1 Cor. i. 26-27.}.

8. Describe how riches, power, learning, nay,
natural goodness, often prejudice [men] against [the] Gospel.

9. On self-sufficient virtue, on putting up our
own feelings, etc., as the rule. These men \emph{complete in themselves}
…

10. Apoc. iii. [vv. 1, 2, 8, 17, etc.], 1 Cor.
iv. [vv. 4, 7, etc.], and not thrown upon God.

11. But I have [not] got at the bottom of the
mystery. I have been speaking only of the \emph{called}, but [there
is] a second [and] wonderful mystery perfectly hid from us—who are
the \emph{chosen}?

12. The visible Church does not stand for the
invisible future elect. Those rich men who are in the Church may
be holier than the poor. So many of the saints [were both] rich and
noble men.

13. [The] moral is, the necessity of waiting on
God's grace, and not quenching it.

\xsection{February 1 (Fourth Epiphany) Present State of Our Oratory}

1. INTROD.—This
day, commencing with this evening, is a great day for our
Congregation, for it is the anniversary of its establishment in
England.

2. This day four years [ago in England], and
again this day three, in Birmingham.

3. The Purification, though not the greatest
feast, [is] a good day, suitable to those who are beginning a work in
an heretical country.—

\begin{quote}

(1) It is a forlorn day in winter.

(2) Christmas gone, Lent coming.

(3) A little child and a poor mother coming to the Temple.

(4) Purification reminds us of necessity of purity of heart.

\end{quote}

4. To me \emph{especially} interesting, for it
has been my great feast-day for thirty years. Thirty years this year
since I was brought under the shadow of our Lady \footnote{Elected Fellow of Oriel
(\emph{the House or Hall of Blessed Mary}) in 1822.}, whom I ever wished to love and honour more and more.
And thus, when I became a Catholic, it was the day of the
Congregation, etc.

5. God has blessed us through her intercession
for three years in this place (Alcester St.). We have gradually
prospered, year after year, and now a more definite establishment at
Edgbaston.

6. Everything has come naturally, like a tree
growing, and we hope it will still [grow].

7. About the Achilli matter. When it first arose,
I said, 'The devil is here. Look not on prosecutor, lawyers, friends,
etc. They are all weapons of the devil.' A NET—pulling
strings close. \emph{Vide} Psalter.

8. Therefore the remedy was prayer. What showed
this more, was the extreme difficulty [of the case].

9. Eph. vi. 12, 'We wrestle,' etc.

10. \emph{Number} of prayers offered.

11. The sequel has shown it—a great noise
ending in nothing, so as to disappoint—first a roaring lion, then a
serpent slinking away; so it is \emph{now}. People will say, 'Oh,
there was no great danger.'

12. If we fail, it will be because we do not pray
enough.

13. Therefore commend ourselves to our Lady.

\xsection{August 15 (Eleventh Pentecost) On Our Lady as in the Body}

1. INTROD.—\emph{Question}.—Whether
this feast, [the Assumption, is] not inconsistent with the Immaculate
Conception; for why should our Lady die if she did not inherit Adam's
sin?

2. \emph{Answer}.—Because she was under the
laws of fallen Nature, and inherited its evils, except so far as sin
[is concerned]. Thus our Blessed Lord [suffered fatigue, pain and
death]. Thus she had not perfect knowledge from the first. She had
need of shelter, clothing, etc., not in a garden [as our first parents
were].

3. Hence, since all men die, she died. Our Lord
died.

4. Yet even as regards the body, our Lord
observed a special dispensation about her. Hence she was not only
protected from diseases, but from torture, wounds, etc.

5. It was becoming that she who was \emph{inviolata},
\emph{intemerata}, should have no wound.

6. The difference between men and women as to
warfare. The women protected and sit at home. How many a wife, or
sister or daughter, suffers in mind, and you hear them say, 'O that I
were a man!' And they suffer in soul, [as the] saints about \emph{the
cross} [who were] not martyrs [suffered].

And hence Mary had a sword through her [heart].
Mental pains, like bodily. And this her pain.

7. And hence she brings before us the remarkable
instance of a soul suffering, yet not the body.

8. She lived therefore to the full age of human
kind. [In this she was] different from our Lord.

9. What a picture this puts before us! Fancy her
thirty, forty, fifty, sixty, looking still so beautiful and young, not
fading, more heavenly every year; so that she grew in beauty, and the
soul always grew in grace and merit.

10. And then, fancy the increased pain at the
absence of Christ, [for she lived] fifteen or sixteen years without
Him!

11. On the long life and waiting of the
antediluvian patriarchs—Jacob's 'I have waited for Thy salvation, O
Lord'; Moses; Daniel; the souls in Limbo Patrum like Mary, though the
time [of her waiting] shorter.

It was like purgatory, waiting for Christ's face;
except with merit and not for sin.

12. Hence [it is] not wonderful [that] it is a
pious belief that she died from love. This alone could kill that body.
It was a \emph{contest} between body and soul. The body so strong, the
soul so desirous to see God. No disease could kill that body. What
killed it? The soul, that it might get to heaven.

13. (1) By languishing; (2) by \emph{striving} to
get loose.

14. Hence [it was] fitting that, when she did get
loose, her Son should not let the body be so overmatched and overcome,
but at once that the soul had got the victory, He raised up the body
without corruption.

15. Our Advocate in heaven.

\xsection{December 8, 1853 On the Peculiarities and Consequent Sufferings of Our Lady's Sanctity}

1. INTROD.—Genesis
iii. We cannot be surprised at our Lady's Immaculate Conception.

2. The reason is so plain that it seems
axiomatic, nor, though it has been a point of controversy, do I think
any holy person in any age has ever really denied it; if they seemed
to do so, it was something else they opposed.

3. Has not God required holiness wherever He has
come?—(1) burning bush \footnote{'Come not nigh hither:
put off the shoes from thy feet, for the place whereon thou standest
is holy ground.'—Exodus iii. 5.}
(2) 'Be ye holy, for,' etc. \footnote{'Sanctify yourselves,
and be ye holy: because I am the Lord your God.'—Levit. xx. 7.}
(3) priests' purifications; (4) consecration of Temple and tabernacle;
(5) without sanctity, no one, etc.; (6) Confession before Communion.
If, then, our Lady was to hold God, etc.

4. Still more, if from her flesh, etc.

5. Hence, though the Church has never proposed it
as a point of faith \footnote{See p. 116, sec. 6.},
it is not difficult to conceive it should be one, and there has been a
growing wish that the Church could find that it was part of the
original dogma. Indeed, it is almost saying what has been said in
other words, for if no venial sin, must there not be Immaculate
Conception?

6. Now to explain what the doctrine is. Eve, as
Adam, had been not only created, but constituted holy, grace
given, etc. Eve was without sin from the first, filled with grace from
the first.

7. When Adam and Eve fell, this grace was
removed; and this constitutes the state of original sin. Describe war
of passions, etc. This is the state into which the soul of man comes
on its creation. Nothing can hinder it but a \emph{return} of the
great gift.

8. Now in the text \emph{she} was to restore, and
more, the age of Paradise. She was promised upon the fall. Eve has
been deceived. \emph{She} was to conquer. How would this be the case,
unless Mary had at least the gifts which Eve had?

9. We believe, then, that Mary had this
sanctifying grace from the moment she began to be.

10. This being the case, I wish you to
contemplate her state. First, her wonderful state \emph{before} her
birth. She had knowledge and the use of reason from the first. This
[was] necessary for love, therefore she had it. What knowledge ?—(1)
supernatural, (2) not physical, (3) of divine objects—[as] the Holy
Trinity, which commonly requires external instruction.

11. Not of sin. Here difference from our Lord, by
way of illustration.

12. Consequences—her idea of disobedience; no
recognition of separate sins. It is only temptation brings this
knowledge home to ordinarily innocent people. She would know she \emph{could}
disobey if she would, but it was like willing to jump down a
precipice; she was sure not.

13. She would not be able to comprehend \emph{how}
people came to sin. And if the supernatural information told her the
fact, she would take it of necessity simply on faith.

14. Let us suppose her passing out of her first
infancy. She is taught external things. She is taught to read. She
learns Scripture. She hears of the sins of her people. She has to take
it on faith.

15. She is a little child, not three years old,
but she cannot pass her mother's threshold but the very scent of the
world overpowers her. It is a bad world: how is she to live in it? She
understands many things: she does not understand it.

16. At length she is taken to the Temple, and
there she lives ten years—what a blessed change!—in the presence
of her God. But even then, though she looks at the priests as God's
ministers, yet, alas, how is she to bear the world, even in its best
shape!

17. Time comes that she must return. Alas! she
has a growing suffering; she is thrown on the world. Do you not see
that there cannot be a more insufferable penance than to be thus
perfectly holy, yet in this unholy world? I know she has full
consolations, but she is in a sinful world, and has the \emph{poena damni}.

18. She looks back on the happy mysterious time
which passed between the creation of her soul and her birth.

19. What a comfort to find herself transferred to
St. Joseph's charge! This is the first alleviation, for a time, which
God gives to her penance.

20. Then the angel Gabriel. Ah! here is an
alleviation indeed. She is no longer desolate for thirty years.

21. Prophecy of Simeon. Loss of Jesus at twelve
years old. His ministry. His crucifixion.

22. O Mary, you were young, now you are
old—old, yet not as other old people, dwindling, but increasing in
grace to the end. But oh what a penance! \emph{O commutationem}! \footnote{See Note
9, p. 337.}

23. And to go about the world! to go to Ephesus!
Oh wonderful! Your journey to St. Elizabeth, to Bethlehem, was with
your Son. Now you journey further without Him.

24. CONCLUSION.—The
holier we are, the less of this world [can we endure].

25. Fitting to be the Feast of the Congregation
[of the Oratory] since, especially in a country like this, we must
begin with holiness.

\xsection{July 23, 1854 (Seventh Pentecost) [Nature and Grace]}

1. INTROD.—Text:
'Jesus loved him.'

2. Explain the circumstances. And then we come to
this anomaly—that God loves \emph{for something in them} those who
will not obey His call.

3. Now this is a difficulty surely which we feel
ourselves. People are (1) amiable, (2) conscientious, (3) benevolent;
they do many good actions, but are not Catholics; or not in God's
grace.

4. Explanation. Nature not simply evil. We do not
say that Nature cannot do good actions without God's grace. Far from
it. Instances of great heathens.

5. What we say is that no one can get to \emph{heaven}
without God's grace.

6. Contrast of two states as on two levels: (1)
moral virtues with 'their reward,' industry, etc., has a reward in \emph{this}
life.

7. (2) Spiritual state of grace. It has all these
virtues and a \emph{good} deal more, and especially faith.

8. This is why faith is so necessary. Explain
what faith is, as a door. It is a sight, [power of vision]. It is
looking up to God. When we \emph{pray}, we have faith, etc., etc.

9. Now what an awful thought this is when you
look at the world—if something \emph{more than} Nature is necessary
for salvation.

10. People say, 'If I do my duty'—'He was such
a good father'; 'He was upright,' etc., etc. All this is good, but by
itself will not bring a man to heaven.

11. When you think what heaven is, is it
wonderful? Think of our sins. Is it wonderful God does not give
forgiveness to Nature?

12. Is it wonderful that grace alone can get
repentance?

13. Let us turn this [over] in our hearts.

\xsection{August 6 (Ninth Pentecost) 'No One can Come to Me except the Father,' etc.}

1. INTROD.—I
said, a fortnight ago, that when we saw what is good in those who are
external to the Church, we must say that it is from Nature, and did
not prove that such persons were in God's favour.

2. This is true, but you may insist that
Protestants, [as well as] those who do not believe that Christ is God,
etc., etc., have an \emph{appearance of religion}; that you cannot
deny your senses; that as you believe them in other things, \emph{e.g}.
that they are honest, so you must here; that they must have \emph{grace}
if they have \emph{faith} and \emph{love}, and therefore must be in
God's favour and in the way to heaven.

3. I am going, then, to give a further answer.
First, I \emph{grant} they show often real faith, real hope, real
love, and that it comes from grace, and that while they obey that
grace, etc., they are in a certain sense in the way to heaven; but
still this is quite consistent with what I have said.

4. All men in God's wrath. How are they brought
out of it? By God's grace coming like a robe (the ordinary way in
baptism, and afterwards by penance) and making them pleasing to Him. \emph{Few}
are in this state. It is called \emph{the state of grace}, and it is
the state to die in, and since we may die any moment, the state to
live in, if we would be safe.

5. And though few are in this state, it is the
state in which God wills all to be in, for Christ died for \emph{all}.

6. As He sends out \emph{preachers} all over the
earth, and as still more, guardian angels, so \emph{graces}.

7. To all He gives grace, even to those who are
not yet in His favour, or \emph{in} grace. He gives them this grace in
order that they may \emph{come} into a state of grace—heathens,
idolaters, Jews, heretics, all who are not Catholics. All have grace
without knowing it—'[even when they are] without God'—while they
are \emph{far} from Him.

8. When you see men, not Catholics, will good
things, acknowledge it, but understand \emph{why} they [these graces]
are given, viz. like preachers, to bring them into the Church; and
they are brought into the Church by obeying them, though not all at
once.

9. Instances. A kindness to a Catholic [or to]
any strangers—generosity—leads to hearing something about
Catholicity. More grace [follows]. [The man] resisting [at first], but
yielding [gradually], etc., etc., till he is brought in.

10. Again, purity may keep a person from bad
company. This throws time on his hands. He passes a Catholic chapel,
he goes in, and he is attracted by a picture of our blessed Lady, etc.

11. All the while these persons may be out of God's
favour, not yet justified, though He has died for them and wishes to
save them, and is gradually drawing them.

How [about] heathen? Sends angels?

12. And thus I answer the question with which I
began.

13. I entreat all those who are in doubt or
inquiring to be faithful to grace, and they will be brought in.

Continue

Back

\xsection{August 20 (Eleventh Pentecost—Octave of the Assumption) [Rejoicing with Mary]}

1. INTROD.—This,
we know, is one of the most joyful weeks of the year. Our Lord's
Resurrection is, of course, pre-eminently [joyful] (and in like manner His Nativity), as He is above all. But this week is unlike most
other feasts connected with Him, and rather stands at the head of the
saints' feasts, and this is its peculiarity. I will explain.

2. The one idea is congratulation. \emph{Congratulamini
mihi, quia cum essem parvula}. Congratulation is a special feeling.
Not in Christmas, or [in] any act of His economy [or of] His Passion,
not in Pentecost [nor] Corpus Christi, nor in the Sacred Heart, [do we
congratulate]. We congratulate when some great good has come to
another. We do not (strictly speaking) congratulate ourselves, though
we may each other. We congratulate martyrs and saints, etc.

3. Now this life tells us what congratulation is.
We congratulate persons on good fortune, which does not concern us
[ourselves], on preferment, on a fortune, on escaping danger, on
marriages and births, on honours, etc.

4. On Catholicity only [\emph{i.e}. alone] \emph{realising
unseen} things and carrying human feelings into the supernatural
world. Hence care of those who [have] departed—purgatory—heaven.

5. Now consider St. Paul's words. \emph{Gaudere cum
gaudentibus, flere cum flentibus}—congratulation and compassion,
or pity [opposed to] two bad states of mind, [\emph{epichairokakia}]
and envy. Congratulation and compassion both disinterested and
unselfish, but congratulation the more. What is so beautiful as to see
in the case of brothers and sisters, (\emph{e.g}.) where a younger
rejoices in the gain of an elder, etc.

6. Now we congratulate Mary at this time of year, after her long waiting—sixty years. What a purgatory! This
very circumstance that all her life was God's, made the trial longer.
But now, as Christ ascended, so has she.

7. But again, even this congratulation has often
something selfish in it; men hope to get something for themselves
through their promoted friend. This is true also in the supernatural
order, but with this difference, that the one desire is good, the
other evil.

8. We cannot \emph{covet} unseen good. Again, we
do not deprive another of it.

9. Hence we \emph{may rejoice} selfishly in Mary's
triumph.

10. We have a friend in court. She is the great
work of God's love.

11. Foolish objection, as if [we asserted] she
were more loving than God—a ring, \emph{e.g}. a pledge of favour to
a person, any favours will be granted.

12. Conclusion.

\xsection{September 3 (Thirteenth Pentecost) [Disease the Type of Sin]}

1. INTROD.—About
the ten lepers in the Gospel.

2. Description of leprosy as a disease. What it
was.

3. It made the person (1)
deformed—(describe)—swollen and disgusting; (2) it was lasting,
not like a fever; (3) incurable.

4. Lepers were driven out of society, they were
so loathsome; and they became like beasts. Travellers describe them
now as outside the cities in troops.

5. Now all this is \emph{sin}. Go through the
particulars, as the angels see it. Describe our souls.

6. Since we are one and all sinners, we do not
understand it. But the angels must revolt from us, but for their love.
We are an exception to the intellectual creation—except the devils.

7. Parallels: (1) a person with a bad temper; (2)
a vulgar person—we shrink from them.

8. Yet our Saviour loved us, in spite of all
this.

9. Enlarge on this. Take the cases of saints: (1)
tending the leper; (2) sucking sores; (3) Father Claver with the
Blacks; yet all this is nothing to Christ['s charity to us].

10. Here, to say nothing else, [is] difference
from our Lady. She had never seen heaven.

But He came [from heaven] among us, and now gives
Himself to us in the Holy Eucharist. You know how we shrink from dirt,
etc.

11. Thus we have at once two thoughts—humility
and thankfulness. How can we be proud of anything we are? How can we
not love Christ?

\xsection{December 25 [Christmas Joy]}

1. On the special \emph{beauty} of the narrations
of the Gospel, especially as regards our Lord's birth, and of these
Luke ii. So much so, that unbelievers have called them myths.

2. Luke ii. Describe the scene. It sends us back
to Paradise and to Adam and Eve, and to the Canticles.

3. We might fancy [there had been] no fall. [We
see] Christ, as if He did not come to die, and His immaculate Mother;
the angels; the animals, as in Paradise, obeying man.

4. We all seem caught and transformed in its
beauty—'from glory to glory'—as St. Joseph.

5. But, many Christmases as there have been, this
has something peculiar. A crown given to Mary. The Feast of the
Conception ever precedes Christmas, but this year something has been
done.

6. This year, as you know, the Pope, in the midst
of the bishops of the world, has defined the Immaculate Conception,
viz. that Mary had nothing to do with sin.

7. We were sure that it was so. We could not
believe it was not. We could not believe it had not been revealed. We
thought it had, but the Church did not say it was, etc.

8. Not out of place here. As we sing to Mary when
the Blessed Sacrament is exposed, so now on today.

9. And we of The Oratory have a special interest
in it. For our Church is raised under the invocation of Mary
Immaculate; and, as queens give largesses on their great days, so now
that this crown is put on her head, she has, we think, shown us
especial favour.

10. You recollect, some of you, three years ago,
our trials: the world flourishing (Achilli matter) my going to
Ireland; Lady Olivia Acheson's illness and death; and the illness of
three intimately connected with us. All this weighed us down. \emph{The
Christmas midnight Mass three years ago}.

11. The contrast now: benefactions to our house
and Church.

12. Our state in the University through your
prayers.

13. We may expect trouble again—joy as if
sorrowing, sorrowing as rejoicing. But God is all-sufficient.

\xsection{August 19, 1855 Our Lady the Fulfilling of the Revealed Doctrine of Prayer}

(\emph{Vide
above}, p. [21], \emph{August} 11, 1850)

1. INTROD.—In
this week we especially consider our Lady as rising to her doctrinal
position in the Church. Her first feast and this. The Immaculate
Conception and the Assumption, both doctrines.

2. She is the great advocate of the Church. By
which is not meant Atonement, of course. We know perfectly that she
was saved by her Son. But she is His greatest work, and He has exalted
her to this special office.

3. Hence from the first, \emph{advocata nostra}.
St. Irenaeus, and pictures at Rome in St. Agnese, etc.

Now to understand this, we must throw ourselves
back into the world as it is by nature. Everything goes by \emph{law}.
This order is the most beautiful proof of God, but it is turned
against Him, as if it could support itself.

Hence Revelation is an interruption and
contravention—all of it miraculous.

4. Now here we have a most wonderful doctrine of
Revelation brought before us in its fulness, viz. the efficacy of
prayer.

5. Nature uniform. \emph{How} has prayer its
power? Worship [we understand to be] right, and adoration and
thanksgiving; but how petitioning and supplication?

6. This then is the marvel, and the comfort which
Revelation gives us, viz. that God has broken through His own
laws—nay, does continually.

7. This so much that prayer is called \emph{omnipotent}.

8. Even Protestants grant all this. (Quote \emph{Thomas
Scott}.)

9. Now our Lady has the gift in fulness; not
different from us except in degree and perfection. This is her feast.

10. Hence it is that the more we can go to her in
simplicity, the more we shall get.

\xsection{August 26 (Thirteenth Pentecost) Thankfulness and Thanksgiving}

1. The gospel of the day. Were not ten cleansed?
etc.

2. Does not this event seem strange? Yet how
thankless we are. We have all to condemn ourselves. There is nothing
in which our guilt comes more home to us.

3. How we pray beforehand; how we petition again
and again. Do we return thanks \emph{even once}?

4. I think this feeling comes upon men, that they
are not equal [to the task]; that \emph{words} will not do; and so
they do nothing from being overpowered. And this grows into a habit;
and thus, when we gain our object, we suddenly leave off our prayers
and coldly accept the favour. But still we may show our gratitude by \emph{deeds}
and by recurrent remembrance. We might remember the \emph{day}; we
might perpetuate our gratitude.

5. 'Where are the nine?' and he, the tenth, was a
Samaritan! (Other instances—woman at the well; good Samaritan.) It
is a paradox which is fulfilled, that the less a man has the more he
does. The centurion and the Syrophoenician.

6. When we have a number of blessings, we take
them as our due. We do not consider that they are so many accumulated
mercies. Thus the Jews especially, etc.

7. Now let us think what we can claim of God, and
what He has done. Preservation perhaps implied \emph{de congruo} in
creation. But how much He has done for us! for each one in his own
way—yet so much to every one, that every one is specially favoured—favoured
as no one else.

8. Survey your life, and you will find it a mass
of mercies.

9. Hence the saints, three especially—Jacob,
David, St. Paul—are instances [of thanksgiving].

10. Close connection with hope and love. This
gratitude is the greatest support of hope, and hence those saints who
have been patterns of gratitude were patterns of hope.

11. On setting up memorials.

12. Gratitude is even a kind of love, and leads
to love. Against \emph{hard} thoughts of God. Not [being] too proud to
admit to \emph{ourselves}, 'At least He is good to ME.'

\xsection{September 2 (Fourteenth Pentecost) Service of God Contrasted with Service of Satan}

\footnote{'Not preached.'}

1. No man can serve two masters.

2. This is true, even \emph{because} they are
two, but much more if [they are] opposed. In all things we must throw \emph{our
heart} into our work. It is the only way in which any work is done
well. This is how men succeed in any line.

3. Yet, though this is certain, men forget it as
to religion. They think to serve God without taking His service \emph{exclusively}.

4. What is meant by \emph{exclusive} service? Is
it going out of the world? No. There are persons so called—but it is
not that.

5. But [it is] subordinating all things to God's
service. Whether we eat or drink, etc.

Parallel of worldly matters. A worldly man
carries his aim into all things. He is thinking of his business
wherever he is.

6. So in religion. And this is what is meant by
loving \emph{God above all} things. And this is why such love alone
keeps us in God's favour.

7. To be religious, then, is not merely to have a
\emph{respect} for religion, to do some of its \emph{duties}, to
defend it, to profess it, but

8. It is to live in God's presence; to know the
whole economy of redemption.

9. Hence the necessity of meditation.

10. Warning, because the world is likely to \emph{crush
out} our religion.

\xsection{September 9 (Fifteenth Pentecost) Life of the Soul}

1. INTROD.—Gospel
[Luke vii. 11-16—raising to life of the son of the widow of Naim].

Our Lord's miracles are especially typical—(1)
leprosy—heresy; (2) demoniac—cleansing the soul from the evil
spirit; (3) blind—John ix.; (4) loaves—so this.

2. It brings before us the natural state of
man—state of the whole world [typified in it].

3. What is meant is, \emph{not} that man may not
have natural powers, but [being lacking in] spiritual, that left to
himself, he will know nothing of the unseen world. In one sense, then,
the world is alive, in another dead.

4. It is in this sense that the soul is dead. Now
if dead, observe the greatness of that death. (1) Dead men are without
sense or feeling: so the soul as to heavenly things, motives, objects,
etc. (2) [A dead body provokes] fear and odiousness: so the [dead]
soul in the sight of angels and Almighty God.

5. (3) As to the outward form [of the dead] it is
the same [as the living], and this suggests much. (i)
Imitation—Christianity in the world. (ii) Simulation, because they \emph{know}
more than they \emph{do}, and pretend from shame. (iii) [Souls that
are dead may still have] actual grace, [and] habits formed under it.

6. Yet in God's sight [they are] dead. Now
consider Eph. ii. [see vv. 4 and 5]. \footnote{'But God … even when we
were dead in sin, hath quickened us together in Christ.'}

7. Now reflect on all this—the terrible state
of the world—in detail; here, there and everywhere. Yet, as dead men
do not know they are dead, neither does the world.

8. On Christ, the sole source of life, from today's
gospel—Gal. ii.

9. On the love which life implies.

\xsection{September 16 (Sixteenth Pentecost) Septem Dolorum—Election}

1. INTROD.—Nothing
is, of course, so awful as the question of election, about which so
much is said in Scripture. It is not to be supposed that I am going
into any depths here.

2. The doctrine, as I shall take it, is this, and
most practical; and I will first illustrate it.

3. Take the case of some large and new
institution in a nation, which requires a great many new hands, \emph{e.g}.
a new department of revenue, a new commission, some speculation
abroad, the post office, railroads, the war.

4. Such an institution, especially if a
speculation or expedition (1) promises great rewards to those who take
part in it; (2) it is not for every one to get, but he must make
interest; (3) no one will get part in, or receive the rewards of, if
he does not join it.

5. Enlarge. As a question of justice. Suppose a
man who went on with his own trade, etc., complaining that he had no
part of the receipts of a speculation in which he took no part, etc.

6. Apply. Draw out the state of this world—its
trades, occupations, aims; its science, literature, politics, etc.
People may acquit themselves well, and get the reward of their
occupation, which is the reward of \emph{this world}, \emph{e.g}. such
as wealth, fame, etc., etc.

7. But a new system comes in. Almighty God
proclaims a different reward, viz. eternal life to those who take part
in His objects, etc. You see it is quite distinct from Nature.

8. Enlarge on the \emph{interest made} to get a
place—no claim because [a] good father, a good subject, etc., etc.

9. Here, then, we have the election. If we want
to take part in it, we must join it.

10. The cross of Christ puts a different
complexion on the \emph{whole} of life. If a man takes up any new
course, his old ways are flat in comparison.

11. \emph{Septem dolorum} in connection—we must
take part with her.

\xsection{September 23 (Seventeenth Pentecost) Love of God}

1. INTROD.—The
gospel is the second which we have lately had on the precept of the
love of God.

2. Nature tells us we should love God. Nay, a
natural inclination and leaning to the love of God.

3. Still, it never will lead us to love. It fails
for want of strength, and the feeling comes to nothing and dwindles,
as a tree of the south planted in the north. Grace essential.

4. On pure love of God—illustrate—single,
real, for Himself, \emph{e.g}. we are to love men \emph{propter Deum},
thus \emph{not propter} [\emph{seipsos}], etc., which is Nature. If,
then, we love God by association [sic], or merely for His benefits,
etc., it is not enough.

Love delights in the name of God, likes to hear
of Him, likes to think of Him, likes to act for Him, [is] zealous for
His honour and a champion for His cause.

5. But this is not all. It is not merely looking
at what does not notice us, as the Pantheists say. It is a friendship.
Three things are necessary for friendship: (1) mutual love; (2) mutual
consciousness and sympathy; (3) mutual intimacy—intercourse.
Companions, walking with God, Luke xxiv \footnote{The journey to Emmaus.}? Apply to confidence in God's loving us.

6. But this is not all—\emph{dilectio}: choice.
And no common choice, but \emph{above} all things.

7. Thus it is pure, amicable, mutual and
sovereign.

8. Now to see what it is, we may see what it is
not; and parallel it to worldly principles. Take the course of
men.

9. (1) They begin with self-indulgence and
self-gratification. Here is something which is not love, yet acts as
love does.

10. (2) Perhaps ambition, martial spirit. This
possesses them—\emph{this} not love.

11. (3) Love of home: [a man is] a good father, a
good son, [devotes himself to such duty with] concentration of mind \footnote{\emph{I.e}. without
reference to God: not for His greater honour and glory.}—\emph{this} not love.

12. (4) He gets wealthy, and is tempted to make
wealth his \emph{enough}—\emph{this} not love.

13. (5) Love of consistency, character; self his
centre—\emph{this} not love.

14. (6) Ease and comfort in old age—\emph{this}
not love.

15. \emph{How} are we to gain love? By reading of
our Lord in the Gospels.

\xsection{December 25 Christmas Day}

1. INTROD.—Today
a change in the history of mankind. Many important eras and
seasons—this the most important. And it is described in various
terms in the services [of the feast]. \emph{Melliflui facti sunt coeli},
etc.

2. All things created good; but man is fallen.

3. Man fell, and the angels fell before him; but
the case of the two is different. The angels were pure spirits, and
have but one nature; man has two natures. Angels are spirit, but
man is made up of soul and body. An angel is good or bad: if good,
[there is] nothing to resist the good; if bad, nothing to resist the
bad. If they fell, they fell once for all. If man fell, there is a
contest between the flesh and spirit, reason and passion.

4. Man not simple [in his nature]; [he is made up
of] two principles. Illustration of these two formidable principles in
man, by comparison to man and beast. In each heart of man there is
what may be called man's true nature and beast's nature. Power of wild
animals. The wild principle of man has carried him away.

5. It first showed itself in the fall
itself—passion—then Cain and Abel. Thence it swept over the world.
Wars, murder, injustice, sensuality, crimes of all sorts.

6. Thus things [went] continually from bad to
worse, and did Almighty God suffer it, there is no depth to which man
would not descend. In 1600 years [he had become] so bad that [God
sent] the deluge.

7. The earth restored; but how vainly! Man soon
got almost as bad as before. He cast off God; he set up idols; he
tyrannised over others. He went on to found states, and he impressed
sin upon them—idolatry mixed up with politics, with all the usages
of society—marriages, business contracts, births, deaths and
burials, recreations and institutions. And the raising temples and
stamping it on the great cities, and then misusing and devoting the
creature to idolatry, Rom. viii. [20] \footnote{'For the creature was made
subject to vanity.'}. And thus sin got so established as to exert a tyranny
over each individual. Those who would have been better [were victims
of] bad education, ridicule, persecution.

8. And then the struggles of the unregenerate and
remorseful.

9. Who can estimate the entire establishment of
evil? In vain judgments, God's pleadings, etc.

10. Now this being so, was it not plain that if
there was to be a change, God alone could do it? If a Redeemer, He
must be God. So this is the great event now [beginning].

11. Yet not at once a bloody combat, but a little
child.

\xsection{May 25, 1856 (Sunday within the Octave of Corpus Christi) [Devotion to the Holy Eucharist]}

1. INTROD.—There
is no feast, no season in the whole year which is so intimately
connected with our religious life, or shows more wonderfully what
Christianity is, as that which we are now celebrating. There is a
point of view in which this doctrine [of the Body and Blood of Christ]
is nearer to our religious life than any other. And now I will explain
what I mean.

The Holy Trinity unseen. The Nativity, Easter,
etc., past. But this is the record of a present miracle, a present
dispensation of God towards us.

2. [In devotion there is always] one difficulty
to counteract. Our Lord came 1800 years ago. How shall we feel
reverence of what took place 1800 years ago? We are touched [with]
pity, gratitude, love, by what we see. None of us have seen or heard
even those, who saw those who saw those who saw Him.

How shall we learn to live under the eye of God?
Now we know how difficult it is to keep up the memory of things. Then
again, books, how little can they do for us! It is a great thing to be
moved [even] once in a way by a book, but we cannot count upon their
moving us habitually. Accordingly an historical religion, as it is
called, is a very poor and inefficacious [\emph{a word illegible}]. We
see it in the case of Protestants. Their religion is historical, in
consequence they speak of Christ as a mere historical personage—the
titles they give Him, etc., etc.—there is a want of reality, etc.

This is one difficulty in the way of practical
devotion.

3. A second difficulty. The world is in
wickedness. Satan is god of the world; unbelief rules. Now this
opposition to us has a tendency to weigh us down, to dispirit us, to
dull our apprehensions, etc.

These are two extreme difficulties in the way of
religion.

Now observe,

4. How almighty love and wisdom has met this. He
has met this by living among us with a continual presence. He is not
past, He is present now. And though He is not seen, He is here. The
same God who walked the water, who did miracles, etc., is in the
Tabernacle. We come before Him, we speak to Him just as He was spoken
to 1800 years ago, etc.

5. Nay, further, He [does] not [merely] present
Himself before us as the object of worship, but God actually gives
Himself to us to be received into our breasts. Wonderful communion.
Texts—Father, Son and Holy Ghost.

6. This [is] how He counteracts time and the
world. It [the Blessed Sacrament] is not past, it is not away. It is
this that makes devotion in lives. It is the life of our religion. We
are brought into the unseen world.

7. These thoughts are fitly entertained, and
themselves increased at this season, when St. Philip's day comes.
Quote passages in his life to show his delight in the Blessed
Sacrament. He has died on this day. We cannot have a better
preparation for his, than this, feast.

8. Let us rejoice in Jesus, Mary, Philip.

\xsection{July 27 (Eleventh Pentecost) [On the Healing of the Deaf and Dumb Man]}

1. INTROD.—We
read these words in today's gospel, 'They bring unto Him,' etc. \footnote{'They bring unto him one
deaf and dumb; and they besought him that he would lay his hand upon
him.'—Mark vii. 32.}

2. The man is cured, and two things go to his
cure—Christ's word and act, and His disciples bring him to
Him. Christ does not heal without His disciples, and they cannot heal
except as bringing to Him.

3. So it is now—the great ordained
system—Christ the Author of Grace, and His friends whom He brings
round Him, and makes His family, the step towards obtaining grace \emph{by
prayer}.

4. Christ can do all things. He created, He
redeemed without any one else; but He saved [saves?] through the
co-operation of others—by the saints above and the Church below.

5. Christ can do all things—He \emph{gives grace}
too, and it is only by His ordained system—merit a promise—a
contract, etc., etc.

6. Christ can do all things, and He does not
confine Himself to [co-operation of] others, so far as this, that all
over the earth, external to His Church, He hears those who call on
Him. He has many ways. Every one has a guardian angel. Case of Hagar.

7. \emph{But} He does this to bring them on into
His Church, that \emph{they} too may become His friends.

8. And it must be recollected that the Holy
Church Universal is praying everywhere [for them]. Mass [continually
offered].

9. Abraham and Moses. God reveals that His
friends \emph{may} pray, 'I say not that I will ask the Father,' etc.
\footnote{'In that day you shall ask
in My Name: and I say not to you, that I will ask the Father for you:
For the Father himself loveth you.'—John xvi. 26-27.}

Therefore it is that we call our Lady our
advocate, and the saints intercessors; for our Lord has made over this lower office to them, and stands in the higher, of the Giver
of grace.

10. Thus the salvation of the world is in our
hands, [\emph{e.g}. of]

11. England—Birmingham.

12. Therefore let us \emph{pray}.

\xsection{August 29 (Fifteenth Pentecost) [The Raising to Life of the Son of the Widow of Naim—Luke vii. 11-16]}

1. INTROD.—The
Holy Fathers are accustomed to derive a spiritual lesson from the
miracle recorded in the gospel of this day. It was a miracle exercised
on one, but it was a sort of specimen of what takes place by God's
love so often. It was done once, but it images what occurs
continually.

2. This was a young man borne out to his burial,
and his mother is weeping over him.

The mother is the Church, who has born him in
baptism, when he was born again and became her child.

He has fallen away, and is dead in sin. He is
here carried on his way, like Dives, to be buried in hell.

3. How awfully he is carried forth! Slowly, but
sure, as the course of a funeral.

Describe his odiousness—death so fearful, every
one shrinks from the sight. Children in the streets turn away. Those
only bear it who love the corpse, or have duties towards it. So \emph{with
the soul}. How angels must shrink from the dead soul!—the
guardian angel bears it. How horrible it looks even [if in] venial
sin, much more in mortal!

The mother bears it—the Church does not
excommunicate.

4. Its bearers are four: (1) pride, (2)
sensuality, (3) unbelief, (4) ignorance. We see these from Adam's
original sin, and they are in every sinner, though perhaps in a
different order in different persons. There are those who go on,
through God's mercy, in the right way. But I am speaking of cases of
sin.

5. Now I believe generally \emph{pride} comes
first—obstinacy of children; disobedience; quarrelling; refusing to
say prayers; avoiding holy places, etc. Thus the soul being left open
to the evil one, he proceeds to assault it with sensuality.

6. Sensuality. A person does not know when he is
proud, but this [sensuality] need not be described, for every one who
yields to it knows what it is. God has set a mark upon it, the mark of
sting of conscience, because it is so pleasant; whereas pride is
unpleasant to the person who exercises it.

7. Thirdly, unbelief. Pride and sensuality give
birth to unbelief. A man begins to doubt and disbelieve.

8. Fourth, ignorance. At last he does not know
right from wrong.

9. And thus a soul is led out to be buried, to be
buried \emph{in hell}. And how many reach that eternal tomb!

10. Wonderful electing grace of God, choosing
one and not another, coming without merit—the Church cannot do
it.

11. We all have received it [this electing grace]
without merit. Let us prize it when we have it.

\xsection{September 7 (Seventeenth Pentecost) [Love of Our Neighbour]}

1. INTROD.—Sometimes
it is said that there is one, sometimes two, great commandments.
Charity is the great commandment. Though properly the love of God, it
involves love of neighbour.

2. We have not seen God. How are we to ascertain
that we love Him? Feelings are deceptive. Thus, as by a test, by
loving others, by love of man. And so St. John says, 1 John iv. [12].

3. First, we should love man merely as the work
of God. If we love God, we shall love all His works. Undevout men walk
about, and look round, and they never associate what they see with
God; but everything is the work of God. And though we should not be
superstitious, we should destroy nothing without a reason. Cruelty to
animals [is] as if we did not love God, their Maker; nay, wanton
destruction of plants.

4. Thus, even if mankind were of a different
species, as fellow-beings [they would have] a relationship [to us].

5. But they are \emph{of our blood}, Acts xiv.
Adam [our one father].

6. And all involved in Adam's sin—the sympathy
of sin, as all in sin, in misery and transgression, and in
danger of ruin.

7. Hence Gen. xviii. \footnote{Abraham interceding for
Sodom and Gomorrah.}, Ps. cxviii. 139 \footnote{'My zeal hath made me pine
away: because my enemies forgot thy words.'}.
St. Paul, Rom. ix. [3] \footnote{'I wished myself to be
an anathema from Christ for my brethren.'},
Acts xvii. [26] \footnote{'He hath made of one all
mankind.'}. Our
Lord weeping over Jerusalem. Missionaries to heathen countries, as St.
Augustine who came here.

8. Of course, zeal for God also [moved these to
heroism], but the sight of souls dying [more directly].

9. Still more if Christians, for then we are
brought near to God. He who dwelt in solitary light once, now has
round Him a circle of holy beings, so that we cannot love Him without
loving them. Hence the glory paid to saints, as His garment.

10. Besides, we love the divine attributes and
character in the saints: 'He who loveth God, loveth His brother also.'

11. This the condemnation of those who oppose the
Church.

12. On the other hand—love of saints—love of
our Lady as God's mother, [a] sign of predestination. She the great
work [\emph{i.e}. the greatest of God's works] and the glory of our
race. Let us at this season beg her to make us full of that love of
herself, and of all those who have God's grace, and of all whom God
has made.

\xsection{September 28 (Twentieth Pentecost) Seven Dolours}

1. INTROD.—The
most soothing of all the feasts of Mary. What a contrast the first
portion of the Blessed Virgin's history is to the latter! We sinners
have no sympathy with the first part of her life. She had nothing but
joy, increasing up to that day which heralded its reverse. It was at
the height of her earthly joy that the reverse began—her seven
dolours.

We say seven, but that is a perfect number only;
her woes were continuous.

2. Go through her life—Presentation,
Annunciation, Visitation, the Nativity, Shepherds, Magi,
Purification—and then we hear of a sword. And the flight into Egypt;
avoiding Herod; loss of our Lord in the Temple; death of St. Joseph;
[our Lord] leaving her to preach; [His] crucifixion and [her]
bereavement.

3. Parallels of Moses, Deut. xxviii. Solomon at
dedication; and Transfiguration with prophecy of suffering and so
riding in triumphantly into Jerusalem, and 'Crucify Him!'

4. Yet in truth it would seem that she knew it
all from the first, though we don't know when it was told her. This is
something which equalises the two [portions of her life]—the
knowledge beforehand [of her woes]. And it is this which gives a
character to her whole life. All through that first calm time, she
knew it was but the stillness before the storm, and she could
not enjoy what was so joyful. All along there was the vision of One
lifted on the cross, and the sword pierces her heart.

5. Describe the cross—and she by it! This is
the key of her life on earth.

6. Ignorance is bliss—animals, men [even] do
not know what is to happen to them.

7. And this was the peculiarity of her life.
[Bodily] pain, trouble, etc., come at fixed times and go, but it is
otherwise with mental: foresight and memory make them continuous. This
is the sword in Mary's heart, the peculiarity of it being that it is
mental.

8. And again she did nothing—only
suffered—did nothing indeed, except in internal acts. A champion
acts, and a martyr acts. Hers was mere suffering.

9. And especially the sight of the suffering of
another which we cannot help. A mother seeing her child suffer. Case
of Hagar.

10. Many a wife, many a mother stands by and
says, 'O that I could take a part!' Martyrs declaring themselves, and
suffering because others were. [Yet Mary suffered] not like Hagar \footnote{'I will not see the boy
die.'—Gen. xxi. 16.}, but like the brave mother in the Maccabees.

11. This is the \emph{compassio} of Mary.

12. Suitable to us, most soothing of feasts: for
mental pain more widely spread than bodily, in this age especially.
Care, anxiety from difficulty of livelihood,—[those terrors of an]
intellectual age—madness and heartache; remorse at sin. In all Mary
is our sympathy and comfort, etc., etc.

\xsection{October 12 (Twenty-second Pentecost) The Maternity of Mary}

1. INTROD.—There
is no feast of our Lady which comprehends so much as this. It is a
sort of central feast. It connects all that is taught about her in
one.

2. A number of feasts look \emph{towards}
it—the [Immaculate] Conception, Birth, Purification, Visitation,
Nativity. Her becoming a mother is the scope in which they end. For
this all her graces, etc., \emph{because} she was to be the Mother of
God, and a temple set apart for Him.

3. What is meant by being the Mother of God?
Mother of the Person of the Son—God's blood—God's flesh, etc., and
so God's Mother.

4. So high an office \emph{required} a due
preparation, as St. John the Baptist or the apostles, but much more.

5. And the reward and power [were in] proportion.
\emph{Monstra te esse Matrem}.

6. And thus we are brought to that other set of
doctrines included in the Maternity. For she is our mother as well as
God's. And thus this feast becomes not only one of the most wonderful,
but of the most soothing.

7. Two natures in Christ—so she was mother of
Him who was God as well as man. 'Behold I and my children,' etc., Heb.
ii. 13.

8. Hence, 'Behold thy Son—[Behold] thy Mother,'
John xx.

9. Here is its connection with the seven dolours.
Her first birth without pain; her birth of us with pain.

10. It became her who was to be a mother to us,
to be so far like other mothers as to have pain.

11. On the constant, unwearied affection of a
mother's love; (on many not having experienced it) but nothing
extinguishes it. The father gives up the son, brothers despair of him,
but she remains faithful to the end, hopes against hope, does not mind
slights, ingratitudes, etc.

12. Here you have the maternity of Mary. You
cannot weary her, she never reproaches, etc. Therefore do we pray her
to help us in the \emph{hour of death}, for she will not leave us.

13. Especially as men get old and lose their
earthly relations and those who knew them when young.

14. Who are our constant friends but our guardian
angel, who has been with us since our youth, and Mary, who will be
with us to the end?

\xsection{October 19 (Twenty-third Pentecost) Purity of Mary}

1. INTROD.—If
there is one thing more than another which marks Christianity, it is
the honour given to virginity. We, who have ever heard the doctrine,
cannot fancy how it must come upon the heathen at the beginning by the
contrast.

2. And indeed the Holy Fathers appeal to it from
the first as a great miracle. When we consider the state of the
heathen, etc. So wonderful that numbers of persons should be
found who were willing to debar themselves even of the marriage state,
living in chastity.

3. Moses, Aaron, the Priests, the Prophets.

4. Nay, the Jews—hardness of the heart,
divorce, polygamy.

5. Nay, celibacy was not held in honour even from
a religious reason. They each wished to be mother of the Messias.

6. Hence the force of the prophecy, 'A virgin
shall conceive.' And when the time came, St. John the Baptist went
before Him a virgin. He Himself, the Messias, pre-eminently such; and
His Virgin Mother, and His favourite disciple, the other John, St.
Paul, and all of them, either gave up their wives or had none.

7. Hence we see the force of the doctrine of the
Immaculate Conception. A new thing was coming upon the earth. It was
fitting that it should begin with a new beginning, as Adam's at the
first—of grace before sin.

8. A new thing, though Joshua \footnote{See Note
10, p. 338.}, Elias, Eliseus.

9. The heathen philosophers, stern, proud, etc.,
whereas, St. Gregory insists, humility must be with chastity, and our
Lady a special instance of humility.

10. But further, the celibacy of false religions
has been negative—the \emph{absence of love}.

11. This indeed is what is imputed to
us—blighted affections. The peculiarity of Christian celibacy is
that it is from love to God—'and followed Thee.' St. Jerome in
Breviary.

12. The more we love God, the more we are drawn
off from \emph{earth}.

13. The Blessed Virgin's Purity arose from the
excess of her love.

\xsection{December 25 (Christmas Day) Omnipotence in Subjection}

1. INTROD.—(1)
They say that love does not reason, \emph{i.e}. so intent on [its]
object that it does not regard itself or its own feelings; and so of
adoration and praise. Thus Christ was born in silence; not a word from
our Lady or St. Joseph, or the shepherds or the magi. The angels
indeed, but very briefly.

And thus, I suppose, we all feel little disposed
to speak today, as interfering with enjoyment. (2) A second reason is,
because love has so \emph{many} thoughts which reason cannot draw out
fully and do justice to. Or, if we preach, we do it for the honour of
the day.

2. If we speak, the first natural thought [is]
that every feast, as it comes, is the best. Nothing like Christmas.
But really we have reason to say so. Easter is the higher, but the
sufferings of Christ, which we contemplated, are a shock which sudden
reversal to good does not remove. And our own sin and penance [have
preceded it] \footnote{See Note
11, p. 338.}. But
Christmas [is] as if we had never sinned. Some divines think that
Christ would have come into the world though man had not sinned. Thus
this feast has not necessarily the idea of sin in it, though in fact
Christ came for our sin. Seeing the end from the beginning, as
Moses seeing the [promised] land, through a valley of conflicts.

3. But if, for the honour of the day, I must take
one thought or lesson to put before you, it shall be the adorable
marvelousness of what may be called the humiliation of the Divine
Being, as at this time of year. (1) Omnipotent—what He can
do—create and destroy worlds—He can do what He will, therefore it
would seem that God \emph{could not} humble Himself. (2) Idea that God
is so high that He cannot listen to man. (3) For consider who He is.
[He has] no [obligation of] justice towards us, as none on our part
towards beasts. (4) If He only \emph{attended} to us (texts to the
contrary, Isa. lvii.—'Inhabiteth eternity'; 'Shall God dwell on
earth?' \footnote{'Is it then to be
thought that God should indeed dwell upon earth?'—3 Kings viii. 27.} If 'Emmanuel'
only meant this). (5) But He has taken our nature.

4. Now observe particulars. (1) Nine months in
His mother's womb; (2) swathing bands; (3) infant—carried about,
etc.: Simeon—Egypt; (4) subject to them, when He even displayed what
He really was; (5) worked at trade; (6) laid hold of, as beside
Himself; (7) His Passion; (8) His Crucifixion; (9) now in Tabernacle;
(10) in our breasts.

5. Example to us. We are most of us in
subjection; why not sanctify it?

6. This St. Philip [did] by sacraments, humility,
detachment, purity, and joy or peace, and cheerfulness.

\xsection{December 28 (Sunday—Holy Innocents) Suffering}

1. INTROD.—The
three feasts about Christmas, as if to tame down its joy, bring before
us suffering.

2. And so the events about our Lord's Nativity:
(1) Circumcision, (2) Purification, '[a] sword,' etc., (3)
Epiphany—massacre of infants.

3. Remarkable that the children should suffer,
because it is the age of innocence.

4. It suggests to us the doctrine of original
sin—that man has fallen. Pain would not be, with man upright. Here
then we have a proof that man is under God's displeasure—\emph{pain}
not death.

5. Sufferings of children: (1) from illness, (2)
from cruel parents, etc. \emph{Nothing worse than} to see a helpless
child in great pain.

6. But, however, the Holy Innocents were
otherwise circumstanced. This martyrdom was an [entrance] into the
Church. Their sufferings meritorious.

7. St. Rose and other holy women, who inflicted
on themselves penances extraordinary.

8. The Church like a joint-stock (all who share
it must be cleansed).

9. Let us rejoice in this feast then;
particularly it is for mothers whose children suffer. \emph{All the
sufferings of baptized children merit}, and all innocent profit in
suffering.

10. Let us thank Him who turned sufferings of
children to account.

11. The merits of saints ever growing, of
martyrs, and souls going from purgatory to heaven; \emph{of children
suffering and dying in infancy}.

\xsection{January 4, 1857 (Octave of Holy Innocents) Passage of Time}

1. INTROD.—All
times, all days are the beginning of a year, but especially when the
date changes.

2. Time, as present, is momentary, as future, is
unknown, as past, is irrevocable.

3. As present, momentary. No standing still.
While we speak, it goes. We are all older when we leave this church
than when we enter it. Whether it be joy or sorrow, it goes. We look
forward to a great day; we keep a great festival. It comes once in a
year. [As] grains in an hour-glass, it is gone ere it is well come.

4. And on what road is this swift time driving?
On a road of darkness. We are every moment entering and driving along
an unknown future—on a steam-engine on a railroad in the dark.
Accidents may happen any moment. Unseen dangers waiting for us. Balaam
and the angel. Hence Jacob asking God's blessing on his journey. St.
Raphael. We are not merely journeying, we are rushing forward, and to
what?

5. To judgment. On the importance of time.

6. Thirdly, the past is irrevocable. What would
we give to wipe out much!

7. On the necessity of taking good heed how we
spend time. Counsel of perfection never to misuse time. Vow by some
saints.

8. \emph{Desideria efficacia et sterilia}.

9. Let us begin the new year well.

\xsection{April 5 (Palm Sunday) Falling Away}

1. INTROD.—Too
awful a subject commonly, as leading [men] to despond; yet useful
sometimes, and natural at this season.

2. Now first let us lay down about nature and
grace—[that] nature can do many things, but cannot bring to heaven.
Grace is like a new nature, and joins us to the heavenly family; and
they are saved who die with this grace; those lost who are without it.

3. This answers the question: Will good departed
from avail? As some Protestants say, 'Look how a man lives, not how he
dies'—(explain).

4. Proof, Ezech. xviii. [24]. And rightly, for
the sovereign Lord of heaven can prescribe His terms.

5. Now this chapter leads to a further thought,
viz. that much as is said to encourage repentance, as much perhaps is
said to warn against falling, as if the prospect, or chance, or issue
on the whole were equal.

6. \emph{E.g}. our Lord, 'I came not to call.'
But on the other hand, recollect the number of passages such as 'Two
shall be in the field'; 'Ten virgins'; 'He that persevereth,' etc.;
'Many that are first,' etc.

7. So St. Paul, preacher of repentance: but Heb.
vi. [4-6] \footnote{'For it is impossible
for those who were once illuminated, … and are fallen away, to be
renewed again to penance.'}.

8. So holy Simeon, 'This child [for the fall,
and for the resurrection of many in Israel,' Luke ii. 34].

9. This text of holy Simeon especially fulfilled
at Passion, when two special examples.

10. Multitude on Palm Sunday, \emph{vide} their
being in grace [implied] in the prayers [second and last] in the
Blessing of Palms. Cp. our Lord's weeping—disappointment of the
foolish virgins.

11. Judas. Our Lord chose him when he was in
grace—trace about him—'the ten indignant,' Mark x. 32, etc. \footnote{At the request of the
sons of Zebedee.}

12. Some fall away at one age, some at another.
Go through this.

13. On natural habits produced by supernatural
acts deceiving the old.

14. Our Lady. Prayer—pray \emph{lest} we fall,
if we fall, and for others.

Continue

Back

\xsection{May 31 Whitsunday}

1. INTROD.—If
we get to heaven, the wonderful peace will be our great blessedness,
the blessedness of the end having come. And not only the end, but the
consummation, for not only will labour be over, but our reward will
have come.

2. Today is the nearest approach we have to such
a consummation. Describe 'in one place,' etc. \footnote{'And when the days of Pentecost were
accomplished, they were altogether in one place.'—Acts ii. 1.} Today we celebrate the day when Almighty God exhausted His
gifts upon us. They were long promised, but on this day they
were all poured out. He emptied the fulness of His mercy on us this
day.

3. Long have we been following the course, from
Christmas to Easter, etc.; but now the end has come. We are called
upon to be thankful for, and enjoy what we have been anticipating.

4. So was it at the first Pentecost, and more
strikingly still. For ages had the Church been expecting this day,
etc.

5. And again so different from their
expectations. This happens to \emph{us}. We pray, and we do not get
our answer. Yet we do in a higher way. So the apostles. (1) They did
not know their Lord was to suffer; (2) that He was to go. Yet He said,
'If I go not away,' etc. \footnote{'If I go not, the
Paraclete will not come to you; but if I go, I will send him to
you.'—John xvi. 7.}
Yet the fulness, when it came, did not disappoint them.

6. All our infirmities, sins, etc., are reversed
in the coming of the Holy Ghost. Sin is gone, fear is gone, etc.

7. All that we have of good comes from this day.
All the sacraments from this day. If baptism gives, etc., it is from
the day of Pentecost. If confirmation, if penance, etc. If [the Holy]
Eucharist. If faith, if hope, etc. If chastity.

8. It is life for death. If we are dry, if cold,
if deified, if sickly, if wounded, etc. It is the sweet refreshing
breath. If we are to overcome the foe, etc. If we have a refrigerium
in purgatory. If at last we mount to heaven.

9. This day especially St. Philip's feast. He
preached for fallen Christendom. We too.

\xsection{October 11 During Exposition for Troubles in India}

1. When our Lord was on the Cross, He said,
'Father, forgive them.'

2. Now we know well what the obvious lesson of
that prayer is. It shows the love of the Creator in compassionating
His children when they were sinning; nay, sinning with the most awful
intensity of outrage, for they were wounding, torturing, putting to
death that nature which He had assumed.

3. But there is a further reflection perhaps, not
so obvious, to be deduced from it, and which is very much to our
purpose to consider, in reference to the great calamities which have
on this day been so solemnly brought before our consideration by our
bishop, and in consequence by us before the throne of grace.

4. He Himself, though ineffably holy in His human
nature, still had that very same nature which in them who assailed Him
was capable of such sin. If they were of His nature, on the other
hand, He was of theirs. If they had the guilt of being His brethren,
He had the shame of being theirs. We know how ashamed men commonly are
when any one connected with them does anything wrong. The bad deed of
any one of our blood is in a certain sense our own bad deed, and is an
humiliation. Now our Lord, in His own proper nature as God, is
infinitely separate from all beings whatever, but He took on Him
a created, a human, a frail nature, when He came on earth. He became a
child of Adam. He took on Him that fallen nature which He had made
perfect at the time that He created [it], but which had lost its
perfection, and which anyhow was always [in] its own essence and by
itself frail. If angelic natures have, separate from the grace of God,
imperfection, much more has man's nature. Our Lord took on Him a
nature which in any other (except His mother) but Him would be sure to
sin. He took on Himself a nature which nothing but the grace of God
could save from running into sin, from that inherent imperfection
which attaches to the creature. He [His human nature] could not sin,
but the reason why it could not was not because it was intrinsically
higher or better than the nature of any other son of fallen Adam, but
because the presence of Himself in it, of Himself who was God,
rendered it utterly removed from sin and incompatible with it. Still,
His human nature was such that, had it not been His, it might have
sinned. But it never was by itself, it never had been without Him.
From the first moment of its existence He had taken it up into
Himself; He had created it for Himself, and thus it was absolutely and
eternally secured from all sin.

5. But still He would know and understand,
infinitely more than we can, the shame of having a nature which was in
itself peccable. And therefore the sins of all His brethren weighed on
Him, and were in one sense His, because He partook their nature, had a
share in a common possession which was a very shameful possession. In
this sense, though most pure, He bore Him a body of death and
the sins of the whole world.

6. At various times He shows this feeling: when
He sighed and said 'Ephpheta'; when He wept at Lazarus's grave. These
were the signs of the burden He was bearing, who, in partaking our
nature, had \emph{in solidum} the sins of that nature on Him. And so
in the Garden, when He sweated drops of blood, it was the weight of
that fallen nature which He had assumed which made Him weary even unto
death.

7. And when He was on the cross, He had this
additional woe, that the evil of human nature now showed itself in a
new way, as rising up against Him who bore it. It was the climax of
its depravity, that it turned against Him who for its sake had
voluntarily put it on. And, while He felt its ingratitude, He felt
perhaps equally its shame, as the father of a family is ashamed of his
Sons' acts against him, and though he feel them, dare not mention
them, because they fall back upon himself.

8. And therefore, when He was lifted up upon the
cross, He would not be angry with His torturers, lest it would seem as
if it were His own act, for it was the act of that very nature in
others which He bore Himself, and, as when we have a hand or a foot in
pain, we are not angry with it, but feel a tenderness towards it, so
He felt a tenderness to that fallen nature which was showing itself so
awfully devilish in His persecutors, for it was His own.

9. My brethren, I do not know whether you see
whither I am leading you by this train of thought. We are on
this day engaged, in obedience to the call of our bishop, etc.

10. Now I suppose most of us have heard something
or other of those indescribable horrors which have been perpetrated by
the revolted Hindoos and Mahomedans in the instance of our dear
countrymen and countrywomen in India.

11. (Go through them.)

12. We are not only horrified, but angry. Dost
thou well to be angry? It is very horrible, but let us not 'think it
strange,' 1 Peter iv.

13. Now I don't doubt the right of the Ruler.
'Vengeance is mine,' 'not the sword in vain.' But Catholics, of all
people, have nothing to do with rule or responsibility in India.

14. What I am impressing on you is that these
enormities belong to our nature, and that we ought to consider that we
are of one blood with [those who did them], have one nature, and that
that nature is such as might cover it. There is not any one of us but
might in other circumstances have committed the same. (Hazael \footnote{'Their strong cities thou
wilt burn with fire, and their young men thou wilt kill with the
sword, and thou wilt dash their children, and rip up their pregnant
women. And Hazael said, But what, am I thy servant a dog, that I
should do this great thing?'—4 Kings viii. 12-13.}, 'Thou art the man,' 2 Kings xii \footnote{'And Nathan said to David,
Thou art the man.'—2 Kings xii. 7.}. I assure you, my brethren, I speak in earnest when I say that,
much as you pity the persecuted, yet should [you] pity the persecutors
more.)

15. Instances to show it in every age and
country.

16. Therefore I say to any person who indulges in
any bitter feeling about these dreadful transactions, that the
question is whether that feeling is not the same in kind, though
different in degree, with that which at this minute is making our
soldiers in India, according to the confession of their officers, \emph{demons}.
How do we make matters better by sharing and propagating the
savageness of human nature? 'Thou art the man.' Hazael.

17. I turn to a truer view. These sufferings,
certainly of children, are martyrdoms [\emph{vide} sermon, December
28, 1856]. And how many brought to repentance. The long suspense
described by a lady in a letter from Cawnpore. Her little child
restless and nervous, etc., etc. Priests and nuns have suffered.

18. May Our Blessed Lady, whose Maternity this
day is, protect them.

19. Let us pray that it may all be overruled to
our country's good.

\xsubsection{Instances.}

1. The savage conquerors in the East—Zingis,
Timour, etc. 'Zingis depopulated the whole country from the Danube to
the Baltic in a season, and the ruins of the cities and churches were
strewed with the bones of the inhabitants. He allured the fugitives
from the wood under a promise of pardon, got them to gather in the
corn and grapes, and then put them all to death. At one place he put
to death 300 noble ladies in his presence. He divided cities into
three parts. He left the infirm and old, enlisted all the young
men into his army, made the women, the rich, and the artisans his
slaves. Almost fabulous his slaughters: at Maru 1,300,000, at Herat
1,600,000, at Neisabour 4,647,000.' 'Timour at Delhi massacred 100,000
prisoners, because some of them showed exultation when the army of
their countrymen came into sight. At Ispahan 70,000 human skulls, at
Baghdad 90,000.'

2. The Persian shah about fifty years ago, in
Morier's \emph{Zohrab}: three bushels of eyes at Asturabad.

3. Herod with the innocents. Persians and Rome
savage persecutors of Christians. Tortures.

4. Middle Ages. (From the newspapers of the last
week.) When the Norman barons conquered England, they persecuted the
poor English in the most savage ways to gain money or their
submission. The \emph{Saxon Chronicle} says: 'The men were hung by the
feet and by the thumbs, and thus smoked over a smouldering fire.
Knotted strings about their heads and pulled tight till they pierced
the brain. They were put into dungeons with adders and toads; they
were put into chests too short and narrow to hold them, and thus were
crushed. They were attached by sharp iron collars to a beam, so that
they could in no ways sit, nor lie, nor sleep, but they must bear the
iron. Thus many thousands were exhausted by hunger.'

5. Guicciardini and Muratore.

6. The German and Spanish soldiers of Charles V.
in Rome, just before St. Philip's time, on getting possession of Rome,
put to death 4000 soldiers and inhabitants. I will not stop to speak
of their plunder of the great riches which it contained, but I
am speaking of their savage cruelties. They got all the cardinals,
bishops, prelates and nobles they could, and put them to the torture
to extract money from them. (St. Caietan.) They carried off the noble
Roman ladies, and the nuns from all the convents, and treated them to
the most horrible outrages. Their shrieks resounded on all sides. And
this went on for days. Many died in their torments, and soon
afterwards. It was worse than the Goths. Only a few years ago, as many
as sixty priests are said to have been shot by the triumvirs at Rome \footnote{In 1848.}.

7. Not much more than a hundred years ago, there
was an attempt in Scotland to place the old royal family upon the
throne, and a rising of the people. After the English had got a
victory, they put the wounded Highlanders to death in cold blood. They
dragged them out of the huts or thickets and shot them, or dispatched
them with the stocks of their guns. One farm-house with twenty wounded
men in it, they burned to the ground; and all this though the
Highlanders had been behaving in the most noble and generous and
merciful way to the wounded English. Moreover, they brought together
into heaps all the wounded from the field of battle, ran them through
with their swords, and cut the throats of those who were found sick in
bed.

8. And a century before that, Cromwell and his
English soldiers had committed as great, or greater, atrocities in
Ireland. When he took Drogheda, he offered quarter to all who would
lay down their arms. And when they did so, he broke his word,
and began an indiscriminate slaughter. The massacre lasted for five
days, so that the streets literally ran with blood. They killed not
only the soldiers who were in arms against him, but the inhabitants of
the town. A thousand fled for safety to the church, and he killed them
all in it. He then marched to Wexford, and did the same. He killed
here, too, the inhabitants as well as the soldiers. Three hundred
women gathered round the great cross of the place; they were all put
to the sword. Some writers say that the slaughter at Wexford amounted
to 5000, and Cromwell himself confessed that it was as much as 2000.

9. Nay, what is the state of feeling of our own
soldiers at this minute in India? You will say they have reason; but
who has not a reason? The Hindoos have thought they had. Now what do
we read? One of the officers before Delhi writes: 'We must have blood.
The streets of Delhi will be a fearful sight. Our men are mad for
revenge.' Another says: 'I only trust all the women and children will
have been removed (by the time we take the place), for when we are
once inside few will be spared.' Another says: 'Our men cannot be
restrained, and they are \emph{like demons} let loose.' Another says:
'I believe the city will be given up to three days' plunder. I fear it
will make our Europeans very undisciplined. Heaven knows, they are
hard enough to control here; but when they have once gone in like
bloodhounds, and been allowed to plunder, they will be downright
demons.' These things are perhaps going on now.

\xsection{April 11, 1858 (Low Sunday) [The Church]}

1. INTROD.—Last
week I spoke of one of those great and august works with which our
Lord followed up the great Act of Sacrifice, viz. the foundation of
His Church. Nor can there be a more suitable time than this season to
speak of it, considering it was the chief concern, as far as we know,
of the forty days; vide the gospel of this day.

2. Now in this we differ from all other religions
about us. They all profess to have the truth as well as we profess it,
but there is one thing they do not profess, viz. that their religious
society is founded by Almighty God. We do of ours.

3. And since they do not profess it, they will
not let \emph{us} have what they have not themselves.

4. State the doctrine. We profess, not only our
religion, but our society to come from Almighty God; we profess it to
be divine. We profess it to have a multitude of privileges, etc.

5. Now you may be asked sometimes by a serious
objector, sometimes by an inquirer, how it is that we know that the
Church comes from God? I answer that it bears the proof of it to all
serious men on its very face, if they will but be patient to examine;
and I will say how.

6. I said on Good Friday concerning the world
that its strength is in the \emph{look} of things. Men associate
together, say the same thing, and seem strong. They keep up
appearances. But there is an \emph{inside} to things as well as an \emph{outside}.
And here is the weakness of the world as a prophet, that it does
not touch the inside.

7. Men cannot live for ever on externals. They
have heart, affections and aspirations, and the world cannot satisfy
these. They have a conscience; they sin, and need direction.

8. Now this is what our Saviour, when on earth,
did for His disciples; and thus He attached them to Him. He was a
living object of worship—(1) He gave pardon; (2) He gave direction.

9. When He went, He said He would not leave them
orphans.

10. This was fulfilled in the Church: (1) pardon,
(2) direction, (3) presence—(enlarge).

11. Hence suited to our \emph{need}—(enlarge).

12. Faith only requisite.

\xsection{April 25 (Third Easter) [The Holy Eucharist]}

1. INTROD.—Passage
in today's gospel: 'Yet a little while.' \footnote{'A little while, and now
you shall not see me: and again, a little while, you shall see
me.'—John xvi. 16.} \emph{Vide} also John xiv.

2. This brings before us the thought of the Holy
Eucharist, in which it is so wonderfully fulfilled.

3. (This great marvel or miracle, describe
generally. Some remarks on it.)

4. Now first, the doctrine of the Resurrection.
How wonderful this is. Describe it. It is as stupendous a miracle as
any, though luckily Protestants retain it.

5. Another marvel. The risen Body shall ascend to
heaven and live there for ever and ever.

6. Now see what this implies. We cannot suppose
that our present gross bodies shall be in God's presence for eternity.
Accordingly St. Paul says that 'we shall all be \emph{changed}' \footnote{1 Cor. xv. 51, 52.}—'animal body and spiritual body.' \footnote{'Seminatur corpus animale,
surget corpus spiritale.'—1 Cor. xv. 44.}

7. Now our Lord as the first-fruits is already
gone to heaven, therefore His Body would be altered.

8. Now we know in a way what great changes matter
goes through—ice, melting iron, gas. On the butterfly as an emblem.

9. Four great properties [of a glorified
body]—impassibility, activity, brightness, subtlety. Again (1) doors
being shut (subtlety), (2) appearing and disappearing (activity):
another form [of activity], ascending up on high, (3) brightness, Apoc.
i. [14-15].

10. Two others: (1) in many places at once, (2)
at a point. As to the second, how did He get through doors? Child
increasing, growing, etc. As to the first, God a spirit. His appearing
to St. Paul.

11. These two in the Holy Eucharist.

12. Let us adore.

13. How wonderful, by making it a miracle, He has
kept it secret, for the \emph{world will not believe}, John xiv., xvii.
And so when He was on earth, John i. [5].

\xsection{May 30 (Trinity Sunday) [The Blessed Trinity]}

1. INTROD.—This
is the everlasting mystery. All other mysteries arise in \emph{time},
this in eternity.

2. No words of man can explain it. Three are one.
Now soul and body are one, all the faithful are one. The soul, when
enveloped in the Divine Essence, hereafter is one with it. But no
illustration of earth can give the faintest shadow of the truth.

3. Yet though so difficult for the reason, it is
not difficult for devotion, and thus is fulfilled the saying, 'Thou
hast hid these things from the wise and prudent, and revealed them to
Thy little ones.'

4. For each truth concerning the Holy Trinity is
easy. It is their \emph{combination} is the mystery. Exemplify. That
there should be one God—that God should be called or be the
Father—that He should be the Son—that He should be the Holy
Ghost—but that all three propositions [four] should be true \footnote{See Note
12, p. 339.}.

5. Explain more fully. The Father is absolutely
the One God, as if no Son and Spirit, etc. All the attributes, etc.
belong to the \emph{Father}, and to the Son, etc. Whatever the Father
does or is, that the Son does or is—not as two, but One.

6. But they divide offices in mercy to our
infirmity.

7. Now for devotion. The Father, the Creator, the
Preserver, Governor, Judge. Source of all good; the harbour of our
rest; heaven, etc.

8. The Son has taken on Him our nature, etc.

9. Holy Ghost Sanctifier, etc.

10. Now while we address each in \emph{devotion}
as the One God, we may leave it to the next world how Each of
Three can be the One God.

11. The joy of heaven, when all mysteries will be
removed.

\xsection{January 30, 1859 (Fourth Epiphany) Blessed Sebastian Valfre}

1. INTROD.—The
day. Blessed Sebastian born about thirty years after St. Philip's
death.

2. General resemblance between the two: (1) St.
Philip's early devotion to God; no mortal sin; hard life, long life,
and hidden life. (2) Blessed Sebastian the same. More is known of his
boyhood—(give instances). Like St. Philip, he labours till the day
of his death. Circumstances of his death.

3. Differences from St. Philip: (1) St. Philip
without object [in life at first]. \emph{Others} go to Rome for
preferment. He did not aim at being priest; he did not aim at founding
a Congregation—like Benedictines, no great \emph{work} \footnote{'These early Religious [\emph{i.e}.
the first Benedictines] ... had little or nothing to do with
ecclesiastical matters or secular politics; they had no large plan of
action for religious ends; they let each day do its work as it
came.'—'The Mission of St. Philip,' \emph{Sermons on Various Occasions},
p. 225.}, but [like] St. Vincent of Paul, e.g. \emph{works}. (2)
Blessed Sebastian had the definite object of being a priest.

4. Hence Blessed Sebastian's particular
character—of priestly, pastoral work of every kind—(go into
details). \emph{He} differed from other priests [of his time] in his \emph{incessant}
work.

5. Well is it that his feast is this year on a
Sunday, for we have just set up an altar to him.

6. And well is his feast at this time of the
year, for it is the time we came to Birmingham.

7. We have been now ten years in Birmingham, and
when I thank God for what He has allowed us to do, I can suitably do
it, for I have had less to do with it than others. You know how many
of us have devoted ourselves to missionary work.

8. Pray then for us—(details).

9. And pray for The Oratory at Turin—(details).
Subscribing to the Achilli fund; our going there—[their] simplicity
of life, etc. Their present troubles. Little to choose between one
country and another. All good men persecuted.

\xsection{August 19, 1860 (Twelfth Pentecost) God the Stay of Eternity}

1. INTROD.—The
gospel says, 'What shall I do to inherit?' etc. Here this man,
whatever his own character, asks an all-important question.

2. He implies the soul will live for ever.

3. What is eternity? Why, it is awful. I cannot
call it good in itself. Some good and wise people have said so, but
for me it is the most awful thought in the world.

Consider it. Time breaks to pieces everything;
much more does eternity. Our soul can never die, but it can get older
and older. Fancy this—older and older, colder and colder, so that
the longer we lived the more miserable [we should become]. Therefore,
when I look at eternity itself, it is a sort of living death to
creatures such as man, and no good. Who can bear the weight of eternal
years?

4. The scribe, then, does not ask for 'living for
ever,' but for 'eternal life.' Life is something more than living; it
is to live vigorously, to be always young, etc., etc. Many have no
youth, as some years have no spring. It is therefore to be happy, and
happier and happier as time goes on.

5. This being the case, it is plain also that
nothing but what is infinite can sustain eternity. We read in romances
of two persons determining to die, and die together, and care for
nothing else, not even God—vain thought! We want something more than
ourselves, something more than the creature. We must be associated
then, and one with the Creator.

6. God then, the Almighty and the Infinite, is
the only stay of eternity.

7. Now then we see the meaning of our Lord's
answer to the scribe, of loving God, for He alone is eternal, and
unless we are conformed to Him, we shall be miserable in eternity.

8. Let us learn to love. We know what it is on
earth to love a person. Signs of love—liking the presence, speech,
etc., of the loved person; taking up his opinions, etc., etc.

\xsection{September 2 (Fourteenth Pentecost) [The Holy Angels—I]}

1. INTROD.—This
month leads us to think of the holy angels. It is a far larger subject
than I can get through this evening. There are two points of view in
which they are to be considered—in nature, and in grace. And this
evening I will speak of what they are in their nature.

2. God created them in the beginning of all
things, with all other things which He created. When He created the
heavens He created them, and He created them in the heavens. Here is
the vast difference from earth; for man was created on earth, in order
that in time he might attain to the heavens.

3. Simple spirits—hence no form—angels \emph{with
wings}. Mere appearances—as in the Holy Eucharist.

4. No shackle of body. We too are spirits, but in
bodies (bodies \emph{part} of us, disembodied saints \emph{desire}
their bodies). Hence we are \emph{sluggish, passionate}, etc. Hence \emph{we
sleep}, not they. We cannot move about quickly; they in the
twinkling of an eye from heaven to earth.

5. \emph{Most perfect of creatures}—the \emph{image
of God's attributes}.

6. Their knowledge most comprehensive. They do
not \emph{learn}, they do not discover, but at once from their nature
they know intuitively all things of the world; whereas the greatest
philosophers with pains only knew a little.

7. They know God and His attributes by nature,
even without grace. They understand His attributes, etc. They see God
in all things, never being seduced by the creature, as separate from
Him.

8. They have a natural love of God, from the
perfection of their \emph{reason}. They love Him above all things.

9. They love each other—and each order of
angels, in its own degree, fittingly.

10. Three points which they have not by nature:
(1) knowledge of the future; (2) of the heart; (3) of the mysteries of
grace.

CONCLUSION.—Many
wonderful things in this world, but an angel more wonderful than all.
If a creature so wonderful, what the Creator?

\xsection{September 9 (Fifteenth Pentecost) The Holy Angels—II}

1. INTROD.—Recapitulate.
The Creator might make ten thousand worlds, each more perfect than the
preceding, all more perfect than this. We know of but one besides
this—the universe of angels. This may be otherwise. The angelic
world differs from this, in that each part is perfect and independent
of any other part.

2. Differences, but they \emph{all} excel in two
things naturally—(1) strength, (2) purity.

3. Purity—'As the angels' \footnote{'In the resurrection
they shall neither marry, nor be married, but shall be as the angels
in heaven.'—Matt. xxii. 30.}—no bodies. Strength—Exod. xii. \footnote{The destroying
angel—death of the firstborn in Egypt.}; 2 Kings xxiv. \footnote{The angel of the
pestilence whom David saw by the thrashing-floor of Areuna the
Jebusite.};
4 Kings xix. \footnote{The angel that slew the
host of Sennacherib.}. Their
voice—Apoc. x. [3] \footnote{'He [the angel] cried
with a loud voice, as when a lion roareth: and when he had cried,
seven thunders uttered their voices.'};
1 Thess. iv. [16] \footnote{'The Lord himself shall
come down from heaven ... with the voice of an archangel ... '}.
Number—count the lowest—everywhere guardian angels—one to every
man, though at one time a thousand millions of men.

4. Their differences. Some think no two [are]
alike, but differ specifically, as eagle, dove and nightingale.
Indeed, it is difficult, as I have ever thought, to consider pure
spirits other than specifically distinct, because since no parts or
whole in the angelic world, there are no logical laws in it (except
virtue). But, leaving this question, [there are] nine orders in three
hierarchies—(enumerate) \footnote{See Note
13, p. 340.}.

5. Such by \emph{nature}, now by \emph{grace}.
From the first instant of creation endowed with
grace—habitual—faith, hope, charity. Knowledge of Holy Trinity,
etc.

6. All of them holy, but in \emph{proportion} to
their nature. All have all virtues, but each order rises, having not
only \emph{all} virtues in greater perfection, but a characteristic
virtue.

7. (1) Angels—contentment; (2)
archangels—imitation of the perfection of all the other
orders—absence of all pride and rivalry; (3)
principalities—simplicity of intention.

8. (4) Powers—tenderness and sweetness; (5)
virtues—courage; (6) dominations—zeal.

9. (7) Thrones—submission and resignation; (8)
cherubim—knowledge; (9) seraphim—love.

10. Honour due to the angels, Exod. xxiii., Josh.
v., Judges vi., xiii., Daniel x. Explain Apoc. xix. 10, that St. John
was so great that he was not to adore.

11. Let us honour them in the best way, but
imitating, like the archangels, the virtues of each order.

\xsection{September 16 (Sixteenth Pentecost) The [Holy] Angels—III}

1. INTROD.—Recapitulate.

2. The angels were all created perfect and gifted
with supernatural holiness. Even Lucifer, etc.

3. They were first of all put on their trial.
They did not see the face of God.

4. The time of this trial—no natural term like \emph{death}—shorter
than men['s] because of their spiritual nature, as it was so
penetrating, etc., might stand or fall for good in a short time. (Why
not in an instant?)

5. Who fell? [Some] out of all the orders.
Lucifer a seraph.

6. The numbers. Some think a third—Apoc. xii.
[4].

7. The sin of the angels, one and the same in
all, from imitation. Lucifer led them.

8. \emph{What} [was] the sin? All [sins] in one
doubtless, but especially pride. What \emph{kind} of pride? Obstinacy,
ambition, disobedience, arrogance?—all doubtless, but especially and
initially reliance [on] and contentment in \emph{natural} gifts, with
despising \emph{supernatural}.

9. Additions to this pride (1) a sort of \emph{sensual}
love of self; (2) presumption, ambition, hatred of God; (3) jealousy
of man who was to be created.

10. Battle in heaven. (Michael—'Who as God?')
Each party trying to convert the other to its own side.

11. Cast into hell—fire in their
spirit—though they are now out of it [till the day of judgment].

12. Allusion to matters going on in Italy. Good
and bad not so keenly divided as in angels, but still it is the devil
against Michael.

\xsection{September 30 (Eighteenth Pentecost) The Holy Angels—[IV]}

1. INTROD.—About
guardian angel.

2. The different works of angels. The word [\emph{angelos}]
denotes work and service.

3. What orders of angels have to do with this
universe? The lowest, \emph{i.e}. the \emph{angels}, are the
ministers. Mundane or exterior, and heavenly or domestic works.
Extraordinary missions—the cherubim of Eden [Gen. iii. 24]—the
seraph [in] Isaias [vi. 6, purifying the prophet's lips with living
coal from the altar]—Gabriel and Mary [the Annunciation]. One
[angel] making charge over to another to execute.

4. First work—'rolling the heavens' \footnote{Probably a quotation
from some poet.} [\emph{i.e}. directing the movements of the heavenly
bodies]—science need not [be supposed to have] superseded this—see
my sermon, Parochial, etc., vol. ii. [The Powers of Nature. Feast of
St. Michael, etc.]—John v. [Pool of Bethsaida].

5. Second work—guardians of nations, provinces,
cities, bishopricks, churches. 'Let us depart hence.'

6. Of individuals. Every one from the time of the
soul's creation to death. And every one. Judas, Antichrist.

7. St. Frances of Rome.

8. (1) Odiousness of the charge, \emph{e.g}. St.
Paul linked to a soldier; (2) condescension, etc.; (3) encouragement
to us, and comfort.

\xsection{October 7 (Nineteenth Pentecost) Cardinal Virtues—Prudence}

1. INTROD.—Apparently
the greatest. It is called prudence, wisdom, judgment or discretion.
'Be ye prudent as serpents.'

2. For to resist self (or temperance), the world
(or fortitude), and to be in grace (justice), is obviously necessary
for all, but why prudence?

3. Again, a man on looking back will often say,
'By the grace of God I overcame myself—the world—and I generally
served God and my neighbour, but alas! all my troubles have come from
want of prudence.'

4. St. Anthony (in Cassian), when all the monks
assigned different virtues as the greatest, said prudence—because it
hinders virtues from becoming vices. This partly lets us in to what
prudence is. Let us take different instances of imprudence.

5. (1) Virtue being in a man, prudence is the
directing principle; what is virtue in one man is not in another.

6. (2) On \emph{turning-points} in life—as men
mistake their way in the mountains and come to precipices.

7. (3) On avoiding occasions of sin—temptation
nearly always comes before sin, as bad food, air, lodging, etc.,
before illness. If we avoided temptation, how little sin we
should do; but prudence is the directing principle.

8. (4) Command of tongue—sudden words, what
harm they do! Our Lord when they attempted to entrap
Him—Joseph—David.

9. (5) On avoiding scandals. We ought ever to
have our eyes about us lest we do others harm.

10. But how are uneducated men to do what seems a
virtue of the \emph{perfect}? At least they may take advice.

11. And they may pray. Two passages in Scripture:
Proverbs ii. 3-5 \footnote{'If thou shalt call for
wisdom, and incline thy heart to prudence; If thou shalt seek her as
money, and shalt dig for her as for a treasure; Then shalt thou
understand the fear of the Lord, and shalt find the knowledge of God.'},
Ecclus. li. 11 \footnote{'I remembered thy mercy,
O Lord, and thy works, which are from the beginning of the world. How
thou deliverest them that wait for thee, O Lord, and savest them out
of the hands of the nations. Thou hast exalted my dwelling-place upon
the earth, and I have prayed for death to pass away. I called upon the
Lord, the father of my Lord, that he would not leave me in the day of
my trouble, and in the time of the proud without help.'}.

\xsection{October 14 (Twentieth Pentecost) Cardinal Virtues—Justice}

1. INTROD.—Justice
a name for \emph{all} virtue. The robe of justice—justification. How
great then must be the virtue proper so-called.

2. And so the beatitude: 'Who hunger and thirst
after justice.'

3. The attribute of God enhances this, for the \emph{first}
attribute we know Him by is justice; viz. in conscience—\emph{before}
experience, before the knowledge of providence, before we look out
into the visible framework of the world. Justice and all-knowledge the
two; and in Christianity it is the two, love and justice. And where
should we be without Christ's justice? Merits of saints founded on the
covenant.

4. What \emph{is} justice? Giving to all their \emph{due};
text in the Romans, 'Honour to whom honour,' etc. Hence it is
synonymous with the habit of 'doing one's duty,' whether to God or our
neighbour. To God adoration, devotion, etc., and to the holy angels,
etc., but I shall not \emph{insist} on this part of the subject.

5. To man it is summed up in the maxim, 'Do as
you would be done by.' This is placing conduct on the basis of
justice. This basis of justice, for not 'as others would \emph{like}
you to do,' but 'ought to wish you to do.' And so 'forgive us our
trespasses, as,' etc.; Matt. xviii. 23, parable \footnote{'The servant who owed
ten thousand talents.'}; 'If I have washed your feet,' John xiii. 14.

6. Parts of justice: (1) truth, (2)
honesty—restitution; the terrible onus of restitution shows how
important a virtue justice is.

7. (3) Faithfulness, and (4) gratitude, \emph{e.g}.
to parents.

8. (5) Liberality—detachment from money as
being the opposite to rapacity and avarice.

9. (6) Courtesy in manner and act.

10. (7) Equity, consideration, kindness in
judging, putting oneself into other person's situation. Not
'swift to wrath,' James i. 19; Ephes. iv., last verses \footnote{'Let all bitterness and
anger ... be put away from you,' etc.}.

11. Application on the contrary—our only notion
commonly of justice, is justice to ourselves, hence anger [\emph{dia t\={e}n
phainomen\={e}n hubrin}], etc., etc.

12. CONCLUSION.—At
the judgment this is the attribute God will exercise. Our justice will
then have a peculiar claim, while we are invoking God's promises.

\xsection{October 21 (Twenty-first Pentecost) Cardinal Virtues—Fortitude}

1. INTROD.—Fortitude
and temperance (unlike prudence and justice), and fortitude
especially, virtues of warfare in a fallen world. Cowardice the
opposite. We know about bravery and cowardice in human matters. How
our warfare spiritual, Eph. vi. 12 \footnote{'For our wrestling is
not against flesh and blood, but against principalities and powers,
against the rulers of the darkness of this world, against the spirits
of wickedness in the high places.'}.

2. 'Overcometh the world,' 1 John v. 4 \footnote{'Whatsoever is born of
God overcometh the world.'}; overcometh the devil, Apoc. xii. 10-11 \footnote{'The accuser of our
brethren is cast forth, who accused them before our God day and night.
And they overcame him by the blood of the Lamb, and by the word of the
testimony; and they loved not their lives unto death.'}.

3. Hence the Old Testament puts it forth as the
characteristic virtue. The spies of the Lord, Deut. xxxi. 7 \footnote{'And Moses called Josua,
and said to him before all Israel, Take courage and be valiant: for
thou shalt bring this people into the land which the Lord swore he
would give to their fathers.'}; Josh. i. 6, 7, 9 \footnote{'Take courage and be
strong ... Take courage and be very valiant ... Behold I command thee,
take courage and be strong. Fear not and be not dismayed, because the
Lord thy God is with thee in all things whatsoever thou shalt go to.'}—Gideon,
David; Aggeus ii. 4 \footnote{'Yet now take courage, O
Zorobabel, saith the Lord; and take courage, O Josua, son of Josedec,
the high priest; and take courage, all ye people of the land, saith
the Lord of hosts, and perform: (for I am with you, saith the Lord of
hosts).'}.

4. In the new covenant, martyrs, active courage
as well as passive—St. Ignatius. St. Barlaam—his hand burnt off.
All the children of the city coming to the governor saying, 'Kill us,'
and he saying: 'O cacodaemons, have you not precipices and halters?'

5. This is how Christianity was set up—a whole
epistle, the Hebrews, not to say 1st of St. Peter, on the duty and
virtue.

6. But you will say this is beyond us. How is it
a cardinal natural virtue? Well, I can give instances, \emph{e.g}.
'because iniquity shall abound,' etc. \footnote{'Because iniquity shall
abound, the love of many shall grow cold.'—Matt. xxiv. 12.}; cowardice—'lest they be discouraged' (\emph{ut non pusillo
animo fiant}) \footnote{Coloss. iii. 21.}.

7. Cowardice in telling the truth.

8. Cowardice in resisting evil, in not going
after the way of sinners in act and deed.

9. Impatience of ill-usage from others.

10. Impatience at continued evils;
disgust—giving up.

11. This brings me to perseverance. It is
difficult to persevere in any course, though no positive obstacles or
opposition. How great this cardinal virtue then, as connected with the
end of life.

12. If the merits of the martyrs are to assist
us, let us merit that assistance by some portion of their bravery.

\xsection{October 28 (Twenty-second Pentecost) Cardinal Virtues—Temperance}

1. INTROD.—Temperance
contrasted with fortitude, as within with without, and the pleasant
with the painful.

2. Now to explain it. Our soul may be said to
have in it two natures, and at variance, and so opposed that peace and
unity implies the subjection of one to the other—as two combatants
will fight till one or other is thrown.

3. Reason and passions—grief, joy, anger,
desire of having, fear—all going into extreme manifestations, and
needing a controller. We see it in brute animals. When they cease [?]
it is not that reason governs them, but the object [that excites them]
is removed.

4. Comparison of a child on horseback. On the
other hand, a rider who has perfect command—the Tartars, who live on
horseback.

5. Now in the case of the warfare of the soul the
struggle more serious and the dangers greater, because (1) the
passions have \emph{instruments}, as being united to the body; (2)
objects sensible; whereas the object of the reason and
conscience, Almighty God, is \emph{unseen}.

6. Therefore a certainty of the subjection of the
soul, unless for a remarkable virtue, viz. temperance, or
self-government, or control. It is the very critical, or cardinal, or
most essential and directing virtue.

7. This self-rule is what \emph{makes a man};
without it a man is a slave, etc.—laments and curses himself, etc.

8. Hence not heroic, but we see what it is in the
saints. It is the characteristic of the saints, and thus is inflicted
[\emph{sic}] on us, that in its degree it is the characteristic of a
man. You may have wondered why a saint is characteristically
mortified.

9. In saints we specially see how it subserves
the soul; their fastings, etc., etc., are to make them pray better,
etc., etc.

10. I need not give instances as in the former
virtues, but I will mention specially—

11. The necessity of temperance in \emph{thoughts}
and in \emph{words}.

12. If we would have the saints assist us, let us
cultivate that virtue which was their distinction.

\xsection{February 24, 1861 (Second Lent) State of Innocence}

1. INTROD.—State
of our first parents. Image and likeness \footnote{'Let us make man to our
image and likeness.'—Gen. i. 26.}.

2. Image in nature. All things are in a way in
the image of God, as being His creatures. Man in a special way.

Soul (1) a spirit; (2) immortality; (3) knowledge
illimitable \footnote{Note he does not say \emph{infinite}.}; (4) free
will; (5) Godlikeness, as Satan said \footnote{'You shall be as
gods.'—Gen. ii. 5.}.

Body—beautifulness, perfection of form, etc.,
etc. but still defects.

Soul—passions against reason; body—mortality.

Thirdly, war of soul with body, as having
different ends. Strange the body cannot be without the soul, nor the
soul well without the body; yet they cannot agree together.

3. Almighty God knows what He has created, and
therefore He did not leave man thus, but gave him a supernatural gift.

4. Likeness—(explain)—sanctity. In fact
[also] health and strength, viz. three subjections: (1) Soul to God,
(2) passions to reason, and (3) body to soul \footnote{These 'subjections'
pertain to the gift of \emph{Integrity}.}.

5. Hence a knowledge of mysteries. An absence of
passions, only good affections.

6. Yet we need not lament paradise—on account
of the future glory promised to us.

7. '\emph{Lest he take of the tree of life}.'

\xsection{March 3 (Third Lent) State of Original Sin}

1. INTROD.—State
of original sin. Deprivation of grace; consequences, the 'wounds.'

2. Stripped of God's supernatural gifts or grace.
He might never have given it, and then no punishment. But since
we were intended for heaven, it was a great punishment. Take the case
of a person born to wealth, etc., of cultivated mind, etc., banished
to a desert island.

3. And as a spendthrift involves all his
descendants, so here.

4. But worse. I described last week the triple
subjection. Well, when man cast off God, his passions and affections
rebelled against his reason, and his body against his soul. Case of a
strong man or child on horseback. Daniel in the lions' den.

5. This great calamity constitutes the wounds of
human nature—parable of the Good Samaritan—\emph{first} stripped, \emph{then}
wounded.

6. Now here I shall view them as three—the
absence of three subjections—sloth, selfishness, sensuality.

7. Describe sloth—that deadness, blindness of
soul, dislike of prayer, disgust at religion, liking to ridicule it.
Dislike of ruling our mind and heart, etc., etc.

8. Therefore selfishness—making self the
centre, etc.

9. And therefore sensuality—idolising the
creature.

10. And then these may be considered sins against
God, our neighbour and ourselves. Contrasts—love of God, love of our
neighbour, and self-command.

11. They branch into the seven deadly sins: (1)
sloth; (2) pride, avarice, anger, envy; (3) gluttony and luxury.

12. They tend to utter death. Are you to live
after this life? What is your state then, and is God in heaven?
Again I say, what is your state then, with the world swept away?

13. What is your duty? As a ruined man might try
to repair his fortunes. Children of this world labouring to regain an
ancestral estate.

14. Two great graces: illumination and
excitation.

15. 'Now is the acceptable time.'

\xsection{March 10 (Fourth Lent) [Restoration]}

1. INTROD.—All
rational creatures find a need in their nature, and are insufficient
for themselves.

2. But man especially. We see this contrasting
him with the inferior animals.

3. He has a body as they have, but observe the
difference. They are suited to their \emph{habitat} by nature, he not.

4. Brutes need no dwellings, no clothes, no
prepared food, no cooking. Bread and wine both manufactures, etc.,
etc. Caves, skins and furs, teeth and claws, stomachs. Armour too and
arms; flight.

5. Again, if they need anything, Nature supplies
them with instinct, \emph{e.g}. changing colour in the north, etc.,
nests, holes, care of young, etc.

6. But consider man. He is not adapted to this
world in which he finds himself, and would die if left to himself.
Arts are necessary, dwellings, etc., etc.

7. Hence again he has to live in a society—for
each needs the aid of many others—so many trades, etc. Then again
language. This not by nature, but by imitation. Education, books, etc.

8. In like manner, as he is thus dependent in
body and mind, so he is in \emph{soul}, in religion. Insist on the
analogy, but with this difference, that for body and mind he can get
help from other men, as by education, but this change (\emph{al}.
addition) in the spirit comes \emph{from God alone}. When man fell
everything went, but with this difference, that he had powers within
himself in course of time to remedy the evil in all matters of this
world, but \emph{not of his soul}.

9. In paradise he was healthy, etc., etc. He
needed no habitation, etc. However, in process of time he could by his
own powers find out food, medicine, dwellings, etc., hence societies,
kingdoms; \emph{but not religion}.

10. Nay, worse still, the very advancement in
society, in civilisation, is antagonistic to religion. Society viewed
on its religious side is the world, one of our three enemies.

11. (Analyse.) Each man condemns himself, but if
another does the \emph{same}, the example is a safeguard, defence, and
excuse. How when a multitude [does the same], the world becomes a
prophet antagonist to conscience.

12. Also the pomp, glory, etc., of the world
becomes an idol.

13. And the more society grows, the worse the
world.

14. I said last Sunday the man got worse and
worse as time went on, much more society.

15. Explain 'progress.' \footnote{See Note
14, p. 341.}

16. Yes, in worldly matters.

17. But in religious, not.

18. Worse and worse—judgments, flood, etc.

19. Only true progress in the individual, in the
heart.

20. Let us at this season follow it out.

Continue

Back

\xsection{November 2, 1862 (Twenty-first Pentecost) On the Gospel of the Day— [The Parable of the Servant who Owed Ten Thousand Talents—Matt. xviii.]}

I consider this parable, and the other passages
of our Lord's teaching which are parallel to it, of a very awful
character. I think all of us will say so who seriously turn their
minds to consider them. (Go through it.)

It is introduced by a question of St. Peter,
which itself may be viewed in connection with another declaration of
our Lord's on the same subject, which is recorded in the 17th of St.
Luke, vv. 3-5 \footnote{'Take heed to yourselves: If thy brother sin
against thee, reprove him; and if he do penance, forgive him. And if
he sin against thee seven times, and seven times a day be converted
unto thee, saying, I repent; forgive him. And the apostles said to the
Lord, Increase our faith.'} (Quote.)
Apparently in allusion to this, or in some connection with it, St.
Peter asked: 'Lord, how often,' etc. Matt. xviii. 21-22 \footnote{'Then came Peter unto him,
and said, Lord, how often shall my brother offend against me, and I
forgive him? till seven times? Jesus saith to him, I say not to thee,
Till seven times: but, Till seventy times seven times.' And thus our
Lord takes the opportunity to follow out and complete the great
evangelical doctrine which He had begun to declare in the passage
recorded in St. Luke. St. Peter asked if seven times would be enough,
and our Lord answered, 'I say not,' etc., etc.}.

In the same way in the sermon on the mount, Matt.
v. 22-24 \footnote{'I say to you, Whosoever
is angry with his brother shall be in danger of the judgment: and
whosoever shall say to his brother, Raca, shall be in danger of the
council: and whosoever shall say, Thou fool, shall be in danger of
hell fire. If therefore thou offer thy gift at the altar, and there
thou remember that thy brother hath anything against thee; Leave there
thy offering before the altar, and go first to be reconciled to thy
brother, and then coming thou shalt offer thy gift.'}. And He has
introduced it as one of the seven petitions of His own prayer, which
is the first element and type of all our devotions, and which we say
every day. Forgiveness of injuries then bound up in the very idea of
prayer in the evangelical law; and our Lord in a passage in St. Mark
seems distinctly to say so; for after speaking of the faith which will
move mountains, He proceeds, Mark xi. 25-26, 'And when you shall stand
to pray, forgive, if you have ought against any man: that your Father
also who is in heaven may forgive you your sins. But if you will not
forgive, neither will your Father that is in heaven forgive you your
sins.'

Now this great Christian precept is often
expressed in these two words: viz. that when injury is done to us, it
is our duty to forgive and forget. Let us dwell upon these.

Now, at first sight, we shall all of us allow
that it is a very beautiful precept, especially when we are young,
when our hearts are light and open, and our tempers generous; we shall
on the one hand think it admirable and great, and, not having had to
practise it in fact, we shall be drawn to it, think it easy, and
resolve to observe it as life goes on. I can fancy young people
drawing before their minds pictures of injuries done them, of their
forgiving the injuries, and returning good for evil. And when they
read accounts of men who have done so, and instances of generosity,
magnanimity, patience and nobleness in this respect, they are greatly
moved and filled with a love of the virtue. Nor is it only a beautiful
precept, it is of a most useful and expedient character too. Every one
must confess who turns his mind to the subject, that the world would
go on far better, that all men would be happier, if this precept was
universally observed. For what is a greater or wider scourge of man
than war, dissensions and litigation? and though these miseries arise
in a great measure from covetousness (James iv.), they arise still
more from passion, from a sense of injuries, from a fierce
determination to retaliate, from a thirst for revenge. James iv. 1-2,
'From whence are wars and contentions among you? ... you covet, and
have not; you kill, and envy, and cannot obtain.'

To forgive and forget, then, is (1) at first
sight a beautiful, an admirable precept, and (2) one which on long
experience leads to the greatest benefit to mankind. All men are
interested in its recognition and observance; yet it will be found not
at all easy in fact, but a very difficult precept, one which is but
rarely obeyed and very partially, where it is not altogether
neglected; and further, one to which many plausible objections may be
made, and many arguments in favour of a contrary course, which become
formidable when they are brought to defend that unwillingness to obey
it; and the difficulty of obeying it, which in matter of fact
will be found in human nature.

Now I will first set down what I conceive the
precept to be, and next consider how the objection to it arises.

(\emph{I was interrupted, or I meant to have written
a sketch of a whole sermon. I have forgotten now my arrangement. I put
down some isolated} [\emph{topoi}].) (1) Not to forgive is even
contrary to justice, a \emph{higher} kind of justice than natural
justice, for we should do as it has been done by Almighty God to us.
(2) Forgetting, yes, as God forgets, for He forgets by putting aside,
behind His back, our sins. (3) We should put aside also, for a reason
special \emph{to us}, for the thinking of injury is a temptation to
avenge it. (On distrust necessarily remaining after forgiveness \footnote{These words were added in
pencil. It is not clear whether they belong to (3) or to (4).}.) (4) Mere emotion is not revenge. (5) Though we must put aside
the injury, we must not put aside the injurer, for that would be
hatred—this the \emph{cardo} of the difficulty of the precept. (6)
On being obliged to \emph{speak} to persons with whom we have
quarrelled. This has \emph{exceptions}, \emph{e.g}. if they are likely
to tempt us to sin, which perhaps \emph{was} the injury; but such
exceptions must be determined by a \emph{director}. (7) It seems to be
contrary to justice if injuries are not punished. This is true, but we
must not judge in our own case. (8) Contrary to nature to forgive.
Yes, but sin and redemption (see above, 2). (9) This is what this age \emph{forgets}
when it speaks in favour of revenge. (10) Men do not believe in
redemption, nor that they are sinners. Hence Mahomet. (11) Do I
put forgiveness [merely] as a \emph{condition} [of obtaining
forgiveness for ourselves]? No, one who believes in what Christ has
done has no heart for revenge. (12) Onesimus—Christ says 'forgive
me' by the lips of the fellow-servant. (13) Man's duty to pray for
injurers. (14) Pray to meet them in heaven, (15) when all angularities
will be rubbed off, and we shall be able truly to love them. (16) We
and they are sinners; let us help each other.

\xsection{April 5, 1863 (Easter Day) [Silent Joy]}

'I sat down under His shadow whom I desired, and
His fruit was sweet to my mouth (\emph{gutturi}).' \footnote{Cant. ii. 3.} So says the Spouse in the Canticle.

Aggeus ii. 8, 'The desired of all nations shall
come, and I will fill this house with glory.'

Malachias iii. 1, 'And presently the Lord whom
you seek, and the angel of the testament whom you desire, shall come
to his temple.'

And therefore the Bride sat down, \emph{as Mary at
His feet}. And so it is, whether in great joy or sorrow, we are
silent. Each emotion, when profound, produces a calmness. Thus Job's
friends, Job ii. 13, 'And they sat with him on the ground seven days
and seven nights, and no man spoke to him a word: for they saw that
his grief was very great.'

Thus in Christ's death and resurrection.

Look in the gospel and you will see that it is no
peculiarity of ours, of our race, of these times, but it is a
deep characteristic of our nature. Who is it that speaks in the
gospel, from the time when Mary poured the alabaster box to the time
when Jesus ascended into heaven? \footnote{The rest of the sermon is
from a slip of paper pasted in the book, apparently the notes which
the preacher took with him into the pulpit.}

And so as regards our Lord—read the whole
account—it is not the disciples who speak, it is Jesus—Jesus in
the supper, \emph{the Last Supper}, the garden, the Passion, the
Cross, the Resurrection.

It is the wicked who speak—they speak who speak
to sin—Chronista, Christus, Synagoga \footnote{The three deacons of the
Passion on Palm Sunday.}—all [speakers except Christ] grouped under that word [Synagoga].
Judas, chief priest, false witnesses, Pilate, the multitude.

Exceptions: Pilate's wife, centurion.

Even apostles—as Peter who denied Him.

Two speakers alone—( 1) when proper? (2) who is
it? \footnote{The meaning of this is not
clear.}

Thus fulfilled, Psalms xiii. 3 \footnote{'Their throat is an open
sepulchre: with their tongues they acted deceitfully; the poison of
asps is under their lips. Their mouth is full of cursing and
bitterness.'}.

Nay, a time came when even the wicked were
silent—'that every mouth should be stopped'—and the Lord alone
spoke. He spoke His seven last words amid the silence: only one
exception—the penitent thief.

Thoughts good and bad ... His resurrection. 'My
Lord and my God.' 'Lord, what shall this man do?' And Acts i.

\xsection{April 12 (Low Sunday) [Faith]}

1. INTROD.—In
today's epistle and gospel which I have just read, we have brought
before us what is one of the great lessons of this sacred season, viz.
the necessity of \emph{faith} as the foundation of the Christian life.
Last week we considered the Passion and Resurrection, today the faith
by which we receive them as by a channel.

2. Contrast faith and reason. No great work done
by mere reason, even in this world.

Not that they are opposed, but faith has the
power of anticipating, and arrives at first at what reason scarcely
guesses at at last—St. Peter and St. John in Keble's poem \footnote{See Note
15, p. 341.}.

3. Therefore, since Almighty God works by human
means, He chose faith as the faculty which does great things.

4. Natural and supernatural faith. It is often
what we mean by genius in Nature, which sees what others see not, etc.
We believe in the existence of God, though it can be proved also; and
so of Christ, etc.

5. Hence the multitudes converted, etc., etc.

6. Quote passages from epistle 1 John v. 4-9 \footnote{'And this is our victory
which overcometh the world, our faith. Who is he that overcometh the
world, but he that believeth that Jesus is the Son of God? ... If we
receive the testimony of men, the testimony of God is greater.'} Cf. 1 Peter ii., then the gospel of yesterday [John xx. 1-9;
St. Peter and St. John coming to the Sepulchre]; then gospel of
the day, John xx. 19-31 \footnote{'Jesus saith to him,
because thou hast seen me, Thomas, thou hast believed: blessed are
they that have not seen, and have believed.'},
in which faith and reason are contrasted and the superiority given to
faith.

7. And lest faith should be confused with
enthusiasm, general notes are given to steer us, viz. the very effect
of faith, 'Who is he that overcometh the world?'—and so, 'Who is he
that overcometh the flesh?' 'Who is he that overcometh the devil?'
etc.

8. This \emph{first, before} faith we see on a
large scale in the world.

9. After faith we see the effects in our
heart—the justifying by reason what we have done by faith.

\xsection{January 3, 1864 St. John the Saint of the Time, of the New Year, etc.— Old Year Ending, New Year Beginning, etc.}

1. 'Canst thou drink of the chalice?' \footnote{'Then came to him the
mother of the sons of Zebedee with her sons adoring, and asking
something of him. Who said to her, What wilt thou? She saith to him,
That these my two sons may sit the one on thy right hand, the other on
thy left hand, in thy kingdom. And Jesus answering said, You know not
what you ask. Can you drink the chalice that I shall drink? They say,
We can. And he saith to them, My chalice indeed you shall drink: but
to sit on my right or left hand is not mine to give, but to them for
whom it is prepared by my Father.'—Matt. xx. 20-23.} etc. 'What shall this man do?' \footnote{'Peter, turning about,
saw that disciple whom Jesus loved following; who also lent on his
breast at supper, and said, Lord, who is it that shall betray thee?
Him therefore when Peter had seen he saith to Jesus, Lord, what shall
this man do? Jesus saith to him, So I will have him remain till I
come, what is it to thee? follow thou me.'—John xxi. 20-22.}

You see how little they knew, as we. The prospect
quite dark.

2. So sanguine and eager to serve his Lord, yet
how differently from what he thought! The throne on right hand and
left were deferred till the next world. So we, and we shall find at
the end that God is faithful and gives us our wish, but how
differently from what we expect! 'Commit thyself to God, and He will
give thee the desires of thy heart,' yet in His way, not ours. Thus
Jacob who said 'few and evil,' \footnote{'The days of my
pilgrimage: ... few and evil.'—Gen. xlvii. 9.} yet speaks of the Lord who had been with him from the
beginning to this day, etc. \footnote{'The God that feedeth me
from my youth until this day.'—Gen. xlviii. 15.}
So Solomon naming God's mercies, 'Thou hast not failed,' yet how
differently!

3. He is the saint of the longest lived; he
covers all length of life. He is the saint of the young, the
middle-aged and the old. Hence the appropriate addresses, 'Little
children,' etc.: 1 John ii. 1, 'My little children, these things I
write'; \emph{ib}. 18, 'Little children, it is the last hour'; \emph{ib}.
iii. 7, 'Little children, let no man deceive you'; \emph{ib}. 18, 'My
little children, let us not love in word, nor in tongue; but in deed
and in truth,' etc. And what is his experience? 'The world lieth in
wickedness.' 1 John ii. 15-16, 'Love not the world, nor the things
which are in the world. If any man love the world, the charity of the
Father is not in him. For all that is in the world is the
concupiscence of flesh, and the concupiscence of the eyes, and the
pride of life.'

\xsection{July 24 (Tenth Pentecost) [The Pharisee and the Publican]}

1. INTROD.—'God
be merciful to me a sinner.'

2. In these words is contained the essence of
true religion.

3. \emph{Why}? Because they refer to conscience
as leading the mind \emph{to God}.

4. All men have a conscience of right and wrong,
Rom. ii. 14.15 \footnote{'For when the Gentiles,
who have not the law, do by nature those things that are of the law,
these, having not the law, are a law to themselves: who shew the work
of the law written in their hearts, their conscience bearing witness
to them, and their thoughts between themselves accusing or also
defending one another.'}.—\emph{the
conscience accusing}, etc. But it does not lead them, when they
transgress it, to \emph{God}. They are angry with themselves. They
know they are wrong, and are distressed, but it does not lead them to
religion; at the utmost it leads them to understand a sin against
their neighbours—as cruelty, etc. But when it leads the soul to
think of God, then that soul may be very sinful, but at least it has
something \emph{of true religion} in it.

5. 2 Cor. vii. 10-11 \footnote{'For the sorrow which is
according to God worketh penance steadfast unto salvation: but the
sorrow of the world worketh death. For behold the selfsame thing, that
you were made sorrowful according to God, how great carefulness it
worketh in you, yea defence, yea indignation, yea fear, yea desire,
yea zeal, yea revenge! In all things you have shewed yourselves
undefiled in the matter.'}. And so 'to Thee then only have I sinned, and done evil before
Thee.' Ps. 1.

6. Hence in the text the reason why the publican
was \emph{more justified}, because he understood that his offenses
were against God.

7. But see what comes from this. Directly a man
realises that what he does wrong is against God, then he feels how
much more extensive it is, viz. of the thoughts.

8. And how much more intensive, viz. as against
the Highest. He calls it \emph{sin}.

9. Then he grows in his notions. As blows don't
pain at first, so sin may pain hereafter.

10. Thus he sees it is an offence against the
moral nature of God.

11. Hence all diseases are but types of sin.

12. Hence idea of guilt.

13. Hence need of a cleansing.

\xsection{July 31 (Eleventh Pentecost) [On the Gospel of the Sunday— The Healing of the Man Deaf and Dumb]}

Various maladies which our Lord cured, typical of
various sins.

1. Blind.—Those that have not faith, and do not
apprehend doctrine.

2. Deaf.—Those that are without devotion and
cannot hear the songs of angels.

3. Dumb.—Those who through cowardice or pride
do not confess the Gospel, though they believe in it.

4. Without taste (and smell).—Those who have a
dull, unsensitive conscience.

5. Lame.—Those who are slothful.

\xsection{August 7 (Twelfth Pentecost) [Love of God]}

By contrast—love of God. Luke x. 27, 'Thou
shalt love the Lord thy God with thy whole heart, and with thy whole
soul, and with all thy strength, and with all thy mind.'

Love of God—emotion not necessary; straw burns
out quicker than iron.

1. Desire of His exaltation.

2. For His own sake, not ours: (1) \emph{complacentia},
(2) \emph{ad majorem} [\emph{Dei}] \emph{gloriam}.

3. Yet with our own personal interest in it, as a
mother or sister follows the history of a son or brother with
sympathy, though without \emph{personal} gain, etc.

Loving God for His own sake. None can be saved
without love.

\xsection{August 7 The One Sacrifice}

1. Christ a sacrifice. We keep the feast at
Easter: 'Christ our Pasch,' etc.

2. What is meant by sacrifice?—offering,
killing, eating. Objects—(1) worship, (2) thanksgiving, (3)
propitiation, (4) impetration.

3. Heathen sacrifices, 1 Cor. x. 20, 'The things
which the heathen sacrifice, they sacrifice to devils, and not to
God.' \emph{Taurobolium}, etc.; human sacrifices, etc.

4. Jewish sacrifices, Lev. ix., x., xvi.

5. Fulfilled in Christ. Two things—death and
intercession. Therefore the Aaronic priesthood and bloody sacrifice
fulfilled in one act and time by our Lord in the flesh. Scripture
speaks of our Lord as not only fulfilling a bloody sacrifice (which is
once for all), but besides this of a sacrifice according to
Melchisedech, that is, bread and wine.

6. Superiority—once for all; all sin.

7. His mercy in continuing the sacrifice, that it
might not be a mere matter of history.

8. This is the Mass; which is His sacrifice
reiterated, represented, applied, as He continues it in heaven. Mal. i.
11 \footnote{'For from the rising of
the sun even to the going down, my name is great among the Gentiles,
and in every place there is sacrifice, and there is offered to my name
a clean oblation: for my name is great among the Gentiles, saith the
Lord of hosts.'}.

9. Order of Melchisedech.

10. One priest, one victim, one sacrificing,
everywhere.

11. From the first—various liturgies—points
in common. Detail of rites as in the bloody [sacrifice on Mount
Calvary]—as our Lord taken before Pilate, etc., etc.

12. Arguments for apostolicity of the Mass.

\xsection{August 14 (Thirteenth Pentecost) [Love of God]}

1. I\textsc{ntrod}.—I said last
week that no one can be saved without \emph{love of God}. This the
awful truth.

2. In fact this is plain, but considering the
state of the case—the immortal soul—how tired it will get of
everything in eternity, except of something which is infinite. God in
Himself a world; His attributes infinite.

3. \emph{Yet how can we love Him}? See how much
against our nature it is. We take delight in things of the world,
etc., etc., in science, in literature, etc. These are our aims; but to
love God is \emph{an aim} above our nature.

4. Granted \emph{it is so}. However, God does not
command impossibilities.

5. Therefore He gives us grace to raise us above
our nature. Even angels need grace.

6. What is grace? and what does it do for us?

7. Let us pray God for it.

\xsection{August 21 (Fourteenth Pentecost) [The Life of Grace]}

1. Nothing is more common than to think that
natural virtue, what we do by nature, is sufficient for our salvation.

The state of most men is sin, but as to those who
go the highest [it is] natural virtue. Put it in the way of an
objection. Why is not this enough? Two things confused with each
other—the improvement of things in this world, which natural virtue
can do, and the salvation of the soul by grace.

2. What most men consider enough is this—if
they follow what they think right, if they do the duties of
their station, if they do what their conscience tells them, and so
live and die. As to prayer, the best prayer is to do their duty here;
they think the next world may take its chance.

3. Now most men do not get so far as this. They
live in sin; but the utmost they think of is to be saved mainly by
their own strength, and by doing the common duties of life without
thinking of religion, though they may acknowledge that on great
occasions God helps them, within or without, but is it when \emph{dignus
vindice nodus} \footnote{'Nec Deus intersit, nisi
dignus vindice nodus, Inciderit.' Horace, \emph{Ars Poet}., 191-2.}.
They do not see the necessity of thinking of God, but they say that
the best service is to do those duties which come before them.

4. Particularly at this day. When men think that
religion is unnecessary, that the world will advance merely by its own
powers.

5. On the other hand, \emph{the life of grace}—virtues
through grace.

6. Natural virtues bring on the world—doubt1ess
social science, political economy, science of government, etc.,
etc.—but \emph{I want to be saved}.

\xsection{September 18 [The Mass]}

The Mass is to be viewed in two aspects—(1) as
it regards our Lord; (2) as it regards us. As it regards Him, it is
the great act of sacrificial atonement. As regards us, the great act
of intercession. Texts, Rom. viii. 32 \footnote{'He that spared not even
His own Son, but delivered him up for us all, how hath he not also
with him given us all things?'}; Heb. vii. 22-25 \footnote{'By so much is Jesus
made a surety of a better testament. And others indeed were made many
priests, because by reason of death they were not suffered to
continue: But this, for that he continueth for ever, hath an
everlasting priesthood, whereby he is able to save for ever them that
come to God by him: always living to make intercession for us.'},
ix. 15 \footnote{'And therefore he is the
mediator of the new testament, that by means of his death, for the
redemption of those transgressions which were under the former
testament, they that are called may receive the promise of eternal
inheritance.'} and \emph{passim};
1 Tim. ii. 5-6 \footnote{'For there is one God,
and one mediator of God and men, the man Christ Jesus; who gave
himself a redemption for all.' ...}.

1. He is the great High Priest who is ever
offering up His meritorious sacrifice, and the Mass is but the earthly
presence of it.

2. While He offers it above, the whole Church
intercedes. (1) Mary on high, and the saints with her. Thus a heavenly
Mass is now going on above. (2) Below—not a light benefit that we
may intercede.

We have indeed a hope within us that God will
hear us for ourselves, but will He hear us for others? It is only
through His wonderful meritorious sacrifice that we have this power,
and therefore fitly in the Mass is the intercessory gift exercised.
Therefore the very privilege of Catholics above others is intercessory
prayer; it is the imputation and the imparting to their prayers the
merit of the sacrifice. Therefore St. Paul says, 'Pray without
ceasing.' St. James, etc. All intercessory prayer all over the whole
world, \emph{e.g}. litanies, the priest's office, the breviary,
is as it were in presence of the Mass. It is the great act of
communion, etc.

\xsection{January 1, 1865 [Eternity]}

1. All days are the beginning of new years, but
we have especially reason to place the first day in this time, for the
season in which it comes is the beginning of a new year, because it is
the beginning of a new revolution in this world's course. The earth is
asleep, and I may say dead; and as man's extremity was God's
opportunity, when things are at their worst they begin to mend. The
sun stays in its downward course—it turns back, the days become
longer, etc., etc. The year awakens, and human thought and activity
with it—the farmer because of the ground; the navigator looks for
favourable weather and the right wind; the warrior opens his campaign;
parliaments, etc., etc.

2. MOTION.—Such
is this wonderful world, in which all is motion—begins, goes on,
increases, and dies again, year after year, and man in detail, day
after day, goes on to his work and his labour till the evening.

3. CHANGE.—Such
it is with us, and with an end. What does it end in? We pass in the
course of 365 days the day of our death—like walking over our
gravestone. What does it end in?—a state in which time ceases, or
rather time, it may be said, stops. Time in this world is marked by
motion. Motion, or what is commonly called change, is the very
fulfilment of this state of things.

4. END OF
CHANGE.—But the day will come
when time brings with it no changes—(past, present, and future
because [there is] change)—when all is the same—day after day, age
after age—in short, when time stops—\emph{an eternal now}. This we
call eternity.

5. TIME WITHOUT
CHANGE IS
ETERNITY.—Properly time cannot
stop; it runs on as I am speaking. There is nothing to end it; but as
soon as there is no change in it, it is eternity. All our thoughts,
ideas, etc., will stop: they will be fixed and one and the same. As
they are good or bad, it will be heaven or hell.

\xsection{March 5 (First Lent) [Sin]}

1. INTROD.—At
this time of year, 'Come let us reason together; argue with me, saith
the Lord.'

EXPOSTULATION.—Isa.
i. 2 \footnote{'Hear, O ye heavens, and
give ear, O earth: for the Lord hath spoken, I have brought up
children and exalted them, but they have despised me.'} and xliii. 21-26
\footnote{'This people have I
formed for myself; they shall shew forth my praise. But thou hast not
called upon me; neither heat thou laboured about me, O Israel ... I am
he that blot out thy iniquities for my own sake, and I will not
remember thy sins. Put me in remembrance, and let us plead together:
tell me if thou hast any thing to justify thyself.'}, Mic. vi. 1-2 \footnote{'Hear ye what the Lord
saith; Arise, contend thou in judgment against the mountains, and let
the hills hear thy voice. Let the mountains hear the judgment of the
Lord, and the strong foundations of the earth: for the Lord will enter
into judgment with his people, and he will plead against Israel.'}.

God, most blessed from eternity, created us, not
for any good that we could do to Him. He would not be happier,
stronger, etc., by creating us. On the other hand we are wholly
dependent on Him. The axe does not depend on the carpenter for
beginning of [existence], nor the son on the father for continuance of
life, but God made the dust, out of which we are, out of nothing, etc.
He sustains us, etc. We are entirely His work and property, and should
do \emph{His service}. Yet we have \emph{cast off His yoke}.

2. But again, He made us in order \emph{to bless}
us. He knows of what we are made. He knows what will make us happy.
Yet we have refused to be blessed; we have sought our own happiness.

3. Two claims—duty and interest. Let us
confess. We have preferred to be our own masters; we have refused to
believe that sin is an evil. We will not believe what an evil sin is;
we have no loathing or horror of it.

4. But now consider what sin is. God is infinite.
It is the one thing which may be said to be of an infinite \emph{nature}
besides God. It is inexhaustible, irremediable; it is greater than
angel or archangel, a rival infinity to God—'against thee only have
I sinned.' According to the person injuring [injured?], so is the
injury, \emph{e.g}. insulting a superior. Sin is the lifting up the
hand against the infinite benefactor.

5. He will leave me to myself. What will become
of me?

6. Save me from myself.

\xsection{December 2, 1866 [Omniscience of God]}

1. INTROD.—Omniscience
and omnipresence of God—\emph{knowing the heart}; incomprehensible;
millions of men, yet He knows all that goes on in the heart [of
each one] and \emph{remembers}.

2. [Incomprehensible] yet familiar to children \footnote{See Note
16, p. 342.}.

3. Scripture—1 Kings xvi. 7 \footnote{'And the Lord said to
Samuel, Look not on his countenance, nor on the height of his stature;
because I have rejected him: nor do I judge according to the look of
man; for man seeth those things which appear, but the Lord beholdeth
the heart.'}, 1 Paralip. xxviii. 9 \footnote{'And thou, my son
Solomon, know the God of thy father, and serve him with a perfect
heart and a willing mind: for the Lord searcheth all hearts, and
understandeth all the thoughts of minds: if thou seek him, thou shalt
find him; but if thou forsake him, he shall cast thee off for ever.'}, 2 Paralip. vi. 30 \footnote{'Hear thou from heaven,
from thy high dwelling place, and forgive, and render to every one
according to his ways, which thou knowest him to have in his heart;
(for thou only knowest the hearts of the children of men).'},
Jeremias xvii. 10 \footnote{'I am the Lord who
search the heart, and prove the reins; who give to every one according
to his ways, and according to the fruit of his devices.'},
Apoc. ii. 23 \footnote{'I am he that searcheth
the reins and hearts: and I will give to every one of you according to
his works.'}. Future
judgment—Rom. xiv. 10 \footnote{'For we shall all stand
before the judgment seat of Christ. Therefore every one of us shall
render an account to God of himself.'},
1 Cor. iv. 4-5 \footnote{'For I am not conscious
to myself of anything; yet am I not thereby justified: but he that
judgeth me is the Lord. Therefore judge not before the time, until the
Lord come, who will both bring to light the hidden things of darkness,
and will make manifest the counsels of the hearts.'}, Heb.
iv. 12-13 \footnote{'For the word of God is
living, and effectual, and more piercing than any twoedged sword, and
reaching unto the division of the soul and the spirit, of the joints
also and the marrow, and is a discerner of the thoughts and intents of
the heart. Neither is there any creature invisible in his sight: but
all things are naked and open to his eyes.'}.

4. Suitable to this time of year—the particular
and general judgment.

5. The keenness of the judgment—as above, Heb.
iv.—magnifying-glass, the wonders of the microscope, a new world,
diseases. Hence we must feel we do not know ourselves. \emph{Therefore}
1 Cor. i., 'judge nothing before the time.'

6. Most awful, but different way in which good
and bad take it.

7. The bad dread it. Adam and Eve in the garden.
'And then shall they say to the rocks, Fall upon us,' etc., etc., Luke
xxiii. [30] \footnote{'And they shall begin to
say to the mountains, Fall upon us; and to the hills, Cover us.'}, Apoc. vi.
[16] \footnote{'And they say to the
mountains and the rocks, Fall upon us, and hide us from him that
sitteth on the throne, and from the wrath of the Lamb.'}.

8. The good desire it—to be known to God, Ps.
cxxxviii \footnote{'\emph{Domine, probasti me},
etc.}.
Purgatory—willing victims.

9. This is one test whether we can bare our
hearts before God.

\xsection{March 20, 1870 On the Gospel of Third Lent}

1. INTROD.—The
diseases which our Lord cured were typical of sins. The dumb spirit,
who is he? One who will not go to confession, or who cannot, who has
not the opportunity. I wish I could describe him and his misery.

2. Time was, before the Gospel, there was no
personal individual confession. It is one of the great gifts of Christ's
coming.

3. Christ came to fulfil all the needs of
man—to give him hope, peace, strength, joy, and all virtues and
blessings. Now let us see what is one special need of his nature.

4. Man is a social being. The instrument of
society is the great gift of speech.

\emph{Begin thus} [\emph{2nd scheme}].

\begin{quote}

(1) Man is a social being.

(2) Speech the great tie and bond of society; dumbness and deafness
generally go together. It is said that blind men are more cheerful
than deaf and dumb, because society is \emph{a truer world than the
physical}.

(3) No man is sufficient for himself—the voice is an outlet. No
greater misery than to be shut up in oneself—\emph{speech} is the \emph{great
relief}. How dull it is to see beautiful things without
companions to speak to. We must say all that is in our heart. As the
pleasant things, so also the painful. Difficulty of keeping a
secret, or of not speaking to others when we have been ill-treated.

(4) Nay, Almighty God not by Himself, but with His Son and Spirit.
From eternity love, and not power.

\end{quote}

5. The devil alone is solitary—and evil
spirits—this the worst misery of hell.

\emph{Begin thus} [\emph{3rd scheme}].

\begin{quote}

(1) From this dumbness we may gain a great
spiritual lesson.

(2) Social nature. Whatever we feel we bring out. Praise and prayer.

(3) And so all angels. One society in heaven. Praise and prayer.

(4) And so God Himself.

(5) Evil spirits and evil men on the contrary.

(6) On all our affections and passions relieved by words.

(7) Keeping secrets, etc.

(8) Confession one kind of speech.

(9) Those who from want of opportunity, from pride, from despair, do
not confess.

(10) Comfort of confession.

(11) Those who don't are like the evil spirits.

(12) Happy all Catholics, if they knew their happiness.

\end{quote}

\xsection{April 3 Gospel for Passion Sunday}

1. INTROD.—We
veil our crosses. On the various gospels descriptive of our Lord's
hiding Himself.

2. He hid Himself from the Jews because they had
refused the light.

3. He is the light of the world and the light of
the soul.

4. Abraham had first 'seen' Him—on Moriah—and
the other prophets, as if mounted on high. And all the Jews, though
they had not seen Him, had heard of Him and expected Him. He was the
'expectation of the nations.'

5. At length 'He came unto His own,' etc. 'The
light shineth in the darkness.'

6. A warning to all of us lest we receive the
grace of God in vain. A \emph{yearly} warning.

7. We cannot be as others. We have had great
opportunities. We mix with Protestants. They have their own
views. They argue and conclude on their own basis. They are sharp and
clever men of business; good politicians; on \emph{their own}
principles right. No wonder they think so differently, for the great
bulk of them have not seen what we have seen. But Luke x. 23-24 \footnote{'Blessed are the eyes
which see the things which you see: for I say to you, that many
prophets and kings have desired to see the things which you see, and
have not seen them; and to hear the things which you hear, and have
not heard them.}.

8. O let us beware lest we ever get blinded. Isa.
vi. 9-10 \footnote{'Go, and thou shalt say
to this people, Hearing, hear and understand not; and see the vision,
and know it not. Blind the heart of this people, and make their ears
heavy, and shut their eyes; lest they see with their eyes, and hear
with their ears, and understand with their heart, and be converted,
and I heal them.'}.

9. 'Strive to enter the strait gate: for many, I
say to you, shall seek to enter, and shall not be able,' Luke xiii.
24.

\xsection{April 24 (Low Sunday) [Faith]}

1. INTROD.—Hope—Christmas;
love—Pentecost; faith—Easter.

2.Because there was at first so much doubt, etc.
(1) The first blow was that our Lord should die—this seemed
impossible. (2) That He should rise from the dead.

3. Hence on Low Sunday epistle and gospel.

Now under these circumstances it seemed
reasonable that our Lord should give them the testimony of
sight, touch, etc., for, unless some one saw Him again, how were the
apostles, how was the world to know it.

4. But a deeper lesson. Sight could not be given
to all, because our Lord was going to heaven, and those who did not
see must believe on the witness of others. Now the Gospel was to last
to the end of the world. Therefore He in His love determined that one
of the apostles should be away and not see Him.

5. This was Thomas, who, being in the state of
confusion which they all were in before they saw Him, persisted in
that unbelief which at first they all had. When the women testified,
the apostles would not believe. When the apostles testified, Thomas
would not believe.

6. We all know what happened. Our Lord graciously
granted, etc., but He said: 'Blessed are they that have not seen and
have believed.'

7. This is one lesson. Our Lord speaks to us.
Thomas thought it hard he had not the evidence the rest had. \emph{He was
not content with what was sufficient}. This \emph{the great lesson}.
Doubtless sight is more than the witness of other men.

8. Let us take the Gospel of St. John. There are
miracles more wonderful than in the other gospels, \emph{i.e}. those
addressed to the intellect, not the imagination, etc.; and he
testifies to the truth, and so do the Christians around him. John xxi.
24, 'This is that disciple who giveth testimony of these things, and
hath written these things that we may know that his testimony is
true.' The early Christians had no greater evidence than we have,
but they believed it more vigorously; hence they went through so
much.

\xsection{May 8 Patronage of St. Joseph}

1. INTROD.—'Yet
a little time, and ye will not see me,' etc.

2. That time when Christ came to each apostle,
was at the death of each. He says the time was short of this
life—and though He was going they would go soon; and He only went
before them to prepare a place.

3. Yet though their life was short, how long it
seemed by being so full of suffering.

4. St. Paul's sufferings (though greater,
perhaps): 2 Cor. xi. 24-28 \footnote{'Of the Jews five times
did I receive forty stripes save one. Thrice was I beaten with rods,
once was I stoned, thrice I suffered shipwreck, a night and a day I
was in the depth of the sea; In journeying often, in perils of waters,
in perils of robbers, in perils from my own nation, in perils from the
Gentiles, in perils in the city, in perils in the wilderness, in
perils in the sea, in perils from false brethren; In labour and
painfulness, in much watchings, in hunger and thirst, in fastings
often, in cold and nakedness. Besides these things that are without,
my daily instance, the solicitude for all the churches.'},
Acts xx. 22-23 \footnote{'And now, behold, being
bound in spirit, I go to Jerusalem, not knowing the things that shall
befall me there: Save that the Holy Ghost in every city witnesseth to
me, saying that bonds and afflictions wait for me at Jerusalem.'}, 1 Cor.
iv. 11-13 \footnote{'Even unto this hour we
both hunger, and thirst, and are naked, and are buffeted, and have no
fixed abode; And we labour, working with our hands: we are reviled and
we bless; we are persecuted and we suffer it: We are blasphemed and we
intreat: we are made as the refuse of the world, the offscouring of
all until now.'}.

5. And so of all Christians then. They were tried
with long unsettlement and uncertainty—their lives in their hands;
persecution any day—inscription in the catacombs—'O wretched we,'
etc.—'If in this life only.' \footnote{'If in this life only we
have hope in Christ, we are of all men most miserable.'—1 Cor. xv.
19.}

6. And so of all the saints—confessors,
ascetics, etc., etc.—they are all in trouble; and when we think of
them we think of pain, penance, etc.

7. This is what supported them, hope, viz. that
Christ comes again, that their sufferings would end with this life;
that they would be rewarded by being with Him.

8. Hence heaven was their \emph{patria}, their HOME—Family,
Father, peace—all was trouble here.

9. There is but one saint who typifies to us the
next world, and that is St. Joseph. He is the type of rest, repose,
peace. He is the saint and patron of home, in death as well as in
life.

10. Let us put ourselves under his protection.

\xsection{August 7 (Ninth Pentecost) [The Omnipotence of God and Man's Free-Will]}

1. God is almighty, but still this does not mean
that He can do everything whatever, for if so He could do
contradictions. There are some things, of course, which are impossible
to Him because the very thought of them is an absurdity, \emph{e.g}.
He can never cease to be holy; He can never wish to cease to be holy,
etc., etc.

2. And so again, much more when He created, He Himself, as it were, put obstacles in the way of the exercise of
His omnipotence—things which once were possible ceased to be
possible. He made a sort of covenant with creation in creating. He
forthwith made Himself a minister to His creation, which could not
stand of itself.

3. And much more when He created rational beings,
who can exercise a will of their own, and do right or wrong, He can't
do what He would. We say, 'Thy will be done.' It is difficult to
conceive how. When God had once created a being who could do right and
wrong, He suspended His own prerogative of 'His will being done.'

4. Especially when He makes a covenant, for then
He is bound by its terms. And further, such beings bring His
attributes into operation and they seem to contradict each other—as
justice and love.

5. I come to this conclusion: that men who rely
on the boundless mercy of God do not understand how the matter stands.
He has other attributes, and they act according to the case—'Let me
alone,' \footnote{'Let me alone, that my
wrath may be kindled against them … And Moses besought the Lord,'
etc.—Exod. xxxii. 10-11.} power of
intercession. God chose the Jews, etc. They are an example of what I
mean. He willed their salvation. He did all things He could \footnote{This 'could' is a
preacher's word; it must not be theologically pressed. See Note
17, p.
342.} for them, and He cast them off. How awful is this—His will
was not done! By creating beings who could have a will of their own,
He circumscribed His own power.

6. Now I am led to these thoughts by the epistle
and gospel of the day.

7. Go through the epistle \footnote{1 Cor. x. 6-13.} and gospel \footnote{Luke xix. 41-47.}.
Isa. vi. 10, 'Blind the heart of this people, and make their ears
heavy, and shut their eyes; lest they see with their eyes, and hear
with their ears, and understand with their heart, and be converted,
and I heal them.' Matt. xiii. 14 [above passage quoted], John xii. 40
[same passage quoted], Acts xxviii. 26 [same quotation].

8. Our last [end]—our 'enemies may come about
us,' etc. God forbid.

\xsection{August 28 (Twelfth Pentecost) [God Our Stay in Eternity]}

1. The lawyer asks, 'What shall I do?' etc. Our
Lord refers him to the duty of loving God.

2. Now that this is our plain duty is clear. It
is the \emph{condition} of heaven. But it is more than that—I wish
you to see that it is the nature of things.

3. We all wish to live a long life. We are all
fears and awe at death. Why? Well, partly because it is the loss of
life; but more, because it is leaving what we know and going to what
we do not know, and for judgment. This world our \emph{home}—it is
going into a strange country. This is the main reason.

4. Men (especially Protestants) talk vaguely of
going to glory, etc. Now let us contemplate what it is, going into a
strange country. What will be our happiness there? Let us look at it
in a common-sense way. What is to constitute our happiness? What is to
occupy us in eternity? Why, even of \emph{this} world men get tired!
You hear of old people who are ready to die, not because they
like death but because they are tired. Now if many men are tired of
eighty years, supposing they were to live on till two hundred, would
not they be \emph{tired} then? The world, too, would get more and more
strange to them—solitary. Much more eternity.

5. Some men say, 'We shall see the wonders of the
universe'—curiosity gratified. And they will take us a long time
certainly, and memory will fail, so that we may begin again when we
forget. But how soon we get tired of sight-seeing! We long to get
home.

6. Home, that is it; what is our home?

7. God and the love of God.

8. Thus \emph{necessitate medii}.

\xsection{Christmas Day [The Advent of Christ Foretold]}

1. INTROD.—There
are many subjects in which we have nothing in common and cannot
sympathise with each other. But if there is a day which puts us, high
and low, rich and poor, on a level, it is this. Angels at Nativity,
Resurrection, and Ascension, are above us even in their nature and
speech.

2. Fall—the Evil One getting usurped possession
of the earth.

3. But deliverer promised from the first, and
even the time of His coming, though long after, determined.

4. Therefore expectation of freedom all through
the East, and in the West.

5. No event thus known beforehand. God's
providences in the natural world are generally sudden—a great man
arises accidentally; great discoveries—and wars, as \emph{the present}
\footnote{The Franco-German War.}. How sudden—and so
end of the world.

6. But this contrasted to them. It was as well
known beforehand as many of the calculations of science, like the
eclipse we had a few days ago—all upon deep principles of law.

7. And so now that He is come, though the time of
His second coming is not determined as the first, let us be sure that
all is decreed, and goes upon fixed laws in its season, though His
coming is put off again and again, and we are deceived. Thus in the
physical, terrestrial world all is confusion at first sight—the
earth rises and falls, water rushes in, the face of the land changes,
but all on law.

8. So, whether the temporal power is established
or falls—

9. Only let us be ready for His coming.

\xsection{June, 1871 (Trinity Sunday) [Mysteries]}

1. INTROD.—Our
happiness consists in loving God. And we cannot love Him without
knowing about Him. And we cannot know about Him, ever so little,
without seeing that He is beyond our understanding, \emph{i.e}.
mysterious. These are thoughts for today.

2. What do we mean when we speak of God? The Creator. Well, how could He make all things out of nothing?

3. Or again, our Judge, who speaks in our \emph{conscience};
and yet, how can He read our heart?

4. Or again, Providence. Yet how can He, in spite
of the laws of Nature, and the separate wills of ten thousand minds,
turn everything to good for each of us?

5. Union of justice and sanctity with mercy;
power with skill.

6. Thus to be religious at all, to know and
believe anything of God, we must believe what we cannot understand, \emph{i.e}.
mysteries. It is as our Creator, Judge, Providence, having being, and
upholding good that we love Him.

7. And so of revealed religion. The Holy Trinity,
the Incarnation, why are these revealed?—to give us reason for
loving God.

8. Show how they lead us to love God.

9. But why need we love God?—because we are to
live after death. And then, where shall we find ourselves if we have
not love of God?

10. Those things here—(1) sensible comforts,
(2) activity, (3) affections.—where are they then?

\xsection{July 2 (At St. Peter's) [The Visible Temple]}

'Whether you eat, or drink, or whatsoever you do,
do all to the glory of God,' 1 Cor. x. 31.

1. What do these words mean? What do they enjoin
upon us?

2. We have our duty towards God, our neighbour,
and ourselves. Now we may in a certain way fulfil these duties without
doing them to God's glory.

3. \emph{E.g}. we may do our duty to God from
mere fear, or from habit, or from human respect; from expedience, \emph{e.g}.
going to Communion once a year, saying prayers, keeping from
particular sins—being respectable—this right, but not enough. To
our neighbour from pity, from benevolence, from family
affections—this too, good, but not enough. And so to ourselves. We
may be virtuous, and proud or self-conceited. That is, we may do
things good, and in a certain sense be good in doing them, \emph{yet not
to the glory of God}, \emph{i.e}. because \emph{not from love}.
This is one thing, then, that is meant by the text.

4. Then again, what is meant by doing \emph{all
things}? We have only rare opportunities of doing our duty. How can
we eat and drink to [the glory of God]?

5. (1) Eating and drinking. (2) Use of the
tongue—bad conversation. (3) Reading—curiosity. (4) Amusements in
kind and in reason \footnote{Perhaps 'season.'}.
(5) Work—idleness, justice. (6) Sickness. (7) Punishments and
penances.

6. Thus the whole day—'Pray without
ceasing'—Matt. v. 16 \footnote{'So let your light shine
before men, that they may see your good works, and glorify your Father
who is in heaven.'},
Phil. iv. 8 \footnote{'For the rest, brethren,
whatsoever things are true, whatsoever modest, whatsoever just,
whatsoever holy, whatsoever lovely, whatsoever of good fame; if there
be any virtue, if any praise of discipline, think on these things.'}.

And so especially the worship of God. God has
told us to pray. Now let us apply this to the service of God. To
pray together, and publicly. This implies, of course, rites of
religion, and buildings to perform them in. How can these be done to
God's glory? Now, I can understand men saying, 'No religious rites, no
common worship; religion is private and personal.' But I cannot
understand [them] saying, 'It is common and public, it has rites, it
has houses,' and not to bring those houses under the commandment of
glorifying God, being edifying, etc.

7. Now how do we glorify God in religious houses
or churches? In making them \emph{devotional}. No matter what
architecture, etc., \emph{devotional} is the end, towards God and
towards men.

8. And costly ('of that which cost me nothing,'
etc. \footnote{'Neither will I offer
burnt offerings unto the Lord my God which cost me nothing'—2 Kings
xxiv. 24.}) as a means of
expressing devotion—Aggeus i. \footnote{Where the people are
reproved for neglecting to build the Temple.}, Isa. lx. [13] \footnote{'The glory of Libanus
shall come to thee ... to beautify the place of my sanctuary; and I
will glorify the place of my feet.'},
Apoc. xxi. \footnote{Description of the New
Jerusalem.}. Hence
David, 1 Par. xvii. \footnote{David's purpose to build
the Temple.}—Ps.
cxxxi. (\emph{memento Domine}, David); 1 Par. xxix \footnote{David, by word and
example, encourageth the princes to contribute liberally to the
building of the Temple.}.—[his] zeal for the house of God; his singers, his psalms—i
Par. xxv \footnote{The number and divisions
of the musicians.}. This made
him \emph{according to His own heart} [1 Kings xiii. 14].

9. Now you know what this tends to. Why is it
that I come before you today? It is because I felt a profound
appreciation of the work in which your priest was engaged, and a
true sympathy in his exertions. I recognised in him a zeal for the
honour of God's house such as that of David, whose spirit was troubled
that his God had no abode fit for Him. I knew that for years and years
his spirit chafed within him that he could not perfect in this place
that idea of solemnity and beautifulness in the visible temple which
he had in his mind. Twenty years and more, to my knowledge, has this
idea occupied his mind. Then, too, he honoured me by asking me to take
here some part in promoting his work, which he has committed to me
now. Then he did a part—and now, by his persevering zeal, and the
munificence of pious men, he has been able to do more; and he urges
you, through me, to take part in, and to complete his service of zeal
and love. And in the next place he calls [you] to a religious act in a
religious way. He appeals to you on a Sunday, not on Monday, Tuesday,
etc. He has taken the legitimate ecclesiastical means of asking for
your contributions, which is possible on a Sunday. He does not take
means of raising money which are not possible on a Sunday; he does a
sacred work on a sacred day.

[Further], his object has special claims from the
circumstances of this church. It is the mother church of Birmingham.
It is dedicated to St. Peter. In subscribing to it you are testifying
your loyalty to the Holy See in its troubles. Lastly, on the feast of
the Visitation, when all Nature rejoices and Mary sings the \emph{Magnificat}—2
Cor. viii. 7 \footnote{'That, as in all things
you abound in faith, and word, and knowledge, and all carefulness;
moreover also in your charity towards us, so in this grace also you
may abound.'}.

\xsection{July 30 (Ninth Pentecost) The Jews—[Christ Weeping over Jerusalem—Luke xix. 41-47]}

1. Only one nation thus selected.

2. And that from its very root.

3. Two thousand years before our Lord, \emph{i.e}.
four thousand nearly from this time.

4. This people has had records, not traditions
only.

5. It is a specimen of God's governance, in the
midst of prevailing confusion, all over the earth.

6. Two cautions: (1) Children suffering for their
parents.

7. (2) Tower of Siloam.

8. Mercies, rejection of mercies, punishment.

\xsection{August 6 (Tenth Pentecost) The Divine Judgments}

I said last Sunday, 'Jews suffering for the sins
of fathers.' Is not this condemned by today's parable, in which
Christians will be behaving like the Pharisee? \emph{Answer}. (1) Not
by private decision. (2) Not individually, but nationally.

1. I mentioned last week the subject of the Jews,
but I could not continue without explaining clearly about judgments.
To continue, first there is judgment in the \emph{next} world—yes,
but in this also. In one sense \emph{all} suffering is a judgment of
sin—in one sense consequence of \emph{Adam's} sin; (1) individuals,
(2) nations.

2. But it does not therefore follow that \emph{we}
can say what are judgments and what not.

3. This is what religious men are very apt to do
by their private judgment. Irreligious men scout the idea [of divine
judgments].

4. Some indeed force themselves upon us, because \emph{all}
feel this, \emph{e.g}. (1) if a man were struck dead for lying, [or]
(2) if he committed sacrilege against the Holy Sacrament, stealing,
etc.

5. But (1) if in party matters, in which good men
are on both sides, if in political, he uses his private judgment, he
is wrong.

6. Yet how often this is done! a death, a
misfortune is interpreted our own way.

7. Again, (2) national judgments. First, this
does not show that the suffering nation is worse than others. Tower of
Siloam—Pharisee and publican.

8. Nay, nor that the people of that time are
worse, for they \emph{may} be suffering for the sins of the fathers.

9. Thus we come again to the Jews. They may be in
judicial blindness, but not by the fault of this generation.

10. They were taken without the merit of
individuals into covenant, and now they are put out. And since no one
can say, 'Jesus is the Lord but by the Holy Spirit,' therefore, as
Protestants are blind but without their fault, so the Jews.

\emph{Rather thus}:

INTROD.—The
Pharisee judged the publican. Thus I am led to the subject I touched
on last week. Instances: to say a man is wicked because
unfortunate—Job's friends; to say sudden death is a sign, etc.; to
take party or political views; to say nations are special
sinners who suffer—as France—Tower of Siloam.

Censoriousness is judging by our private opinion.
But certainly there are judgments. What is on record? What God says,
either by revelation, or the voice of mankind, or by the Church. \emph{E.g}.
a case of lying followed by death—for \emph{vox populi}, etc.; a
case of blasphemy or insulting the Blessed Sacrament. Then as to
nations, there is only one case revealed, the Jews, and even in that
case we do not judge individuals—(explain).

Charity thinketh not evil—quote 1 Cor. xiii.

\xsection{August 13 (Eleventh Pentecost) Continuing the Subject [of Divine Judgments]}

1. INTROD \footnote{Marginal note against
Introduction: 'This should be at the end.'}.—We have in the Old Testament, but we have nowhere else, an
unveiling of God's providence. It is not so now. The fortunes of the
Church and of the Holy See are not commented on now unerringly. Then
there were inspired prophets and inspired books: but there are none
such in these times. Thus Scripture is once for all.

2. Interpretation of the history of the
Israelites and the nations around, especially Israel—\emph{always
against their own will}, Ezek. xx. 32 \footnote{'Neither shall the
thought of your mind come to pass by, which you say, We will be as the
Gentiles, and as the families of the earth, to worship stocks and
stones.'}. No nation on earth has so great a history as the Jews; none
has so great a future.

3. Worldly prosperity does not go with true
religion now, as it did then. Then, in order to show \emph{that there was
a God}, He wrought in a special way—[also] \emph{in order to show
that He did work}; and it is our evidence of a Providence till the
end of time \footnote{There is a note which
was apparently intended for insertion in this paragraph, enumerating
the different kinds of inspired prophetical writings—history,
psalmody, ethics, predictions.}.

4. God has given us the greatest evidences in the
fact of the Jewish people.

5. Three great visitations: in Egypt; in Babylon,
taken from their land; and now in the world at large.

6. Moses' prophecy: 'Ye shall not be as the
nations.' They were \emph{so unwilling} to be a special people.

\xsection{August 20 (Twelfth Pentecost) [Divine Judgments Continued]}

1. INTROD.—Epistle
and gospel are on formality of Jews. This brings me to the subject
which I wish to continue.

2. Prophecy, if disobedient, idolatrous, to be
scattered, Lev. xxvi., Deut. iv.

3. This fulfilled in the first captivity and
dispersion—few returned, etc.

4. But return they did. And then a second and a
worse dispersion to this day.

5. Now why? For they boasted to \emph{keep} the
law; no idolatry.

6. It is clear that they must have committed a
grave fault; and it was this—they kept the law only in the letter,
not in the spirit.

7. 'Neighbour' in today's gospel, and so \emph{external}
purity, etc. 'I fast twice a week,' etc.

8. The prophets had warned them in vain. 'I will
have mercy and not sacrifice.'

9. Consequently they understood no part of the
true meaning of their Scriptures. As they made precepts formal, so
they made prophecies of the Messiah carnal.

10. And it ended in their being \emph{blind}, and
rejecting the Saviour when He came.

11. This then the great sin—greater than any
former—the crucifixion of our Lord.

12. 'His blood be upon us and on our children.'
So unto this day.

13. This then why they are without homes, without
honour, and without spiritual light.—from that curse which they
invoked upon themselves.

14. You will say that no one can really suffer
for the sins of others. True, they will be judged according to their
light. But the reason they have no more light is because their fathers
sinned.

15. Let us beware, for we at least can ruin
ourselves.

\xsection{August 27 (Thirteenth Pentecost) [Divine Judgments Continued]}

1. INTROD.—One
out of ten lepers returned thanks, and he a Samaritan; an election.
Thus reprobate were the Jews.

2. In consequence they cast off their Saviour,
and were in consequence cast off by God.

3. Now it will be observed they have from the
first been wanderers more than other people—in Abraham's time, in
Egypt, in the wilderness—but still unsettled. In \emph{unsettled times}
some stay necessary.

4. Hence \emph{their Temple} as a pledge, 2 Sam.
vii. 10 \footnote{'And I will appoint a
place for my people Israel, and I will plant them, and they shall
dwell therein, and shall be disturbed no more.'}. A pledge of
the \emph{gathering together} of the people.

5. Hence it was so beautiful, etc.

6. While it remained, they remained. When it
fell, they were scattered.

7. Hence in early times holy men believed and
predicted that it never could be rebuilt.

8. Hence Julian attempted to rebuild. Who Julian
was.

9. The more wonderful because it was the notion
of the Fathers that Antichrist will rebuild Jerusalem.

10. What happened.

11. They never will be able to rebuild this
temple till they get back into their land—never to get back till
they become Christians—and then it will be a Christian and not a
Jewish temple which they will build.

12. I end as I began when I spoke on this subject
first. It is a \emph{wonderful proof of the providence of God}. And He
will not desert His Church or frustrate His word now, though perhaps
not by miracle.

\xsection{September 17 (Sixteenth Pentecost) [Divine Judgments Continued]}

1. INTROD.—I
have lately been speaking of the wonderful history of the Jews, which
bears so much on the conviction which we have of the truth of
Christianity. We read in today's gospel of the Jews, and so
continually, and we know our Lord was a Jew, our Lady a Jewess, etc.
Yet how little do we know about them, etc., etc.

2. \emph{The Jewish history is the beginning of
Christianity and of its evidences}. The mustard seed. Abraham the
mustard seed, the father of the faithful. God has founded one church,
and that from the beginning. Slow, \emph{as geological formations}. As
we cherish a plant—in \emph{the hothouse}, etc.

3. It was the divine purpose that that seed, as
existing in Abraham, should fill the earth. He meant gradually to
train the people, his descendants, till at length the Christ or
Messias should be born among them, and in His name they should [go]
forth, etc.

4. He did not use them in order to cast them off.
The gifts and callings, etc. \footnote{'As concerning the
gospel, indeed, they are enemies for your sake; but as touching the
election, they are most dear for the sake of the fathers. \emph{For the
gifts and calling of God are without repentance}.'—Rom. xi.
28-29.} Jerusalem instead of Rome, etc. \footnote{If the Jews had not
rejected Christ, Jerusalem would have remained the Holy City.}

5. But when the time came, they would not—they
thought God could not do without them. 'Stones—children of Abraham.'
'Many shall come,' etc. Parable of the great supper and the
vineyard. 'Lo we turn to the Gentiles.'

6. This is a warning to all Christians. It is a
warning to the Roman people who seem to have cast off the Holy See,
for it is not certain that the Pope might not change St. Peter's see,
and it is quite certain that he might simply leave Rome as Jerusalem
was left.

7. It is a warning to each of us.

Continue

Back

\xsection{September 24 (Seventeenth Pentecost) [The Old and New Testaments]}

1. INTROD.—On
the gospel.

2. 'The Lord said to my Lord.' \footnote{'The Lord said to my Lord, Sit on my right
hand … If David then call him Lord, how is he his son?'—Matt. xxii.
44-45.}

3. Great truths put obscurely in the law. Both as
regards prophecy and religion and morals.

4. The law of Moses and the Old Testament like a
bud, and the new law the open flower, \emph{e.g}.—

5. The first commandment.

6. The thoughts of the heart—'with all thy
heart,' Matt. v.; 'If thy eye is single.'

7. Impurity, Matt. v.—divorce and polygamy.

8. The second great commandment—parable of the
good Samaritan.

9. But something the same—\emph{faith}, the \emph{Church},
the order of ministers, and rites.

10. But all these might be \emph{dead} without
love of God, etc.

11. Let us beware lest we are outside Christians.

\xsection{March 31, 1872 (Easter Day) [Victory of Good over Evil]}

1. This day commemorates the victory of truth
over falsehood, of good over evil, of Almighty God over Satan—quote
Matt. xxviii. 1.

2. Not a recent event, the existence of
evil—millions of ages ago, a revolt in heaven—rebel angels; thus
Satan the god of this world. And the conflict began first in
heaven—'Michael and his angels.' Then the devil was cast out, and
came down to the earth. Then it went on to the greater conflict with
the Son of God.

3. Wonderful there should have been such a
conflict and such a victory.

4. (1) No evil without His permission. This is
one wonder.

5. (2) Then when permitted, He might have
destroyed it by a word; but He suffered it.

6. (3) He might have let it run its course, and
die as a conflagration dies out.

7. (4) But He determined on a conflict and a
victory.

8. (5) And a victory of apparent weakness over
force.

9. This was His will, and since He chose this
way, we believe it to be the best way.

10. This has been the character of the conflict
ever since. There has been a conflict, and a victory of weakness over
might. Martyrs.

11. Holy See.

12. Comfort of this time.

13. We do not know what is coming, but we do know
that we shall conquer.

\xsection{April 7 (Low Sunday) [Faith Conquering the World]}

1. Faith is inculcated on us both by the epistle
and gospel of this day.

2. What is faith? Why it is that secret inward
sense in our conscience and our heart that God speaks to us,
accompanied by a sense of the duty to obey Him \footnote{See Note
18, pp. 342-3.}—a sort of voice or command bidding us to believe, telling us
to yield ourselves to Him.

3. Thus, if we hear any one scoff at religion,
speak against God, or against our Lord, or the Blessed Virgin, the
saints or truths of the Gospel, or at the Church, we are spontaneously
shocked and turn away. And if unhappily we listen or read, a feeling
of remorse and distress and sorrow comes upon us.

4. Faith not opposed to reason, but anticipates
it. It is a \emph{short} cut.

5. It is (1) evidently the beginning of religion.
And (2) it was a new thing when our Lord came (except among the Jews).
(3) It 'overcame the world.'

6. It overcame the world. St. John prophesied
when he said it should 'overcome.' How would Christianity have
progressed without it?

7. It overcame the world—by contrast, 'When the
Son of man cometh,' etc. \footnote{' … shall he find, think
you, faith upon earth?'—Luke xviii. 8.}

8. We need not take this to mean there will be \emph{no}
faith, but observe a contrast.

9. We do not know when this time will be, but we
understand from what we see that a time \emph{will} come. The unbelief
now is dreadful, and should remind us of that time.

10. Let us pray that when He comes we may be
found watching.

11. The trial came on the apostles suddenly,
their faith failed.

\xsection{April 14 (Second Easter) [Faith Failing]}

1. INTROD.—The
good Pastor hardly made Himself known to His disciples than He went to
heaven. \emph{He went away before men believed in Him}.

2. This was His will—'Not to all the people.'
Enumerate how few—the most five hundred brethren at once—but then,
as it seems, 'some doubted.'

3. For it was His will that 'the just should live
by faith,' Hab. ii. 4—and then thrice in St. Paul \footnote{Rom. i. 17, Gal. iii.11,
Heb. x. 38.}.

4. Accordingly elsewhere He says, 'We live by
faith, not by sight' \footnote{2 Cor. v. 7.}—Abraham's
faith. So our Lord's miracles. And He said, Mark xi. 22-23, 'Have
faith in God. Amen I say to you, That whosoever shall say to this
mountain, Be thou removed, and cast into the sea; and shall not
stagger in his heart, but believe that whatsoever he saith shall be
done unto him; it shall be done.'

5. Thus the FOUNDATION
of the Church is faith, Matt. xvi. 13-18, and when faith goes the
Church goes. The angels: 'Ye men of Galilee ... so also will he come
again.' \footnote{'Ye men of Galilee, why
stand you looking up to heaven? this Jesus who is taken up from you
into heaven, shall so come as you have seen him going into
heaven.'—Acts i. 11.}

6. \emph{The Church cannot go till faith goes};
and as the Church will last as long as the world, therefore when faith
dies out the world will come to an end.

7. I repeat few had faith when our Lord went, and
few will have faith when He comes again. The foundation of the Church.

8. Hence the words, Luke—'Shall He find faith
on earth?'

9. All this makes us look to the future,
especially when there is a failure of faith.

10. The prophecies distinctly declare a failing
of faith.

11. On listening to prophecies in circulation \footnote{On the number of these
prophecies and their character, see Poulain, \emph{The Graces of Interior
Prayer}, p. 346 [English translation].}—not to be trusted.

12. Of course I am not denying that holy people,
nuns, etc., sometimes prophesy, but Scripture is surer, 2 Peter \footnote{The reference may be to 2
Peter iii. 9, 'The Lord delayeth not his promise,' etc.}. Of course it requires an interpreter, but still there is
something to guide us in the literal text.

13. The awful future—'of that day and hour
knoweth no one'—but it is profitable to read the words of Scripture,
though we but partially understand them.

\xsection{April 21 (Third Easter) [The Second Coming]}

1. INTROD.—'Modicum,'
etc. 'A little while' \footnote{'A little while, and now
you shall not see me; and again a little while, and you shall see
me.'—John xvi. 16 (opening words of the gospel of the Sunday).}—the
disciples were perplexed.

2. Our Lord spoke as if He were to come again
soon. And certainly many of His disciples thought He would. They
thought not exactly that He would end the world, but that He would
come to end the present state of it, to judge the wicked and introduce
a holier world. Nay, at one time even the apostles.

3. But no one knows when, not even the angels.

4. It seems to have been our Lord's wish that His
coming should always appear near.

5. He gave indeed signs of His coming, but every
age of the world has those signs in a measure.

6. The signs were the falling away and the coming
of some great enemy of the Truth called Antichrist, who should bind
together all the powers of the world; that as there was war between
the good and bad angels in heaven, so between the servants of Christ
and Antichrist on earth.

7. This then is our state. In every age things
are so like the last day as to remind us that perhaps it is coming;
but still not so like that we know.

8. Every age is a semblance, a type in part of
what then at last will be in fulness.

\xsubsection{(Same as last; another scheme.)}

1. 'Modicum.'

2. So they would explain the angels' words, '\emph{Viri
Galilaei}.' \footnote{See footnote 1 \textbf{[}Note
6\textbf{]},
p. 224.}

3. And the mass of the disciples did think our
Lord would come back soon.

4. Our Lord said 'no,' 'but of that day,'
etc.—'like a thief in the night,' etc.—2 Peter iii. [10] \footnote{'But the day of the Lord
shall come like a thief … '}, 2 Thess. ii. 3 \footnote{'Let no man deceive you:
... for unless there come a revolt (\emph{discessio}) first, and the
man of sin be revealed,' etc.}.

5. Here, then, signs mentioned—viz. \emph{discessio}
[a revolt] and Antichrist.

6. Before the end a great spiritual war between
Antichrist and the remnant of believers in the world.

7. You may say, 'Then the time of Christ's coming
is known.'

8. No, for this reason—every time is such as to
be like, and to remind us of the last day.

9. True, always [cause for] fear—the world
always seems ending.

10. It is the great mercy of God, and the power
[of prayer that delays the end].

11. However, at length the time will come. [Some
alterations or additions were made in secs. 7-9, which it has been
found impossible to embody in the text given above. Their \emph{placing}
must be left to the ingenuity of the reader.]

Sec. 7. (a) 'In spite of this, in every age
almost, Christians have thought the end coming.' (b) 'Moreover,
though its fulness at the end, always in the world; many Antichrists.'
(c) 'But still it is our duty ever to look out for Him.'

This last, viz. (c), is followed by—'8. Hence
He has made the end \emph{always} seem near.'

Another addition to or substitution for sec. 8
is, 'At this time things very [?] like the end.'

Added in pencil as a substitution or addition to
sec. 9: 'It will also keep us from being over frightened now about
present signs.'

\xsection{April 28 (Fourth Easter) [Prophecy]}

1. INTROD.—I
have said that as our Lord went away suddenly, so will He come again.
Next, that there will be a great token of His coming, viz. a falling
away. Thirdly, that it will still be sudden, because that falling away
is in almost every age, or, at least, again and again.

2. Now some passages in Scripture about the
falling away: 1 Tim. iv. [1] \footnote{'Now the Spirit
manifestly saith, that in the last times some shall depart from the
faith, giving heed to spirits of error, and doctrines of devils.'}, 2 Tim. iii. [1-5], \emph{ib}. iv. \footnote{'Know also this, that in
the last days shall come dangerous times. Men shall be lovers of
themselves ... having an appearance indeed of godliness, but denying
the power thereof.'} [3] \footnote{'For there shall be a
time when they will not endure sound doctrine,' etc.}, 2 Peter
iii. [3-4] \footnote{'In the last days there
shall come deceitful scoffers, walking after their own lusts, saying,
Where is the promise of his coming?'}.

3. About the suddenness, Matt. xxiv. 27, 'For as
lightning cometh out of the east, and appeareth even unto the
west, so shall the coming of the Son of man be.'

4. Infallible word, commented on by theologians
through ages: 'Blessed are those who hear,' etc.

5. Of course at all times there is a spirit of
prophecy in the Church, and there are holy men and women, though there
is no proof of this [in the stories now abroad].

6. It seems to me a great pity that Catholics
leave Scripture prophecy, which is the infallible word, for rumours
and stories about prophecies without foundation, \emph{e.g}. at this
very time.

7. Orval coming up again \footnote{See Poulain, \emph{The
Grace of Interior Prayer}, p. 345 [English translation].} (vide \emph{Rambler}, vol. iv. p. 73).

8. Shifting according to circumstances—instance
of 1748.

9. But still the word of God stands sure and
cannot be superseded. If they are true, they co-operate with Scripture
and do not oppose it.

10. Now this great contrast between these
so-called prophecies and Scripture, the one prophecy good, the other
evil.

11. Those who are always looking for good, are
always disappointed; but it is our comfort and glory to know that the
Church always triumphs, though it seems always failing.

12. Hence two lessons: (1) The bad state of
things is to remind us of His coming and its suddenness. (2) We are
encouraged under it by the feeling it is our \emph{special portion} to
be in trouble, 2 Cor. iv. 8-9 \footnote{'In all things we suffer
tribulation, but are not distressed; we are straightened, but not
destitute; We suffer persecution, but are not forsaken; we are cast
down, but we perish not … '}.

13. Three lessons: (1) To remind and warn. (2) To
calm us, because in every age. (3) To give us faith and hope, from the
sight of the Church's continual victory.

\xsection{May 5 (Fifth Easter) [Holy Scripture]}

1. We are so ignorant, and the world so confused,
that there is a natural desire to know the future (trust in
superstitions, fortune-tellers, etc., etc.). A future must come, and
we know nothing about it, and desire to know it.

2. Fortune-tellers about ourselves and public
affairs—almanacs; and so Catholics have their prophets.

3. All those reports such as Catholics are apt to
be beguiled with, have two tokens of error. They do not appeal to or
carry on Scripture prophecy; next, they are different from Scripture
prophecy, as I said last week.

4. Now, though there was no direct comfort and
instruction to be gained from Scripture prophecy, it would be a duty
to keep it in view, because it \emph{is} in Scripture—because it is
the inspired word of God.

5. And this great evil arises from neglecting it,
because Protestants take it up and interpret it wrongly; they
interpret it against \emph{us}—\emph{our} Scripture becomes a weapon
in their hands because we have relinquished it to them.

6. But great \emph{edification} does come from
reading Scripture prophecy; and a blessing is promised on those who
read. Very little is told us about the future; nothing to
gratify our curiosity, but with it real edification.

7. The Apocalypse brings before us the conflict
between Christ and the world.

8. And so of other Scriptures—the Gospel the
best spiritual book—St. Paul's epistles, the Psalms.

9. Pius VI.'s
declaration \footnote{Letter to Martini,
Archbishop of Florence, 'on his translation of the Bible into Italian,
showing the benefit which the faithful may reap from their having the
Holy Scriptures in the Vulgar Tongue.'}. This why
so many French and Italians have become infidels.

10. To know Christ is to know Scripture—an
anchor.

\xsection{May 12 (After Ascension) [The Wonderful Spread of Christianity]}

1. INTROD.—On
the wonderful beginning and spread of Christianity.

2. (Describe it.) Twelve men, etc., etc.

3. So they went on gaining ground for centuries,
till at length, etc.

4. Then how great their greatness! Think of the
Pope, etc., etc.

5. Yet which was the more wonderful of the two?
Why it is not wonderful that a temporal power should have temporal
strength.

6. Another most remarkable thing is that while it
was gaining ground, it all along thought that it was failing, and the
end was coming.

7. They did \emph{not} think so in the time of
its great prosperity, when it really \emph{was} failing.

8. (Now this presentiment of failure is to show
the strength of the Almighty. We have this grace in earthly
vessels.)

9. It arises from the prophecies. We know evil in
this world, not good, is promised us.

10. Again, it is a type to bring before us the
last age when it will fail (God so contrives the events of this world
that, etc.), and when Christ will come from heaven at the last moment
to save.

11. (Horsley's letter \footnote{See \emph{Discussions and
Arguments}, pp. 107-8 where the letter is quoted.}.)

12. Passages from Malvenda about Rome.

13. The wisdom of God is stronger, etc.

\xsection{June 9 (Third Pentecost) [The Fall of Man]}

1. INTROD.—The
ninety-nine are the angels, the one is man.

2. Man is one because perhaps there are
indefinitely more angels than men; and next, because Adam was one
head, the head of our race. We all sinned in Adam, but each angel who
fell sinned in himself.

3. The account of Adam's fall.

4. Now, to understand how great it was, we must
consider Adam's high gifts. It was a miracle almost, a violation of
his nature and state, that he fell, for he had so many gifts.

5. Had he been like us we could understand it;
but he was not like us. But on his falling he \emph{lost} those gifts,
and became what men are now, and that we can understand.

6. He came under God's anger—he was prone to
sin; he was under captivity of the devil. The whole face of the world
external was changed, as winter instead of summer—that \emph{world},
I may say, deprived of angels, of God's countenance, and full of the
devil; even innocent things became infected and means of temptation.

7. He lost those gifts, and therefore, when he
had offspring, he transmitted to them that nature which he had; but he
could not transmit those gifts which he had forfeited.

8. Such, then, is our state as children of Adam.
We are what he was after sinning—in precisely the same state—and
that state is called 'original sin.' \emph{We have not} the advantage
which Adam had.

9. Now, if a man says this is mysterious, hardly
consistent with justice, I answer: (1) The whole of revelation must be
mysterious, we do not know enough to defend it, because it is part of
a whole system.

10. (2) God is not bound to give us high gifts
such as He gave Adam. It is sufficient that He gives us such grace
that it is our fault if we do not go right.

11. (3) But, again, Christ came to set all right.

\xsection{June 16 (Fourth Pentecost) [The World, the Flesh, and the Devil]}

1. INTROD.—The
whole creation travaileth.

2. All creatures must be imperfect and tend to
corruption if left to themselves. All creation \emph{which we see}—the
visible world.

3. The visible world requires a support to its
laws; they cannot support themselves.

4. And still clearer as regards separate beings.
All things in \emph{fluxu et transitu}.

5. Brute animal passion—but without sin—but
no brute passion but exists in man.

6. Such excesses the gift of reason is to hinder
and subdue; and therefore sin in not doing so.

7. But the conflict so strong that it requires
the grace of God.

8. Now we see the state to which original sin,
the sin of Adam, has reduced us. It has rendered us like the brutes,
because it has deprived us of grace, yet left us in sin.

9. This stripped human nature is called in
Scripture \emph{the flesh}—(Cain's fratricide, the flood,
destruction of Sodom, state of things when our Lord came)—

10. And is our second giant enemy. Our first
enemy is the devil.

11. Now trace the effects of the flesh—the
growth of evil in individuals, in bodies; the power of
example—encouraging each other, appealing to each other; false
maxims—affecting to \emph{teach}.

12. This \emph{the world}, a creation of the
flesh—our third great enemy.

13. Thus fallen man has to fight against three
great enemies.

14. Let us \emph{never forget} we are servants
and soldiers of Christ, Eph. vi. 11-17 \footnote{'Put you on the armour
of God, that you may be able to stand against the deceits of the
devil. For our wrestling is not against flesh and blood, but against
principalities and powers, against the rulers of the world of this
darkness, against the spirits of wickedness in the high places.
Therefore take unto you the armour of God, that you may be able to
resist in the evil day, and to stand in all things perfect. Stand
therefore, having your loins girt about with truth, and having on the
breastplate of justice; And your feet shod with the preparation of the
gospel of peace; In all things taking the shield of faith, wherewith
you may be able to extinguish all the fiery darts of the most wicked
one. And take unto you the helmet of salvation, and the sword of the
Spirit (which is the word of God).'}.

\xsection{June 23 (Fifth Pentecost) [The World Rejecting God]}

1. All men like to be independent and have their
own way, and in many things they can profitably be so and get on more
to their advantage than when they are under rule, but—

2. In one thing they cannot—in religion and
duty.

3. And for this reason: because we are made up of
two principles which war against each other. One or other must be the
master.

4. Satan knew this, both man's desire to be
independent and the impossibility of it. He knew that man must either
be God's servant or his own, and that he, man, did not know this. So
he tempted him with, 'Ye shall be as gods,' and waited securely for
his consequent falling under his own power.

5. Therefore man, rejecting his true Lord,
admitted a usurper. This brings in atheism, \emph{i.e}. idolatry with
immorality. And therefore he always tends to get worse and worse, and
unless God interfered he would become unbearable.

6. But God has always pleaded with man ('My
spirit shall not always,' etc., and 'The Spirit intercedes'), and thus
reserved a remnant. This remnant has pleaded for the world and saved
it. It is the salt of the earth.

7. The deluge—till only eight persons. Earth
filled with violence. 'They ate and drank,' etc.

8. Sodom. 'If ten persons.'

9. When our Lord came. Rom. i.

10. Ever since, it has been the elect few who
have saved the world and the Church.

11. When at length 'He shall not find faith on
the earth,' He 'cometh.'

12. On what in this age takes the place of
professed idolatry, and is really atheism.

\xsection{August 4 (Eleventh Pentecost) [Miracles—I]}

1. INTROD.—The
gospel miracle; other miracles.

2. People say, Why are not miracles now? (1) in
complaint; (2) in unbelief. We know there are \emph{not such} nor \emph{so
many} as once.

3. But let us consider why miracles were
necessary in the beginning—the then state of the world. Even if the
great powers of the world had been inspired to enforce Christianity,
how would that prove it \emph{true}?

4. Mere men as the preachers, so \emph{weak},
they would need something to give them \emph{authority} and \emph{weight}.

5. (1) The world had to be startled and awed,
which weak preachers could not do;

6. (2) Secondly, to be convinced, which worldly,
powerful preachers could not do.

7. A miracle when \emph{real} is what man cannot
do.

8. It was just suited to the case. Common sense
tells us it is just what \emph{would} convince us.

9. Why not now then? It was necessary, especially
in the beginning.

10. And hence it is still accorded by God in
converting the heathen—St. Gregory Thaumaturgus, St. Martin, St.
Augustine, St. Patrick, St. Boniface, St. Francis Xavier.

11. But though we have not miracles as in the
beginning, (1) dealings of God with the human soul are like miracles.

12. And (2) so are providences and answers to
prayer. (Not miracles now, because want of faith. \emph{Vide} the
gospel).

\xsection{August 11 (Twelfth Pentecost) [Miracles—II]}

1. INTROD.—Why
we do not see miracles.

2. We believe that miracles \emph{are} wrought
now, though they are few.

3. I have spoken of miracles wrought by apostles
of countries.

4. And so of saints. If I am asked why miracles
scarce, I answer, Saints are scarce. We cannot conceive common men
doing miracles.

5. You will ask, \emph{Why} are saints scarce
now? It has ever been that times vary. There are sometimes bursts of
supernatural power and greatness.

6. So the Psalms, xliii. \footnote{\emph{Deus auribus nostris}.
In which the Church commemorates former favours and present
afflictions.}, lxxiii. \footnote{\emph{Ut quid Deus}. A
prayer of the Church under grievous persecutions.},
lxxxviii. (finis) \footnote{'Be mindful, O Lord, of
the reproach of thy servants,' etc.}, and
Isaias li \footnote{An exhortation to trust
in Christ. He shall protect the children of His Church.}.

7. But when there \emph{are} saints there are
great miracles. St. Philip.

8. But you will say, If there are few saints on
earth, yet there are many in heaven; why do they not do miracles from
heaven, as St. Philip used to do, as we read in the accounts appended
to his life?

9. Because we have not \emph{faith}—not \emph{individuals}
merely, but the population. (Enlarge on this.)

10. \emph{Vide} Luke xix. 26, Matt. xxi. 27, Mark
ix. 23, Mark vi. 5.

11. Because men say, 'Unless we see signs and
wonders,' etc., in a haughty way.

12. Miracles now come as a reward to faith, in
those who do not look out for them. Not denied then.

\xsection{August 18 (Thirteenth Pentecost) [Christ's Presence in the World]}

1. INTROD.—We
have read, Sunday after Sunday, as today, of our Lord's miracles; but
did we see Him, I do not think that [the miracles] would most strike
and subdue us.

2. Not His works, \emph{but Himself}.

3. But here I explain something. Strange to say,
it was His will that, seen by casual spectators, He should seem like
another man, Isa. liii. 3 \footnote{'His look was as it were
hidden and despised, whereupon we esteemed him not.'};
and hence John i. 5, 10 \footnote{'The light shined in the
darkness; and the darkness did not comprehend it.' \emph{Ib}. 10, 'He
was in the world, and the world was made by him, and the world knew
him not.'},
and Mark vi. 3 \footnote{'Is not this the
carpenter, the son of Mary, the brother of James, and Joseph, and
Jude, and Simon? are not also his sisters here with us? And they were
scandalised in regard of him.'}. And
the Samaritan woman, John iv. And this specially so in the case
of bad men, Luke xxiii. 11 \footnote{'And Herod and his army
set him at nought, and mocked him.'},
John xix. 9 \footnote{'And Pilate said to
Jesus, Whence art thou?'}.

4. When we had seen Him two or three times, if we
were not utterly dead to truth we should find that He had made a deep
impression on us, on looking back, though we did not perceive it at
the time, Luke xxiv. (Emmaus).

5. Next, supposing we could stay and gaze on Him,
then what would first strike us would be His awful infinite repose,
the absence of all excitement, etc., etc. All that is told us of Him,
all His words and works, brings out this—and doubtless His aspect.

6. Next, if we could still look on, if we could
see His eyes, two things would strike us; first, His seeing us through
and through. Hence He is often said to 'look.' Mark iii. 5, 'And
looking round about on them with anger'; \emph{ib}. viii. 33, 'Who,
turning about and seeing his disciples, threatened Peter,' etc.; \emph{ib}.
xi. 11, 'And he entered into Jerusalem, and having viewed all things
round about.'

7. Secondly, compassion. Mark x. 21, 'And Jesus
looking on him loved him'; Luke xxii. 61, 'And the Lord turning,
looked on Peter: and Peter remembered the word of the Lord.'

8. And then when He began to speak! the tones of
His voice! John vii. 46, 'The ministers answered, Never did man speak
like this man'; Matt. vii. 28, 'And it came to pass, when Jesus had
fully ended these words, the people were in admiration at his
doctrine: For he was teaching them as one having power.'

9. Hence He \emph{draws} men. Matt. ix. 9, 'And
He saw a man sitting in the custom house, named Matthew: and he said
to him, Follow me. And he rose up, and followed him.' Virtue going out
of Him. Mark v. 30, 'And immediately, Jesus knowing in himself the
virtue that had proceeded from him, turning to the multitude, said,
Who hath touched me ?' \emph{ib}. vi. 56, 'And whithersoever he
entered, into towns, or into villages, or cities, they laid the sick
in the streets, and besought him that they might touch the hem of his
garment: and as many as touched him were made whole.'

10. All this, even though He did no miracle.

11. This is what we must look for in heaven.

12. And yearn for it \footnote{These last words are
barely, if at all, legible.} in the Blessed Sacrament.

\xsection{August 25 (Fourteenth Pentecost) [The 'Two Masters']}

1. INTROD.—Two
masters. \emph{Why} cannot we serve two masters? Most men wish to
serve God and the world.

2. What is it to have a master? what is meant by
it?

3. Not merely an employer; this not enough.

4. A master is one who has some hold over us. In
old times slaves, but now it is by compact. If I promise, if I take
wages, I willingly take a master. As children are naturally subject to
parents, so, by free will, servants to masters. They may \emph{change},
but while they have a master they are \emph{bound}.

5. Now on \emph{serving} a master. Consider St.
Paul, Eph. vi. 5-6 \footnote{'Servants, be obedient
to them that are your lords according to the flesh, with fear and
trembling, in the simplicity of your hearts, as to Christ; Not serving
to the eye, as it were pleasing men; but as servants of Christ, doing
the will of God from your heart.'}.

6. And if so of all masters, so especially of the
good—idea of a \emph{household}.

7. Now we see what in religion is meant by God
being our master. (1) He has created and \emph{bought} us. (2) We have
made an everlasting contract with Him. (3) It is not a contract in
this or that— as \emph{employers}—but we are of His household and
family. (4) We are one of His, and must study His interests. (5) He is
a good master.

8. \emph{Hence}, if our Lord is our master, we
can have no other master, and we must be full of zeal and love.

9. He has given Himself \emph{wholly} to \emph{us}.

10. The other—\emph{Mammon}! So not only we can't
have two; we must have \emph{one}.

11. Now let us ask ourselves: \emph{Is} in fact
God our master? Do not we follow our own will, taking one day one
master, another another.

12. There would not be all this variety of
religions, and this infidelity in the world, if men really made God
their master. They would soon agree together.

On men of no party. Apoc. iii.—Laodiceans \footnote{It is not clear where
these words were intended to come.}.

\xsection{September 1 (Fifteenth Pentecost) [Miracles]}

1. INTROD.—Miracle
on the widow's son at Naim.

2. Open, public—so on Lazarus, John xii. Matt.
ix. [6], 'that you may know.' \footnote{'But that you may know
that the Son of man hath power on earth to forgive sins, (then said he
to the man sick of the palsy,) Arise,' etc.} So Acts iii. \footnote{The miracle upon the
lame man.};
so Elias, 3 Kings xviii \footnote{Elias cometh before
Achab. He convinceth the false prophets by bringing fire from heaven;
he obtaineth rain by his prayer.}.

3. But many others our Lord \emph{forbids} the
proclaiming. Thus He takes the blind men into a house, or charges
them, etc., Matt. ix. 27 ff. \footnote{'And as Jesus passed
from thence, there followed him two blind men, crying out and saying,
Have mercy on us, O Son of David. And when he was come to the house,
the blind men came to him.'} And again, still more remarkably Jairus's daughter, Mark v. 43
\footnote{'And he charged them
strictly that no man should know it.'}, Luke viii. 56 \footnote{'Whom he charged to tell
no man what was done.'}.

4. Now this will tell us how to answer the
question about miracles now. There are miracles now, but not such
miracles as in the beginning—not public ones. They were in order to
establish the religion—but now the religion established.

5. Then they might be wrought by bad men—by
Judas—Matt. vii. 22. But now they are marks of sanctity of the
persons or the things by which they are wrought.

6. Hence (1) the workers do not \emph{proclaim}
them.

7. (2) Not so marked, by running into [\emph{i.e}.
not easy to be distinguished from] providences.

8. (3) Not so discernible—to one, and not to
another.

9. (4) No \emph{necessity} to believe them, for \emph{the
Church does not propose them}.

10. (5) Not to be urged on unbelievers.

11. As I have said before, the miracles of the
Catholic Church are those which are personal to ourselves. (1)
Goodness of God to us in the course of life. (2) His grace given to
our souls.

\xsection{September 8 (Sixteenth Pentecost) ['The Riches of His Glory']}

1. INTROD.—The
epistle [Ephesians iii. 13-21].

2. Do you understand St. Paul's words, 'The
riches of his glory,' etc.?

3. We have here a glimpse of what heaven is. 'Eye
hath not seen,' etc.

4. It was the support of St. Paul against the
world.

5. The world has \emph{its} 'depth' and 'height,'
etc. Illustrate deep science, high power, glory, etc.

6. It is this which makes the world the false
prophet; it preaches and seduces us with false maxims.

7. It is grievous to say, but it must be said,
that almost all we read, the periodical press, is in this respect a
false prophet.

8. The devil said he had 'all the kingdoms of the
earth.' Things good under bondage of evil.

9. Therefore God gave us the Church—as the true
Prophet to bring the glories of heaven before us.

10. All sacraments, etc., with this object.

11. So Scripture a \emph{revelation} of the next
world—especially our \emph{Lord's person}.

12. And so the saints and their history—a whole
family round our Lord.

13. This is the reason why we are allowed to
think so much of our Lady; why she is \emph{given} us to employ our
thoughts. Protestants say we make too much of her. Now which is best,
to think too much of her, or of the world?

\xsection{September 22 (Eighteenth Pentecost) [Disease a Type of Sin]}

1. INTROD.—The
paralytic in the gospel.

2. The cures to typify the spiritual disorders
and diseases of mankind. This \emph{one} reason of the special
character of our Lord's miracles.

3. And it is well to consider the \emph{variety}
of bodily diseases with this view. They are horrible, but we may be
sure that the various spiritual maladies are far more horrible.

4. And the \emph{least} sin, for its quality is
so bad—in this it goes beyond bodily diseases, for bodily
infirmities admit of degrees much more. And it goes beyond the analogy
of disease in these respects: (1) because universal to the race; (2)
because so intense.

5. But the case of sin may be likened to the
analogy of offences against the senses, as to which the least
imperfection is \emph{destructive}; \emph{e.g}. the sweetest nosegay
spoiled by one bad scent of one dead leaf. One drop of bitter in the
most pleasant drink. And so of hearing, one discordant note. And so in
the sciences—in astronomy the slightest motion [vibration in
an observatory]—or in the mirror or glass the slightest dimness; and
in chemistry, poisons; and in medicine, etc.

6. The whole creation marred. Then why did God
allow it? I answer—

7. That is a question not for the present time.
You don't inquire \emph{how} a fire arose before you have extinguished
it.

8. Next, our Lord came to \emph{destroy} sin.
This the characteristic over all other religions ('not the righteous,'
\footnote{'I am not come to call
the righteous, but sinners to repentance.'—Mark ii. 17.} 'repent ye,' \footnote{'Jesus preached: Repent
ye, and believe the gospel.'— Mark i. 15.} the lost sheep); they [other religions] acknowledge sin, but
they cannot cure it.

9. He takes away the guilt, and the power [of
sin].

10. This by His death and passion.

11. This the fundamental doctrine—texts.

\xsection{October 13 (Twenty-first Pentecost) [Forgiveness of Injuries]}

1. INTROD.—Today's
gospel [the king taking an account of his servants].

2. Parallel passages; Luke xvii. 3-4 \footnote{'Take heed to
yourselves: If thy brother sin against thee, reprove him; and if he do
penance, forgive him. And if he sin against thee seven times in a day,
and seven times in a day be converted to thee, saying, I repent;
forgive him.'}.

3. So far easy, for it is scarcely supposable
that one should have so little generosity as to refuse
forgiveness to one who confessed himself wrong and asked to be
forgiven.

4. But when he does not ask to be forgiven; if he
persists in opposition and injury, and goes on doing harm, and takes a
wrong course. Yet this commanded too. The Lord's Prayer—Matt. vi.
14-15 \footnote{'And if you will forgive
men their offences, your heavenly Father will forgive you also your
offences: But if you will not forgive men, neither will your Father
forgive you your offences.'}, Mark xi. 25-26
\footnote{'And when you shall
stand to pray, forgive, if you have ought against any man: that your
Father also who is in heaven may forgive you your sins. But if you
will not forgive, neither will your Father that is in heaven forgive
you your sins.'}, Rom. xii. [18-20] \footnote{'If it be possible, as
much as is in you, having peace with all men. Not revenging
yourselves, my dearly beloved, but give place unto wrath: for it is
written, Revenge to me; I will repay. But if thy enemy be hungry, give
him to eat; if he thirst, give him to drink: for doing this thou shalt
heap coals of fire upon his head.'}.

5. Or again, supposing he does not do so, asks to
make it up, still there may be, you may say, such difficulties as
these: I may wish to keep at a distance, for:

6. (1) \emph{E.g}. I cannot trust him; he is a
dangerous man.

7. (2) He is likely to do me spiritual harm.

8. (3) The sight of him is a temptation, an
irritation to me; we shall be best friends at a distance.

9. (4) I shall be a hypocrite if I make it up,
for I don't like his doings.

10. (5) I ought to protest against him.

11. ANSWER.—'If
you \emph{in your hearts forgive} not every one his brother'
[Matt. xviii. 35]. You must love him. Col. iii. 12-13 \footnote{'Put ye on therefore, as
the elect of God, holy and beloved, the bowels of mercy, benignity,
humility, modesty, patience; Bearing with one another, and forgiving
one another, if any have a complaint against one another: even as the
Lord has forgiven you, so do you also.'}; Matt. v. 44-47 \footnote{'Love your enemies do
good to them that hate you, and pray for them that persecute and
calumniate you; That you may be the children of your Father who is in
heaven, who maketh his sun to rise upon the good and bad, and raineth
upon the just and the unjust. For if you love them that love you, what
reward shall you have? do not even the publicans this? And if you
salute your brethren only, what do you more? do not also the heathens
this?}.

12. OBJECTION.—'But
I do not like him.' How can I love him? \emph{This} is a fundamental
difficulty.

13. ANSWER.—Can
you pray that you may meet him and love him in heaven? You and he are
both far from what you should be; and each has to \emph{change}. Look
on the \emph{best} part of his character—learn sympathy with him.
Think how he suffers. Purgatory useful for this—to bring you and him
nearer to each other.

\xsection{October 20 (Twenty-second Pentecost) [Final Perseverance]}

1. INTROD.—Epistle
for the day, perseverance in grace.

2. Two things plain: (1) perseverance necessary,
Matt. xxiv. 13 \footnote{'He that shall persevere
to the end shall be saved.'}, Ezech.
xxxiii. 18 \footnote{'For when the just shall
depart from his justice, and commit iniquities, he shall die in them.'}.

3. (2) Not in our power, but a special gift of
God. We cannot merit it.

4. Now what is merit? (Explain.) By ourselves not
only not perseverance, but nothing can we merit.

5. Because (1) by ourselves we can do nothing
pleasing to God, because of our sinfulness; and (2) because how can
anything we do be worthy of \emph{heaven}? what proportion? Luke xvii.
7-10 \footnote{'But which of you,
having a servant ploughing or feeding cattle, will say to him, when he
is come from the field, Immediately go sit down to meat? And will not
rather say to him, Make ready my supper, and gird thyself, and serve
me, whilst I eat and drink; and afterwards thou shalt eat and drink?
Doth he thank that servant for doing the things which he commanded
him? I think not. So you also, when you shall have done all these
things that are commanded you, say, We are unprofitable servants: we
have done that which we ought to do.'}.

6. (1) Therefore the grace of God, and (2) His
promise; thus we can be said, first, to please God, and secondly, to
merit.

7. And these two by the merits of our Lord and
Saviour.

8. But there are two things we cannot merit—the
first grace and the last.

9. As to the first grace, it is plainly God's
free bounty which has made us Christians.

10. As to the last, it is God's free bounty, in
spite of the accumulation of merits. No extent of merit is sufficient
to gain perseverance—the just may fall, however holy, etc. Think of
Solomon; think of Judas. It is a special gift to die in grace.

11. Two conclusions. First let us continually
pray that God would give us this special gift of dying in grace.

12. This may comfort us when we lose our friends,
that God may in His mercy have taken them then, when they were in
grace.

\xsection{January 12, 1873 (Sunday in Epiphany) [Manifestation of the Kingdom of Christ]}

1. INTROD.—The
Magi.

2. They were a prophecy and anticipation of what
was coming.

3. We know the kingdoms of this world became the
kingdom of Christ.

4. Two things wonderful: (1) that such a conquest
should be made; (2) that it should be prophesied.

5. That kingdom is passed by, three hundred years
ago. We have, however, the \emph{remains}—cathedrals, ruins of
abbeys—the usages of society, etc.

6. So that we are known as the 'old religion.'
And what is old comes from our Lord, and what is new came from men.

7. This, then, is the wonderful manifestation of
past times.

8. But now it is wellnigh past—while it lasted
it was comparatively easy to believe when there was only \emph{one}
religion.

9. But now Satan, who has his instruments in
every age, says: 'There are so many religions, none is true; they are
all false.'

10. Those who live will find a \emph{wave} of
infidelity overspread the land. What they are to do.

11. There is another \emph{manifestation} \footnote{The other \emph{manifestation}
seems to be the Church with her notes. The claims of this witness to
be interrogated are put off with 'There are so many religions,' etc.}: 'Come and see'—whereas men keep you from coming and seeing.
'A witness in court'—we say, 'Let us actually see him.' But
no—here it is 'so many religions, etc., Catholicism is going down,'
or 'Catholics are a fallen race,' etc., etc., instead of above, 'Come
and see.' Reading the Gospels. John ix., Luke xvii. By 'religious men'
is meant those who have experiences.

\xsection{January 27 (Third Epiphany) [Men of Good Will]}

1. INTROD.—The
centurion in the gospel of the day. Account in St. Matthew, in St.
Luke.

2. He was a heathen, etc. This is how our Lord
began the Church, when as yet there was none, and addressed Himself to
those who were well inclined, and gained them.

3. This is what is meant by men of good will in
the angels' song.

4. Instances: Nicodemus, John viii.; Gamaliel in
Acts v.; Luke ix. [49], 'who followeth not us' \footnote{'And John, answering,
said, Master, we saw a certain man casting out devils in thy name; and
we forbade him, because he followeth not with us.'}; Syrophoenician [woman], Matt. xv., Mark vii.

5. And so now. We must not repel them or treat
them harshly, or laugh at them, etc.

6. They illustrate the secret work of
grace—from grace to grace.

7. Tests of being \emph{bonae voluntatis} [of
good will]—\emph{not} justice, sense of fairness, and benevolence,
though these are praiseworthy—and we must be grateful to such men.

8. But (1) humility from sense of sin. 'Lord, I
am not worthy.' 'Even the dogs,' etc.

9. (2) Sense of duty. 'I am a man under
authority.'

10. (3) Devotion. 'He has built us a synagogue,'
Luke vii.

11. Let us beware lest those who have less
advantages than we have outstrip us. 'Many shall come from the east
and the west,' etc.

\xsection{March 2 (First Lent) [God Our Stay in Eternity]}

1. INTROD.—We
must draw near to God.

2. This means to contemplate, to recognise, to
fear, to love. Now let us see the necessity of this.

3. Here we are tempted to make the world our God,
because we see it, and do not see God.

4. But consider what our state is when we are
dead; our senses then are all gone.

5. Consider this: we have five senses, and we
know what a deprivation the loss of any one—sight or hearing or
touch or feeling—any one.

6. But in death they all go together. See what we
are reduced to. It is true we cannot have any bodily pain—and that
is what people are apt to say, 'All his pain is over.'

7. True, but is there no pain of the mind? Do
we know how acute pain of the mind is?—surely we know it even
in this life.

8. Let us consider our being suddenly cut off
from all intercourse except with ourselves—a truly solitary
confinement; worse, for that here is only loss of hearing, \emph{i.e}.
conversation.

9. Supposing in addition it comes on us that we
should not be thus, except for our own fault!

10. Now it is clear that we should have no remedy
unless God visited us and gave us light.

11. The light of glory, the light of heaven, the
only thing.

12. But suppose we have no desire for it, no love
of it. Suppose we look back in fond \emph{regret} to this world.

13. Therefore the love of God is the \emph{only}
way in which we can be happy.

\xsection{June 22 (Third Pentecost) [The Lost Sheep the Type of Fallen Man]}

1. INTROD.—Gospel,
one sheep in wilderness, man; the ninety-nine, angels.

2. Contrast between angels and man. They so
great, we so low.

3. Yet Psalm viii. 4-5, 'What is man,' etc. \footnote{'What is man, that thou
art mindful of him? or the son of man, that thou visitest him? Thou
hast made him little less than the angels,' etc.}; 'out of weakness were made strong' \footnote{Heb. xi. 34.}; 'when I am weak, then I am strong' \footnote{2 Cor. xii. 10.}; 'these things the angels desire,' etc. \footnote{'Which things the angels
desire to look into.'—1 Peter i. 12.}; 'joy among the angels of God.' \footnote{' ... upon one sinner
doing penance.'—Luke xv. 10.}

4. For, see the difference. Even angels fell; and
even for angels no restitution. You would think they were more
convertible—they had no powers to return.

5. Could, then, any being return, if not angels?

6. Man not only in the image of God, but of the
beasts of the field.

7. Incarnation.

8. All things possible with God.

9. Hence a great multitude.

10. Hence \emph{saints}.

11. Woe is us, if \emph{elect}, yet such as we
are.

\xsection{January 4, 1874 [The New Year]}

1. Difference of feelings of young and old
towards a new year.

2. The young with hope and expectation; the
mature with anxiety.

3. The young look forward first for a
change—each year brings changes. And to them they are changes, as
they think, for the better; they are older, stronger, more their own
masters, etc.

4. And secondly, the future is unknown, and
excites their curiosity and expectation.

5. It is different with them who have some
experience of life. They look (1) on \emph{change} as no great good;
they get attached to things as they are, etc.

6. But (secondly) the ignorance of the future, so
far from being good, is painful—in truth it is one of our four
wounds. Ignorance of all things, especially of the future—of what a
day may bring forth—of suffering, bereavement, etc.

7. Thus, like railway train, bowling away into
the darkness.

8. Ignorance what sufferings and bereavements are
in store—of death—of the day of death. We \emph{walk over our own
dying day}, year by year, little thinking.

9. It may be a work-day, or holiday, or a 'many
happy returns' [day].

10. All things make us serious. \emph{This we know},
that death is certain; and then the time comes when there will be \emph{no}
change—for time is change—and no ignorance.

\xsection{June 28 (Fifth Pentecost) [The Jewish and the Christian Church]}

1. 'Unless your justice [exceed that of the
scribes and Pharisees, you shall not enter into the kingdom of
heaven,' Matt. v. 20].

2. The Jews, then, God's people, and their Church
God's Church. It was the Ark. The world lay in wickedness, and in the
wrath of God, except that holy Church which God founded by Moses. The
Pharisees its rulers.

3. It was salvation, for 'salvation of the Jews.'
So now.

4. It taught God's law. 'Moses' seat.' So now.

5. Indefectible, never to end. You will say it
ended. No, it changed into the Christian Church.

6. But though Jewish Church could not fall away,
its members could. And \emph{so now}. Even its rulers could fall
away, though they taught what was right—Moses' seat; and so could
the body of its people, and \emph{so it did}. They relied \emph{on their
privileges}, and were cast off. St. John the Baptist said, 'Flee
from the wrath to come,' to Pharisees coming to His baptism. And so
Christians may [fall away]. This is a warning to us, and St. Paul so
makes it, Rom. ii.

7. Therefore whatever is said to or about the
Jews is a warning to us.

8. Thus what is said about the Samaritans. (Who
were the Samaritans?) Many are singled out as better than the Jews.
(1) The good Samaritan; (2) the grateful Samaritan \footnote{The leper who returned
to give thanks.}. They are like Protestants. So Protestants may be better than
we in spite of 'salvation from the Jews.'

9. Nay, heathen were better than the Jews, \emph{e.g}.
centurion—'Many shall come,' etc., Matt. viii. 11-12 \footnote{'Many shall come from
the east and the west, and shall sit down with Abraham, and Isaac, and
Jacob, in the kingdom of heaven. But the children of the kingdom shall
be cast out.'}; Tyre, Sidon, Sodom \footnote{See p. 47 (footnote).}
[Matt. xi. 21-23].

10. Thus at present countries on the
Continent—they may be cast off as the Jews were. Protestants in
England may be better.

11. But we must look to ourselves. Many are
called. \emph{Strive} and \emph{seek} [Luke xiii. 24].

12. All those who are in earnest, though they
know their imperfections, must not fear.

\xsection{August 2 (Tenth Pentecost) [Revelation—Word of God (I)]}

1. INTROD.—I
have been reading from Scripture, viz. an epistle and gospel. Why?

2. What is meant by Scripture, Scriptures?
Writings, \emph{the Word of God}, or \emph{revelation}—through \emph{different}
ages.

3. Why has God given us a 'Word'? Because we are
so ignorant.

4. Two Testaments. First with one nation (Old),
then with people of \emph{all} nations (New).

5. The Bible, the Book.

6. By the by, why are Catholics said to \emph{burn}
the Bible? They never do, or have (unless they committed an act of
sin); but what they \emph{burned} was not the Bible but a Protestant
translation. (Also without comment.)

7. The Church \emph{comments} and \emph{explains}.
Now as to the Old Testament, or the Word of God to the Jews,

8. Law and prophets,

9. Till our Lord came.

\xsection{August 9 (Eleventh Pentecost) Revelation—Word of God [II]}

1. INTROD.—Recapitulation.
Scripture—Scriptures—two Testaments—Bible—teaching, and
therefore inspiration.

2. When I say 'inspired'—not in science or art,
etc.

3. Difference of Old and New—Old imperfect, and
through so many ages; New perfect, and once for all in one age.

4. Here I shall speak of the New. The
Apostles—inspired—our Lord God. Heb. i. 1-2 \footnote{'God, who at sundry
times and in divers manners spoke in times past to the fathers by the
prophets, Last of all in these days has spoken to us by his Son, whom
he hath appointed heir of all things, by whom also he made the world.'}.

5. OBJECTION.—Why
not their words [inspired], if their writings? Why not their speeches?
Why not their conversation? Of course it was. All they said about RELIGION
was. They might not know about the earth going round the sun, etc.

6. But it might be objected, on the other hand,
that such sayings were not recollected. But some might be.

7. This is what Catholics called 'tradition,' and
in which we differ from Protestants. Meaning of the word 'tradition.' \emph{Vide}
epistle for this Sunday \footnote{1 Cor. xv. 1-10. 'For I
delivered [\emph{tradidi}] unto you first of all which I also
received.'}.

8. Things we know by tradition: (1) that
Scripture is the inspired Word of God; (2) what books the Bible
consists of—for these Protestants need tradition; (3) the Mass, etc.

9. And so natural. Every school, every set of
workmen, go by tradition—'common law' is tradition.

10. Hence we say there are two parts of the Word
of God, written and unwritten.

11. But still, surely tradition \emph{may} go
wrong. Yes, and Scripture may be wrongly interpreted.

12. Therefore \emph{the Church} decides, as being
infallible.

13. Hence there may be mistaken reports of
miracles, prophecies, etc., but we must see what the Church says about
them.

Continue

Back

\xsection{March 28, 1875 (Easter Day) [The Seen and the Unseen Worlds]}

1. INTROD.—This
is the greatest day of the year, because it is the day on which our
Lord rose from the dead. He said while He died on the Cross, '\emph{Consummatum
est}'—'It is finished,' and in the Resurrection we behold the
fruit of His 'finishing.'

2. (1) The miracle itself—and this special,
because without corruption; others [who were miraculously raised to
life] died again. (2) Next it is a \emph{reverse} after a sorrowful
week. (3) It was a conquest of the foe. (4) It was the exaltation of
our nature.

3. But further, it gives us a great lesson, never
to despair. There are two worlds, and nothing which we see in this
world images to us what is going on in the next.

4. This world runs by laws. All things go on as
at the beginning of the Creation. The sun rises and sets; and so human
affairs. They thought by killing Him to have stamped out religion.

5. And so now—the political world, commercial,
scientific; telescopes, calculations, ships, etc., etc.—but another
world going on \emph{too}. This world a veil.

6. The mass of men only see this world. Each
man enters in this world with hopes for a career, etc., etc. Not
wrong in minding this world, but in not minding the next. Their view
of life.

7. The flood; 'marrying and giving in
marriage'—\emph{but} another world \emph{too}.

8. We walk by faith, not by sight.

9. Therefore acts of faith, hope and charity.

10. Therefore attending Mass, in which the whole
mystery of redemption, atonement, resurrection, etc., [is set forth].

11. Let us thank God for giving us eyes, and pray
Him to give others eyes too, for He died for all.

Perhaps no laws in heaven, but every act from God's
personality; and each perfect in itself, so that we could not reason
from one to the other \footnote{This paragraph is placed at the end. It was
written on an empty space, and it is not clear where it was meant to
be.}.

\xsection{June 6 (Third Pentecost) The Sacred Heart}

1. INTROD.—Many
devotions in Holy Church. This is one which has spread of late years,
more, I may say, than any other. Today is the special feast of it; and
[this] leads me to explain in what it consists.

2. Our Lord is One. He is the one God. He took on
Him a manhood, a body and soul; that body from Mary. Still, He was
one, not two—one, as each of us is one.

3. We too, in our way, are each of us one, though
we are two—soul and body—and the body has parts; [nevertheless] each of us is \emph{one}. This is what is meant by
speaking of our Lord's [oneness] as we speak of our own.

4. And though each of us is thus composite, we
can love each other as one, though of so many parts. And in like
manner, though our Lord is God and man, with a soul [and body], we can
contemplate Him as one, and worship, love Him as one.

5. Further, if I said I loved the face, or the
smile, or liked to take the hand of my father or mother, it would be
because I loved them. And so, when I speak of the separate portions of
our Lord's human frame, I really am worshipping Him. So in the Blessed
Sacrament we do not conceive of His Body and Blood as separate from
Him.

6. Devotions at various times [and ages]—the
Wounds, the Blood, the Face—and in like manner the Heart. We \emph{worship}
[\emph{each}] \emph{as Him}, as that One Person who is God and man; we
worship [\emph{Him}] by the memento, the pledge of His Heart.

7. Why? The Heart a \emph{symbol}—so the Wounds
and the Blood. [In contrast with these] a symbol is sometimes that
which [only] expresses and reminds—thus water, oil, wine, bread.

8. What is the Heart the symbol of?—of His
love, His affection for us, so that He suffered for us—the agony in
the garden.

9. Moreover, of His love in the Holy Eucharist.

10. The Heart was the seat, first, of His love
for us; secondly, of His many griefs and sorrows.

[\emph{The following were appended, apparently as
alternatives}:]

7. Of two things especially to remind us now,
when the world is so strong—His power and His love. He will overcome
by love.

8. The Heart is the emblem of His love—in
worshipping It we worship Him.

\xsection{Christmas Day [Christmas Joy]}

1. Almighty God condescends to be represented in
human language as hoping, fearing, suffering disappointment,
repenting, feeling anger, etc. But there are two human feelings and
affections which may be predicated of Him, not in the way of figure,
but \emph{proprie}—joy and love; of course I mean as being perfectly
free from human passion.

2. Of the two great festive seasons of Easter and
Pentecost [on the one hand], and Christmas [on the other], Easter,
with the fifty days of Pentecost, is the season of love, Christmas of
joy.

3. I need hardly say so—our churches, our
altars are dressed up as token of our joy—and our houses, according
to our opportunity—and we meet together for social enjoyment, and to
provide festive meals and entertainments of various kinds for the poor
and for children. And all this, of course, is right, and is the proper
token of our faith and hope, of our Christian joy. It is, I say, the \emph{season}
of joy, and therefore it is fitting that we should exhibit these signs
of our being full of joy.

4. But a warning needs to be given. It is quite
possible, rather, alas! it is not uncommon, for men to stop at,
to get no further than the outward signs. Nay, I may say that it is
quite plain in a country like this, in a place like this we may—nay,
we do—content ourselves with rejoicings which are temporal, earthly,
visible, without going at all to what is the real reason, after all,
for those external signs. All around us men are doing so, and thus we
are led to do as they do.

5. Therefore most necessary to recollect that
there are two sorts of joy, earthly or outward, heavenly or inward,
and that we may easily place the former in the place of the latter.

6. We shall understand best what the true joy is
by what is told us in Scripture of the first Christmas. If that joy
really consisted in anything external, joy could not have been there.
(1) Go through the journey to Bethlehem—winter—slow
journey—lagging behind—no room in the inn, as the inn was \emph{full},
doubtless the stable also—caravanserai—the stable, etc. State of
the stable—not like our stables, neat and clean. This is what the
shepherds found. They themselves had a hard time of it, watching their
flock by night, but they came to a worse place; not so cold, but less
like a home. Yet, I say, they rejoiced. Contrast of Herod in his
palace close by. (2) The angels sang 'Glory to God in the highest; and
on earth peace to men of good will'; and the angel said to the
shepherds, 'Behold, I bring you tidings of great joy,' and the
shepherds returned 'rejoicing': and Mary [in] the \emph{Magnificat} at
an earlier time shows her thoughts. Yet what were the outward
circumstances?

7. Paragraph in newspaper: 'Somehow Christmas
festivities fall flat when one has grown up.' The shepherds, etc.,
rejoiced even in [the midst of] their own outward discomfort and
hardships; but men of the world cannot lastingly rejoice even in the
midst of their good things. Oh how this shames our delicacy, our
desire of comfort, etc.! Of course we may thankfully take what God
gives us; but at least, while we rejoice in these gifts, let us not
forget to let our inward spiritual rejoicing keep pace with our
external.

8. Let us recollect the apostle's words. [\emph{Perhaps
the preacher quoted here from the epistle read at the first mass of
Christmas}: 'The grace of God our Saviour hath appeared to all men,
instructing us that, denying ungodliness and worldly desires, we
should live soberly, and justly, and godly, in this world; looking for
the blessed hope, and coming of the glory of the great God and our
Saviour Jesus Christ.'—Titus ii. 11-13.]

\xsection{December 26 (St. Stephen) [The Martyrs]}

1. INTROD.—The
first martyr—what meant by a martyr—witness for the truth. Christ
the first martyr \footnote{'Thou sayest that I am a
king. For this was I born, and for this came I into the world, that I
should give \emph{testimony to the truth}.'—John xviii. 37.}, but He \emph{more}
than a martyr.

2. There is one God, but He was forgotten by all
the earth (except the Jews).

3. If one God, only one religion. But every
nation had its own god or gods, and they never thought of interfering
with each other.

4. God suffered this for a long while (Acts xiv.
16 \footnote{'Who in times past
suffered all nations to walk in their own ways.'}), but at length, etc.

5. Hence preachers and evangelists,
apostles—but men did not like to be interfered with. It was a new
thing; hence persecutions, and the preachers became martyrs.

6. In Jerusalem they first suffered, because they
came in collision with the prejudices of the Jews—such St.
Stephen—then a wider range—the apostles all martyrs.

7. The \emph{Te Deum} calls them an
army—(enlarge upon this). Contrast to Mahometans—nay, to
Protestantism, which spreads, not indeed by persecution, but by the
patronage, etc., of the State. What can be more wonderful than an army
conquering by being beaten?

8. The most horrible deaths; stoning is bad
enough, but it [was] only one way—St. Andrew, St. Bartholomew, St.
Peter, St. James, St. Paul, St. John (oil)—the young, the old, the
weak, the strong, men and women.

9. FIRST REFLECTION.—\emph{Thus}
are \emph{we} Christians. What, under God, do we not owe to them?

10. SECOND
REFLECTION.—How \emph{comfortable}
our lives are! The thought of the martyrs should support those too who
are in pain, etc., and those who see their friends in pain.

11. All this should humble us.

\xsection{January 9, 1876 [Life a Pilgrimage]}

1. Life a journey or pilgrimage, Gen. xlvii. 9,
(Jacob) \footnote{'The days of my pilgrimage
are a hundred and thirty years: few and evil, and they have not come
up to the days of the pilgrimage of my fathers.'}, Luke xiii. 33 \footnote{'Nevertheless I must walk
today and tomorrow.'}, John xi. 9 \footnote{'Are there not twelve
hours in the day? If a man walk in the day, he stumbleth not.'}. This
is a thought befitting the beginning of the year.

2. In a journey we have a start and a goal. So
life. Again, in a journey, obstacles—rivers, mountains, etc. So in
life, temptations.

3. Now journeys have different lengths—so
different lives have different lengths; one man dies old, another
young—each life is \emph{long enough} to reach the goal. Each length
according to capabilities—one can go three miles an hour, another
four, etc.

4. The length of each is determined by the length
of light. No one has to travel in the dark, John ix. 4 \footnote{'I must work the works of
him that sent me, whilst it is day: the night cometh, when no man can
work.'} and xi. 9-10 \footnote{'If a man walk in the day,
he stumbleth not, because he seeth the light of this world. But if he
walk in the night, he stumbleth, because the light is not in him.'}—as
one man's journey might be near the pole, another's far
south—different times of year.

5. If we linger or deviate on a journey, the
light goes.

6. Have \emph{we} not lingered or gone out of the
road? Double loss of time if we have to get \emph{back}; and then how
to find the way?

7. We shall have to run.

8. Perhaps a carriage—increase of grace.

9. Now let us think of the past year and the year
to come.

\xsection{February 20 (Sexagesima) [Christ our Fellow-Sufferer]}

1. INTROD.—We
have to labour and suffer, as I said last week \footnote{There are no notes of the
sermon here alluded to.}, but we have this support and consolation, that Christ labours
and suffers with us. This a great subject.

2. Adam fell. God never puts on us more than we
can do; He gives grace sufficient.

3. But it is much more than this. He might have
forgiven and restored us \emph{without} Christ's death; but He has
done so in a more excellent way.

4. The Prince of Wales going into a labour
prison, putting on dress of convicts, having his hair cut, all for the
sake of converting convicts. So—

5. Christ has sought us—but more, for He has
wrought and suffered for and instead of us.

6. Still more; not only He has taken \emph{ours},
but has given us His—the vine and branches—one body, He the head,
Rom. xii. [5] \footnote{'So we, being many, are
one body in Christ.'}; 'Why
persecutest thou me?' [Acts ix. 4].

7. We are all [that] He is—sons of God—full
of grace—heirs of heaven.

8. Is not this sufficient to sweeten labour?

\xsection{February 27 (Quinquagesima) [Communion with God]}

1. God the Creator of all—all things \emph{depend}
on Him.

2. But the happiness of intellectual beings is
not only [in] dependence [upon], but in \emph{union} with Him.

3. This union shows itself in communion—that
is, a fellowship—intercourse of thought, or a spiritual
conversation.

4. The fall of Adam has placed a huge obstacle,
as a wall or a mountain, between us and God, and Christ has broken it
down. He has opened the kingdom of heaven to all who believe. This is
why He took flesh and came on earth.

5. Now this communion requires love and grace on
the part of God, and faith and prayer on the part of man.

6. In His part God is not wanting. His love is as
expansive, as diffusing, as powerfully and constantly overflowing as
the sea, or as the wind, or as the flame, and whereas their expansion
is for evil, that of the Divine Attributes is for good.

7. Now we have instances of this communion
between God and man in Scripture clearly defined.

8. Enoch 'walked with God'—and Noe. What is
meant by 'walking with God' is plain; men who are companions on a
journey \emph{talk} while they \emph{walk}. The two journeying to
Emmaus; our Lord joined them—this was communion.

9. Another image is that of \emph{friend}.
Abraham is \emph{friend} of God, 2 Paralip. xx. 7 \footnote{'Didst not thou our God
kill all the inhabitants of this land … and gavest it to the seed of
Abraham \emph{thy friend} for ever?'}, Isa. xli. 8 \footnote{' … the seed of
Abraham my friend.'},
James ii. 23 \footnote{'And the scripture [Gen.
xv. 6] was fulfilled, saying, Abraham believed God: ... and he was
called the Friend of God.'}. Friends
are in possession of each other's confidence; and Gen. xviii. [17],
'Shall I hide from Abraham what I am about to do?'—and [Abraham's]
intercession for the cities.

10. And so Moses—Num. xvi. 18 and Ex. xxxiii.
11.

11. What was, then, the privilege of the few, for
the Jews were 'servants' in Judaism, is the right of all Christians. \emph{Vide}
Luke xxii. [4], 'friends,' and John xv. 15, when our Lord speaks of 'omnia
quaecunque,' etc.

12. CONCLUSION.—(1)
Those who make friends of the world cannot have this Divine
friendship; (2) Those who have that Divine friendship have a disgust
of worldly friends. 'Jesus Christ the same yesterday, and today, and
for ever.'

\xsection{First Lent [1876] [Sin]}

1. INTROD.—This
is a time when the Church calls upon us to put off sin—both the \emph{reatus}
[guilt] and punishment of it—and the indwelling power.

2. That there is such a thing as sin—something
distinct from everything else—I shall take for granted; it is an
offence against God: 'against Thee only have I sinned.'

3. It is different from a mistake, intellectual,
social, or practical; different from an offence against
beautifulness, etc.

4. I shall take this for granted. And the
question is about its \emph{gravity}—is it of little consequence or
great?—\emph{this} is my point.

5. Now here I shall dwell on one consideration,
viz. of what Almighty God thinks of it, by what He has done in
consequence of it.

6. The GOSPEL.
What is meant by Gospel, and why?

7. It is the coming of God in our flesh, thus to
humiliate Himself—to suffer, to preach, etc., etc. Can it be a light
thing to bring about this? This is indubitable if Christianity is
true—there cannot be two opinions.

8. Nay, He had to suffer—to be tortured—to be
crucified. Texts.

9. His terror: it was more then than the outward
appearance, John xii. [27] \footnote{'Now is my soul
troubled; and what shall I say? Father, save me from this hour: but
for this cause I came unto this hour.'},
Luke xxii. \footnote{The agony at sweat of
blood.}, Heb. v.
[7] \footnote{'Who in the days of his
flesh, with a strong cry and tears offering up prayers to him that was
able to save him,' etc.}.

10. Now what is to be said to those \emph{who ignore
all this}?

\xsection{March 12 (Second Lent) [Hell]}

On future punishment—on hell.

First, we are not fair judges of the malignity of
sin.

1. Because like men who have ever lived in a
mine, and never seen the light of day.

2. Because skill in any art (and so in holiness)
raises the standard.

3. Because culprits make bad judges.

Secondly, God is a consuming fire; sanctity burns
what is not holy.

Thirdly, consider Scripture. St. John Bapt.,
'burn up the chaff'—Matt. xxv., Mark ix., fire—St. Luke, Dives and
Lazarus—St. Paul, St. Peter, St. John, St. Jude.

\xsection{March 19 (Third Lent) [Punishment of Sin]}

1. Recapitulation—God 'a consuming fire' [Deut.
iv. 24]; the nature of things; His nature cannot alter; He cannot
alter it; He cannot make sin a blessedness, or make it other than
antagonist to Him.

2. An atmosphere—we cannot live in water, or
except in air—a sinful soul cannot breathe the atmosphere of the
spiritual world; that atmosphere is fire to it.

3. While in this world punishment is sometimes
delayed; but even here it commonly comes as heathens witnessed, (thus
old men suffer for their youth). But [if not here] certainly
afterwards—a 'fiery indignation' [Heb. x. 27], etc.

4. Not only deliberate sinners, but even God's
own children; and this is what I wish to bring before you,
especially this evening. This is the doctrine of purgatory.

5. Such is the divine law—even after repentance
and reconciliation after a firm faith; yet those whom God loves, who
die in grace, nevertheless suffer for the sins which they did before
their conversion and after it.

6. Matt. v., 'last farthing.'

7. 1 Cor. iii., 'by fire.'

8. (1) Delaying repentance. 'I will repent by and
bye'—but purgatory.

9. (2) Every sin has its punishment.

10. (3) Joy that there is a purgatory.

11. (4) The willing plunge. The content of
purgatory—next to the content of paradise.

\xsection{March 26 (Fourth Lent) How to Escape Purgatory}

1. The first thing to do is to be in constant
union here with our Lord and Saviour; that is, in a state of grace.

2. This we are able to do through the sacraments
and our own care—that care is to avoid mortal sin. Venial sin no one
[can avoid] but by special privilege; but it is our own fault if we
fail into mortal.

3. If we are always in and out of grace (1) shall
we persevere? (2) at least we shall not save ourselves from purgatory.

4. Next, if we are in grace and reunited to
Christ, our penitential works tell—1 Cor. xi., 'judged'; (1)
prayers—'Knock, and it shall be opened'; (2) alms; (3) fastings.

5. Thirdly, indulgences. Still depending on a
state of grace.

6. Fourthly, but better than all, contrition.

7. Two kinds of contrition. With love or with
fear; for God or for self.

8. Anyhow there must be hatred, grief, resolve.

9. But contrition with love does everything; it
saves without the sacraments when they can't be had.

10. St. Vincent Ferrer and his penitent (\emph{Bail},
t. 5).

\xsection{April 16 (Easter Day) [Gifts of the Resurrection]}

1. INTROD.—This
day is the foundation of all our hopes for eternity. Recollect that
our souls will never die; the body dies, but not the soul ever; always
in being. The body dies—rises; but there is no break of continuing
being as regards the soul.

2. What condition are our souls in to \emph{encounter}
immortality? Who is there who, if he can bring himself to think
steadily, would like to live for ever with no better outfit for
eternity than he has? Who would like to go out of this world to
judgment with no clearer conscience than he has?

3. In this tremendous difficulty our Lord came to
be our Saviour. The Son of God came, etc., etc. He died. He took upon
Him all our miseries, and made Himself a sacrifice.

4. And he has gained us great gifts by which to
reverse our state. Salvation consists in five. We have virtues
which God cannot have, \emph{e.g}. faith and hope, but these [five]
gifts are much more; [they make us] 'partakers of the divine nature'
[2 Pet. i. 4]. Five—GRACE and
TRUTH, John i. 16-17 \footnote{'And of his fulness we
have all received, grace for grace … The law was given by Moses;
grace and truth came by Jesus Christ.'}; LIGHT and LIFE;
\emph{ib}. 4, 9 \footnote{'In him was life; and
the life was the light of men ... The true Light, which enlighteneth
every man.'}, and \emph{ib}.
xiv. 6, 'the way,' etc. \footnote{'I am the way, and the
truth, and the life.}
PEACE or JOY,
\emph{ib}. xiv. 27 \footnote{'Peace I leave with you,
my peace I give you ... Let not your heart be troubled.'}. So
in apostolic greetings in the Epistles. All five Divine
Attributes—'partakers of the divine nature.'

5. The first is grace or holiness (instead of
unholiness).

6. The second truth (instead of ignorance),
witness to the truth—faith—hearing.

7. The third is light—\emph{seeing}, \emph{knowledge},
everything clear—as eyes correspond to ears. John viii. 12 \footnote{'I am the light of the
world.'}; \emph{ib}. xii. 35, 36 \footnote{'Walk while you have
light, that the darkness overtake you not ... Whilst you have the
light, believe in the light.'}.

8. Fourth, life, John x. 10 \footnote{'I am come that they may
have life, and have it more abundantly.'}, 28 \footnote{'I give them [my sheep]
life everlasting.'}, \emph{ib}.
iv. 14 \footnote{'The water that I will
give him shall become a fountain of water springing up to life
everlasting.'}, \emph{ib}. v.
24 \footnote{'He who heareth my word
... hath life everlasting.'}.

9. Fifth, peace or joy, John xiv. \footnote{V. \emph{supra}.}, \emph{ib}. xx. Twice \footnote{The salutation 'Peace be
to you,' after the Resurrection.}; joy, \emph{ib}. xv. 11 \footnote{'These things I have
spoken to you, that my joy may be in you, and your joy may be filled.'}, 1 John i. 4 \footnote{'These things we write
to you, that you may rejoice, and your joy may be full.'}.
All five in John xvii. [Christ's prayer for His disciples].

10. Oh how different is all this from the ideas
and language of the world!

11. Let us recollect, to realise it in ourselves
is the only true way of keeping the Resurrection.

\xsection{November 19 (Twenty-fourth Pentecost— Sixth Epiphany) [The First and Second Advents]}

1. The mustard seed.

2. The Church between the two comings of Christ.

3. Those comings both awful, Mal. iii, and Mal.
xxiv. 29, etc.

4. [But] this difference—the first expected;
the latter sudden—Mark xiii.—the ten virgins.

5. The first—Jacob's prophecy [the time fixed, \emph{i.e}.
\emph{when} the sceptre shall have departed from Judah]—Daniel's 70
weeks=490 years.

6. The second [sudden] like the Flood and [the
destruction of] Sodom—Luke xvii. 26-30.

7. Hence 'watch and pray,' Mark xiii.

8. And so St. Paul—a first duty to \emph{wait},
1 Thess. i. 10 ['to wait for his son from heaven'], Rom. xii., and
Heb. x.

9. But it may be said, What difference between
this and waiting for the death of each?

10. Against building and planting—progress.

11. Making this structure and polity of visible
society a god to be worshipped, though the individual dies.

12. No, all we see will come to nought, however
great and beautiful, Isa. ii. \emph{finis}.

\xsection{December 3 (First Advent) [The Second Advent]}

1. The gospel this day a portion of our Lord's
prophecy [of the destruction of Jerusalem and the second coming],
rising out of the apostle's admiration of the beauty of the Temple.

2. It seemed to them too beautiful to be
destroyed.

3. This is the way with men; they wonder at their
own great works; they look up to the great works of their fellows; and
when they are able to trace out the beauty of God's creation, that too
they make an idol, and bow down before their own work, loving the map
[which they have made] of it, and say they have discovered it to be
too beautiful to be broken.

4. It is well for man, an \emph{homuncio}, thus
to think, for he can do a great thing but once; but God destroys His
own works, however beautiful, because He could from His infinite
resources create many worlds each more beautiful [than] those that
were before it.

5. And His own works He regards not, if they have
not that note of \emph{sanctity} which He breathed on them in the
beginning.

6. Therefore when the Jews, His own people, came
to nought, if He did not spare the work of His own hands, much less
[was He likely to spare] Herod's work.

7. Babylon and Nabuchodonosor.

8. Now in this [the present] state of society it
is pride, not open sensuality [which is conspicuous]; \emph{i.e}.
think of the greatness of an army, of a popular assembly, of some
queen's ball. But whenever the world looks imposing and likely to
last, that is the most likely time that it will be brought to an end,
or at least [is the likely time] of some great judgment.

\xsection{December 10 (Second Advent) The Immaculate Conception}

1. The great feast of this church [of the
Birmingham Oratory], being the mystery to which it is dedicated.

2. Describe Adam's state [before he fell], and
original sin. He became a child of wrath—enemy of God—and the
dwelling-place of the evil spirit—[\emph{Lex orandi lex credendi}:
this last point illustrated by the] exorcisms at baptism.

3. Now, could this be the state of the Blessed
Virgin?

4. [It could not be.] How is this proved? By \emph{meditating}
on the Incarnation. (Explain)—deeper and deeper knowledge the saints
have in gazing on the Beatific Vision—and in this life the more we
meditate on divine truths the more we find.

5. She is [the] Mother of God. He not merely
inhabited a man, etc. Doctrine of the Incarnation.

6. Therefore immaculate in [her] conception.

7. Hence the Holy Fathers, etc., and so on to our
times, till Pius ix. has made it an article of faith.

8. It does not, cannot interfere with the supreme
glory of her Son, but subserves it.

\xsection{December 17 (Third Advent) [Signs of the Second Advent]}

1. Sign of our Lord's coming, though we don't
know the day, viz. an apostasy or revolt—[\emph{ho anomos}] [the
lawless one, 2 These. ii. 8]; Antichrist [cf.] [\emph{h\={e} anomia}],
Matt. xxiv. 12.

2. Give circumstances of St. Paul saying
so—belief that our Lord was \emph{then} to come when St. Paul wrote.

3. Our Lord says, Matt. xxiv. [9, 'Then shall
they deliver you up to be afflicted, and shall put you to death: and
you shall be hated by all nations for my name's sake']—contrast of
prosperity, \emph{ib}. 38, ['For as in the days before the flood they
were eating and drinking ... so also shall the coming of the Son of
man be'], and persecution, the greatest persecution, as holy men have
anticipated.

4. But why will they not persevere, as [the]
first Christians [did]? Want of faith (\emph{vide} verse 12). We can
do all things by faith if we have faith; but false reason cuts at the
root (\emph{vide} verse 24, false prophets).

5. Sophistry and false reason—even the elect.

6. This all is opening on us—like the last age.

7. But as sun shining through clouds, or as a
dying man kept alive by prayer, always going and never gone, so (for
the chance of more conversion and more elect), the world ever dying.

8. Alas! the next generation—young people, I
fear for you! \footnote{Father John Pollen, S.J.,
remembers the great impression made on him by this sermon, and
especially by the pity expressed for the rising generation. He heard
it as a boy at the Oratory School.}

9. Let us at this time of year pray that as
Christ on His first coming came with preparation, so we may be
prepared for His second coming.

\xsection{December 24 (Fourth Advent) [The First Advent]}

1. 'A thousand years as one day.' Quote 2 Pet.
iii. 7-8 \footnote{'But the heavens and
earth, which are now, by the same word are kept in store, reserved
unto fire against the day of judgment and perdition of the ungodly
men. But of this one thing be not ignorant, my beloved, that one day
with the Lord is as a thousand years, and a thousand years as one
day.'}.

2. Illustrated in preparation for Christianity.
Abraham two thousand years before His coming.

3. He was told, 'In thee [shall all the kindred
of the earth be blessed,' Gen. xii. 3], but his faith tried by the
delay.

4. Mount Moriah the illustration of that faith
which was against sight, for Isaac was the child of promise.

5. This a type of all holy men; instances in Heb.
xi. 13 \footnote{'All these died
according to faith, not having received the promises, but beholding
them afar off.'}.

6. Not only did not our Lord come, but all events
seemed to look the other way.

7. Prophets said again and again, \emph{Persevere in
faith}; but what chance was there when the people rebelled and
apostatised, and there were captivities, etc., and at last deportation
of the whole people to Babylon?

8. Elias and the seven thousand—a remnant.

9. Solomon, 'There hath not failed [so much as
one word of all the good things that he promised by his servant
Moses'], 3 Kings viii. [56]; Isaias to Hezekiah, Isaias xxxvii. 32 \footnote{'For out of Jerusalem
shall go forth a remnant, and salvation from Mount Sion.'}; Habacuc iii. 18 \footnote{'But I will rejoice in
the Lord; and I will joy in God, my Jesus.'}.

10. At length Herod king; a heathen—but man's
extremity God's opportunity. Christ at length came.

11. Apply this to these times—and personally to
individuals.

[A line along the margin shows that § 8 was to
be inserted in § 7, presumably after 'apostatised.']

\xsection{December 31 [The Past not Dead]}

1. We are accustomed to keep the beginning of the
year, but not the ending; we congratulate each other on the new
year, but we let the old year go its way. Why is this?

2. It is, first, because the old year, whether it
was a happy year or unhappy, suggests thoughts of pain. If unhappy,
that is painful; if happy, it is [painful] because it is gone.

3. Also, because the past year is dead, and
whatever is dead naturally inflicts pain upon us. And therefore we
turn to the \emph{new} year with hope.

4. But is it really dead? No; in one sense it is
awfully alive. All things live to God—the past as well as the
present. No; \emph{to us only} the past year is dead, and to us also
one day it is to [be] alive again. On the last day the books will be
opened, chronicling all events. Quote Apocalypse xx. 12 \footnote{'And I saw the dead,
great and small, standing in the presence of the throne; and the books
were opened: and another book was opened, which is the book of life:
and the dead were judged by those things which were written in the
books, according to their works,'}.

5. Think of the number of events—\emph{e.g}. in
the newspapers—yet they are nothing to the sum-total.

6. Every soul has its history; every soul is
immortal and independent.

7. What are the events of the year but a history
in millions of souls of the unceasing warfare between good and evil.
We talk of battles in the world, etc., but what are they?

8. How many have gone right? how many wrong? how
many turning-points for life? how many have died good and bad, young
and old?

Photographs; light from distant stars not yet
arrived here. How much we need God's protection; the future quite
dark.

\xsection{May 13, 1877 (After Ascension) [Particular Providence]}

1. INTROD.—[The
Ascension] the end of the miraculous series of events which our Lord's
life comprises. From birth to Ascension, as is said in the Creed,

2. Till He shall come to judge.

3. The world goes on by fixed laws, and they are
such as in themselves are good, and subserve and proclaim His General
Providence. All these contrivances and final causes, etc. But nothing
personal in this—a cold system. He does not speak. No encouragement
to us to speak to Him—Acts xiv. 17 \footnote{St. Paul's address to
the Iconians: 'Nevertheless he left not himself without testimony,
doing good from heaven, giving rains and fruitful seasons, filling our
hearts with food and gladness.'} and xvii. 26-27 \footnote{St. Paul's address to
the Athenians: 'And hath made of one all mankind to dwell upon the
whole face of the earth, determining appointed times, and the limits
of their habitation. That they should seek God, if haply they may feel
after him, or find him.'}.

4. But He has been more merciful than this—\emph{the
hairs of our head numbered}. \emph{Particular Providence}
everywhere and always. How was He to show this particular providence?
By suspending the laws [of Nature] by miracles.

5. This He did. But that particular providence He
has mercifully brought [out] in a distinct form, first in the Mosaic,
and then in the Gospel economy. He has suspended His laws.

6. Especially when, after suffering, He ascended
to heaven.

7. He promises us heaven. He has gone first to
prepare a place for us individually there.

8. \emph{Patria}—Heb. xi. [13], pilgrims \footnote{'All these died
according to faith, not having received the promises, but beholding
them afar off, and saluting them, and confessing that they are \emph{pilgrims}
and strangers on the earth.'}; Eph. ii. [19] \footnote{'You are no more
strangers and foreigners.'}.

9. \emph{Sursum corda}—Col. iii. 1 \footnote{'Therefore if you be
risen with Christ, seek the things that are above, where Christ is
sitting at the right hand of God.'}.

10. We think of meeting our \emph{friends} in
heaven; we do not think of Him who is the best of all friends.

11. Don't say we don't know Him; Gospels,
especially Gospel of St. John, bring Him close to us.

12. \emph{Sursum corda}—Col. iii. 1. Pray
without ceasing, after pattern of Luke xxiv. 52 \footnote{'And they adoring went
back into Jerusalem with great joy.'}, Acts i. 14 \footnote{'All these were
persevering with one mind in prayer.'}.

\xsection{May 27 (Trinity Sunday) [The Holy Trinity]}

On this day we close our celebration of the
merciful truths of the Gospel by a solemn commemoration of the Holy
Trinity,

1. Lest we should forget who and what God is.

2. He has so humbled Himself, \emph{e.g}. our
Lord's human life, also the Holy Ghost—\emph{a gift} poured
out—received—quenched—[\emph{al}.] extinguished—grieved.

3. Therefore the Church appoints this feast, that
we may have that holy fear, and awe, and wonder at Him which
becomes His greatness, while we believe in Him and love Him.

4. State the doctrine—that there is a Divine
Trinity or Triad, a Divine Three in heaven. Each is God in the fulness
of Divine Attributes, yet there is only One God, as the Creed
says—this is the beginning and end of the Holy Truth—we cannot say
more or less.

5. Three devotional sentiments towards
it—faith, fear, love—all three feelings together.

6. Faith—that there should be a mystery,
congruous. He might not have told us, but if not, still we might be
sure there was one. Unitarians—\emph{Credo quia impossibile}. I
would not believe in a God who had no mysteries.

7. Awe, as I have said [secs. 1-3].

8. Love; the mystery no difficulty, for \emph{each
part} of it is clear; the most ignorant can confess God the
Father—then [God] the Son, etc. \footnote{See p. 158, footnote.} Though how all three [propositions] are true together, and why
the second does not contradict the first, etc., the wisest cannot
know. And in heaven we shall know all.

9. The Mystery brings before us the peace, as
well as the love which [we] shall have in heaven.

\xsection{June 24 (Fifth Pentecost) [Anger]}

1. INTROD.—Epistle
[1 Pet. iii. 8-10], Gospel [Matt. v. 20-24].

2. The anger thus spoken of by our Lord \footnote{'Whosoever is angry with
his brother, shall be in danger of the judgment.'} is a sin against \emph{charity}; it is something more than
indignation, for it is a \emph{personal} feeling, the consequence of
some slight or injury to oneself.

3. And more than want of self-command. Passion or
irritableness, which is parallel to concupiscence, a sin against self.
Passion or irritableness leads a \emph{child} to beat the ground. To
shout, to abuse, to strike is a relief. All this is sinful, but not [\emph{in
se}] against charity; one may really love the person who provokes
us.

4. Yet, though it [anger] is not this common
affection, in spite of the phenomena being the same, it is a common
sin, and far worse than an affection.

5. It shows itself in this, that it does not go
off or evaporate; it remains as spite, resentment, a grudge, a desire
of revenge—and who will say this is not common?—in a feeling of
repulsion, alienation, hatred.

6. \emph{How} wrong it is we shall see by the
contrast of the \emph{Christian} feeling, as darkness is understood by
light—how beautiful the generosity, the nobleness of returning good
for evil! Joseph and his brethren; David with Saul.

7. I have said it is a personal feeling; but
there is a kind of hatred, which is only partially personal, or only
at first, but when the keen personal feeling is gone, leaves a habit
of hatred and repulsion—a dull negative feeling. This too is against
charity.

8. It may be said, indeed, How can I help it? I
don't like the man, as I don't like a certain taste or sound,
etc.

9. On the other hand, you must feel towards all
men as those you can bear to meet in heaven.

10. This one effect of purgatory, to burn away in
every one of us that in which we differ from each other.

\xsection{August 5 [The End of Man]}

1. Why are we placed here on earth? This is a
question which comes often to children, and is the beginning of their
responsibility.

2. Too often the question ceases to be asked by
the young soul; but it drowns the thought in the levity of the small
world of children around it.

3. But I repeat it, my brethren, think of it
now—\emph{Why} are you here in this world? Were you put here merely
to eat, drink, sleep, etc., for a certain [number] of years, to marry,
to grow old, to die? Were you put here merely to get on in life, to
make a fortune and a name, to gain power, influence, to be in a
position to gratify ambition, etc.? You know you have a higher end
than this.

4. Now consider what the real reason is. You were
put here to prepare yourself for a higher and eternal state; and for
this all the riches, power, name—all the cleverness, sharpness and
knowledge you may have or acquire, nay, I will say, all the industry,
all the affectionateness, all the good-heartedness you have by nature
(though these qualities are entirely good)—will not avail at all.

5. You are come to make the raw material of your
souls into (as I may say) a vessel of honour for the Lord's house
above.

6. Consider the instance of various trades on
earth, and you will understand—bread, pottery, moulding, and the
fine arts—a building, a statue.

7. So there is the raw material of your soul—it
is called in Scripture the flesh; it is human nature in the rude
condition to which Adam's sin has reduced it. Take the instance of the
brute animals—you will in a measure understand what human nature
is—the passions, etc.

8. On bringing the soul into shape—the small
trials of every day.

9. How different good old men are from what they
were when young. So, on the other hand, you can't help moulding
yourself. Woe to you if you mould the wrong way.

\xsection{June 9, 1878 (Pentecost) [The Coming of the Holy Ghost]}

1. INTROD.—The
Holy Ghost, whose feast is today, is God. But God is the object of
every day's devotion, not of festivals. Saints, our Lord's Humanity,
have festivals. What is today? (This is so large a subject that I fear
I shall hardly say all I wish this morning.)

2. God is in Himself all good, and not only all
good, but one characteristic of that goodness is the attribute of
communicating that goodness to all His creatures. He gives Himself to
them in order that they may in their measure partake of that
perfect[ion] which is He.

He is not only the Giver but the gift, and this
day we thus view Him as the gift; for, considered [as] the great and
heavenly gift of all good, He is the Holy Ghost; therefore this day is
the solemn commemoration of that mercy of God by which He has given
Himself to us. The gift—(keep to that—say if we have peace, etc.,
etc.)—and this a feast of St. Philip, etc.; \emph{vide} sermon on
Whitsunday, p. 146.

3. But further. God the Holy Ghost is not [only]
the great gift, but the promised gift. His goodness is especially
shown in His giving us Himself in spite of difficulty and resistance.
This is a world of sin and evil, created good by Him in the beginning,
and therefore His goodness is specially shown in reparation, and in
the triumph of His goodness over evil at the end of a long contest
occasioned by sin. And here is another aspect of this day—it
commemorates the end of a great and long course of providence, the
accomplishment of promise and prophecy. It was not only a gift, but a
promised gift. Hence our Lord speaks of it as 'the promise of the
Father.'

4. A long process involving in it the mission of
the eternal Son to mankind; the world had to be made fit in the
fulness of time—meanwhile the chosen people, etc.

5. Promises and prophets—a long expectation.

6. And when our Lord had come, then, not for a
long time, but still for a time, for forty days, the suspense
continued.

7. The apostles in expectation, 'Lord, wilt thou
at this time,' etc.

8. At last Whitsunday came; wind and fire;
Malachias iii., 'Fire upon earth'—Holy Ghost by fire.

\xsection{July 21 and 28 (Sixth and Seventh Pentecost) On the New Creation}

I

1. We find from St. Paul that \emph{its} life [\emph{i.e}.
the life of the New Creation] is gratitude for our Lord's sufferings
for us, 2 Cor. v.; Gal. ii., sin.

2. Gratitude implies (and requires) (1) sense of
sin; (2) faith.

3. Leads to (3) hope; (4) hope implies fear.

4. And gratitude is a kind of love.

II

But love is charity, such as is necessary to
fulfil the command for 'eternal life,' 'Thou shalt love the Lord thy
God with,' etc.,

Involves

1. Gratitude.

2. Likeness to God. Like loves like, \emph{i.e}.
love of appreciation.

3. The love of friendship—'Abraham the \emph{friend}';
John xv., 'friends'—which involves a mutual consciousness of
love—John xxi., 'Lord, thou knowest,' etc.

4. Companionship—'Walk before me'; 'Walk with
me'—journey to Emmaus.

5. \emph{Dilectio}—preference.

\xsection{(Twentieth Pentecost) [Faith and Thanksgiving]}

'Unless you see signs and wonders, you believe
not,' John iv. 48.

1. INTROD.—This
gospel opens a subject too long to treat well today, especially as
this day has subjects of its own, one of which I cannot pass over.

2. How is it that our Lord seems to accuse the
nobleman of unbelief?

3. He certainly does in some sense, but in what
sense?

4. His disbelief was not as if he did not believe
our Lord's power; he did not try our Lord's power to get \emph{evidence}
(see the narrative), but he was (naturally indeed) so earnestly set
upon his one desire, his son's recovery, that he did not show
resignation. He urged our Lord to make haste.

5. Contrast Martha and Mary—'Lord, he whom thou
lovest is sick' [John xi. 3]. And the event showed how fully they
might trust—for [though] He did not make haste, still he who was
even dead was raised again.

6. Ephesians v. 19-20 \footnote{'Speaking to yourselves
in psalms and spiritual canticles, singing and making melody in your
hearts to the Lord; Giving thanks for all things in the name of our
Lord Jesus Christ to the Father.'}; Phil. iv. 3 \footnote{'In everything, by
prayer and supplication, with thanksgiving, let your petitions be made
known to God.'};
1 Thess. v. 18 \footnote{'In all things giving
thanks; for this is the will of God in Christ Jesus concerning you
all.'}; 2 Cor.
vi. 10 \footnote{'As sorrowful, yet
always rejoicing.'}.

7. I am led to these remarks by the thanksgivings
to be made in this day's service for the answer to the prayers we have
been offering by direction of his Lordship for a blessing on this year's
harvest and the trade and manufactures of this country.

8. The distressed father in the gospel asked for
a miracle; we do not, but our cases are the same—whether by miracle
as he, or by providence as to us, both he and we prayed to God for an
object. We have both severally received what we ask for.

9. His distress and ours.

10. He showed his gratitude, and so, I trust,
shall we.

11. He was admonished—(go through his case).

Continue

\xchapter{Catechetical Instructions}

Back

\xsection{August 28, 1849 The Creed—De Deo—I}

1. INTROD.—On
the \emph{articles} not embracing the whole truth—a number of
others, but we believe whatever the Church teaches; if you are in
difficulty this the great rule—(enlarge—come to be taught, no use
coming else). Why called the Apostles' [Creed]?

2. By 'God' we mean one who is all-perfect—this
is the only idea of God; therefore the heathen gods are called
'not-gods.' \footnote{Jer. ii. 11; ib. v. 7.} They do not
answer the idea.

3. (Not pleasant to inquire into the proofs; it
is an irreverence, therefore pass over it lightly \footnote{He would also have felt it
superfluous in 1849 'to inquire into proofs.' 'Whatever my anxiety may
be about the future generation, I trust I need at present have none in
insisting, before a congregation, however mixed, on the mysteries or
difficulties which attach to the doctrine of God's existence,' etc.—\emph{Discourses
to Mixed Congregations}, p. 265.}.) We are obliged to believe it, for else, was this world for
ever? If not, who created it? and who created its Creator, etc.? Thus
it is simplest, we cannot help believing in a God, nay, believing the
most incomprehensible attribute, viz. that He had no beginning.

4. God must be all-perfect, for (1) He has made
the world, and therefore must be more perfect than that which He made.
He has given to it of His fulness, therefore is there wisdom, power,
beauty, etc., in the world (think of the human soul)—then He much
more so. This thought alone gives us an indefinitely high notion of
God, considering the extent and wonders of creation.

5. Next, think that He created from nothing.

6. Go through His attributes.

7. The one singled out is 'Almighty,' and for two
reasons: (1) Because we are to consider, not the time \emph{before}
creation, but creation and after; (2) because it is the \emph{Creed},
which has reference to God's omnipotence.

8. 'Creator of heaven and earth'; (1) Creating
from nothing; (2) number of attributes discovered from the material
world; if there be beauty, He is beautiful; if Spirit, He Spirit, etc.

9. Creator of angels, men, things inanimate.

10. Proof of all this from above: (1) Two most
incomprehensible [mysteries] from the nature of the
case—eternity \emph{a parte ante} [\emph{i.e}. before creation], and
creation from nothing, and if these, what shall we not give to Him?
(2) Idea in human mind, in conscience, and hence His moral attributes
holy, just, ever vigilant, all-seeing. (3) The visible world and all
the senses. (4) The variety of attributes which the world and the soul
shows—wisdom, greatness, and minuteness, grace and beauty,
comeliness, goodness. (5) Providence—therefore all good. Providence
as shown in general [and] in our own life.

11. Doctrine of the Holy Trinity.

\xsection{[September 4] De Deo—II}

1. Recapitulation—from the conscience and our
personal history we have all the great doctrines about God which so
much concern us, and which would suffice for our believing in Him, \emph{though
there was no external world}, viz. holiness, justice, omnipresence,
ever-watchfulness, mercy, and future retribution.

2. About the argument from the external world,
and why it is dangerous at this day \footnote{See Note
19, pp. 343-4.}; \emph{because it tells us nothing about sin}; it [the world]
was made \emph{before} sin.

3. The attributes the external world
adds—infinite power, wisdom, skill, etc.

4. Additions: (1) From the nature of the case one
perfection implies another. (2) Eternity \emph{a parte ante}, creation
from nothing, etc., from the nature of the case.

5. How is this faith, if it is gained thus by
reasoning?—(explain—how reason goes a certain way to draw a
conclusion)—that is not faith—(explain the process)—\emph{desire of
the truth}.

6. How faith is a completion through grace; this
grace may be given to those who know not revelation; but if revelation
comes in their way it will lead them to it.

7. This is why Catholics hold apostasy so
wretched a thing—it is not the mere change of opinion, but a going
against grace.

\xsection{September 11 De Angelis—III}

1. Angels the first work of God.

2. Different from the human race, as created all
at once; not from a pair—myriads—almost an infinite number like
the stars; nine ranks, yet there may be \emph{more} unknown.

3. Immaterial—their appearing in human
form—Raphael in [book of] Tobias.

4. Incorruptible; immortal in their nature.

5. Their knowledge; perfect from the first, not
learning by \emph{discursus}, first one thing, then another; knowing a
thing wholly at once. They do not know the future—nor men's hearts,
except from divine revelation. God alone knows these—(bad angels
tempt by objecting evil or exciting the thoughts).

6. Created in grace. And first, What is \emph{grace}?—something
beyond \emph{nature}. Nothing can love God really and know Him,
and attain to heaven, without a gift beyond itself.

7. The fall of the angels—\emph{pride}—as sin
of thought, for it was \emph{all} the sin they \emph{could} fall into,
being spiritual natures. Pride is relying on oneself for happiness,
not on God's grace; envy would follow, viz. against men. Their \emph{naturalia
remanserunt integra}—they lost (1) supernatural beatitude,
intuitive vision of God \footnote{\emph{I.e}. in prospect.};
(2) justice and grace; (3) their intellect darkened, and their will
confirmed in evil.

8. They range the earth and dwell in the air.

9. They do not properly possess the mind of the
possessed except indirectly, by raising fantasies, etc. Of sinners and
unbaptized, the \emph{soul} is not possessed but ruled; hence even
innocent persons are sometimes possessed.

10. Good angels were confirmed in grace.

11. Guardian angels for the faithful, and perhaps
for the \emph{infideles et reprobati}.

\xsection{September 25 De Mundo Visibili et Brutis—IV}

1. On the visible world, as described in Genesis;
created by degrees.

2. Created good—in what sense?—relatively
good. A better world conceivable; all creatures, as such,
imperfect—grace perfects. He has not given His grace to all
[creatures].

3. This seen especially in the phenomena of brute
animals; how far they are like men—in structure and
make—remarkably like—one idea \emph{physically}—for man is an
animal and something [besides]. On the brute animals in contrast with
men. They have not souls \footnote{He is using the word soul
in the popular sense.}.

4. First: They do not know they exist; they
cannot \emph{reflect}; they are like our minds in dreams or vacant
vision or hearing; they don't anticipate death; they may see other
brutes killed, but don't fear; no burying bodies—no despair of life,
or suicide.

5. Second: Sensations—nothing more; they hear
sounds, but do not know what is meant by them; they cannot speak.

6. Third: They act upon their sensations in a
mechanical way, as by smell, etc.

7. Four: [They] eat and drink, etc., for the sake
of eating and drinking, etc.—not in excess—still because
mechanical.

8. Five: No abstract ideas—of justice, truth,
etc.—no idea of duty.

9. Six: No governing power in their affections.

10. Seven: No perfectibility.

Hence they cannot sin, though they have impulses,
etc., which in man are sin. Man is apt to argue that a thing is not
sinful if it is natural, whereas it may be sinful \emph{in him}, who
has means to prevent it \footnote{Cp. \emph{Discourses to Mixed
Congregations}, p. 149.}.

\xsection{October 2 De Gratia et Gratia Amissa—V}

1. INTROD.—After
brute animals we come to man.

2. Man a compound being—how soul and body can
be together a mystery \footnote{Cp. Parochial and Plain
Sermons, vol. iv. p. 283.},
and what follows, though strange, not so much so [\emph{i.e}. not so
mysterious]—compounded of body which aims at sensible good, and soul
which aims at spiritual. Hence \emph{grace given} that there might not
be war.

3. This the state of Adam—original justice,
grace—moreover, absence from death and the \emph{vestigia mortis}.

4. Adam sinned; he could not have had a lighter
trial. What was his punishment? A stripping of this grace, etc.

5. Enlarge on this as humbling man, [viz.] as a
mere creation of God he is very imperfect—\emph{Pulvis es}, etc.
('Dust thou art, and to dust thou shalt return']—[but now] he is not
[as] in that state he would have been in had he not had grace.
[Parable of] good Samaritan [the wounded man being typical of our race
after the Fall]—evil of souls, four wounds—of body, death, and
disease. As living things gasp and die under an exhausted receiver,
though air is not part of their nature.

6. This state of man after the Fall—it is
called original sin—and since Adam's sin is imputed, it shows itself
in the above privations.

7. No one but has grace enough to save
him—sufficient grace—but we cannot add up correctly many times
running \footnote{See \emph{Discourses to Mixed
Congregations}, pp. 128-9.}. Things needed
not be yet are—efficacious grace.

8. Process of grace; actual grace, drawing a man
on, one grace improving [improved upon?] leads, \emph{de congruo},
to another—so to contrition with desire of the Sacrament, which
justifies, viz. habitual grace.

9. Children, habitual grace—but still actual
grace is needed besides.

10. Grace of perseverance.

11. Our Lady without original sin.

12. The effect of this to humble us—creatures
imperfect.

\xsection{October 9 De Redemptione—VI}

1. Sin leads to the doctrine about Christ: 'And
in Jesus Christ, His only Son, our Lord … born of the Virgin Mary,' \emph{i.e}.
God coming in the flesh.

2. Consider the \emph{cumulus} of sin—all the
sins of every individual through centuries, and to the end of the
world. The offence to God, how great!—infinite—though the malice
finite.

3. God might have condemned all men—He might
have pardoned all—and that without any satisfaction; but He
determined to take a punishment \emph{equal} to what their sins
deserved. Now man could not pay this, and therefore Christ came, who
was God.

4. God passed over the angels who fell. He looked
lovingly upon man, His youngest creation, and, as great doctors teach
\footnote{Though by no means
generally: see \emph{Sermons to Mixed Congregations}, p. 321 note.}, He would have become
incarnate even if man had not fallen, though He then would not
have suffered, viz. to show the \emph{glory} of His Attributes in a
created nature.

5. Christ's coming prophesied of through so many
ages. Holy men looking out for Him.

6. Why His coming delayed—to show God's free
decrees.

7. Jesus—(explain the word): (1) Saviour, Phil.
ii. [9-10] \footnote{'For which cause God
also hath exalted him, and hath given him a name which is above all
names: that in the name of Jesus every knee should bow,' etc.}, Acts iii.
[6] \footnote{'In the name of Jesus
Christ of Nazareth arise and walk.'},—the own name of
a person always conveys tenderness and familiarity. (Saviour.) (2)
Joshua His type—all saints have delighted in it [viz. the name
Jesus].

8. Christ—Prophet, Priest, and King.

\xsection{October 16 De Trinitate—VII}

1. INTROD.—The
mention of the Redeemer leads us to a great mystery.

2. It is to be expected that there are mysteries;
we cannot tell [beforehand] what—and those which \emph{are} are sure
to surprise us, and we say: 'I expected some, but not this, for this
is so strange.' \footnote{'I consider that this
mysteriousness is, as far as it proves anything, a recommendation of
the doctrine. I do not say that it is true because it is mysterious;
but that if it \emph{be} true, it cannot help being mysterious. It
would be strange, indeed, as has often been urged in argument, if any
doctrine concerning God's infinite and eternal Nature were not
mysterious. It would even be an objection to any professed doctrine
concerning His Nature, if it were not mysterious.'—Parochial
Sermons, vol. vi., 'Faith without Demonstration.'}

3. A mystery is but a mark of infinity. \emph{Vide}
published discourse \footnote{'The outward exhibition
of infinitude is mystery; and the mysteries of nature and grace are
nothing else than the mode in which His infinitude encounters us and
is brought home to our minds,' etc.—\emph{Discourses to Mixed
Congregations}, p. 309.}.

4. We may be sure that every apparent \emph{explanation}
is a mistake, a heresy. We must begin by confessing it unintelligible.

5. Mystery of Trinity, briefly put.

6. On the analogy of the elements—published
discourse.

\xsection{October 23 De Filio Dei—VIII}

1. INTROD.—On
the divinity of our Lord.

2. God from God, Light from Light;
analogies—the word, the sun, and

3. Passages of Scripture—Proverbs; wisdom—Isaias
ix., John i., Phil. ii., Col. i., Heb. i.

4. Mistakes of Protestants \footnote{\emph{Ib}. pp. 346 ff.}.

\xsection{October 30 De Dominio Seu Regno Christi—IX}

1. INTROD.—Our
Lord in the first place God; but also, He has redeemed us with a
price.

2. Hence contrasted to this world, which is the
(usurped) kingdom of Satan—god of this world.

3. Contrast the two, the mediatorial kingdom of
Christ [and the kingdom of Satan] as in the Two Standards, beginning
with John xvii.

4. An empire—(explain what an empire
is)—Psalms ii., xliv., lxxi., lxxxvii., Isaias xliv., liv., lx.,
Apoc. xix.

5. (Contrast the two as in the Two Standards.)
Prophecies—lamb and lion, Isaias xi. 6 \footnote{'The wolf shall dwell
with the lamb, and the leopard shall lie down with the kid; the calf
and the lion and the sheep shall lie down together; and a little child
shall lead them.'}, Isaias ii. 2 \footnote{'And in the last days
the mountain of the house of the Lord shall be prepared on the top of
the mountains,' etc.}.

6. Spreading by meekness—unlike any other
empire: strong in weakness.

7. Exemplified at this moment. State of the Pope.

8. Yet wars, etc. \footnote{Probably in reference to
the crusades, and perhaps to the military defence of the Papal states.} Yes, but the \emph{strength} is not through war, etc. Explain
therefore 'gathering of every kind' in but not of [the world].

9. Contrast the kingdom of Christ and Satan as in
the Two Standards.

\xsection{November 6 De Nativitate Christi ex Virgine—X}

1. INTROD.—'Conceived
by the Holy Ghost, born of the Virgin Mary.'

2. Predestination of the Blessed Virgin, even
before the foresight of the Fall.

3. Gen. iii. [15] \footnote{'I will put enmities
between thee and the woman, and thy seed and her seed; she shall crush
thy head,' etc.}—\emph{inimicitsas} [enmities]; \emph{ipsa} [she].

4. Parallel of Eve and Mary—this idea of a
woman kept up through Scripture in types, though the relationship to
the Shiloh [Gen. xlix. 10] is not always preserved—(1) Sara, Ishmael
scoffing [at]; (2) Mary, Moses' sister; (3) the canticle of canticles.

5. Isaias vii. [14] \footnote{'Behold, a virgin shall
conceive,' etc.}.

6. When our Lord came, perhaps no need of further
notice; yet (for our Lady did not come forward at once into public
view—\emph{e.g}. in Catacombs) in prophecy, Apoc. xii., [she is
prominent] to the end of time.

7. Particulars. Christ might have been born in
the ordinary way, but other way more fitting—Immaculate Conception.
[Her question to the angel], 'How shall this be?' [Her] vow of
virginity—first who did so [\emph{i.e}. made this vow]—why
married; a true marriage.

8. Conceived of the Holy Ghost—all works belong
to the Three Persons [of the Blessed Trinity]—but as wisdom is
attributed to the Son, etc.

9. Christ had \emph{all} grace from first—did
not grow in grace.

10. [Our Lady suffered] no pain in child-bearing;
Eve \emph{é contra}; hence in representations of our Lady [before the
manger] she is made kneeling, etc.

11. Ever virgin.

12. Mother of God—this to secure the doctrine
of the Incarnation. \emph{Vide} published Sermon \footnote{'The Glories of Mary,'
etc., in \emph{Discourses to Mixed Congregations}.}.

\xsection{November 13 Passus et Crucifixus—XI}

1. INTROD.—\emph{Sub
Pontio Pilato}.

2. This is introduced to mark the fulfilment of
prophecy, which fixed the time.

3. Forty weeks—four empires—Genesis xlix.—fortunes
of Judah—Herod—the Romans.

4. Hence the crucifixion instead of stoning. Thus
other prophecies fulfilled.

5. On the 'fulness of time'—one (Roman)
empire—universal peace—Greek philosophy—[heathen] religions worn
out.

6. Another reason for 'Pontius Pilate,' viz. to
prove the historical reality of our Lord's coming against Docetae,
etc. Nay, in last century Dupuis with his three hundred years sooner.

7. On the contrary, God \emph{was} man; God
suffered, died, was buried, etc.; but not as God, but as man—in His
human nature.

8. Still it was His will who is Highest to make
Himself lowest. Indeed, it is in all its parts the most awful of
mysteries—death of the cross: (1) what hanging now is, yet
indefinitely worse—the ignominy of the position; as we fix noxious
birds up; (2) nakedness; (3) contempt and mockery—His disciples
leaving Him; (4) no part of the body without its suffering; (5) His
delicate make more susceptible of pain.

9. Pains of His soul—the bloody sweat—no
support from God or from sense of innocence; feeling of guilt; feeling
of responsibility.

10. Yet the cross our triumph—sanctified by Him
who hung on it; predicted under the [figure of] brazen serpent.
It is now a means of grace.

\xsection{November 20 Mortuus, Sepultus—XII}

1. INTROD.—'Dead
and buried; He descended into hell.'

2. \emph{Dead}. He died for our sins—but also
because He came subject to the laws of our fallen nature; man, though
not naturally immortal, was not to die. 'Dust thou art,' etc.—even
our Blessed Lady, Enoch, Elias, and so Christ.

3. Hypostatic union preserved even in blood,
which after the Resurrection was all gathered up.

4. Buried. This mentioned to show He was dead.

5. Hell eternal prison; purgatory; Limbo Patrum.

6. He went there, not as the others, but to
triumph and take them out.

7. And so on the third day He rose again.

8. Even the instruments of the Passion a triumph;
the cross—its meaning changed: He sanctified it. History of its
finding.

9. The goodness of God not only in saving us, but
in condescending to our weakness in religion; difficult to form \emph{spiritual}
ideas of God—the populace and [the] philosophers. Hence He has
deigned [firstly] to take a body, [and secondly], with [it] a
(personal] history \footnote{'While, then, Natural
Religion was not without provision for all the deepest and truest
religious feelings, yet presenting no \emph{tangible history} of the
Deity, no points of His personal character (if we may so speak without
irreverence), it wanted that most efficient incentive to all action, a
starting or rallying point,—an object on which the affections could
be placed, and the energies concentrated.'—\emph{Oxford University
Sermons}, p. 23. The italics are our own.}.
And this is perpetuated in the two great devotions of the
Blessed Sacrament and of our Lady.

\xsection{November 27 Resurrexit et Ascendit—XIII}

1. INTROD.—These
articles are also mysteries in the Rosary. You should be familiar with
the narrative of them in Scripture—(read it).

2. Our Lord remained forty days on earth—why
forty? the flood ['was forty days upon the earth,' Gen. vii. 17]; [the
Israelites] forty years in the wilderness; Moses (forty days) on the
mountain; Elias ['walked in the strength of that food forty days and
forty nights unto the mount of God,' 3 Kings xix. 8]; our Lord
[fasting forty days] in the desert; hours of His death [forty hours in
the tomb].

3. His state [after the
Resurrection]—wonderful; He came and He went. His dwelling, perhaps
though [\emph{some word has been omitted}] on a terrestrial paradise
where [were] the bodies of the saints who were raised with
Him—especially [also] Enoch and Elias, who for their delay claimed
the sight of Him.

4. What He taught—the constitution of the
Church, the orders of the Hierarchy, the matter and form of the
Sacraments.

5. At length He ascended at midday; He was taken
down from the cross and buried in the evening— rose again in
the morning. This was not a proper place for His glorified body; He
ascended of His own power as God and as man. [See] \emph{Catechismus
Romanus}.

6. '\emph{Sitteth} on the \emph{right}
hand'—(explain).

7. There He sits as our one Mediator and
Intercessor—quote Rom. viii. 34 \footnote{'Christ Jesus that died,
yea, that is risen also again, who is at the right hand of God, who
also maketh intercession for us.'}, Heb. vii. [24-25] \footnote{'But this, for that he
continueth for ever, hath an everlasting priesthood. Whereby he is
also able to save for ever them that come to God by him, always living
to make intercession for us.'}.
In what His intercession consists—in presenting His human nature.

8. Difference between [His and the] Blessed
Virgin's intercession; He is God; she is powerful through prayer;
hence we do not say to Him, \emph{Ora pro nobis}.

9. Also He ascended to fix our minds on heaven in
love; to exercise our faith in One who is absent; to give ground for
our hope.

10. Conclusion of foregoing.

\xsection{December 4 Inde Venturus Est, etc.—XIV}

1. INTROD.—['Thence
He shall come to judge the living and the dead.']

2. Particular judgment—some think the soul is
not taken up to Christ literally, for it [\emph{i.e}. this] is
introducing the wicked into heaven, but only intellectually.

3. The soul of the just goes to purgatory, unless
a Saint; of the sinner to hell.

4. General judgment at the end of the world: it
will come suddenly.

5. Signs previous, though unheeded then.

6. Preaching of Gospel all over the earth; now
this has in great measure been done.

7. Apostasy—love of many waxing cold—1 Tim.
iv. [l] \footnote{'Now the Spirit
manifestly saith, that in the last times some shall depart from the
faith, giving heed to spirits of error,' etc.}, 2 Tim. iii.
[1-2] \footnote{'Know also this, that in
the last days shall come dangerous times. Men shall be lovers of
themselves,' etc.}.

8. Antichrist, 2 Thess. ii. [3] \footnote{'And then that wicked
one shall be revealed,' etc.}.

9. Fire burning up all things, and becoming the
purgatory of the living just—[general] resurrection.

10. The judgment—reasons for it; first, to show
the full consequences of good and evil in individuals; to clear the
just; to bring shame to the wicked. Wisdom iv.

11. Secondly, to justify God—Ps. lxxii. 16-17 \footnote{'I studied that I might
know this thing. It is a labour in my sight; until I go into the
sanctuary of God, and understand concerning their last ends.'}, Ps. xlix \footnote{'Gather ye together his
saints to him ... and the heavens shall declare his justice.'}.

12. Matt. xxv. [32] \footnote{'And all nations shall
be gathered together before him: and he shall separate them one from
another, as the shepherd separateth the sheep from the goats.'}, division of good and bad, as if the separation had \emph{already}
been made.

13. No venial sins, only mortal sins [judged] at
last judgment.

14. Necessity of confessing sins now, that we may
not have to confess them then.

\xsection{December 11 Et in Spiritum Sactum—XV}

On the condescension of the Holy Ghost. Creation
implies ministration, and is the beginning of mysteries. It passes the
line, and other mysteries are but its continuation.

\xsection{December 18 Sanctam Ecclesiam Catholicam—XVI}

On the Church as means of grace—the seven
sacraments, etc. On the principle of communication of merits.

\xsection{January 4, 1850 Remissionem Peccatorum—XVII}

1. Forgiveness of sin the proclamation of the
Gospel—and a new idea. It was \emph{reminding} men of what was
necessary for them, which in the world they forget. Mortal sin, how
great an evil! hence to forgive as great an act as to create or raise
the dead.

2. The great boon—because not everywhere. Grace
everywhere, not forgiveness, though in order to forgiveness:
forgiveness on contrition.

3. Forgiveness through Christ (first of all
created natures) when on earth, through the Church; first
through baptism, next through penance (no sin the Church cannot
remit). These two [Baptism and Penance] the \emph{Sacramenta mortuorum}.

\xsection{January 8 Carnis Resurrectionem—XVIII}

1. On mysteries without number all through
revealed religion—[our] ignorance \emph{how} things are. The Creed
begins with mystery, with mystery it ends. Resurrection of the body.

2. Matter, it would seem, could not be made
spiritual (ancient philosophers) \footnote{'Among the wise men of
the heathen, as I have said, it was usual to speak slightingly and
contemptuously of the mortal body; they knew no better. They thought
it scarcely a part of their real selves, and fancied they should be in
a better condition without it. Nay, they considered it to be the cause
of their sinning; as if the soul of man were pure, and the material
body were gross, and defiled the soul. \emph{We} have been taught the
truth, viz. that sin is a disease of \emph{our minds}, of ourselves;
and that the whole of us, not body alone, but soul and body, is
naturally corrupt, and that Christ has redeemed and cleansed whatever
we are, sinful soul and body. Accordingly \emph{their} chief hope in
death was the notion that they should be rid of the body.'—\emph{Parochial
Sermons}, vol. i. p. 276, 'The Resurrection of the Body.'}—apparently the means of temptation; cause of ignorance, of
death, etc.; retards the soul.

3. Heathen philosophers of old time called the
body a prison, etc., as if the soul was pure; they made much of the
soul; hence to say the body was to be raised, to them a shocking
doctrine, Acts xvii.

4. And again, they thought sins of the flesh no
harm, \emph{because} the flesh; it disgusted them to [be told]
the body should rise again, for it implied the need of self-discipline
and mortification—the body being worth something.

5. But in truth all, soul as well as body,
imperfect—all creation imperfect; the grace which can make the soul
perfect makes the body [perfect] too.

6. Christianity, then, raises the
body—Incarnation—Mary—relics of martyrs.

7. Every one will rise with his own body—the
same body—

8. With all their members perfect, and all
defects removed;

9. Yet so far the same that the martyrs will have
their scars.

10. Immortal.

11. Four properties [of the risen
body]—impassibility, brightness (not the same to all), agility,
subtility.

12. Reflection upon glorified bodies.

\xsection{January 11 Vitam Eternam—XIX}

1. INTROD.—The
Creed begins with God, it ends with ourselves; the last articles have
reference to us.

2. Eternal life. Life means more than existence,
for the lost live.

3. It means blessedness or beatitude; and this is
called life, because there is no word which can fitly describe it; so
we must use such words as occur.

4. By blessedness is meant our greatest good, and
this from the nature of man can be nothing temporal, but must be
something eternal. If a man thought his happiness to end, or were not
sure, he would not be happy.

5. It consists in seeing God; not only seeing His
glory, or a likeness of Him, but Himself. Since it is His nature or
essence which will be seen, no likeness will do, for no likeness is
there of His essence.

6. It is seen by means of the \emph{lumen gloriae},
which raises the soul above itself—'In Thy light shall we see
light.' It is by an immediate union to God, and our intellect is
raised above itself in order to it.

7. This light of glory raises the soul above
itself. It [the soul] is what it is, but it is bathed and flooded with
a heavenly light; it puts on a divine form, so that men are called
gods. A red-hot iron, etc.

8. Such is essential blessedness—consisting in
the possession of God. The soul ever sees God present, wherever it
is—the rapturous nature of this privilege. We (most men) know so
little of intellectual joys here, that few illustrations can be given.
Most intense, yet continuous. (Happiness in itself—happiness of
convalescence—happiness of tears; soothing, etc.—happiness of
coming before the Blessed Sacrament—not happiness merely of success,
etc., as on earth, \emph{i.e}. of having gained, at
possessing—saints' raptures.)

9. So much so that the soul could dispense with
everything else—the blessed would not want friends from the earth.
Each could well bear to lose the memory of everything else for God.

10. But God has added these additions: all the
blessed will see each other, and rejoice in each other's glory,

11. And the honour of each other.

12. The glories of the heavenly palace.

13. Let this thought comfort us in the troubles
of this life, and the prospect of purgatory.

\xsection{January 3, 1858 (Octave of St. John) [Christian Knowledge]}

1. Christian knowledge is made up of four
parts—of the Creed, John iii. 11, of the doctrine of the Sacraments,
Acts i. 3, of the Ten Commandments, Matt. xxii. 37, and of prayer,
Luke xi. 1. And these four make up the teaching which is called the
Catechism.

2. We are accustomed to think the Catechism
belongs to children only, and think it does not concern grown men and
women—Catechumens—but this is not so. The Council of Trent
appointed a \emph{Catechismus ad Parochos}.

3. This is shown in the very name catechism, from
[\emph{kat\={e}chein}], to ring in the ears again and again.

4. There is great danger of our knowing only part
of what God has revealed, and danger of our forgetting what we know.
Therefore it is necessary, again and again.

5. Feeling all this deeply, I have resolved, \emph{opitulante
Deo}, to begin a set of strictly catechetical lectures in the Mass
on Sundays, viz. \emph{the four parts above mentioned}.

6. I know how careful our Fathers are in bringing
before you the truths of revelation. I know how you have profited by
their teaching. But there is the danger that some or other of the
truths should be omitted unless there is some system of catechising,
catechetical instruction, at least from time to time. You may hear one
thing three or four times and \emph{another not even once}. You may
know a great deal on some subjects, more than Catholics ordinarily do,
and not enough on others.

7. St. Paul said he had taught \emph{all the}
counsel of God. It is as necessary to know all our duty as to practise
it all; and we cannot practise it all without knowing it. What we
should aim at is knowing all and doing all.

8. In this consists \emph{perfection}. Perfection
does not lie in heroic deeds, or in great fervour, or in anything
extraordinary—many, even good men, are unequal—but in consistency.
This is what old Catholics have when good, in opposition to converts,
and therefore this congregation needs it especially.

9. And so Christian knowledge is not the knowing
about saints or about devotions and the like, though all this is
excellent, but in knowing the four things above.

10. On the Catechism as an instrument of con­verting
Protestants.

\xsection{January 10 [The Creed]}

1. The Creed begins with the word from which the
name creed is taken, \emph{credo}, I believe. Observe it is 'I
believe,' not 'I conjecture,' 'I am of opinion,' 'I know,' but 'I
believe.'

2. There are many things which we 'conjecture,'
'expect,' 'reckon,' or 'guess,' \emph{e.g}. the future generally, the
weather, the state of trade, our own prospects, health, fortune, etc.,
and we have surmises and suspicions about who are our friends, who our
enemies, and we put no great confidence in such guesses, because we
find we are often wrong. This is not what we mean when in the Creed we
say 'I believe.'

3. There are many things we have opinions about,
and strong opinions, and with very little doubt or fear, \emph{e.g}.
what we have learned by long experience of life. Experience is all in
all to many men—to the farmer, to the physician, to the navigator,
the politician, and to all men—hence it is that we trust the old;
they may not be always right, but still their experience gives them a
right to speak. Opinions may be trusted, but yet they are not
infallibly certain, because sometimes men change their opinions. This,
then, is not what is meant by 'I believe.'

4. There are many things we know, \emph{e.g}. by
our senses, by common sense, by reason. Thus we know what we see with
our eyes—all that is round about us, the world, the sky, the earth,
etc. And by common sense and reason that two and two make four, all
the points of moral conduct, the difference of virtue and vice,
conscience, etc.; [in such matters] there is no doubt or fear, but
certainty. But this is not 'I believe.'

5. What, then, is 'I believe'? It is at first
sight as uncertain, doubtful, as any of them, a conjecture, but
really more certain than [mere] knowledge. To believe is to accept as
true what we are told. What more weak, for people are continually
taken in? Why then strong? Because it is God's word—(enlarge on
this: 'Let God be true, and every man a liar' \footnote{Rom. iii. 4.})—through the Church.

6. It comes from a divine grace.

7. Exhortation to believe against appearances,
and to pray for faith.

\xsection{January 17 [Revelation—I]}

1. Last week I spoke of faith as being acceptance
of the word of God as declared by the Church, and since God is not
seen, it is by grace; hence Eph. ii. 8, 'For by grace you are saved
through faith, and that not of yourselves, for it is the gift of God';
and Heb. xi. 6, 'But without faith it is impossible to please God. For
he that cometh to God must believe that he is, and is a rewarder to
them that seek him.'

2. Now it may be objected that these two—God
and recompense—can be known by nature without grace—see Romans
vii. I answer (1) With great difficulty and many obstacles from
passion, etc. (2) They are not enough—St. Peter in Acts iv. 10-12 \footnote{'Be it known to you all
... that by the name of our Lord Jesus Christ ... even by him this man
standeth before you whole ... There is no other name under heaven
given to men, whereby we must be saved.'}; our Lord Himself, John xvii. 3 \footnote{'Now this is eternal
life, that they may know thee the only true God, and Jesus Christ,
whom thou hast sent.'}; or John iii. 16 \footnote{'God so loved the world,
as to give his only begotten Son, that whosoever believeth in him may
not perish, but may have life everlasting.'};
the Baptist, John iii. 36 \footnote{'He that believeth in
the Son hath life everlasting: but he that believeth not the Son shall
not see life.'};
St. John the evangelist, John i. 10-12 \footnote{'He was in the world,
and the world knew him not ... But as many as received him, he gave
them power to be made the sons of God, to them that believe in his
name.'}; St. Paul, Acts xvi. 31 \footnote{'But they [Paul and
Silas] said, Believe in the Lord Jesus, and thou shalt be saved, and
thy house.'}; grace, 1 Cor. xii. 3 \footnote{'No one can say the Lord
Jesus, but by the Holy Ghost.'}.

3. Now these truths [spoken of in the texts just
referred to] are not known by nature, but by God's
word—(illustrate). By natural reason we know many wonderful
things—the sciences—all those wonderful things of this
day—inventions, historical researches, antiquities dug up from the
earth, knowledge of the stars, etc., etc.—but not \emph{all} the
knowledge of men could bring us one step nearer to the knowledge of
those things which concern our salvation. These are \emph{only} known
by revelation, or by the express word of God.

4. This is what is called revelation, because the
veil taken off, or, in other words, by the express word of God.

5. The word of God—(explain)—two kinds;
Scripture and Tradition.

6. This why faith is necessary. And even what we
could know by nature (God and recompense), we must receive on faith.

7. EXHORTATION.—'Let
God be true, and every man a liar.' One man thinks one thing
difficult, another another.

\xsection{January 24 [Revelation—II]}

1. I used last week the word revelation in
connection with faith. And now I am going to explain more exactly what
it is, and why it is necessary.

2. Revelation is necessary, faith is necessary,
on account of our ignorance, which is one of the four wounds of human
nature.

3. Now does it not seem wonderful—but so it
is—that we may know so much of so many things, but so little of the
things of God? Contrast human knowledge and religious ignorance—so
many different opinions.

4. Hence, if we are to know anything of God, it
must be quite in a different way, viz. by His expressing—speaking to
us, or by His word—the word of God—and His word is revelation. For
revelation means taking a veil away. (Illustrate Isa. xxvi., 2 Cor.
iii.—veil over their heads; veil on Moses' face.) \emph{Still} a
veil, as in Blessed Sacrament—mysteries. Still, whatever we know of
unseen things is by the revealing word of God. Hence by faith, not by
sight, hearing by the word of God.

5. Faith, then, receives the revealing word of
God through the Church.

6. Now what does the word of God say? through
what and when does it speak?—through the Church, in two
ways—written Scripture, unwritten tradition.

7. Enter into Scripture and tradition; (1) parts
of Scripture; (2) parts of tradition.

8. All things we receive by faith: not only
Scripture, or only tradition, but both, for there is not a word of
Scripture, whether of prophet, evangelist, etc., nor of tradition,
which is not the voice or word of the Church—Heb. i. 1.

9. Embrace whatever God reveals as soon as you
know that it is revealed.

\xsection{January 31 [Faith]}

1. I have now explained what is meant by the word
of God, by revelation, and by faith, and why they are necessary.

2. There is great correspondence between things
of the body and of the soul. We cannot see without light; and even
with light we need eyes, and in the dark we grope our way. Now by
nature our souls are in darkness, ignorance, etc. Thus you see how it
is there is need of God's word, revelation, and faith.

3. And here you see the reason of a solemn
declaration, 'Without which there is no one can be saved.' We are
going a journey, etc.

4. Our Lord's words, John iii. 18 \footnote{'He that believeth in
him is not judged. But he that doth not believe is already judged:
because he believeth not in the name of the only begotten Son of God.'}.

5. And still more if they refuse light, John iii.
19 \footnote{'And this is the
judgment: because the light is come into the world, and men loved
darkness rather than the light: for their works were evil.'}.

6. This is one great reason why the light of
faith is necessary, because we are so ignorant.

7. Now you will say, 'Is ignorance the fault of
men in general? if so, how? if not, why are they punished with the
loss of salvation?'

8. No one is punished except for his own fault.
No one is punished except for rejecting light. God gives light all
over the earth—enough to make men advance forward.

9. Explain: from one grace to another, from one
step to another—\emph{prayer}.

10. And thus those who are in a great deal of
ignorance may be saved if they are doing their best, and their
ignorance invincible.

11. Heathen, heretics (material), may have divine
faith.

12. Who these are is secret. All we know is about
ourselves. Application to ourselves.

\xsection{February 7 [Apostles' Creed—I]}

1. Now I have explained what is meant by the word
of God, by revelation, by faith. Now after the preliminaries we come
to the Creed. You must not mind my saying the same thing over and over
again—[\emph{kat\={e}ch\={e}sis}].

2. The Creed, then, since it is received by
faith, must be revealed doctrine; so it is.

3. Next it is called Apostles'. Why? force of the
word—because nothing can be of faith, nothing is revealed, except
what comes from His apostles. No revelation since—once for all—as
sacrifice, etc.

4. Not the whole of the apostles' doctrine, but a
certain portion.

5. Why not the whole? Because it is impossible;
the Church alone can tell us the whole; it is [\emph{an illegible word
here}]. We do not for certain know till the Church tells, \emph{e.g}.
Immaculate Conception.

6. Are we not bound to believe the whole? Yes,
with implicit faith—(explain).

7. But what is put down in the Creed is definite
and simple, \emph{e.g}. first into three, then into twelve.

8. So much, because fundamental; for teaching.

9. Because easy of memory.

10. Because intelligible for strangers who ask
about it as a mark of unity.

11. So the cross—Jesus. Simple and
intelligible. 'Christian is my name, Catholic my surname.'

\xsection{February 21 [Creed—II]}

1. I said last that the Creed did not contain all
that we had to believe, but certain portions, and this is put into our
hands for various reasons.

2. First, as a badge of what our religion is:
'Christian is my name, Catholic my surname.' The sign of the
cross—so the Creed.

3. Next, as what is fundamental—which 'infants
in grace know,' 'other foundation,' etc. 'No one can say Jesus is the
Lord,' etc. \emph{Disciplina arcani}.

4. Thirdly, as being easy of memory, being only a
few clauses, a few words in each.

5. Three chief parts—twelve articles—(go
through them).

6. As to the twelve articles, there was a belief
that each apostle gave an article—thence called Apostles'
Creed; but not so, but, as I said last time, because it contains
apostolic doctrine.

7. And hence there were originally lesser
variations in the Creed in various parts of the Church, in various
countries. Rites and ceremonies vary, and though the faith never
varies, the expression of it may. We have an instance of this in the \emph{Creed
of the Mass}. (Exemplify.)

8. Each Church, then, had its own Creed, the same
except in few words, or a few articles put in or out. The Creed which
has remained and which we use as the Apostles' Creed is the Creed
always used in Rome. Saying it at the Confessional of St. Peter.

9. Another thing to be said about the Apostles'
Creed: the Nicene has additions because of heresies—\emph{Consubstantial},
etc.; the Apostles' Creed that of the Church of Rome, where heresy
never was.

10. But lest the Creed should grow too large, the
Council of Ephesus determined that it should not be added to, though
heresies arose; hence Theotokos not introduced.

\xsection{February 28 [First Article of the Creed]}

1. I have been many weeks engaged in explaining
what the Creed is, what is the need of it, and similar questions. Now
then at length we proceed to consider what it contains.

2. The first article begins, 'I believe in God,'
or, as in the Nicene, 'I believe in one God.'

3. Explain what we mean by God, viz. the one
being of beings, self-dependent, etc., all-powerful, could create
infinite worlds, each more beautiful than the one before, with all
other infinite attributes; yet what we know of Him is infinitely less
than what we do not know.

4. Here, then, you see the Creed opens in
mystery, and what is remarkable, though it is a point of faith, it is
also a point of reason. Hence the heathen philosopher asked one, two,
four, eight, etc., days to determine about God. Then in the Mass—\emph{tremunt
potestates}—the name of God not pronounced by the Jews. It is
what every child understands, who prays to God, as far as the highest
intellect.

5. Here you see what is meant by saying that
faith is against reason, viz. above, because reason itself comes to
truths which it cannot comprehend.

6. There is no article in the whole Catholic
faith more mysterious than this, which is the elementary one—nay,
which is the belief of nature, too, without grace, which Protestants
hold.

\xsection{March 7 [Creed, First Article Continued]}

1. On the awful and incomprehensible nature of
Almighty God. The sun a poor type of it, which we cannot gaze on.

2. Hence the Jews never named the name of God.
Hence in the Mass it is said \emph{tremunt potestates}.

3. Hence the Seraphim. Moses at the burning bush;
Apoc. i. 13, 17.

4. This will lead us to show how difficult it is
to speak of Him without contradictions. He is full of mystery.

5. Enumerate the mysteries contained in it. (1)
No beginning; (2) eternity by Himself; (3) then a Creator after an
eternity; (4) out of nothing; (5) ever working though ever at rest;
(6) everywhere as fully as if in one place, yet without parts; (7)
prescience; (8) knowledge of our hearts; (9) infinitely merciful and
just; (10) all-powerful yet blasphemed, etc.; (11) all-loving and
good, yet allowing sin; (12) infinite, yet personal.

6. This prepares us for those mysteries which are
of faith—the Holy Trinity.

7. As attributes in the divine nature which are
all separate, yet all one, so there is a greater and higher mystery
still, viz. three persons.

8. The Father, the Son, the Holy Ghost, each
entirely the one selfsame God, as if the others were not.

9. Still more mysterious, because we have nothing
like it on earth.

10. We should glory, not stumble, at mysteries.
All religions profess to believe in, to meditate on, Almighty God,
but,

11. How few do so, else the whole world would
become Catholic. The world generally, though they say they believe in
God, as little believes as it believes in Catholicism.

This also gave matter for a lecture for March 21:

thus—

1. The world full of mystery.

2. Much more the Maker of the World.

3. Much more the God not of reason merely, but of
revelation ... Holy Trinity.

\xsection{May 8, 1859 (Sunday afternoon lectures)}

INTROD.—I
cannot determine what I shall lecture on till I know who will come,
for the speaker speaks according to the hearers; to speak for speaking'
sake is mere human eloquence, and not practical, and this St. Philip
opposed especially. His Fathers only converse, not preach.

However, so much is certain, that all hearers
come to learn; learning implies knowledge as its object. There are two
kinds of knowledge, natural and supernatural. I shall be sure to
lecture on either natural or supernatural knowledge. One word more,
and that for the sake of spiritual profit, not mere curiosity.

\xsection{May 15 [Faith—I]}

1. INTROD.—I
said that there were two kinds of religious knowledge.

2. And each is gained in its own way. Natural by
sight and reason, supernatural by faith.

3. Reason of distinction, because we \emph{cannot}
learn what is \emph{above} nature except by faith. On natural religion
by sight and reason.

4. Natural religion is from God, sight and reason
are from God. They are good as far as they go. They do for this world,
but they never can get us to heaven.

5. Now the great bulk of mankind live merely by
sight and reason; their religion is natural religion.

6. (Here we have the comments of fact upon the
'narrow way.')

7. Describe how men live by sense and reason,
natural good feeling, good sense, honesty, uprightness, manliness. If
asked their opinion of any thing, act, or opinion or event, they will
judge merely by their common sense; then they go on to \emph{supernatural}
religion, and they still judge by reason, and so differ from each
other.

8. And they will agree in some points with the
others, not in others, hence \emph{private judgment}.

9. Now I appeal to any one, if this is not the
religion of most people. They would profess they only go by reason.
The bulk of men live and die without faith. The notion of going by
simple faith does not enter into their mind.

10. Nay, though they profess to go by Scripture
yet when there is anything they don't like, they explain it away.

11. Contrast faith and everything by faith.

12. Hence faith the foundation.

\xsection{May 22 [Faith—II]}

1. INTROD.—I
shall make this a recapitulation of the last. It is that the great
majority of revealed religion is faith, whilst other religions, the
religions of man, go by reason and conscience only.

2. It must be so, for faith is the correlative of
revelation, faith in God's words and promises.

3. Natural men may be good fathers, gentle,
simple, etc., and good soldiers, good citizens, great and good
statesmen, good kings.

4. Now first, the religion of nature, or of good
persons, who are not Catholics. Contrasted with faith, they are
benevolent, \emph{e.g}., but not simply because God tells them, but
because their disposition carries them that way; they don't think of
getting a reward. Now contrast Tobias—faith in God's word.

5. Great patriot and soldier—Nehemias iv. 9, v.
19, xiii. 14.

6. Even though they are religious men, their
belief is only a matter of opinion. Thus Protestants, saying that they
may hold what they please; they are amazed when you say that you are \emph{certain};
thus they have not the first principle of \emph{faith}.

7. Now every one who lives with no higher
religion than this comes short of eternal life. We may as well fly up
to the sky as expect by these natural powers and exercises to get to
heaven, because faith is away.

8. Now all acceptable religion is \emph{because God
has revealed this or that}. We are all apt to reason, and there is
nothing wrong in reason, so that we do not oppose faith; but the great
thing is to make an act of faith, whatever we do; to say, I believe
this or that or the other on God's word—even in those things which
we might know by nature,

9. Though not denying that those who are not
Catholics may have this divine faith; but it is only as they have it
that they have any chance of salvation.

\xsubsection{Or rather thus:}

1. INTROD.—Importance
of making act of faith. There are two things in religion—doctrines
to be accepted and commands to be obeyed; doctrines may be taken
by \emph{reason}, commands by conscience. We must take both not by
reason or conscience, but by FAITH.

2. Because faith must be the foundation of
everything, and unless we begin with it, nothing is acceptable.

3. \emph{God has spoken}—Rom. x. 17; 1 Thess.
ii. 13. There are many things which we know by nature. God has said
these over again, these and many new things in Revelation, but in
order that they should be acceptable, we must accept on faith even
those things which we know by nature. Whether reason is for, or
scruples at doctrines, we must take them on faith.

4. \emph{E.g}. the being of a God, immortality of
the soul, future judgment, etc.

5. And thus we learn to take others also on
faith, as the word of God. God has spoken.

6. INSTANCES.—Tobias,
not rich, a benevolent man, but with faith.

7. Job, rich, abundant alms, etc.

8. Nehemias, a statesman, patriot, commander.

9. Esau, the instance of a man without faith,
contrasted with Jacob, who had.

10. OBJECTION.—Protestants
often good and religious, and seem really to live by faith.

11. \emph{Distinguo}. Do they really go by faith,
not by private judgment? Do they really believe God has spoken this or
that definite doctrine or command, or do they believe doctrines merely
so far as reason teaches them, and commands as far as conscience?

12. But if so, very well, invincible ignorance
(draw out).

\xsection{May 29 [Faith—III]}

1. I have said, nothing without faith as its
foundation—Heb. xi. 6 \footnote{'But without faith it is
impossible to please God.'},
Eph. ii. 8 \footnote{'For by grace you are
saved through faith.'}; faith
implies an external message—Rom. x. 14 \footnote{'How then shall they
call on him in whom they have not believed? or how shall they believe
him of whom they have not heard? and how shall they bear without a
preacher?'}, 1 Thess. ii. 13 \footnote{'When you received of us
the word of the hearing of God, you received it not as the word of
men,' etc.}.

2. Yet, as I said, it is impossible to go into
the world without seeing that the idea of taking one's doctrine from
an external authority does not enter into their minds. It is always 'I
think.' This is what is meant by private judgment—though Scripture,
yet they put their own sense on Scripture; they take these books,
reject those, etc.

3. This is a most fearful consideration,
considering we are saved by faith. And observe, it is quite
independent of the question of what \emph{is} the true doctrine, what \emph{is}
the true Church. You see most men do not GO
THE RIGHT WAY. It is a previous question. They don't go the way
of faith. From this it is plain, to go no further, that none but
Catholics are in the right way, because \emph{they alone} go by faith.

4. Now in this awful prospect the question
arises, Does no one else go by faith? Does no Protestant go by faith?

5. We can only answer by what we see. Well, they \emph{profess}
not to go by faith. If they \emph{do} go by faith, at least they do
not know it. Alas, it does seem as if we must say that the majority do
not go by faith.

6. Do any? I trust they do. I trust there is a
remnant all over the world who \emph{do} go by faith, and who so far
are in the way of salvation, or rather, towards salvation. And in
explanation how this is, I shall clear my meaning up more fully.

7. But first I shall answer an objection, viz. If
they go by faith, why do they not join the Catholic Church, in which
alone God speaks? It is said, 'My sheep hear My voice.' I answer, they
are out of the hearing of the Catholic Church, and therefore are in
what is called invincible ignorance.

8. Now I will describe the state of such persons
all the world over. Our Lord died for all, grace is given to all. Most
men seem to profit nothing at all by it, but there are those who
profit, \emph{e.g}.

9. Conscience—there are two ways of regarding
conscience; \emph{one} as a mere sort of sense of propriety, a taste
teaching us to do this or that, the \emph{other} as the echo of God's
voice. Now all depends on this distinction—the first way is not of
faith, and the second is of faith.

10. Characteristics of the first way—connected
with pride. The proud will call the other kind superstitious. A person
makes himself his own centre. He says, I shall hold just what seems to
me, what my moral sense tells me, etc. The other considers it the
voice of God, obeys it as such, a call to look out for more light.

11. Development of the idea of God, of faith in
God, and of the feeling of the necessity of God's speaking in order to
their salvation.

12. Hence to the evidences—the visible things
of God, etc., etc.—history, providences, experiences.

13. A person may be a heathen—Mahometan, etc.,
etc.—and yet have this real faith in God, and so far he is on the
way towards salvation. Their ignorance is involuntary and invincible, \emph{i.e}.
not their fault.

14. On being led on. Thus heathens, like the wise
men, led into truth—the Ethiopians to Judaism; Abimelech's dream;
Pharaoh's dreams. Job iv., spirit \footnote{'And when a spirit
passed before me,' etc.}; Job a just man. Magi, Ethiopians, the [\emph{eusebeis}] in
the Acts. Means which the Almighty used before the coming of Christ.

\emph{N.B}.—As just men existed before Christ
came, why not at a distance from the Church? for what the former is of
time, so just men among the \emph{heathen} is of space.

15. And thus St. Thomas said, 'An angel will
speak from heaven rather than a soul fail.' On Christ's sheep hearing
His voice; thus it is a test whether persons come on towards the
Church when they know about it; invincible ignorance the only excuse.
On the difference between death overtaking the shily-shallying, who
are not seeking, and on the earnest inquirer in invincible ignorance
dying before he is a Catholic.

16. Case of Dr. A—in his dreamy state, before
death learning the truth.

\xsubsection{N.B.—Principles in the above.}

1. It is part of the same mystery why death comes
before inquirers are led into the truth, now as of old, for both
depend on the parallel mystery why (1) our Lord did not come from the
beginning of the world; (2) why the Gospel is not preached all over
the earth. There is [\emph{i.e}. would have been] nothing stranger in
the Ethiopian dying before Philip came to him, than in a Jew of
Solomon's day dying before the Gospel was preached; neither had
baptism.

2. What is faith before the revealed dogma is
known, is superstition after, for God has now superseded the natural
ways of seeking Him.

3. 'My sheep hear Me.' Therefore it is only when
there is invincible ignorance that this can be—being led on into the
Church is the test.

\xsection{July 3 [Faith—IV]}

1. INTROD.—What
I have been saying is this, that even heathen (all men) to enter upon
the road that leads to heaven, must live by faith; by faith even as to
those things which they know by reason.

2. And this is so, because God would take us out
of ourselves and make us depend on Him—not make ourselves our own
centre; we must make God our centre.

3. Analogy of Nature—solar system, monarchy in
society, etc. Yet this difference, that in this world there are many
ranks, etc., intermediate between the centre and ourselves, but
as to religion, we every one depend on one centre alone—God.

4. The wishing to have ourselves our own centre
is pride, the sin of Satan. (Enlarge on it.)

5. Instances of faith, inquiry, doubt, a want of
faith. (1) Inquiry—the child Samuel. (2) Nobleman in 4 Kings vii.,
who would not believe Eliseus; unbelief of St. Thomas. (3) Zachary;
doubt, and Nicodemus. (4) Our Lady's and St. Paul's faith. A fifth
state, weak faith of Gideon: 'Lord, I believe.' \footnote{This section might be
rewritten thus:—Instances of (1) inquiry—the child Samuel ['Speak
Lord,' etc.]; (2) want of faith—the nobleman who would not believe
Eliseus, 4 Kings vii. ['If the Lord should make flood-gates in heaven,
can that possibly be which thou sayest'], and the unbelief of St.
Thomas; (3) doubt—Zachary [Luke i. 18], and Nicodemus [John iii.];
(4) faith—our Lady's and St. Paul's faith. A fifth state, weak
faith, as in the case of Gideon [Judges vi. 36-40]. and the father of
the boy possessed by the dumb spirit—'Lord I do believe; help my
unbelief' [Mark ix. 24].}.

\xsection{August 14 On Love}

1. INTROD.—(1)
On Love as not external to the Church (\emph{i.e}. state of grace) as
faith and hope may be. (2) As not in those who fall from grace, while
faith and hope remain. As not the love of concupiscence or hope, nor
gratitude, its object being the beauty of God.

Now I shall show how love comes after faith,
through a distinct grace. Younger sons in Scripture—Jacob, not
Esau—\emph{vide} St. Francis de Sales.

2. On love. Remains of love in nature, as shown
by the drawings of the heart, by people seeking comfort in religion,
in trouble; and this left in order that grace may work \emph{with}
nature, not against nature, \emph{when} it works. Also as a sort of
claim of God upon us, as if He marked us as His property.

On love being produced from faith through
meditation.

3. By nature [man] has far more to do with other
attributes of God—fear, etc., etc. For instance, I have said that
conscience is the means of faith, but it teaches justice principally,
which is the object of fear, not of love. Again, though there is great
goodness in God's providence, yet in action the marks of particular
providence not so obvious.

4. But it is the objects of Christian faith which
cause love. Go through them minutely. At first sight original sin
might be a doctrine which drives us from God. \emph{O felix culpa},
atonement—much more still particular election. Then Mass and the
Holy Eucharist, etc. Faith, you see, is \emph{the} thing, or the \emph{only}
thing necessary as a means of love.

5. Faith leads to love \emph{through} meditation.
The things of faith, \emph{e.g}. the whole doctrine of election is
hard, but when once embraced it has its reward in its powers of
kindling love.

Doctrine of election: (1) This little globe out
of the whole world; (2) not angels but men chosen; (3) Old Testament
elections, the younger for the older; (4) we chosen out of the world
to be Catholics.

6. And though this in the first place gratitude,
not pure love, yet, since all God's dealings to us are so
admirable and glorious, love is kindled at the same time.

7. All this is summed up, especially in the
Gospels, so that as the Bible is the instrument of hope, so is the
Gospels of love.

\xsection{June 3, 1860 Confirmation}

I conceive that St. Thomas means that the \emph{dona
Spiritus Sancti} have the same relation to the motions of grace
towards the supernatural end of man which moral virtue has to the
motions of reason towards the natural end. They \emph{dispose} the
mind that it may be well moved by the Holy Ghost, as the virtues
perfect the appetite that it may be well moved by reason. Also that he
means by wisdom and understanding two speculative gifts, the latter
apprehension, the former judgment, speculative apprehension or
understanding being what Doyle's Catechism calls knowledge
(comprehension) of mysteries, and speculative understanding or wisdom,
the perfect appreciation (comprehension, grasp) of all subjects of
religion, though Doyle calls it a gift of directing our whole lives
and actions to God's honour.

By knowledge and counsel, two practical
gifts—knowledge being a \emph{judicium}, viz. an insight into duty
generally, though Doyle says, 'a gift by which we know, etc., the will
of God'; and counsel being an \emph{apprehensio}, or what Doyle calls
'a gift by which we discover the grounds of the duty,' etc.

I should say, by which we take in all things
correctly as they are.

Wisdom—the estimation of all things rightly
with reference to \emph{our end}.

Understanding—knowledge, mysteries.

Knowledge—knowledge of \emph{means} for our
spiritual welfare.

Council—or prudence, judging of things rightly.

Piety or godliness—directs \emph{appetitus in
alterum}.

Fortitude—against fear of world, etc.

Fear—against concupiscence.

\xchapter{Notes}

\xsection{Notes}




1. Often from weakness of the will, which in
matters of practical conduct makes a man scrupulous, unable to brush
aside futile difficulties, and the like—victims, as they would be
called, of idle fears.

Return to text


2. In the evidences for
religion 'non excluditur omnis dubitandi possibilitas, sed solum
dubitandi prudenter, et ideo mens sub voluntatis motione relinquitur,
ut deliberata formidine deposita, firmiter adhaereat ei quod videt non
posse nisi irrationabiliter in dubium revocari.'

Return to text


3. 'Brother, a virtue of
charity sets at rest our will, which makes us wish that only which we
have, and lets us not thirst for aught else. If we desired to be more
on high, our desires would be out of harmony with the will of Him who
distributes us here.'—\emph{Paradiso}, iii. 70-76; Butler's
translation.

Return to text
